% =====================================================================
% CORRECTION RECORD (source comments; deliberately not in the body)
% Claims made during drafting and later withdrawn.  The body states the
% corrected positions directly rather than narrating the corrections,
% since drafting history is not the reader's concern; the record is kept
% here so that it is not lost.
%
%  1. rho = 1 was said to be the threshold at which the background drops
%     out at negative x.  FALSE -- the crest is bounded, so any rho > 0
%     gives trough dominance there.  rho > 1 is what makes the crest
%     dominate at large POSITIVE x, i.e. Constraint C1.  Body: Rem.
%     "What rho = 1 is, and is not".
%  2. The exponent 1/4 was stated to require rho > 1.  It holds for all
%     rho > 0; the dominant channel is exponential in x either way.
%  3. [body-removed] The polar decomposition x = sigma t^{3/2} R cos(Theta) was adopted
%     with Theta named "technoanatomy".  Withdrawn: R, Theta are
%     Rayleigh/uniform/independent only because the normalised pair is
%     isotropic Gaussian, which fails under a transient driver, under
%     rho < 0, and under conditioning.  Body now forecloses it as an
%     alternative rather than reporting the withdrawal.  SUPERSEDED:
%     the passage was cut entirely.  It had no independent use --- the
%     exact-half theorem is stated on x directly --- so introducing the
%     decomposition solely to reject it was drafting history in
%     disguise, and the framing that it is what a reader would first
%     try was an unsupported claim about reader psychology.
%  4. The KL comparison 1.16 vs 0.94 nats was offered as strength versus
%     anatomy.  Those are two halves of one Gaussian coordinate.
%     Replaced by sign-of-x versus rho.
%  5. "No terminal stage exists" was asserted as the strongest
%     anti-Kardashev result.  True of the driftless model only; the
%     two-condition theorem gives a gated terminal stage.
%  6. Reciprocity (E3) was reported emergent from E1 and E2, on runs
%     that had included it.  Re-tested without: asymmetries reach the
%     maximum 2.0.  It is imposed, not emergent.
%  7. Mutual singularity of Q and P was listed as unconditional.  It is
%     critical-regime only.
%  8. It was conceded that under rho < 0 the extinctionist is closer to
%     right.  FALSE: halting at x* dominates x = 0 for every rho < 0,
%     and both carry P(viable) = 0 in any case.
%  9. The Lindy result was presented as novel.  It is established prior
%     work (Balbi & Cirkovic; Balbi & Grimaldi; Kipping et al.; Ord);
%     only the derived exponent is added.  NOTE: this one is a citation
%     obligation, not a drafting retraction, and REMAINS in the body.
% 10. Two-boundary physics was proposed as a neighbour for
%     Past-Alignment.  Withdrawn as a false analogy -- those condition
%     on complete physical states with dynamics between them.  The
%     indiscriminate time principle is the correct internal answer.
% =====================================================================

% =====================================================================
%  GOLDEN REASONING -- MASTER DOCUMENT
%  Sections 1-2 assembled 2026-08-01.
%
%  SOURCES.  This file is the concatenation of two standalone section
%  files, each of which carries its own full provenance and correction
%  ledger in its header:
%      section-soft-horizons.tex   -> Section 1
%      section-two-channel.tex     -> Section 2 (revision 2)
%  Those ledgers are NOT reproduced here.  Keep the standalone files:
%  they are the authoritative record of who originated what, of the
%  correction log, and of citation-verification status.  Consult them
%  before posting.
%
%  BIBLIOGRAPHY.  The two sections shared four references (AS15, GJW99,
%  VMS79, VY85); the merged list is their union, 35 entries, deduped.
%  The file is SELF-SWITCHING: if biblatex is installed it is used,
%  with the .bib written at compile time by filecontents; otherwise a
%  manual thebibliography is used instead.  Both are generated from the
%  same 35-key list and their key sets were checked equal.  To force
%  one path, edit the \IfFileExists test below.
%
%  NOTE ON VERIFICATION.  The fallback (thebibliography) path is
%  compile-verified.  The biblatex path is NOT -- biblatex and biber
%  were unavailable in the environment where this file was assembled.
%  If biber errors, the usual causes are (a) needing a biber run
%  between pdflatex runs, or (b) backend=biber vs bibtex mismatch.
%
%  CROSS-REFERENCES.  Section 2's references to "Section One" have been
%  converted to \ref{sh:sec} in this merged file.  Label prefixes are
%  disjoint: sh: for Section 1, tc: for Section 2.
% =====================================================================

\documentclass[11pt]{article}

\usepackage[margin=1.1in]{geometry}
\usepackage{amsmath,amssymb,amsthm,mathtools}
\usepackage{booktabs}
\usepackage{array}
\usepackage{enumitem}
\usepackage{xcolor}
\definecolor{grlink}{rgb}{0.10,0.20,0.55}
\usepackage[colorlinks=true,linkcolor=grlink,citecolor=grlink,urlcolor=grlink]{hyperref}

% ---- bibliography database (written at compile time) ----------------
\begin{filecontents*}[overwrite]{golden-reasoning.bib}
@article{AD13, author={Aurzada, Frank and Dereich, Steffen}, title={Universality of the asymptotics of the one-sided exit problem for integrated processes}, journal={Ann. Inst. Henri Poincar\'e Probab. Stat.}, volume={49}, year={2013}, pages={236--251}}
@incollection{AS15, author={Aurzada, Frank and Simon, Thomas}, title={Persistence probabilities and exponents}, booktitle={L\'evy Matters V}, series={Lecture Notes in Mathematics}, volume={2149}, publisher={Springer}, address={Cham}, year={2015}, pages={183--224}}
@article{BBM05, author={Ben Arous, G\'erard and Bogachev, Leonid V. and Molchanov, Stanislav A.}, title={Limit theorems for sums of random exponentials}, journal={Probab. Theory Relat. Fields}, volume={132}, number={4}, year={2005}, pages={579--612}}
@article{BMS13, author={Bray, Alan J. and Majumdar, Satya N. and Schehr, Gr\'egory}, title={Persistence and first-passage properties in nonequilibrium systems}, journal={Adv. Phys.}, volume={62}, number={3}, year={2013}, pages={225--361}}
@article{Bha82, author={Bhattacharya, Rabi N.}, title={On the functional central limit theorem and the law of the iterated logarithm for Markov processes}, journal={Z. Wahrscheinlichkeitstheorie verw. Gebiete}, volume={60}, year={1982}, pages={185--201}}
@book{Bov06, author={Bovier, Anton}, title={Statistical Mechanics of Disordered Systems: A Mathematical Perspective}, publisher={Cambridge University Press}, year={2006}}
@book{CL67, author={Cram\'er, Harald and Leadbetter, M. Ross}, title={Stationary and Related Stochastic Processes}, publisher={Wiley}, address={New York}, year={1967}}
@article{DW15, author={Denisov, Denis and Wachtel, Vitali}, title={Exit times for integrated random walks}, journal={Ann. Inst. Henri Poincar\'e Probab. Stat.}, volume={51}, number={1}, year={2015}, pages={167--193}}
@article{Der80, author={Derrida, Bernard}, title={Random-energy model: limit of a family of disordered models}, journal={Phys. Rev. Lett.}, volume={45}, year={1980}, pages={79--82}}
@article{Der81, author={Derrida, Bernard}, title={Random-energy model: an exactly solvable model of disordered systems}, journal={Phys. Rev. B}, volume={24}, number={5}, year={1981}, pages={2613--2626}}
@book{Dur19, author={Durrett, Rick}, title={Probability: Theory and Examples}, edition={5}, publisher={Cambridge University Press}, year={2019}}
@book{FW12, author={Freidlin, Mark I. and Wentzell, Alexander D.}, title={Random Perturbations of Dynamical Systems}, edition={3}, publisher={Springer}, year={2012}}
@article{GJW99, author={Groeneboom, Piet and Jongbloed, Geurt and Wellner, Jon A.}, title={Integrated Brownian motion, conditioned to be positive}, journal={Ann. Probab.}, volume={27}, number={3}, year={1999}, pages={1283--1303}}
@article{Gol71, author={Goldman, Malcolm}, title={On the first passage of the integrated Wiener process}, journal={Ann. Math. Statist.}, volume={42}, year={1971}, pages={2150--2155}}
@article{Gom25, author={Gompertz, Benjamin}, title={On the nature of the function expressive of the law of human mortality, and on a new mode of determining the value of life contingencies}, journal={Philos. Trans. R. Soc. Lond.}, volume={115}, year={1825}, pages={513--583}}
@article{Hor67, author={H\"ormander, Lars}, title={Hypoelliptic second order differential equations}, journal={Acta Math.}, volume={119}, year={1967}, pages={147--171}}
@article{IW94, author={Isozaki, Yasuki and Watanabe, Shinzo}, title={An asymptotic formula for the Kolmogorov diffusion and a refinement of Sinai's estimates for the integral of Brownian motion}, journal={Proc. Japan Acad. Ser. A}, volume={70}, number={9}, year={1994}, pages={271--276}}
@book{KT81, author={Karlin, Samuel and Taylor, Howard M.}, title={A Second Course in Stochastic Processes}, publisher={Academic Press}, year={1981}}
@article{Kol34, author={Kolmogorov, Andrei N.}, title={Zuf\"allige Bewegungen (zur Theorie der Brownschen Bewegung)}, journal={Ann. of Math. (2)}, volume={35}, year={1934}, pages={116--117}}
@article{Lac91, author={Lachal, Aim\'e}, title={Sur le premier instant de passage de l'int\'egrale du mouvement brownien}, journal={Ann. Inst. H. Poincar\'e Probab. Statist.}, volume={27}, year={1991}, pages={385--405}}
@book{Lan04, author={Lando, David}, title={Credit Risk Modeling: Theory and Applications}, publisher={Princeton University Press}, year={2004}}
@article{MK04, author={Molchan, George and Khokhlov, Andrei}, title={Small values of the maximum for the integral of fractional Brownian motion}, journal={J. Stat. Phys.}, volume={114}, year={2004}, pages={923--945}}
@article{MP01, author={Milevsky, Moshe A. and Promislow, S. David}, title={Mortality derivatives and the option to annuitise}, journal={Insurance Math. Econom.}, volume={29}, number={3}, year={2001}, pages={299--318}}
@article{MV12, author={M\'el\'eard, Sylvie and Villemonais, Denis}, title={Quasi-stationary distributions and population processes}, journal={Probab. Surv.}, volume={9}, year={2012}, pages={340--410}}
@article{Mar49, author={Maruyama, Gisiro}, title={The harmonic analysis of stationary stochastic processes}, journal={Mem. Fac. Sci. Ky\=ush\=u Univ. Ser. A}, volume={4}, year={1949}, pages={45--106}}
@article{McK63, author={McKean, Jr., Henry P.}, title={A winding problem for a resonator driven by a white noise}, journal={J. Math. Kyoto Univ.}, volume={2}, year={1963}, pages={227--235}}
@article{Mol99, author={Molchan, George M.}, title={Maximum of a fractional Brownian motion: probabilities of small values}, journal={Comm. Math. Phys.}, volume={205}, year={1999}, pages={97--111}}
@article{RVY06, author={Roynette, Bernard and Vallois, Pierre and Yor, Marc}, title={Some penalisations of the Wiener measure}, journal={Jpn. J. Math.}, volume={1}, year={2006}, pages={263--290}}
@book{RY99, author={Revuz, Daniel and Yor, Marc}, title={Continuous Martingales and Brownian Motion}, edition={3}, publisher={Springer}, year={1999}}
@article{Ric44, author={Rice, Stephen O.}, title={Mathematical analysis of random noise}, journal={Bell System Tech. J.}, volume={23}, year={1944}, pages={282--332}}
@article{Ric45, author={Rice, Stephen O.}, title={Mathematical analysis of random noise (concluded)}, journal={Bell System Tech. J.}, volume={24}, year={1945}, pages={46--156}}
@article{Sin92, author={Sinai, Yakov G.}, title={Distribution of some functionals of the integral of a random walk}, journal={Theoret. and Math. Phys.}, volume={90}, year={1992}, pages={219--241}, note={translated from Teoret. Mat. Fiz. \textbf{90} (1992), no. 3, 323--353}}
@article{VMS79, author={Vaupel, James W. and Manton, Kenneth G. and Stallard, Eric}, title={The impact of heterogeneity in individual frailty on the dynamics of mortality}, journal={Demography}, volume={16}, year={1979}, pages={439--454}}
@article{VY85, author={Vaupel, James W. and Yashin, Anatoli I.}, title={Heterogeneity's ruses: some surprising effects of selection on population dynamics}, journal={Amer. Statist.}, volume={39}, year={1985}, pages={176--185}}
@article{Wah78, author={Wahba, Grace}, title={Improper priors, spline smoothing and the problem of guarding against model errors in regression}, journal={J. Roy. Statist. Soc. Ser. B}, volume={40}, year={1978}, pages={364--372}}
@book{Bos02, author={Bostrom, Nick}, title={Anthropic Bias: Observation Selection Effects in Science and Philosophy}, publisher={Routledge}, address={New York}, year={2002}}
@article{Car83, author={Carter, Brandon}, title={The anthropic principle and its implications for biological evolution}, journal={Philos. Trans. R. Soc. Lond. A}, volume={310}, year={1983}, pages={347--363}}
@article{Dys60, author={Dyson, Freeman J.}, title={Search for artificial stellar sources of infrared radiation}, journal={Science}, volume={131}, number={3414}, year={1960}, pages={1667--1668}}
@article{Got93, author={Gott III, J. Richard}, title={Implications of the Copernican principle for our future prospects}, journal={Nature}, volume={363}, number={6427}, year={1993}, pages={315--319}}
@misc{Han98, author={Hanson, Robin}, title={The Great Filter --- Are We Almost Past It?}, howpublished={manuscript}, year={1998}}
@article{Kar64, author={Kardashev, Nikolai S.}, title={Transmission of information by extraterrestrial civilizations}, journal={Soviet Astronomy}, volume={8}, year={1964}, pages={217--221}, note={Russian original, Astron. Zh. \textbf{41} (1964), 282--287}}
@misc{SDO18, author={Sandberg, Anders and Drexler, Eric and Ord, Toby}, title={Dissolving the Fermi paradox}, howpublished={preprint}, year={2018}}
@book{Tal12, author={Taleb, Nassim Nicholas}, title={Antifragile: Things That Gain from Disorder}, publisher={Random House}, address={New York}, year={2012}}
@book{WB00, author={Ward, Peter D. and Brownlee, Donald}, title={Rare Earth: Why Complex Life Is Uncommon in the Universe}, publisher={Copernicus}, address={New York}, year={2000}}
@article{Nexus, author={Benander, Ajax}, title={Nexic Reasoning: Defining a Generalized Calculus Over Anthropic Parameters}, year={2024}, month={December}, doi={10.31219/osf.io/jvkcp}, url={https://philpapers.org/rec/BENNRD}, note={Also at osf.io/jvkcp}}
@misc{Ben26Found, author={Benander, Ajax}, title={The Copernican Method Against the Copernican Principle: Anthropic Typicality over Partial Information States}, howpublished={Preprint}, year={2026}}
@misc{Ben26Optimal, author={Benander, Ajax}, title={The Optimal Earth: A Paradigmatic Re-Rendition of Natural History from Coupled Improbability}, howpublished={Working draft}, year={2026}}
@article{BC21, author={Balbi, Amedeo and {\'C}irkovi{\'c}, Milan M.}, title={Longevity is the key factor in the search for technosignatures}, journal={Astron. J.}, volume={161}, year={2021}, pages={222}, doi={10.3847/1538-3881/abec48}}
@article{BG24, author={Balbi, Amedeo and Grimaldi, Claudio}, title={Technosignatures longevity and Lindy's law}, journal={Astron. J.}, volume={167}, number={3}, year={2024}, pages={119}, doi={10.3847/1538-3881/ad217d}}
@misc{Kip20, author={Kipping, David and Frank, Charles and Scarlata, Claudia}, title={Contact inequality: first contact will likely be with an older civilization}, howpublished={preprint}, year={2020}, note={arXiv:2010.12358}}
@misc{Ord23, author={Ord, Toby}, title={The Lindy Effect}, howpublished={preprint}, year={2023}, note={arXiv:2308.09045}}
@article{Crary23, author={Crary, Alice}, title={The toxic ideology of longtermism}, journal={Radical Philosophy}, volume={2}, number={14}, year={2023}, pages={49--57}}
@misc{GM21, author={Greaves, Hilary and MacAskill, William}, title={The case for strong longtermism}, howpublished={Global Priorities Institute Working Paper 5-2021, University of Oxford}, year={2021}}
@misc{Torres21, author={Torres, {\'E}mile P.}, title={Against longtermism}, howpublished={Aeon}, year={2021}}
@article{Bostrom19, author={Bostrom, Nick}, title={The vulnerable world hypothesis}, journal={Global Policy}, volume={10}, number={4}, year={2019}, pages={455--476}, doi={10.1111/1758-5899.12718}}
@book{Tainter88, author={Tainter, Joseph A.}, title={The Collapse of Complex Societies}, publisher={Cambridge University Press}, year={1988}}
@article{EF24, author={Engler, John-Oliver and Fischer, Jan Niklas}, title={Drawing blanks and winning: quantifying global catastrophic risk associated with human ingenuity}, journal={Sustainable Futures}, volume={7}, year={2024}, pages={100165}}
@book{Cooper19, author={Cooper, Keith}, title={The Contact Paradox: Challenging Our Assumptions in the Search for Extraterrestrial Intelligence}, publisher={Bloomsbury Sigma}, address={London}, year={2019}}
@incollection{Plutarch, author={Plutarch}, title={Whether Land or Sea Animals Are Cleverer ({D}e sollertia animalium)}, booktitle={Moralia}, note={\S7; the saying is attributed there to Bion of Borysthenes}}
@article{Hei46, author={Heider, Fritz}, title={Attitudes and cognitive organization}, journal={J. Psychol.}, volume={21}, year={1946}, pages={107--112}}
@article{CH56, author={Cartwright, Dorwin and Harary, Frank}, title={Structural balance: a generalization of Heider's theory}, journal={Psychol. Rev.}, volume={63}, year={1956}, pages={277--293}}
@article{MKKS11, author={Marvel, Seth A. and Kleinberg, Jon and Kleinberg, Robert D. and Strogatz, Steven H.}, title={Continuous-time model of structural balance}, journal={Proc. Natl. Acad. Sci. USA}, volume={108}, year={2011}, pages={1771--1776}}
@article{AKR05, author={Antal, Tibor and Krapivsky, Paul L. and Redner, Sidney}, title={Dynamics of social balance on networks}, journal={Phys. Rev. E}, volume={72}, year={2005}, pages={036121}}
@book{Wal79, author={Waltz, Kenneth N.}, title={Theory of International Politics}, publisher={McGraw-Hill}, year={1979}}
@book{Wal87, author={Walt, Stephen M.}, title={The Origins of Alliances}, publisher={Cornell University Press}, year={1987}}
@article{Aethus, author={Benander, Ajax}, title={Aethic Reasoning: A Comprehensive Solution to the Quantum Measurement Problem}, year={2024}, month={November}, url={https://philpapers.org/rec/BENARA-5}, note={PhilSci-Archive eprint 25713}}
@misc{AethusBrief, author={Benander, Ajax}, title={Aethic Reasoning in Brief: From a Counterfactual Semantics to the Quantum Measurement Problem}, howpublished={Preprint}, year={2026}, url={https://philpapers.org/rec/BENARI-3}, note={The compressed companion; PhilPapers record BENARI-3}}
@misc{BenNewcomb, author={Benander, Ajax}, title={Determinacy without tense: observer-indexed determinacy and Newcomb's problem}, howpublished={companion manuscript}}
@misc{GoldenScience22, author={Benander, Ajax}, title={Golden Science}, howpublished={unpublished manuscript}, year={2022}, note={dated August 2022; archived copy November 2022}}
@book{kripke1980naming, title={Naming and Necessity}, author={Kripke, Saul A.}, year={1980}, publisher={Harvard University Press}, address={Cambridge, Massachusetts}}
@book{Pomeranz00, author={Pomeranz, Kenneth}, title={The Great Divergence: China, Europe, and the Making of the Modern World Economy}, publisher={Princeton University Press}, year={2000}}
@book{Allen09, author={Allen, Robert C.}, title={The British Industrial Revolution in Global Perspective}, publisher={Cambridge University Press}, year={2009}}
@book{Diamond97, author={Diamond, Jared}, title={Guns, Germs, and Steel: The Fates of Human Societies}, publisher={W. W. Norton}, year={1997}}
@book{Arrian, author={Arrian}, title={Anabasis of Alexander}, note={1.15.8}}
@incollection{PlutarchAlex, author={Plutarch}, title={Life of Alexander}, booktitle={Parallel Lives}, note={\S16}}
\end{filecontents*}
\begin{filecontents*}[overwrite]{gr-manualbib.tex}
\begin{thebibliography}{99}\small

\bibitem[AD13]{AD13} F.~Aurzada and S.~Dereich,
\emph{Universality of the asymptotics of the one-sided exit problem for
integrated processes}, Ann.\ Inst.\ Henri Poincar\'e Probab.\ Stat.\
\textbf{49} (2013), 236--251.

\bibitem[Bri]{Britannica_coal} \emph{Coal: world distribution of coal},
Encyclop\ae dia Britannica, consulted 2026.
\bibitem[Mag]{Magoon_coal} M.~Magoon, \emph{The geography of coal
outside Europe}, From Poverty to Progress, 2026.
\bibitem[Aet24]{Aethus} A.~Benander, \emph{Aethic Reasoning: A
Comprehensive Solution to the Quantum Measurement Problem}, November 2024;
PhilPapers BENARA-5, PhilSci-Archive eprint 25713.

\bibitem[AetB]{AethusBrief} A.~Benander, \emph{Aethic Reasoning in
Brief: From a Counterfactual Semantics to the Quantum Measurement
Problem}, preprint, 2026, PhilSci-Archive eprint 30439.

\bibitem[AKR05]{AKR05} T.~Antal, P.~L.~Krapivsky, and S.~Redner, \emph{Dynamics of social balance on networks}, Phys.\ Rev.\ E \textbf{72} (2005), 036121.

\bibitem[AS15]{AS15} F.~Aurzada and T.~Simon,
\emph{Persistence probabilities and exponents}, in: L\'evy Matters V,
Lecture Notes in Mathematics 2149, Springer, Cham, 2015, pp.~183--224.

\bibitem[BBM05]{BBM05} G.~Ben~Arous, L.~V.~Bogachev, and
S.~A.~Molchanov, \emph{Limit theorems for sums of random exponentials},
Probab.\ Theory Relat.\ Fields \textbf{132} (2005), no.~4, 579--612.

\bibitem[BC21]{BC21} A.~Balbi and M.~M.~\'Cirkovi\'c, \emph{Longevity is
the key factor in the search for technosignatures}, Astron.\ J.\
\textbf{161} (2021), 222.

\bibitem[Ben24]{Nexus} A.~Benander, \emph{Nexic Reasoning: Defining
a Generalized Calculus Over Anthropic Parameters}, 2024,
\texttt{doi:10.31219/osf.io/jvkcp}.

\bibitem[Ben26a]{Ben26Found} A.~Benander, \emph{The Copernican Method
Against the Copernican Principle: Anthropic Typicality over Partial
Information States}, preprint, 2026.

\bibitem[Ben26b]{Ben26Optimal} A.~Benander, \emph{The Optimal Earth: A
Paradigmatic Re-Rendition of Natural History from Coupled
Improbability}, working draft, 2026.

\bibitem[BenN]{BenNewcomb} A.~Benander, \emph{Determinacy without
tense: observer-indexed determinacy and Newcomb's problem}, companion
manuscript.

\bibitem[BG24]{BG24} A.~Balbi and C.~Grimaldi, \emph{Technosignatures
longevity and Lindy's law}, Astron.\ J.\ \textbf{167} (2024), no.~3,
119.

\bibitem[Bha82]{Bha82} R.~N.~Bhattacharya,
\emph{On the functional central limit theorem and the law of the iterated
logarithm for Markov processes},
Z.~Wahrscheinlichkeitstheorie verw.\ Gebiete \textbf{60} (1982), 185--201.

\bibitem[BMS13]{BMS13} A.~J.~Bray, S.~N.~Majumdar, and G.~Schehr,
\emph{Persistence and first-passage properties in nonequilibrium
systems}, Adv.\ Phys.\ \textbf{62} (2013), 225--361.

\bibitem[Bos02]{Bos02} N.~Bostrom, \emph{Anthropic Bias: Observation
Selection Effects in Science and Philosophy}, Routledge, New York, 2002.

\bibitem[Bos19]{Bostrom19} N.~Bostrom, \emph{The vulnerable world
hypothesis}, Global Policy \textbf{10} (2019), no.~4, 455--476.

\bibitem[Bov06]{Bov06} A.~Bovier,
\emph{Statistical Mechanics of Disordered Systems: A Mathematical
Perspective}, Cambridge University Press, 2006.

\bibitem[Car83]{Car83} B.~Carter, \emph{The anthropic principle and its
implications for biological evolution}, Philos.\ Trans.\ R.\ Soc.\ Lond.\
A \textbf{310} (1983), 347--363.

\bibitem[CH56]{CH56} D.~Cartwright and F.~Harary, \emph{Structural balance: a generalization of Heider's theory}, Psychol.\ Rev.\ \textbf{63} (1956), 277--293.

\bibitem[CL67]{CL67} H.~Cram\'er and M.~R.~Leadbetter,
\emph{Stationary and Related Stochastic Processes}, Wiley, New York,
1967.

\bibitem[Coo19]{Cooper19} K.~Cooper, \emph{The Contact Paradox:
Challenging Our Assumptions in the Search for Extraterrestrial
Intelligence}, Bloomsbury Sigma, London, 2019.

\bibitem[Cra23]{Crary23} A.~Crary, \emph{The toxic ideology of
longtermism}, Radical Philosophy \textbf{2} (2023), no.~14, 49--57.

\bibitem[Der80]{Der80} B.~Derrida,
\emph{Random-energy model: limit of a family of disordered models},
Phys.\ Rev.\ Lett.\ \textbf{45} (1980), 79--82.

\bibitem[Der81]{Der81} B.~Derrida,
\emph{Random-energy model: an exactly solvable model of disordered
systems}, Phys.\ Rev.\ B \textbf{24} (1981), no.~5, 2613--2626.

\bibitem[Dur19]{Dur19} R.~Durrett,
\emph{Probability: Theory and Examples}, 5th ed., Cambridge University
Press, 2019.

\bibitem[DW15]{DW15} D.~Denisov and V.~Wachtel,
\emph{Exit times for integrated random walks}, Ann.\ Inst.\ Henri
Poincar\'e Probab.\ Stat.\ \textbf{51} (2015), 167--193.

\bibitem[Dys60]{Dys60} F.~J.~Dyson, \emph{Search for artificial stellar
sources of infrared radiation}, Science \textbf{131} (1960), no.~3414,
1667--1668.

\bibitem[EF24]{EF24} J.-O.~Engler and J.~N.~Fischer, \emph{Drawing
blanks and winning: quantifying global catastrophic risk associated with
human ingenuity}, Sustainable Futures \textbf{7} (2024), 100165.

\bibitem[FW12]{FW12} M.~I.~Freidlin and A.~D.~Wentzell,
\emph{Random Perturbations of Dynamical Systems}, 3rd ed., Springer, 2012.

\bibitem[GJW99]{GJW99} P.~Groeneboom, G.~Jongbloed, and J.~A.~Wellner,
\emph{Integrated Brownian motion, conditioned to be positive},
Ann.\ Probab.\ \textbf{27} (1999), 1283--1303.

\bibitem[GM21]{GM21} H.~Greaves and W.~MacAskill, \emph{The case for
strong longtermism}, Global Priorities Institute Working Paper 5-2021,
University of Oxford, 2021.

\bibitem[Gol71]{Gol71} M.~Goldman,
\emph{On the first passage of the integrated Wiener process},
Ann.\ Math.\ Statist.\ \textbf{42} (1971), 2150--2155.

\bibitem[GS22]{GoldenScience22} A.~Benander, \emph{Golden Science},
unpublished manuscript, August 2022; archived copy dated November 2022.

\bibitem[Gom25]{Gom25} B.~Gompertz,
\emph{On the nature of the function expressive of the law of human
mortality, and on a new mode of determining the value of life
contingencies}, Philos.\ Trans.\ R.\ Soc.\ Lond.\ \textbf{115} (1825),
513--583.

\bibitem[Got93]{Got93} J.~R.~Gott~III, \emph{Implications of the
Copernican principle for our future prospects}, Nature \textbf{363}
(1993), no.~6427, 315--319.

\bibitem[Han98]{Han98} R.~Hanson, \emph{The Great Filter --- Are We
Almost Past It?}, manuscript, 1998.

\bibitem[Hei46]{Hei46} F.~Heider, \emph{Attitudes and cognitive organization}, J.~Psychol.\ \textbf{21} (1946), 107--112.

\bibitem[H\"or67]{Hor67} L.~H\"ormander,
\emph{Hypoelliptic second order differential equations},
Acta Math.\ \textbf{119} (1967), 147--171.

\bibitem[IW94]{IW94} Y.~Isozaki and S.~Watanabe,
\emph{An asymptotic formula for the Kolmogorov diffusion and a refinement
of Sinai's estimates for the integral of Brownian motion},
Proc.\ Japan Acad.\ Ser.\ A \textbf{70} (1994), 271--276.

\bibitem[Kar64]{Kar64} N.~S.~Kardashev, \emph{Transmission of
information by extraterrestrial civilizations}, Soviet Astronomy
\textbf{8} (1964), 217--221; Russian original, Astron.\ Zh.\
\textbf{41} (1964), 282--287.

\bibitem[Kip20]{Kip20} D.~Kipping, C.~Frank, and C.~Scarlata,
\emph{Contact inequality: first contact will likely be with an older
civilization}, preprint, 2020, \texttt{arXiv:2010.12358}.

\bibitem[Kol34]{Kol34} A.~N.~Kolmogorov,
\emph{Zuf\"allige Bewegungen (zur Theorie der Brownschen Bewegung)},
Ann.\ of Math.\ (2) \textbf{35} (1934), 116--117.

\bibitem[Kri80]{kripke1980naming} S.~A.~Kripke, \emph{Naming and
Necessity}, Harvard University Press, Cambridge, Massachusetts, 1980.

\bibitem[KT81]{KT81} S.~Karlin and H.~M.~Taylor,
\emph{A Second Course in Stochastic Processes}, Academic Press, 1981.

\bibitem[Lac91]{Lac91} A.~Lachal,
\emph{Sur le premier instant de passage de l'int\'egrale du mouvement
brownien}, Ann.\ Inst.\ H.\ Poincar\'e Probab.\ Statist.\ \textbf{27}
(1991), 385--405.

\bibitem[Lan04]{Lan04} D.~Lando,
\emph{Credit Risk Modeling: Theory and Applications}, Princeton
University Press, 2004.

\bibitem[Mar49]{Mar49} G.~Maruyama,
\emph{The harmonic analysis of stationary stochastic processes},
Mem.\ Fac.\ Sci.\ Ky\=ush\=u Univ.\ Ser.\ A \textbf{4} (1949), 45--106.

\bibitem[McK63]{McK63} H.~P.~McKean, Jr.,
\emph{A winding problem for a resonator driven by a white noise},
J.\ Math.\ Kyoto Univ.\ \textbf{2} (1963), 227--235.

\bibitem[MK04]{MK04} G.~Molchan and A.~Khokhlov,
\emph{Small values of the maximum for the integral of fractional
Brownian motion}, J.\ Stat.\ Phys.\ \textbf{114} (2004), 923--945.

\bibitem[MKKS11]{MKKS11} S.~A.~Marvel, J.~Kleinberg, R.~D.~Kleinberg, and S.~H.~Strogatz, \emph{Continuous-time model of structural balance}, Proc.\ Natl.\ Acad.\ Sci.\ USA \textbf{108} (2011), 1771--1776.

\bibitem[Mol99]{Mol99} G.~M.~Molchan,
\emph{Maximum of a fractional Brownian motion: probabilities of small
values}, Comm.\ Math.\ Phys.\ \textbf{205} (1999), 97--111.

\bibitem[MP01]{MP01} M.~A.~Milevsky and S.~D.~Promislow,
\emph{Mortality derivatives and the option to annuitise},
Insurance Math.\ Econom.\ \textbf{29} (2001), 299--318.

\bibitem[MV12]{MV12} S.~M\'el\'eard and D.~Villemonais,
\emph{Quasi-stationary distributions and population processes},
Probab.\ Surv.\ \textbf{9} (2012), 340--410.

\bibitem[Ord23]{Ord23} T.~Ord, \emph{The Lindy effect}, preprint, 2023,
\texttt{arXiv:2308.09045}.

\bibitem[All09]{Allen09} R.~C.~Allen, \emph{The British Industrial
Revolution in Global Perspective}, Cambridge University Press, 2009.

\bibitem[Arr]{Arrian} Arrian, \emph{Anabasis of Alexander}, 1.15.8.

\bibitem[Dia97]{Diamond97} J.~Diamond, \emph{Guns, Germs, and Steel: The
Fates of Human Societies}, W.~W.~Norton, 1997.

\bibitem[Pom00]{Pomeranz00} K.~Pomeranz, \emph{The Great Divergence:
China, Europe, and the Making of the Modern World Economy}, Princeton
University Press, 2000.

\bibitem[PlAlex]{PlutarchAlex} Plutarch, \emph{Life of Alexander},
\S16, in \emph{Parallel Lives}.

\bibitem[Plu]{Plutarch} Plutarch, \emph{Whether Land or Sea Animals Are
Cleverer} (\emph{De sollertia animalium}), \S7, in \emph{Moralia};
the saying is attributed there to Bion of Borysthenes.

\bibitem[Ric44]{Ric44} S.~O.~Rice,
\emph{Mathematical analysis of random noise}, Bell System Tech.\ J.\
\textbf{23} (1944), 282--332.

\bibitem[Ric45]{Ric45} S.~O.~Rice,
\emph{Mathematical analysis of random noise (concluded)}, Bell System
Tech.\ J.\ \textbf{24} (1945), 46--156.

\bibitem[RVY06]{RVY06} B.~Roynette, P.~Vallois, and M.~Yor,
\emph{Some penalisations of the Wiener measure}, Jpn.\ J.\ Math.\
\textbf{1} (2006), 263--290.

\bibitem[RY99]{RY99} D.~Revuz and M.~Yor,
\emph{Continuous Martingales and Brownian Motion}, 3rd ed., Springer,
1999.

\bibitem[SDO18]{SDO18} A.~Sandberg, E.~Drexler, and T.~Ord,
\emph{Dissolving the Fermi paradox}, preprint, 2018.

\bibitem[Sin92]{Sin92} Ya.~G.~Sinai,
\emph{Distribution of some functionals of the integral of a random walk},
Theoret.\ and Math.\ Phys.\ \textbf{90} (1992), 219--241; translated
from Teoret.\ Mat.\ Fiz.\ \textbf{90} (1992), no.~3, 323--353.

\bibitem[Tai88]{Tainter88} J.~A.~Tainter, \emph{The Collapse of Complex
Societies}, Cambridge University Press, 1988.

\bibitem[Tal12]{Tal12} N.~N.~Taleb, \emph{Antifragile: Things That Gain
from Disorder}, Random House, New York, 2012.

\bibitem[Tor21]{Torres21} \'E.~P.~Torres, \emph{Against longtermism},
Aeon, 2021.

\bibitem[Ben06]{Benatar06} D.~Benatar, \emph{Better Never to Have Been:
The Harm of Coming into Existence}, Oxford University Press, 2006.
\bibitem[VHEMT]{VHEMT} L.~U.~Knight, \emph{The Voluntary Human
Extinction Movement}, founded 1991; \texttt{vhemt.org}.
\bibitem[VMS79]{VMS79} J.~W.~Vaupel, K.~G.~Manton, and E.~Stallard,
\emph{The impact of heterogeneity in individual frailty on the dynamics
of mortality}, Demography \textbf{16} (1979), 439--454.

\bibitem[VY85]{VY85} J.~W.~Vaupel and A.~I.~Yashin,
\emph{Heterogeneity's ruses: some surprising effects of selection on
population dynamics}, Amer.\ Statist.\ \textbf{39} (1985), 176--185.

\bibitem[Wah78]{Wah78} G.~Wahba,
\emph{Improper priors, spline smoothing and the problem of guarding
against model errors in regression},
J.\ Roy.\ Statist.\ Soc.\ Ser.\ B \textbf{40} (1978), 364--372.

\bibitem[Wz79]{Wal79} K.~N.~Waltz, \emph{Theory of International Politics}, McGraw-Hill, 1979.

\bibitem[Wt87]{Wal87} S.~M.~Walt, \emph{The Origins of Alliances}, Cornell University Press, 1987.

\bibitem[WB00]{WB00} P.~D.~Ward and D.~Brownlee, \emph{Rare Earth: Why
Complex Life Is Uncommon in the Universe}, Copernicus, New York, 2000.

\end{thebibliography}
\end{filecontents*}
\IfFileExists{biblatex.sty}{%
  \usepackage[backend=biber,style=alphabetic,maxbibnames=99,giveninits=true]{biblatex}%
  \addbibresource{golden-reasoning.bib}%
  \newcommand{\PrintRefs}{\printbibliography[title={References}]}%
}{%
  \newcommand{\PrintRefs}{\input{gr-manualbib.tex}}%
}

\numberwithin{equation}{section}

\theoremstyle{plain}
\newtheorem{theorem}{Theorem}[section]
\newtheorem{proposition}[theorem]{Proposition}
\newtheorem{lemma}[theorem]{Lemma}
\newtheorem{corollary}[theorem]{Corollary}
\newtheorem{claim}[theorem]{Claim}

\theoremstyle{definition}
\newtheorem{definition}[theorem]{Definition}
\newtheorem{ansatz}[theorem]{Ansatz}
\newtheorem{convention}[theorem]{Convention}
\newtheorem{constraint}{Constraint}
\newtheorem{principle}{Principle}

\theoremstyle{remark}
\newtheorem{remark}[theorem]{Remark}
\newtheorem{example}[theorem]{Example}

\newcommand{\E}{\mathbb{E}}
\newcommand{\PP}{\mathbb{P}}
\newcommand{\R}{\mathbb{R}}
\DeclareMathOperator{\Var}{Var}
\DeclareMathOperator{\Cov}{Cov}
\DeclareMathOperator{\Leb}{Leb}
\newcommand{\eqd}{\overset{d}{=}}
\newcommand{\one}{\mathbf{1}}
\newcommand{\Imm}{\mathcal{I}}
\newcommand{\Fc}{\mathcal{F}}
\newcommand{\Lc}{\mathcal{L}}
\newcommand{\hm}{\mathfrak{h}}
\newcommand{\Ac}{\mathcal{A}}
\newcommand{\Sc}{\mathsf{S}}

\allowdisplaybreaks

\title{Golden Reasoning\\[4pt]
\large Soft Horizons, Survivorship, and the Two-Channel Muffled Hazard}
\author{}
\date{August 2026}
\newtheorem{premise}{Premise}
\renewcommand{\thepremise}{$\Phi$\arabic{premise}}
\newtheorem{objection}[theorem]{Objection}

\allowdisplaybreaks

\title{Golden Reasoning\\[4pt]
\large Soft Horizons, Survivorship, and the Two-Channel Muffled Hazard}
\author{}
\date{August 2026}
\allowdisplaybreaks

\title{Golden Reasoning\\[4pt]
\large Soft Horizons, Survivorship, and the Two-Channel Muffled Hazard}
\author{}
\date{August 2026}
\allowdisplaybreaks

\title{Golden Reasoning\\[4pt]
\large Soft Horizons, Survivorship, and the Two-Channel Muffled Hazard}
\author{}
\date{August 2026}
\allowdisplaybreaks

\title{Golden Reasoning\\[4pt]
\large Soft Horizons, Survivorship, and the Two-Channel Muffled Hazard}
\author{}
\date{August 2026}
\allowdisplaybreaks

\title{Golden Reasoning\\[4pt]
\large Soft Horizons, Survivorship, and the Two-Channel Muffled Hazard}
\author{Ajax Benander\thanks{Email: \href{mailto:akb9408@rit.edu}{akb9408@rit.edu}.}}
\date{August 2026}
\allowdisplaybreaks

\title{Golden Reasoning\\[4pt]
\large Soft Horizons, Survivorship, and the Two-Channel Muffled Hazard}
\author{Ajax Benander\thanks{Email: \href{mailto:akb9408@rit.edu}{akb9408@rit.edu}.}}
\date{August 2026}
\allowdisplaybreaks

\title{Golden Reasoning\\[4pt]
\large Soft Horizons, Survivorship, and the Two-Channel Muffled Hazard}
\author{Ajax Benander\thanks{Email: \href{mailto:akb9408@rit.edu}{akb9408@rit.edu}.}}
\date{August 2026}
\allowdisplaybreaks

\title{Golden Reasoning\\[4pt]
\large Soft Horizons, Survivorship, and the Two-Channel Muffled Hazard}
\author{Ajax Benander\thanks{Email: \href{mailto:akb9408@rit.edu}{akb9408@rit.edu}.}}
\date{August 2026}

\begin{document}
\maketitle
\tableofcontents
\bigskip

\section{Introduction}

\section{Continuing from Nexic reasoning}\label{gc:sec}

\begin{quote}\small
The companion development derives a typicality constraint --- the
Accordance Principle --- and applies it to the natural-historical record
under a weighting it calls Optimation: the cosmic weighting conditioned
on observer-existence \cite{Nexus, Ben26Found, Ben26Optimal}. That
weighting is there a construction, adopted because observer-existence is
the obvious thing to condition on. This section argues that it is the
wrong thing, that the right thing is the \emph{Golden Point}, and that
Optimation is recovered exactly as a consequence rather than assumed as
a premise.

Two claims carry the argument. The \emph{Past-Alignment Principle}
(\S\ref{gc:sub:pastalign}) states that conditioning the cosmic weighting
on the Golden Point and then discarding every attribute outside one's
own past light cone returns precisely the Nexus of Optimation --- so
that the record we observe is, attribute for attribute, the past-light-cone
restriction of the Golden Nexus. The \emph{Golden-Conditioning
Hypothesis} (\S\ref{gc:sub:conditioning}) argues by contradiction that
this is no coincidence: observer-existence is too weak a condition to
have produced the record, since it survives perturbations that the
record does not. What conditioning on observers cannot explain,
conditioning on the Golden Point does.

The consequence is a reversal of order. Optimation ceases to be the
framework's starting point and becomes its first derived result; the
Golden Biosphere ceases to be an object defined by survival asymptotics
and becomes the conditioning target from which the rest follows. What
the argument does \emph{not} settle is the future: whether the alignment
that holds behind us extends ahead of us is left to the three
hypotheses of \S\ref{gc:sub:extrapolation}, which are stated and not
adjudicated.
\end{quote}

% ---------------------------------------------------------------------
\subsection{Vocabulary and standing conventions}\label{gc:sub:vocab}

\begin{convention}[Terms carried from the companion]\label{gc:conv:vocab}
The argument runs in the companion's vocabulary
\cite{Nexus, Ben26Found}, and the six terms it uses are fixed here.
An \emph{Aethus} (plural \emph{Aethae}) is a partial information state:
some attributes stated, all others genuinely blank, individuated by its
stated content and closed under entailment from it. The \emph{Cosmic
Nexus} is the occurrence-weighted measure over Aethae, with no
conditioning absorbed into it. \emph{Optimation} is that measure
conditioned on the biosphere-level observer attribute, and the
\emph{Nexus of Optimation} the resulting weighting. The
\emph{Aetherian Nexus} is the Aethus of our own natural history ---
the one whose stated attributes we read off the record. \emph{Centric
unfolding} is the accumulation of stated attributes along a
particular reasoner's history. And the \emph{Golden Point} is the event
of a biosphere's passage into the Golden regime, in the sense
\S\ref{gr:sec} makes precise; the \emph{Golden Nexus} is the cosmic
weighting conditioned on that event.
\end{convention}

\begin{convention}[Grades]\label{gc:conv:grades}
As elsewhere: \textbf{[P]} proved, \textbf{[C]} cited, \textbf{[S]}
argued at scaling or schematic level, \textbf{[O]} open, and
\textbf{[$\Phi$]} a philosophical premise carrying its own falsifier.
This section contains no \textbf{[P]}: its two central claims are
\textbf{[$\Phi$]} and \textbf{[S]} respectively, and the extrapolation
hypotheses are \textbf{[O]} by construction. Nothing here is asserted at
a grade the argument does not reach.
\end{convention}

% ---------------------------------------------------------------------
\subsection{The problem: what does Optimation condition on?}
\label{gc:sub:problem}

Optimation is introduced in the companion as the cosmic weighting
conditioned on observer-existence, and the choice is natural: an
observer is what is doing the reasoning, so an observer is what the
reasoning conditions on. The choice is nevertheless a choice, and it
is load-bearing --- every rarity verdict the companion issues is
relative to it.

The difficulty is that observer-existence is \emph{robust} in a way the
record is not. Strip from our history the coal that powered
industrialisation, or the metallurgical sequence that made rocketry
possible, or any of a dozen similar contingencies, and what remains is
still a world with observers in it. Observerhood survives these
perturbations; the specific character of our record does not. A
conditioning target that is indifferent to the difference cannot
explain why we find the difference realised.

That gap is what the rest of this section fills. The proposal is that
the operative condition was never observer-existence but the Golden
Point, and that Optimation is what Golden-conditioning looks like when
restricted to the region of spacetime one can actually see.


\noindent Two questions follow, and the rest of this section answers them
in order. \emph{Which} target should replace observer-existence, and on
what grounds --- \S\ref{gc:sub:continuity}. And what reason is there to
think the replacement is not merely better motivated but actually the
condition under which the record was produced ---
\S\S\ref{gc:sub:pastalign}--\ref{gc:sub:humanhistory}, in two
independent legs.

% ---------------------------------------------------------------------
\subsection{Choosing the target: continuity and non-privilege}\label{gc:sub:continuity}

Before the argument for Golden-conditioning, a prior question: why the
Golden Point rather than some more familiar target? ``Technologically
advanced civilization'' is the obvious candidate and is the wrong one,
for two reasons that compound.

It is not well-defined. Any threshold one names --- writing, metallurgy,
industry, spaceflight --- is a stipulation, and the arbitrariness is not
cosmetic: every rarity verdict issued relative to such a target inherits
it.

And it is \emph{discontinuous with natural history}. A conditioning
target must be applicable across the whole record if the record is to be
read as one object. ``Civilization'' applies to the last few millennia
and to nothing before, so a framework built on it cannot run a single
conditioning across both the end-Cretaceous impact and the Industrial
Revolution. It must instead splice two analyses at a joint it has no
principled way to locate.

\begin{definition}[Departure Point]\label{gc:def:departure}
The \emph{Departure Point} is the event past which \emph{every}
continuation fails to advance --- the first moment at which all
scenarios lead to a failure to advance past that point. The modal
quantifier is the definition and not an embellishment of it: what is
being named is the closure of possibility, not the observation of
non-progress.

\emph{This is what makes it the dual of the Golden Point.} That was
defined as the point past which a biosphere is too capable at self
preservation to regress past it --- i.e.\ the point past which every
continuation \emph{avoids} foreclosure. The two are exact mirrors, and
both are \emph{points} for the same reason: irreversibility. A state
one may still leave is not a point but a phase, and neither concept is
about phases.

Two consequences of stating it this way.

\emph{Biological death is a mechanism, not the definition.} Death is one
way a trajectory is foreclosed, and it happens to be the most visible
one; irreversible loss of the capacity to advance is another, and the
definition is indifferent between them. What is extracted is the
\emph{phenomenology of foreclosure}, taken apart from the particular
mechanism of biosphere-death in which it is most legible --- and once
extracted it applies wherever there is a trajectory to foreclose,
which is why the notion carries across the scales
Principle~\ref{gc:prin:continuity} requires.

\emph{Non-advancement alone does not suffice.} A biosphere that fails
to progress for an age has not thereby reached its Departure Point; it
has had a lull, and a lull is reversible. Only when no continuation
remains that advances has the point been passed. This is the
distinction that the phrase ``failure to advance'' elides and that the
quantifier restores.

Two consequences for how the model of \S\ref{gr1:sec} should be read.
The hazard $h^{\mathrm{tot}}_t$ of \eqref{tc:eq:hazard} is the hazard
\emph{of foreclosure}, not of death; the muffling exponent $x$ is
accumulated against that; and every Imperfection Magnitude Drop is a
movement toward it. And $\PP(\Lambda_\infty<\infty)=0$ is the analytic
form of the modal claim --- it says that \emph{no} continuation
advances, which is precisely the quantifier above. This is why the
trapped regime of Proposition~\ref{gr1:prop:trichotomy}(iii) has
departed: not because a biosphere halting at $x_\ast$ has stopped, but
because with $\rho<0$ the viability probability is identically zero, so
no continuation available to it advances. Halting is not the departure;
having nowhere left to go is.

The term is used rather than ``extinction'' deliberately, and the
reason is the same one that governs the choice of the Golden Point
itself. ``Extinction'' is a biological category, applying to taxa and
carrying the vocabulary of one discipline; the Departure Point must
apply at every scale the record contains --- to a Mesozoic lineage as
readily as to a biosphere tomorrow --- and must name an event in the
Aethic block universe rather than a fact about organisms. It is
therefore defined at the same primitive level as the Golden Point,
which is what allows Principle~\ref{gc:prin:continuity} to be stated at
all.

\emph{Internal and external.} The Departure Point divides by the
locus of its cause. An \emph{internal} Departure Point is brought about
by factors within the biosphere's own homeworld; an \emph{external} one
is inflicted from outside, and is roughly constant in magnitude over
long timescales because the cosmos is roughly constant in makeup. That
division is the ontological origin of the two channels of
\eqref{tc:eq:hazard}: the crest channel carries the external hazard and
is bounded, the trough channel carries the internal and is not.
Technoanatomy (Definition~\ref{gr1:def:anatomy}) is then the statement
of how the second responds to capacity acquired against the
first.\footnote{\emph{Provenance.} The Departure Point, its
internal/external division, the constancy of the external threat, and
the observation that a biosphere may come to draw energy out of itself
are all set out in the author's \emph{Golden Science} of November 2022,
some years before the present formalism and before the assistant's
involvement. That document also states the Golden Point as the point
past which a biosphere cannot regress, derives a positive immortality
probability from an exponentially decaying hazard, argues for
``biosphere'' over ``civilization'' as the unit, and names the Trace
Connection as an open problem --- which \S\ref{gc:sub:pastalign} is the
attempt to close. The formalism here should be read as supplying
machinery for commitments already made there, not as generating
them.}
\end{definition}

\subsubsection{The implication cascade: departure as a least fixed point}
\label{gc:subsub:cascade}

Definition~\ref{gc:def:departure} carries a recursion that is worth
making explicit, because it is already available in the companion
framework in finished form and because it changes what the model's
hazard is a hazard \emph{of}.

The internal and external hazards of \eqref{tc:eq:hazard} \emph{imply}
departure; they are not exhaustive of it. A biosphere may reach a
configuration from which no route avoids those hazards, and by the
modal definition that configuration \emph{is} a Departure Point,
attained before either hazard has fired. The statement ``no matter what
you do, you reach a Departure Point'' is itself a Departure Point ---
and since that statement can again be true one step earlier, the
implication propagates.

\begin{remark}[This is the third postulate's operator, and it is proved]
\label{gc:rem:thirdpostulate}
The companion's third postulate supplies exactly this recursion
\cite{Aethus}. Over states ordered by descent, with $\mathrm{Base}$
collecting the non-recursive modes, it defines a safety operator and
its dual
\begin{align*}
F(Y)&=\bigl\{A\notin\mathrm{Base}:\ \exists B\sqsubseteq A,\
\forall C\sqsubseteq B,\ C\in Y\bigr\},\\
G(X)&=\mathrm{Base}\cup\bigl\{A:\ \forall B\sqsubseteq A,\
\exists C\sqsubseteq B,\ C\in X\bigr\},
\end{align*}
both monotone, so that by Knaster--Tarski the greatest fixed point
$\nu F$ and least fixed point $\mu G$ exist for arbitrary branching and
are complements of one another; for finitely-branching trees the
closure ordinal is at most $\omega$, so the cascade terminates.

Read $\mathrm{Base}$ as departure by internal or external hazard.
Then $G$ is the doom operator in the sense above --- every route has a
doomed continuation --- and the Departure Point is
\begin{equation}\label{gc:eq:departurefix}
\mathcal{D}\;=\;\mu G\;\supsetneq\;\mathrm{Base}.
\end{equation}
Three consequences, each already carried by the companion's
machinery rather than assumed here.

\emph{The implication is one-directional by design.} The companion
states the third postulate as an implication and not a biconditional,
which is precisely the containment in \eqref{gc:eq:departurefix}: the
hazards suffice for departure and are not necessary to it.

\emph{The Golden/Departure duality is a theorem, not a construction.}
The Golden Point is $\nu F$ --- there exists a route all of whose
continuations remain safe --- and the Departure Point is $\mu G$. Their
being exact mirrors is the complementarity of the two fixed points, so
the duality asserted in Definition~\ref{gc:def:departure} need not be
stipulated: it is inherited.

\emph{Human extinction sits in the difference
$\mu G\setminus\mathrm{Base}$.} It is not itself a biosphere-level
departure by internal or external hazard,
since the biosphere persists; but with no muffling agent remaining,
every continuation reaches $\mathrm{Base}$
(Proposition~\ref{hr:prop:extinction}). It therefore enters at the
first step of the cascade. The classical reading --- that human
extinction is a Departure Point because it is only a matter of time ---
is the $n=1$ case of the recursion, and \eqref{gc:eq:departurefix} is
what licenses it.
\end{remark}

\begin{remark}[The model's hazard is a lower bound]
\label{gc:rem:hazardlower}
The consequence for \S\ref{gr1:sec} should be stated plainly, since it
qualifies every rate in this work. The hazard $h^{\mathrm{tot}}_t$ of
\eqref{tc:eq:hazard} is the hazard of $\mathrm{Base}$: it prices the
internal and external mechanisms and nothing else. The hazard of
\emph{departure} is the hazard of $\mu G$, which by
\eqref{gc:eq:departurefix} is at least as large and in general strictly
larger --- because a biosphere departs when its fate is sealed, not
when the sealing is executed, and the interval between those can be
long.

So the model's timescales are conservative in a specific direction:
where it reports a departure time $T$, the Departure Point proper may
have been passed considerably earlier, and the survival curves of
\S\ref{gr1:sub:verdict} are upper bounds on the probability of not
having departed. Nothing in the asymptotic analysis is disturbed by
this --- a cascade that fires earlier by a bounded lead time leaves
exponents untouched --- but the quantitative readings inherit it, and
the anatomy urgency of \S\ref{gr1:subsub:fluct} is sharpened by it: the
window in which one may still act closes at the cascade's rate, not at
the hazard's. Computing that rate is open \textbf{[O]}.
\end{remark}


\begin{principle}[Continuity; \textbf{[$\Phi$]}]\label{gc:prin:continuity}
The struggle our ancestors underwent against the Departure Point in
natural history is continuous, philosophically and ontologically, with a
society's agency-mediated struggle against the Departure Point. They are
not analogues but the same activity under different mediation.
\end{principle}

The Golden Point is the target that respects this, and it does so
because of what kind of definition it has: \emph{survival conditioning
of its type}. That is primitive enough to apply without alteration to a
Mesozoic mammal evading theropod predation and to a modern society
avoiding its own weapons. Both lower a hazard; both prevent an
Imperfection Magnitude Drop; neither requires the vocabulary of the
other. The substance is continuous where the societal construct is not.

\begin{objection}[Agency]\label{gc:obj:agency}
A Mesozoic mammal does not know it is fighting the Departure Point.
There is no agency, no goal, no strategy --- only an animal avoiding a
predator. Calling that continuous with deliberate societal
risk-management equivocates on the word ``fighting''.
\end{objection}

\noindent The objection assumes agency is what individuates the
activity. It is not. By \S\ref{to:sec}, what individuates a task-object
is its \emph{contribution to the root}, and agency is nowhere in that
criterion --- which is precisely why task-objects were required before
this principle could be stated.

What replaces agency is the mediating mechanism, and Optimation
conditions on the mechanism rather than around it. Where a society
lowers its hazard through deliberation, a lineage lowers its through
evolutionary selection; the conditioning does not care which, because
what it fixes is that the hazard be lowered. It therefore reaches
\emph{through} the mechanism to whatever contingencies the mechanism
requires in order to operate.

\begin{example}[Conditioning through the medium: the
$\alpha$-gal case; \textbf{[S]}]\label{gc:ex:medium}
The author's $\alpha$-gal case exhibits the structure exactly, and its
details matter because they show the conditioning having no alternative
route \cite{Nexus}.

Some thirty million years ago the Catarrhine lineage came under
pathogen pressure of a particular kind. The pathogen presented
$\alpha$-gal, a carbohydrate our own cells also produced, so that no
immune response was available which did not attack the host. What
survived carried a mutation disabling \textit{GGTA1}, the gene
responsible for producing it --- a gene retained by every other
mammalian lineage, and lost in ours alone.\footnote{The inactivation of
\textit{GGTA1} in Old World primates is established, as is
pathogen-mediated selection as its leading explanation. The
\emph{severity} of the episode --- whether it approached extinction of
the lineage --- is an inference and not a finding, and the argument does
not rest on it: pressure sufficient to remove a gene universal among
mammals is all that the reasoning requires, and a milder episode
delivering that pressure would serve the example unchanged.}

Now read it under Optimation. The gene was ancient and useful, which is
why it is universal among mammals; and evolutionary selection removes a
gene if and only if there is pressure to remove it. That is the
mechanism's own condition of operation, and it is not negotiable. So
under ordinary conditions the loss was not merely improbable but close
to unavailable: selection had nothing to act on, because the gene was
not costing anything.

A conditioning requiring the loss therefore cannot act on the gene. It
must act on whatever \emph{supplies the pressure} --- and what supplied
it was the contingent existence of a pathogen presenting the very
carbohydrate in question, which converts the gene from an asset into the
one thing standing between the lineage and extinction. Optimation
reweights the pathogen's occurrence to prominence, and the loss follows
\emph{through} selection rather than around it.

This is the general shape. Optimation does not reach past evolutionary
selection to install outcomes; it reweights the contingencies that make
selection deliver them. Agency, where it exists, is one such mechanism
among others --- the deliberative rather than the Darwinian medium ---
and the conditioning is indifferent between them.
\end{example}

\begin{remark}[What the principle buys, and its cost]
\label{gc:rem:continuitycost}
It buys the object the rest of this work analyses. Without continuity,
\S\ref{gr1:sec}'s hazard process would model two different things
spliced at an arbitrary joint; with it, one process runs from
abiogenesis to the present under a single conditioning, and the
technological present is a \emph{position} along it rather than a change
of subject.

The cost is that the principle is a philosophical premise and not a
result, and its scope should be stated at full width rather than left to
be inferred, because it is larger than the conditioning question that
motivated it. Continuity licenses treating a Mesozoic mammal's evasion
of predation, a body's suppression of its own cancers, and an
institution's not turning its capacity against itself as values of
\emph{one variable}. That is what permits \S\ref{ad:sec}'s institutional
machinery to bear on \S\ref{gr1:sec}'s hazard dynamics at all: without
it, the adversivity model would describe the politics of societies and
the survival model would describe the fate of lineages, and no argument
would carry between them.

So the premise is not merely buying a conditioning target. \emph{It is
what makes this one work rather than several.} Should it fail, the
sections do not degrade gracefully --- they separate, and each retains
whatever it had independently while the connections lapse.

\emph{Falsifier:} a demonstration that deliberative and evolutionary
hazard-reduction differ in a respect the conditioning must track ---
that some feature of the record is explicable under one and not the
other. The agency objection is the natural candidate and is answered
above; a better one would consist in exhibiting such a feature rather
than in restating the intuition. A second and more targeted falsifier
follows from the width just claimed: exhibit a hazard-lowering mechanism
at one scale whose governing variable provably fails to correspond to
anything at another --- a property of institutional self-restraint with
no analogue in organismal or lineage-level hazard control, or the
reverse.
\end{remark}

\subsubsection{Why continuity is wanted: the non-privilege of humans}
\label{gc:subsub:nonprivilege}

The principle has now been stated and defended, but not motivated, and
its motivation is not tidiness. Continuity is wanted because it is what
makes the framework non-anthropocentric, and the non-anthropocentrism is
load-bearing rather than decorous.

Begin from what the target is. \emph{The Golden Point is a property of a
biosphere} --- not of a civilization, not of a species. That is why it
applies to an alien biosphere as readily as to ours, and applies without
translation: nothing in ``survival conditioning of its type'' mentions
humans, technology, or society, so nothing must be reinterpreted when
the biosphere in question is not ours.

\begin{proposition}[Recursive anchoring; \textbf{[S]}]
\label{gc:prop:recursive}
Because the target is biosphere-level and the principle is continuous
across mediation, the conditioning may be anchored at \emph{any} Aethus
in the lineage consistent with ours. Take a Mesozoic mammal, its Aethus
a parent of ours rather than an Imperfection Magnitude Drop, and ask
what the continuation of Aethic attributes within \emph{its} future
light cone must look like for its biosphere to reach the Golden Point.
The answer runs through the remainder of the lineage --- and so through
us --- without us having been named anywhere in the question.
\end{proposition}

\noindent Our position is therefore an \emph{output} of the conditioning
rather than a stipulated anchor for it. Humans are not a privileged
split point; the technological present is where the recursion happens to
pass, and its role is emergent from a primitive that mentions no
species. The same question asked of a mammal in the Maastrichtian, of a
lineage that has not yet speciated, or of a biosphere elsewhere entirely,
is one question.

\begin{remark}[Three things a human-centred target would cost]
\label{gc:rem:privilegecost}
The alternative --- conditioning on ``technological civilization'' or on
observers --- is not merely less elegant. It forfeits three things, in
increasing order of seriousness.

\emph{Robustness across biospheres.} An alien case would require
separate treatment, since the target would be defined by features of our
own history. A framework whose central conditioning cannot be applied to
the systems it is meant to reason about has not achieved generality; it
has assumed its instance.

\emph{Robustness across perturbations.} \S\ref{gc:subsub:pastalign-argument}
evaluates histories in which humans do not arise. Under a human-centred
target those histories cannot be evaluated at all --- the target simply
fails to apply --- and precisely the perturbations that carry the
argument become unaskable.

\emph{Non-circularity.} A human-centred target takes \emph{us} as
primitive, so the framework would build its base unit out of the thing
it is supposed to explain. The Continuity Principle is what prevents
this, and that is the deepest reason to want it.
\end{remark}

\begin{remark}[A distinction to keep: definition versus conditioning]
\label{gc:rem:defvscond}
Two claims are easily run together and must not be, since the argument's
validity depends on their separation.

The \emph{definition} of the Golden Point --- a survival property of a
biosphere, in the sense \S\ref{gr1:sec} makes precise --- is
stipulative. Someone who rejects everything else in this section may
adopt it, and the hazard analysis of \S\ref{gr1:sec} stands under it
regardless.

The \emph{conditioning claim} --- that this property is what the
Accordance Principle has been conditioning on --- is not stipulative and
is what \S\S\ref{gc:sub:conditioning}--\ref{gc:sub:humanhistory} argue
for, by reductio on natural history and again on human history. Running
the two together would make the argument circular, defining the
conditioning target into its role and then observing that it occupies
it. Keeping them apart is what leaves the argument something it could
have failed at.
\end{remark}

\begin{remark}[Two non-privilege claims, and their symmetry]
\label{gc:rem:twoprivilege}
The section now makes two claims of this shape, and they are the same
claim along different axes.
Principle~\ref{gc:prin:pastalign} indexes its light-cone restriction to
one's own present \emph{because that is where the observable cone is
anchored, not because the event is distinguished}; the construction at
any other event returns that event's cone.
Proposition~\ref{gc:prop:recursive} says the same of position in a
lineage: the conditioning may be anchored anywhere consistent, and
returns that anchor's continuation. No privileged spacetime event, no
privileged lineage point --- and in both cases the appearance of
privilege is an artefact of which cone the reasoner happens to occupy.
\end{remark}


% ---------------------------------------------------------------------
With the target fixed and its motivation stated, the question becomes
evidential: is there reason to believe the record was in fact produced
under it? The remainder of the section argues that there is, first by
saying what alignment with such a target would look like
(\S\ref{gc:sub:pastalign}), then by two independent reductios --- one on
natural history (\S\ref{gc:sub:conditioning}), one on human history
(\S\ref{gc:sub:humanhistory}) --- of which the second does not require
the premise the first leans on.

% ---------------------------------------------------------------------
\subsection{What alignment with the target would mean}\label{gc:sub:pastalign}

\subsubsection{Statement}\label{gc:subsub:pastalign-statement}

\begin{principle}[Past-Alignment; \textbf{[$\Phi$]}]
\label{gc:prin:pastalign}
Construct the \emph{Golden Nexus} by conditioning the Cosmic Nexus on
the Golden Point. From it, discard every stated attribute whose
referential event lies outside one's own past light cone --- in the
stricter reading, discarding the spacelike attributes as well as the
future-directed ones. What remains is the Nexus of Optimation.

Equivalently: the past light cone of Aethic attributes we actually
observe is the same one as the past light cone through our position in
the Aethic block universe of the Golden Nexus.\footnote{More precisely
than ``our actual light cone'': owing to random contingency in centric
unfolding, our own Aethus is a proper child of the Aetherian Nexus
rather than the Aetherian Nexus entire, so the object being compared is
the past light cone of the Aetherian Nexus's own Aethic block universe.}
\end{principle}

Two features of the statement are deliberate and should not be read past.

\emph{The restriction is to the past cone only.} The principle takes no
position on whether attributes in one's future light cone agree with the
near-optimal part of the Golden Nexus. That silence is not modesty but
scope: one cannot reliably infer future-cone attributes prior to
engaging in Golden reasoning, so a principle that assumed agreement
there would be assuming its own conclusion.
\S\ref{gc:sub:extrapolation} treats the forward question separately and
does not settle it.

\emph{The spacetime point is not privileged.} The principle is indexed
to one's own present because that is where the observable cone is
anchored, not because that event has any distinguished status. The same
construction at any other event returns that event's cone.

\begin{remark}[What the principle claims in the companion's own terms]
\label{gc:rem:oec}
Read against the Optimal-Earth Conjecture \cite{Ben26Optimal}, the
principle sharpens what ``optimal'' is optimal \emph{for}. The
conjecture there holds our record observer-optimal on effectively any
interrogable attribute. Past-Alignment says the optimand is not
vaguely observer-production but specifically the Golden Point: one asks
for the optimal Aethic history \emph{to the Golden Point}, and then
surgically reopens every attribute whose occurrence lies in the future
cone of the indexing event. If the principle holds, cropping the Golden
Nexus to one's past light cone reproduces Optimation exactly --- and it
is the exactness, not the plausibility, that would make this a
first-principles derivation of Optimation rather than a restatement.
\end{remark}

\subsubsection{The argument from prerequisite attributes}
\label{gc:subsub:pastalign-argument}

The supporting argument has two steps.

\emph{Step one: identify the attributes that position us for the Golden
Point.} Our record contains features which, whatever else they do, make
this biosphere unusually able to reach the Golden regime. The capacity
to build rocket vehicles is the cleanest example, since it is what would
allow the biosphere to persist past the eventual death of its star.

\emph{Step two: perturb the record and watch them fail.} Consider
Aethic perturbations in which the dinosaurs never go extinct, in which
animal life never evolves at all, or in which Africa does not exist or
develops very differently. Each comparison presupposes that the
perturbed history and ours share a common term against which a missing
capacity registers as missing, which ordinary object-individuation
cannot supply and \S\ref{to:sec} does: what is compared across these
histories is a Task-Earth and a task-lineage
(Definition~\ref{to:def:taskobject}), not an Earth and a lineage. With
that in place, each perturbation removes the capacity:

\begin{itemize}\itemsep2pt
\item A world in which the dinosaurs persist is one in which, by the
Accordance Principle's own verdict, the lineage does not reach our
level of intelligence, and so does not build rocket vehicles.
\item A world of fungi lacks the animal capacities the construction
requires.
\item A world without the specific African climatic sequence leaves
great apes arboreal, without the bipedal orientation and strategic
cognition the construction requires.\footnote{The climatic sequence in
Africa is what drove our great-ape ancestors to bipedality.}
\end{itemize}

The examples are illustrative rather than exhaustive, and the
generalisation is what matters. By the substitution property of the
Accordance Principle \cite{Nexus}, the argument runs at \emph{any}
rung: a trait of ours may itself be placed in the conditioning seat, and
alternative natural histories are then deduced to fail its occurrence
almost surely. The manoeuvre requires the trait to be identifiable in
histories where it does not occur, which is a task-object
identification rather than an object one
(\S\ref{to:sub:licenses}). So for any Golden-relevant trait one names,
perturbations
of natural history systematically undermine its occurrence probability.

\begin{corollary}[Local optimality of the record; \textbf{[S]}]
\label{gc:cor:localopt}
If for every Golden-relevant trait --- individuated as a task-object,
per \S\ref{to:sec} --- the perturbed histories systematically
depress its probability, then our Aethus locally maximises the
probability of the Golden Point among its neighbours --- and our natural
history is accordingly the past-light-cone restriction of the optimal
Aethus for the Golden Nexus.
\end{corollary}

\begin{remark}[The grade, stated honestly]\label{gc:rem:grade}
Corollary~\ref{gc:cor:localopt} is \textbf{[S]} and not
\textbf{[P]}. It inherits the substitution property's own licence ---
rung-screening, flagged as such at the source \cite{Ben26Optimal} ---
and it establishes \emph{local} optimality among interrogated
perturbations, not global optimality over the space of histories. What
would strengthen it is a perturbation exhibited that leaves a
Golden-relevant trait's probability \emph{undiminished}; what would
break it is a systematic class of such perturbations.
\end{remark}


% ---------------------------------------------------------------------
Past-Alignment as stated records an agreement. It does not yet show the
agreement is more than coincidence, and the two legs that follow are the
argument that it is not.

% ---------------------------------------------------------------------
\subsection{The first leg: natural history}\label{gc:sub:conditioning}

Past-Alignment as stated is an observation about agreement. The
following argues that the agreement is not coincidental --- that
Nexic natural history was never conditioned on Optimation in the first
place, but on the Golden Nexus throughout.

\begin{principle}[Golden-Conditioning; \textbf{[S]}, by reductio]
\label{gc:prin:conditioning}
Suppose, for contradiction, that the Accordance Principle conditions on
personal observer-existence rather than on the Golden Point.

Observerhood is robust under Aethic perturbation. One may strip
probabilistic prerequisites of the Golden Point --- rocket capability,
the presence of industrialisation-driving coal --- and replace them by
their complements without disturbing one's standing as an observer. A
world in which Earth simply lacked the resources to begin an industrial
revolution is a world whose inhabitants are no less observers.

Now select a handful of such attributes: necessary for the Golden Point,
unlikely to survive reopening in the Aetherian Nexus, and admissible in
the sense of Remark~\ref{gc:rem:boundrequires} --- reopening them must
leave the dependent structure intact, so that the resulting Aethae
remain individuated by the same task-objects as ours
(\S\ref{to:sub:licenses}). The totality of
the counterfactual Aethae so obtained carries \emph{substantially more}
Aethic weight than the Aetherian Nexus, while carrying a comparatively
minute probability of the Golden Point.

If Accordance were conditioned on observer-existence alone, it would
therefore place us in that high-weight counterfactual sum rather than in
the Aetherian Nexus. But the Aetherian Nexus is in fact our parent
Aethus. Contradiction. Hence the Accordance Principle is not conditioned
on observers, so long as ``observer'' is robust enough to survive these
perturbations.
\end{principle}

\begin{remark}[The Trace Connection, named in 2022]
\label{gc:rem:trace}
The claim this subsection advances was posed as an open problem, under
its own name, four years before the present argument. The author's
\emph{Golden Science} of August 2022 \cite{GoldenScience22} sets out the
Golden Principle in four numbered points, of which the third states that
our observations are far closer to representing those of an average
Golden-Point biosphere than those of an average position in the universe
--- Past-Alignment in prose --- and the fourth states that some
connection, causal or causal-mimicking, holds between our existence at
this time and place and the Golden Point itself. That paper then names
the connection and declines to supply it: it is to be called \emph{the
Trace Connection} until it is officially derived.

\S\ref{gc:sub:conditioning} is offered as that derivation. What the 2022
statement supplies, and what no later argument could manufacture for
itself, is that the gap was identified, named, and dated before any
candidate filling existed --- so the reductio cannot be read as having
been reverse-engineered to fit a conclusion already in hand. The same
document also carries, in prose and without the stochastic apparatus of
\S\ref{gr1:sec}: the hazard-function framing of biosphere survival; the
observation that an exponentially decaying hazard yields a strictly
positive limiting survival probability; the division of the failure mode
into internal and external channels, together with the insistence that a
Golden biosphere must lessen \emph{both} rather than either; the thesis
that domination is inefficient because it feeds the internal channel,
which is $\Phi3$; the Kardashev critique including the
self-consumption argument; and the definition of the Golden Point as the
point past which a biosphere cannot regress, drawn there against the
Hubble horizon. The formal content of
\S\S\ref{gr1:sec}--\ref{ad:sec} is later; the apparatus it formalises is
not.
\end{remark}

\begin{corollary}[What remains; \textbf{[S]}]\label{gc:cor:remains}
What the reductio leaves standing is the empirical artefact of
Past-Alignment itself: conditioning on the existence of the Golden Point
somewhere in the Aethic block universe reproduces exactly the
inferential form of the Accordance Principle that yields the
stated-attributes of Nexic reasoning. Among the contingent natural
histories one might find oneself at the end of --- a fungal world, a
dinosaurs-survive world, a wholly lemurine one --- the one that must be
conditioned on to reproduce the Nexic state of affairs we actually
observe is the one in which the Golden Point exists.
\end{corollary}

\subsubsection{Bounding the counterfactual weight}
\label{gc:subsub:weightbound}

The reductio of \ref{gc:prin:conditioning} turns on a single inequality:
that the counterfactual Aethae obtained by reopening Golden-Point
prerequisites carry substantially more weight than the Aetherian Nexus.
Stated bare, that is an assertion, and an objector may simply deny it ---
perhaps coal is common given a Carboniferous, perhaps the metallurgical
sequence is near-forced given hands and fire. This subsection replaces
the assertion with a construction, a bound, and a statement of exactly
what empirical input remains outstanding.

\paragraph{The construction.} Fix stated-attributes
$a_1,\dots,a_k$, each a prerequisite of the Golden Point and none a
prerequisite of observerhood. Let $\mathcal A$ denote the Aetherian
Nexus, in which every $a_j$ obtains, and let $\mathcal A_j$ denote the
Aethus obtained by reopening $a_j$ alone to its logical complement,
holding the remaining $k-1$ at their realised values.

\begin{lemma}[Disjointness; \textbf{[P]}]\label{gc:lem:disjoint}
The Aethae $\mathcal A_1,\dots,\mathcal A_k$ are pairwise disjoint, and
each is disjoint from $\mathcal A$. For $j\neq j'$, the pair
$\mathcal A_j$ and $\mathcal A_{j'}$ disagree on \emph{both} slot $j$
and slot $j'$; and $\mathcal A_j$ disagrees with $\mathcal A$ on slot
$j$.
\end{lemma}

\begin{proposition}[Weight ratio; \textbf{[P]}]\label{gc:prop:ratio}
Write $\alpha_j$ for the probability that slot $j$ takes its realised
value, conditional on the remaining slots taking theirs. Then
\begin{equation}\label{gc:eq:ratio}
\frac{W(\mathcal A_j)}{W(\mathcal A)}=\frac{1-\alpha_j}{\alpha_j},
\end{equation}
and by Lemma~\ref{gc:lem:disjoint} the single-flip Aethae may be summed:
\begin{equation}\label{gc:eq:bound}
R \;:=\; \frac{W_\times}{W(\mathcal A)}
\;\ge\; R_1 \;=\; \sum_{j=1}^{k}\frac{1-\alpha_j}{\alpha_j}.
\end{equation}
\end{proposition}

\begin{remark}[What the bound does and does not require]
\label{gc:rem:boundrequires}
Three features of \eqref{gc:eq:bound} are worth isolating, because they
are what make it usable.

\emph{Independence is not required.} Only disjointness is, and that
holds by construction. Dependence among the $a_j$ changes what the
individual ratios \emph{are} --- the $\alpha_j$ become conditional
rather than marginal, as stated --- but not whether the terms may be
added. This matters because the attributes at issue are plainly not
independent, and an argument needing them to be would be unusable.

\emph{It is a lower bound, and deliberately a loose one.} The
single-flip Aethae are a vanishing fraction of the counterfactual set:
the full sum over all Aethae differing in at least one slot is
$\prod_j\alpha_j^{-1}-1$, larger by many orders. But the full sum
requires the joint structure, and the reductio needs only that $R\gg1$.
An undercount that suffices is preferable to an exact figure that is
contestable.

\emph{Admissibility is governed by the Imperfection Magnitude Drop
criterion.} A perturbation enters the sum only if it leaves the
dependent structure intact --- if reopening $a_j$ does not re-randomise
what depends on it. A perturbation that triggers a drop does not yield a
clean single-slot flip, and Lemma~\ref{gc:lem:disjoint} fails for it.
Coal in an inaccessible location is admissible on this criterion: it may
be probabilistically catastrophic without re-randomising the structure
downstream. A relocation of the lineage's speciation to another
continent is not, since the material that later mediates the outcome is
not present there to be held fixed.
\end{remark}

\paragraph{How large must $R$ be, and how improbable must the
attributes be?} The reductio requires only $R\gg1$. Evaluating
\eqref{gc:eq:bound} for $k$ attributes each at probability $\alpha$:

\begin{center}\small
\begin{tabular}{@{}rrrrr@{}}
\toprule
$\alpha$ & $k=1$ & $k=3$ & $k=5$ & $k=8$\\
\midrule
$10^{-1}$ & $9$ & $27$ & $45$ & $72$\\
$10^{-2}$ & $99$ & $297$ & $495$ & $792$\\
$10^{-3}$ & $999$ & $3.0\times10^{3}$ & $5.0\times10^{3}$ & $8.0\times10^{3}$\\
$10^{-4}$ & $1.0\times10^{4}$ & $3.0\times10^{4}$ & $5.0\times10^{4}$ & $8.0\times10^{4}$\\
\bottomrule
\end{tabular}
\end{center}

\noindent The bar is far lower than the informal argument suggested. A
\emph{single} attribute at $\alpha=10^{-3}$ delivers $R\sim10^{3}$
unaided; no portfolio is needed, and the burden reduces to establishing
one prerequisite as genuinely improbable. Correspondingly, the objection
that the counterfactual sum might be comparable requires \emph{every}
candidate prerequisite to have $\alpha$ of order one --- that coal
placement, metallurgical sequence, and the rest are each more likely
than not given the rest of the record --- which is a considerably
stronger claim than the objection is usually asked to bear.

\paragraph{Estimating $\alpha$: the sister-case method.} The remaining
empirical input is supplied by an instrument the companion already
develops, and it should be used rather than reconstructed
\cite{Nexus}.

By the Theorem of Aimless Optimation, Optimation is the Cosmic Nexus
conditioned on the attribute $H_N$ that an Aetherian exists in one's
biosphere; and if an attribute $X$ is statistically independent of
$H_N$, the two Nexae assign it the same weight. Independent attributes
are therefore \emph{controls}: their observed frequency estimates the
Cosmic Nexus probability directly, unlifted by conditioning. The
companion's worked case is the Catarrhini node --- Old World monkeys
furnish over a hundred extant species across twenty-five million years
with no attainment of ape-level intelligence, and their fate being
independent of $H_N$, that frequency transfers as the baseline for our
own lineage at the same node.

This is exactly the quantity \eqref{gc:eq:bound} requires. For an
attribute $a_j$ admitting such a control, $\alpha_j$ is read off the
control's frequency, and its contribution to $R_1$ follows.

\begin{remark}[A corollary on once-only attributes, and the independence
check that governs it]\label{gc:rem:onceonly}
A corollary of the same theorem bears on attributes observed exactly
once, and it is worth stating as an inference rather than a
consistency claim. Once a first occurrence has secured what the Golden
Point requires, a \emph{second} occurrence is statistically independent
of $H_N$ --- it changes nothing about whether an Aetherian exists ---
and by the theorem it therefore retains its Cosmic Nexus weight and is
not lifted. Optimation supplies what is required and does not
overshoot.

The consequence is a discriminating observation. Write $\lambda$ for
the unconditioned expected count across all opportunities. Under the
hypothesis that $\lambda\ll1$, Optimation must lift the attribute for
$H_N$ to obtain, and lifts it to exactly one; \emph{exactly one
occurrence is what that hypothesis predicts}. Under the hypothesis that
no lift was needed, occurrences are Poisson at rate $\lambda$ and
exactly one is a coincidence of probability $\lambda e^{-\lambda}$. The
likelihood ratio between them is $\left(\lambda
e^{-\lambda}\right)^{-1}$:
\begin{center}\small
\begin{tabular}{@{}rrrrrr@{}}
\toprule
$\lambda$ & $1$ & $3$ & $5$ & $10$ & $20$\\
\midrule
ratio favouring lift & $2.7$ & $6.7$ & $30$ & $2.2\times10^{3}$ &
$2.4\times10^{7}$\\
\bottomrule
\end{tabular}
\end{center}
So a single occurrence tells strongly against the configuration having
been \emph{common} --- $\lambda\ge5$ is disfavoured thirtyfold or worse
--- while discriminating only weakly between a genuinely minute
$\lambda$ and one of order unity. Its evidential role is accordingly to
exclude the objection that the prerequisites were unremarkable, not to
supply the magnitude. The magnitude comes from the independent
per-opportunity estimate, and where that gives $\lambda\ll1$ the
attribute contributes approximately $1/\lambda$ to $R_1$.

The corollary is only as good as its individuation, and this is where
the difficulty sits. Rather than proceed directly to an estimate, we
calibrate the instrument on two attributes for which the data are good.
Both return nothing, and the pattern in \emph{how} they return nothing
is what fixes the strategy.

\paragraph{Calibration I: coal.} Britain's coal is the obvious first
candidate, and it fails at the endowment grain. Carboniferous and
younger basins are known on virtually every continent --- the
Appalachians, the Ruhr, the Donets, the Kuznetsk, the Nord-Pas de
Calais--Saar--Lorraine field, and others --- giving eight or more on any
reasonable individuation, and Britain's is not the most favourable, the
Appalachian seams being unusually large and unusually exposed. With
$\lambda\approx8$ we obtain $\alpha=1-e^{-\lambda}\approx0.9997$ and a
contribution to $R_1$ of about $3\times10^{-4}$: nil.

The verdict is confirmed by the author whose thesis is most at stake.
Allen observes that the Ruhr went unexploited not for want of coal but
because Dutch peat held fuel prices down, removing the return on
improving Rhine transport; and he draws the counterfactual explicitly,
that had German coal been developed three centuries earlier the
industrial revolution might have been a Dutch--German achievement rather
than a British one \cite{Allen09}. He identifies Belgium, around
Li\`ege and Mons, as possessing a comparable ratio of wages to energy
cost. So the claim that the path to the twentieth century ran through
northern Britain was always about the \emph{conjunction} of cheap energy
with dear labour and the institutions to exploit both --- never about
the seams, which were widely available. A reader who recalls that claim
should read it that way, and the estimator's verdict agrees with Allen
rather than contradicting him, reached independently and from geology.

\paragraph{Calibration II: domesticable megafauna.} A second candidate,
drawn from a different literature, fails in the same manner. Diamond's
count gives fourteen large mammal species domesticated of some hundred
and forty-eight candidates, with thirteen of the fourteen wild ancestors
confined to Eurasia \cite{Diamond97} --- a striking asymmetry, and the
basis of a well-known thesis. But the relevant question for the
estimator is not how the successes distributed; it is how many
\emph{continents} possessed a workable endowment. Diamond's own
candidate counts are seventy-two for Eurasia against fifty-one for
sub-Saharan Africa, which is not the profile of a rare attribute. Taking
continents as the opportunities gives $\lambda\approx2$ and a
contribution of about $0.16$: nil again.

Diamond's \emph{conjunction} does better --- megafauna together with an
east--west principal axis and the crop package on one landmass gives
$\lambda\approx0.16$ and a contribution near $5.8$ --- which is a real
figure, and still four orders below what a single attribute would need
to carry the reductio alone.

\begin{proposition}[The bound on \emph{parallel} attributes;
\textbf{[S]}]\label{gc:prop:squeeze}
The two failures share a cause, and it is structural rather than
accidental. Call an attribute \emph{parallel} when its opportunities are
a set of co-existing sites --- landmasses, basins, continents --- so
that $\lambda=Np$ with $N$ the site count. For an attribute individuated
at continental scale, $N\approx5$. To obtain
$(1-\alpha)/\alpha\sim10^{3}$ then requires $\lambda\sim10^{-3}$, hence
a per-site probability of order $2\times10^{-4}$ --- a claim that cannot
responsibly be defended for any attribute of the form ``this landmass
possesses feature $X$''. Deep time can deliver minute $\lambda$ for
parallel attributes because its site count is astronomical; recent
history cannot, because its site count is five.

Coal and megafauna are both parallel attributes. That, and not their
recency, is why they failed.
\end{proposition}

\paragraph{Serial attributes escape the bound.} Proposition
\ref{gc:prop:squeeze} is a claim about parallel attributes only, and the
restriction is not a technicality --- it identifies the wrong turn that
coal and megafauna represent. An attribute may instead be
\emph{serial}: a conjunction of successive branch points within a single
trajectory, each of which had to go one way for the next to be reached.

\begin{proposition}[Serial attributes; \textbf{[S]}]
\label{gc:prop:cascade}
Let an attribute consist of $m$ successive branch points, the $i$th
resolved favourably with probability $p_i$ conditional on its
predecessors. Then $\alpha=\prod_{i\le m}p_i$ and the contribution to
$R_1$ is $\alpha^{-1}-1$, which grows \emph{exponentially in $m$}. At
$p=0.75$ throughout: $m=8$ contributes about $9$, $m=16$ about $99$,
and $m=25$ about $1.3\times10^{3}$.
\end{proposition}

\noindent The opportunity structure is what differs. A parallel
attribute is bounded by how many sites exist, and there are five
continents. A serial attribute is bounded by how many branch points the
trajectory contains, and a well-articulated historical episode contains
many. Recent history is therefore not weak evidence; it was being
individuated in the form that makes it weak.

This is how the candidates should be stated. Not ``Alexander existed'',
which is a parallel attribute of no force, but \emph{that he survived
and succeeded through the entire sequence} --- the wound at the
Granicus, Issus, Gaugamela, the near-fatal injury among the Malli, the
absence of a disabling illness before consolidation, and the
Hellenisation surviving the immediate fragmentation of the empire ---
each a branch point with its own $p_i$. Likewise the Roman survival of
the Second Punic War is a sequence, not a fact: the recovery after
Cannae, Hannibal's failure to march, the holding of the Italian allies,
Scipio's African campaign. Likewise the passage of Christianity from a
provincial sect to an imperial religion, and its part in what followed.
Each such episode is one attribute in \eqref{gc:eq:bound}, and each
carries a contribution set by the length of its cascade rather than by
the number of continents.

\begin{remark}[What serial attributes cost]\label{gc:rem:serialcost}
The gain in magnitude is paid for in the independence check, and the
price should be stated plainly, since it is where such an estimate would
fail.

For each branch point one must hold that its failure had no adequate
\emph{substitute} --- that had Alexander died at the Granicus, the
Hellenisation of the East does not simply proceed under another
commander. History is rich in substitutes, and each $p_i$ that admits
one contributes nothing, since the attribute is then not load-bearing
for $H_N$ at that step. The estimator's discipline accordingly transfers
from grain, where it sat for parallel attributes, to substitutability.

Two consequences follow. Cascade length must be counted in
\emph{substitution-free} branch points, not in dramatic incidents, which
will reduce $m$ considerably below the narrative count. And the
per-step probabilities must be conditional on the predecessors, since
the steps are plainly not independent --- a commander who has survived
four battles is not a random draw at the fifth. Both corrections push
$\alpha$ upward, and an honest estimate should expect the eventual
contribution to fall well short of the naive product.
\end{remark}

\begin{remark}[Individuate the branch point by its role, not its
occupant]\label{gc:rem:roleindividuation}
The substitutability audit has a corollary that governs how each branch
point must be \emph{stated}, and getting it wrong is the most likely
route to an inflated estimate.

The counterfactual at each step must be taken over the \emph{role}, not
over the occupant. The question at the relevant branch point is not
whether Christianity in particular arose and spread, but whether
\emph{anything} arose that discharged the role it discharged --- in the
window, with the reach, and with the consequences for what followed. So
individuated, the competing cults of the period are not rivals to the
hypothesis but candidate fillers of the same role, and their availability
raises $\alpha$ at that step rather than lowering it: where a role admits
$\kappa$ independent candidates each of probability $q$, the
role-probability is $1-(1-q)^{\kappa}$, not $q$.

Role individuation is therefore the \emph{conservative} choice, and that
is precisely its merit. It concedes the substitution objection in
advance rather than meeting it afterward: a critic who replies that some
other movement would have served is not refuting the estimate but
restating the quantity already being estimated. What survives such an
audit is a residue of steps at which no filler was available with
comparable probability --- and the claim, where a cascade is offered, is
about that residue.

It also relocates where improbability may legitimately be found. It is
not to be sought in the striking identity of the occupant, which is
where narrative attention naturally goes, but in the narrowness of the
window in which \emph{any} occupant could have appeared. The
transitional periods are the interesting ones for that reason, and the
dramatic incidents largely are not.
\end{remark}

\begin{remark}[The limitation is present ignorance, not the structure]
\label{gc:rem:measurement}
One further honesty is owed about the status of these estimates, and it
cuts in an unusual direction.

The obstacle to evaluating $\prod_i p_i$ for any actual cascade is not
that the quantity is ill-defined or that the method is unsound. It is
that counterfactual history is not yet in a state to supply the
$p_i$ with defensible error bars --- the substitutability audits and the
window widths that Remark~\ref{gc:rem:roleindividuation} requires are
matters on which the discipline has views but not measurements. The
uncertainty is \emph{ours}, and contemporary; it is not a property of
the cascades.

Two things follow. The paper should not, and does not, offer a numerical
cascade estimate: what it offers is the structure
\eqref{gc:eq:bound}, the serial--parallel distinction, and the
individuation discipline, with the arithmetic left for inputs that do
not yet exist. And the expectation that such inputs will exist is a
substantive bet rather than a certainty --- one that quantitative
counterfactual history will mature to the point of supplying audited
branch-point probabilities. If it does, the enumeration becomes
routine and the bound follows. If it does not, the weight claim remains
structurally established and empirically unfilled, which is a weaker
position than we would like and a stronger one than assertion
\textbf{[O]}.
\end{remark}

\begin{remark}[An unnamed significance is a prediction, and it has a
falsification test]\label{gc:rem:unnamed}
Remark~\ref{gc:rem:measurement} leaves the framework in a position that
looks, at first, like an unfalsifiable auxiliary. We do not know what
the establishment of a universalising cult in the late-antique
Mediterranean forced through downstream; the analysis nonetheless
concludes that something Golden-relevant was forced through, on the
strength of \ref{gc:prop:flukes} alone. Stated carelessly this is the
canonical bad move --- every fluke acquires an unnamed significance,
nothing counts against the framework, and the Accordance analysis
becomes a device for dignifying whatever happened.

It is not that, and the reason is that the unnamed significance is a
\emph{prediction} with a definite content and a definite refutation.

\emph{Content.} For a cascade $C$ satisfying the two admissibility
conditions --- genuinely improbable, and load-bearing in the sense of
Remark~\ref{gc:rem:roleindividuation} --- the claim is that there exists
a consequence of $C$ on which the probability of the Golden Point
depends, and that it will be identified. The claim is made before the
consequence is known, which is what makes it a prediction rather than a
description.

\emph{Refutation.} Exhibit a route to the Golden Point that does not
require $C$. If the terminus is reachable without the cascade, then
$C$ was not load-bearing for it, and the predicted consequence does not
exist to be found. This is the substitutability audit of
Remark~\ref{gc:rem:serialcost} run at the \emph{terminus} rather than at
an intermediate --- which is the only place it can be run when the
intermediate is unnamed, and which is a harder audit rather than an
easier one, since the chain to be traced is longer.

\emph{And the criterion is diachronic.} A single unnamed significance
is a debt; a growing ledger of them is a symptom. The framework is
progressive if the ratio of identified to unidentified consequences
improves as counterfactual history matures, and degenerating if the
unnamed list only lengthens while the named one does not. We state this
as the test we intend the programme to be held to: \emph{if the
historiography improves and the significances remain unnamed, that
counts against the framework and is not to be absorbed by it.} The
passage of time is thereby made a test rather than an excuse, which is
what the position in Remark~\ref{gc:rem:measurement} would otherwise
lack.

One asymmetry should be noted, since it limits what confirmation is
available. Failure to exhibit a $C$-free route is weak evidence at best,
because such routes are counterfactual claims subject to the same
immaturity that motivated the remark. The test is therefore one-sided in
the ordinary way: it can refute a particular significance claim and
cannot establish one, and the framework's standing rests on
whether the identifications accumulate.
\end{remark}

\paragraph{The consequence: aggregate, but of serial attributes.} The
reductio requires $R\gg1$, not $R>10^{3}$, and $R_1$ is a \emph{sum}.
Even under the corrections of Remark~\ref{gc:rem:serialcost}, a modest
stockpile suffices: twenty attributes contributing about ten apiece
give $R_1\approx200$, and a single well-attested cascade may supply as
much on its own.

\begin{center}\small
\begin{tabular}{@{}rrr@{}}
\toprule
attributes & each $\approx5$ & each $\approx10$\\
\midrule
$10$ & $50$ & $100$\\
$20$ & $100$ & $200$\\
$40$ & $200$ & $400$\\
\bottomrule
\end{tabular}
\end{center}

\noindent So the programme is an enumeration of serial contingencies,
each articulated as a cascade and each audited for substitutability. The
record supplies candidates readily --- the military sequences above; the
demographic rupture of the Black Death and its effect on the very wage
structure Allen's account requires; and the theses of \cite{Diamond97}
beyond megafauna, several of which admit serial rather than parallel
statement. What the two calibrations established is not that recent
history is barren but that it must be interrogated along the trajectory
rather than across the map.

\begin{remark}[Why the negative calibrations are retained]
\label{gc:rem:whyretain}
Two attributes were examined and both returned nothing. It would be
tidier to omit them, and it would be a mistake.

They demonstrate that the instrument has teeth. A method that never
returns nothing cannot discriminate, and a reader entitled to suspect
that the estimator was constructed to deliver a congenial answer is
answered by two cases where it did not. They also fix the grain
discipline \emph{before} any positive estimate is offered, so that the
eventual stockpile is credible rather than convenient. And they preempt
the two objections a reader arrives with --- that coal was surely the
scarce thing, and that Diamond's asymmetry surely settles the matter ---
with the estimator's verdict rather than with assertion.

Most importantly, the pattern in the two failures is itself the
finding. Had one failed and the other succeeded, the first would read as
a red herring. Both failing, and failing for the same reason, is what
establishes Proposition~\ref{gc:prop:squeeze} and thereby the
aggregation strategy. The negative results are load-bearing.
\end{remark}

This is where the argument's remaining empirical debt sits, and it is
now specific and bounded. What \eqref{gc:eq:bound} establishes is the
structure: given the $\alpha_j$, the bound follows. What
Proposition~\ref{gc:prop:squeeze} establishes is that no single recent
attribute will supply it, so the debt is a stockpile of order twenty
contingencies, each estimated by the sister-case method and each passing
the independence check. The weight claim is accordingly no longer
asserted but \emph{conditionally computed} --- reduced from a
plausibility judgement about an unexamined sum to an enumeration of
bounded terms, each obtained by an instrument calibrated in advance and
each with its required magnitude stated \textbf{[O]}. What remains owed
is the enumeration, and \emph{that} is a research task rather than an
argumentative gap: the terms are individually modest, the sources are
identified, and the standard each must meet is fixed before any is
computed.
\end{remark}

\begin{remark}[The conditioning is not teleological]
\label{gc:rem:teleology}
A reader's first objection will be that conditioning on a future event
to account for present observations is teleology: the Golden Point
reaching backward to arrange its own prerequisites. The framework
answers this at its foundations rather than here, and the answer should
be cited rather than improvised.

The relevant result is the \emph{indiscriminate time principle}, a
corollary of the extrusion principle \cite{Aethus}, whose content is a
pair of mutual non-informations: knowing where the events an attribute's
blanks refer to sit in space and time tells one nothing about when those
blanks resolve for oneself, and knowing when they resolve for oneself
tells one nothing about where they sit. Blankness is a property of the
indexing Aethus, never of a clock reading --- \emph{blanks lack tense},
and an observer's unlearned past is exactly as open as their unformed
future \cite{AethusBrief}.

The objection therefore rests on a premise the framework denies. It
assumes that the Golden Point's futurity makes it a different kind of
object from an unresolved attribute of the past --- that the not-yet-real
is a category. On the Aethic account there is no such category: the
Golden Point is a blank attribute relative to our Aethus, in precisely
the sense that an unread outcome of the deep past is blank, and
conditioning on it is the same operation in both cases. The asymmetry
one feels is real but \emph{entropic} rather than ontological: the
future blocks discriminating information more stringently because it
sits downstream of the entropic arrow, while the past merely happens to
withhold it --- and informational uncertainty is the same thing under
either obstruction \cite{Aethus}.

The framework has already run this argument in the mirror-image case,
which is the strongest evidence that nothing is being improvised for the
present purpose. Newcomb's problem conditions on a \emph{past} event
--- the predictor's sealed placement --- that is blank relative to the
agent's information state, and the companion treatment turns on
declining exactly the assumption at issue here, that a past occurrence
is thereby a settled fact \cite{BenNewcomb}. Past-Alignment conditions
on a \emph{future} event that is blank relative to ours. Same operation,
opposite direction, and by the indiscriminate time principle the
direction is not a difference. A reader who accepts the Newcomb
treatment and rejects this one owes an account of which of the two
mutual non-informations they mean to keep.

What the Golden Point does here is therefore what any conditioning event
does: it selects, and does not act. The Aetherian Nexus is not
\emph{caused} to contain rocketry by a later passage; the weighting
conditioned on that passage assigns overwhelming weight to histories
containing it, and we read the record we are in. Whether anything
stronger than selection holds --- whether the future cone is constrained
rather than merely uninformative --- is exactly what
\S\ref{gc:sub:extrapolation} declines to settle, and the Void hypothesis
is the position on which nothing stronger holds at all.
\end{remark}

\begin{remark}[Why this is not circular, and where it could fail]
\label{gc:rem:circularity}
The obvious objection is circularity: conditioning on the Golden Point
and then finding its prerequisites is no achievement, since the
prerequisites are what the conditioning selects for. The objection would
land against a bare appeal to Past-Alignment, and it is worth being
explicit that it does not land here, because the argument of
\ref{gc:prin:conditioning} is \emph{comparative} rather than
confirmatory. Two candidate conditioning targets are placed side by
side; they issue \emph{different} predictions about the record --- one
indifferent to industrial prerequisites, one not --- and the record
discriminates between them. A comparison that could have gone the other
way is not a circle.

Three ways it could fail, in decreasing order of seriousness. The
weight claim remains the load-bearing one, though
\S\ref{gc:subsub:weightbound} has reduced it from an assertion to a
single measurement: the bound \eqref{gc:eq:bound} is proved, and one
prerequisite at $\alpha\sim10^{-3}$ suffices, so what would remove the
contradiction is now the specific demonstration that \emph{every}
candidate prerequisite has $\alpha$ of order one. Second, the
argument requires ``observer'' to be robust under exactly these
perturbations; a sufficiently demanding notion of observerhood --- one
entailing industrial capacity, say --- would collapse the two
conditioning targets together and dissolve the reductio, though at the
cost of making ``observer'' do work it is not usually asked to do.
Third, the reductio establishes that the target is \emph{not}
observer-existence; that it is specifically the Golden Point is
supported by Past-Alignment rather than forced by the contradiction, and
some third target could in principle serve.
\end{remark}


% ---------------------------------------------------------------------
The first leg is complete but carries one identified weakness, flagged
in Remark~\ref{gc:rem:circularity}: it needs observerhood to be robust
under exactly the perturbations considered, and a sufficiently demanding
notion of observer could deny this. The second leg runs the same
argument on different material, and the weakness does not arise there.

% ---------------------------------------------------------------------
\subsection{The second leg: human history}\label{gc:sub:humanhistory}

The reductio of \S\ref{gc:sub:conditioning} ran on natural history and
carried one identified weakness: it requires observerhood to be robust
under the perturbations considered, and a sufficiently demanding notion
of ``observer'' could deny this
(Remark~\ref{gc:rem:circularity}). Human history supplies a second run
of the same argument on which that weakness does not arise.

\subsubsection{Flukes establish their own significance}
\label{gc:subsub:flukes}

\begin{proposition}[Significance is deduced, not argued; \textbf{[S]}]
\label{gc:prop:flukes}
Let $F$ be an event established as a \emph{fluke} --- of low outlook
probability $p$ --- with no claim yet made about its significance.
Suppose $F$ were statistically independent of the Golden Point. By the
Second Law its Golden Probability equals its outlook probability, $g=p$,
and we should expect with probability $1-p$ to occupy the counterfactual
$\neg F$. We do not. Hence $F$ is not independent of the Golden Point,
and by the Third Law its Golden Linkage exceeds the L-threshold
\cite{Nexus, GoldenScience22}.

The significance of a fluke is therefore \emph{deduced from its
flukiness together with our observing it}, and requires no independent
argument that it mattered. Aggregating over a set of flukes, the union
of their counterfactual Aethae carries a far greater Aethic weight than
our own, and the principle refuses to book our location within the
lesser one as brute luck.
\end{proposition}

\begin{remark}[This settles the contingency-determinism debate, in one
direction]\label{gc:rem:determinism}
Historians divide over whether the path to modernity was contingent or
overdetermined, and the dispute is usually conducted by weighing
narratives. Proposition~\ref{gc:prop:flukes} makes it an inference.

The determinist holds that the apparent flukes were substitutable: no
Alexander, then someone else; no British coal, then some other route.
Substitutability is exactly high $b$ --- a good chance of the Golden
Point without $F$ --- hence low Golden Linkage, hence $g \approx p$,
hence the expectation that we occupy $\neg F$. The Third Law bounds it
quantitatively: $b/a < p/(1-p)$, so the flukier the event, the more
severely substitutability is constrained by our observing it. The
determinist is not refuted by a better narrative; they are constrained
by an inequality whose input is the event's improbability alone.
\end{remark}

\subsubsection{The stockpile}\label{gc:subsub:stockpile}

Two cases, deliberately ordered so that the weaker rhetorically comes
first.

\emph{A single sword-stroke.} At the Granicus in 334 BCE, two months
into an eleven-year campaign, the Persian noble Spithridates brought a
battle-axe down on Alexander's helmet and split its crest; Cleitus the
Black severed his arm before the second blow \cite{Arrian,
PlutarchAlex}. Had the stroke landed, the expedition would have ended
within days of beginning, and with it the Hellenistic world, its
Alexandrias, and the transmission through which much of what followed
ran.

\emph{Where the coal happened to be.} Pomeranz argues that England was
a ``fortunate freak'' whose coal lay close to water and ports, making
steam economically feasible, while China's coal sat in the northwest,
far from the Jiangnan textile economy that might have used it ---
a geological contingency placing coal and the Americas nearer the
western than the eastern end of Eurasia \cite{Pomeranz00}. Allen adds
the labour side: Newcastle carried the highest ratio of labour to energy
costs in the world in the early eighteenth century, and only under that
ratio did mechanisation pay \cite{Allen09}. His summary judgement is
an Imperfection Magnitude Drop stated by an economic historian who has
never encountered one: there was a single path to the twentieth century,
and it ran through northern Britain.

A third case is available but requires reframing. Diamond's account of
Eurasian advantage \cite{Diamond97} is \emph{determinist} at the level
of human history --- geography made the outcome near-inevitable --- and
is criticised on that ground. What the present argument takes from it
lies one level down: the distribution of domesticable species and the
continent's east--west axis are contingent \emph{across possible
Earths} while determining \emph{within} ours. That is precisely the
Accordance profile, low outlook probability with high Golden Linkage,
and Diamond's determinism is at the wrong level rather than mistaken.

\subsubsection{Why this leg is the stronger one}
\label{gc:subsub:strongerleg}

Two reasons, and the second is the more important.

\emph{The robustness premise becomes uncontroversial.} The natural-history
reductio needs it granted that observers survive the removal of the
end-Cretaceous impact, which can be resisted. Nothing comparable is
available here. No one holds that human observers required Alexander's
survival, or Newcastle's coal seams, or Newton. Humans existed for
millennia before each and would have continued after their removal. The
premise that carried the weight and could be denied is here simply not
deniable.

\emph{The conclusion sharpens.} What these contingencies deliver is not
observers but specifically the \emph{technological} condition ---
skyscrapers, the internet, the capacity to leave a dying star.
Observer-conditioning is indifferent to every one of them;
Golden-conditioning is not. So the second leg does not merely repeat
the first at lower risk; it isolates what the conditioning must be
sensitive to.

\begin{remark}[The narrative-bias objection, and the test that answers
it]\label{gc:rem:narrativebias}
The serious objection is selection in the record rather than in the
world. Dramatic near-misses make better stories, are more often written
down, and are more often repeated; a history dense with them may report
a fact about historiography.

The framework predicts an asymmetry that narrative bias does not, and
the prediction is already on record. Attributes with high Golden Linkage
magnitude but \emph{low connectivity} to other attributes should have
low Golden Probability --- lucky avoidance of adversity outranks lucky
survival of it, so isolated close calls should be \emph{rare} in the
record while densely connected contingencies should be
\emph{numerous} \cite{GoldenScience22}. Narrative bias inflates both
alike; the framework inflates one and suppresses the other.

The discrimination favours the framework for a reason worth stating
plainly: \emph{the strongest cases are the dull ones}. Nobody narrates
the Newcastle labour-to-energy price ratio, and no popular account turns
on Jiangnan's distance from northwestern coal seams. Yet that is where
the peer-reviewed economic history converges \cite{Pomeranz00,
Allen09}. Were narrative selection generating the pattern, the
undramatic structural contingencies would be missing --- and they are
the best-attested items in the file. The Granicus is the vivid
illustration; the coal is the evidence, and the ordering of
\S\ref{gc:subsub:stockpile} inverts the rhetorical one deliberately.
\end{remark}


% ---------------------------------------------------------------------
Both legs concern the past cone, and neither speaks to the future. That
silence is deliberate --- \S\ref{gc:subsub:pastalign-statement} recorded
it as scope rather than modesty --- and the three positions available on
the forward question are set out next without adjudication.

% ---------------------------------------------------------------------
\subsection{Extrapolating forward: three hypotheses}
\label{gc:sub:extrapolation}

Past-Alignment is a claim about the past cone and is silent ahead. Three
positions on the forward question are available, and this section states
them without adjudicating.

\begin{definition}[Extrapolation hypotheses; \textbf{[O]}]
\label{gc:def:extrapolation}
\begin{description}\itemsep3pt
\item[Void.] The alignment of our past cone with the Golden Nexus's is
a brute and uninfluential coincidence. Whether or not some connection
obtains, no law or inferential rule constrains our future cone to align
with the Golden Nexus's at the same spacetime position.

\item[Weak.] There is an explicit distribution over futures: the Nexus
describing our future cone assigns each proper child Aethus a weight
proportional to some function of its probability of the Golden Point ---
in the simplest case, that probability itself. Comparing two possible
futures is then possible by an Accordance-like metric, their relative
likelihood tracking their Golden-Point probabilities. Futures in which
the Golden Point does not occur are either incomparable under the
method or, in the strong reading, of probability zero relative to one's
Aethus.

\item[Strong.] The Golden Aethus is already stated in one's own Aethus,
and has been since the beginning of centric unfolding. On this reading
the stated attributes of Nexic reasoning were conditioning on the Golden
Point throughout; we occupy the Aethic block universe with a physical
moving present some centuries before that event, while the conditioning
that it occurs is already fully resolved and supplies exactly the
natural-historical content readable off our past cone. Contingency
remains in \emph{which} route is taken --- which stated attributes we
come to gather, conditioned on their consisting with the Golden Point
--- but outright Aethic disjointness from the Golden Point is
forbidden.
\end{description}
\end{definition}

\begin{remark}[Three questions the Strong hypothesis raises]
\label{gc:rem:strongquestions}
If the Strong hypothesis holds, three metaphysical questions become
pressing and none is answered here \textbf{[O]}: why future-cone Aethic
information should be available at all; how and why that information is
so intimately bound up with the initiation of centric unfolding, being
perhaps its earliest and leading piece of Aethic content; and what the
implications are for the entropic principle of Aethic reasoning and its
account of how the past is distributed between past and future Aethic
cones. Deriving either of the second or third from the other would be a
natural line of attack.
\end{remark}

% ---------------------------------------------------------------------
\subsection{Relation to the rest of this work}\label{gc:sub:relation}

Two compatibility questions arise immediately, and one structural
parallel is worth flagging as an open lead.

\emph{Compatibility with the vantage restriction.}
\S\ref{gr:sec} insists that Golden reasoning is a vantage and not an
obligation --- inferential, licensing conditional expectations rather
than duties, partly on the ground that $Q$ is a measure nothing
occupies. The Void and Weak hypotheses sit comfortably with that
restriction. The Strong hypothesis does not: if the Golden Aethus is
already stated and disjointness is forbidden, then the future-directed
reasoning has a purchase the vantage restriction disclaims. A reader
adopting Strong should understand that it buys something
$\Phi2$ deliberately gave away, and should re-price the argument
accordingly.

\emph{Compatibility with the epistemic ceiling.}
\S\ref{gr:sub:ceiling} holds that no finite observation can certify a
biosphere as Golden, since $Q$ and $\PP$ are locally equivalent and
separate only at infinite horizon. Past-Alignment appears to sit close
to this, since it asserts an agreement between observable content and
the Golden Nexus. The two are reconcilable, and the reconciliation is
that they concern different cones: the ceiling forbids distinguishing
$Q$ from $\PP$ over \emph{future} trajectories, while Past-Alignment
compares \emph{past} attribute content. Nothing in Past-Alignment
certifies that we shall reach the Golden Point; it claims that our
record is what conditioning on that event would have produced, which is
a claim about behind us.

\emph{An open structural parallel.} The trichotomy of
Definition~\ref{gc:def:extrapolation} is suggestively parallel to the
regime trichotomy of \S\ref{gr:subsub:regimes} --- trapped, critical,
viable. In particular the Strong hypothesis resembles the viable
regime, where $\PP(T=\infty)>0$ and $Q\ll\PP$ makes the Golden Biosphere
an ordinary sub-population one may belong to, while Void resembles the
trapped case in which no passage occurs. Should the parallel be exact,
the extrapolation question would not be independent of the anatomy
question: which hypothesis obtains would be settled by $\rho$ and by the
driver's scale function rather than by metaphysics, and
Theorem~\ref{gr1:thm:twocond} would decide it. We flag this as a lead
and do not claim it \textbf{[O]}.

\begin{remark}[The reversal of order]\label{gc:rem:reversal}
If the argument of this section is accepted, the framework's
architecture inverts. Optimation is no longer the starting weighting but
the first derived result --- what Golden-conditioning looks like cropped
to the observable cone. The Golden Biosphere is no longer characterised
by survival asymptotics and then found interesting; it is the
conditioning target, and the survival asymptotics of
\S\ref{gr1:sec} are the account of what such a target consists in. The
sections that follow are then not a sequence of separate results but the
development of a single object from different sides: \S\ref{gr1:sec}
gives its dynamics, \S\ref{gr:sec} its ensemble structure and what may
be inferred from occupying it, and \S\ref{ad:sec} the institutional
machinery that would determine whether a given biosphere reaches it.
\end{remark}

% ---------------------------------------------------------------------
\subsection{Ledger}\label{gc:sub:ledger}

\paragraph{Costs.} \emph{Past-Alignment} (\ref{gc:prin:pastalign}) ---
\textbf{[$\Phi$]}, the section's principal commitment;
\emph{falsifier:} an attribute of the record whose presence
Golden-conditioning does not predict, or a Golden-relevant trait whose
probability is undiminished under perturbation
(Remark~\ref{gc:rem:grade}). \emph{The weight claim} within
\ref{gc:prin:conditioning} --- \textbf{[S]}, asserted rather than
computed, and the reductio's load-bearing step;
\emph{falsifier:} a demonstration that the counterfactual sum does not
outweigh the Aetherian Nexus. \emph{Robustness of observerhood} ---
\textbf{[$\Phi$]}; \emph{falsifier:} a defensible notion of observer
entailing the industrial prerequisites, which would collapse the two
conditioning targets. \emph{Rung-screening}, inherited with the
substitution property \cite{Ben26Optimal}, and carrying its licence
unchanged.

\paragraph{Buys.} Optimation derived rather than posited, and derived
from a comparison that could have failed. A conditioning target that
explains what observer-conditioning cannot --- why the record contains
its industrial and technological prerequisites at all. An explicit
statement of what is \emph{not} claimed about the future, together with
the three hypotheses that exhaust the forward question. And a
reorganisation of the surrounding work into the development of one
object, stated in Remark~\ref{gc:rem:reversal}.

\paragraph{Inherited dependency.} Both arguments now depend on
technoanatomy being a \emph{process} rather than a parameter, for the
reason given in Remark~\ref{hr:rem:forcesprocess}: under a fixed
adverse $\rho$ the cascade pronounces the question closed before it is
asked, and the section's claims are then vacuously true and useless.
The dynamics of $\rho_t$ are open \textbf{[O]}
(\S\ref{gr1:subsub:fluct}), so this section's arguments are conditional
on a piece of the framework that is not yet built. That is a real
exposure and is not reduced by the arguments' internal soundness.

\paragraph{Not claimed.} That the future cone aligns with the Golden
Nexus; that the Golden Point will occur; that any of the three
extrapolation hypotheses is correct; and that the parallel with the
regime trichotomy of \S\ref{gr:subsub:regimes} is more than suggestive.

\section{Task-objects: individuation for a perturbed history}
\label{to:sec}

\begin{quote}\small
Every argument of \S\ref{gc:sec} quantifies over perturbed natural
histories. The dinosaurs never evolve; Africa develops otherwise; a
Golden Trek contains no Earth. Each such argument asks whether some
attribute survives the perturbation, and each therefore presupposes an
answer to a question it never asked: what makes a perturbed Earth still
an Earth, or a perturbed Aristotle still an Aristotle. Under an ontology
of possible worlds the question has a standard answer. Under an ontology
of Aethae it does not, and the perturbation arguments are ill-posed
without a replacement.

This section supplies the replacement, which the author's earliest dated
work already contains: the \emph{task-object}, an object individuated
not by perceived independence from other objects but by the contribution
it makes to a designated root. The concept was arrived at
independently and coincides with Kripke's distinction between rigid and
contingent designators \cite{kripke1980naming}; the departure from Kripke is
not in the distinction but in what is done with it, and in a denial of
the metaphysics his rigid designators presuppose
(\S\ref{to:sub:kripke}). Two liabilities are carried openly: the
correspondence between objects and task-objects is many-to-many in both
directions, so a selection principle is owed and only sketched
(\S\ref{to:sub:selection}); and the reply to the standard identity
objection runs through a commitment about consciousness whose cost is
priced in \S\ref{to:sub:identity}.
\end{quote}

% ---------------------------------------------------------------------
\subsection{The problem: perturbation requires individuation}
\label{to:sub:problem}

Consider the argument of \S\ref{gc:subsub:pastalign-argument} in the
form it was actually run. One perturbs the record --- the dinosaurs
persist --- and observes that the capacity to build rocket vehicles does
not survive. For this to be an argument rather than a stipulation, the
perturbed history must remain comparable to ours: it must still contain
a planet, a biosphere, a lineage, against which the missing capacity can
be registered as missing.

That comparability is what ordinary object-talk cannot supply here.
Perturb the record enough that no informed person would call the
resulting planet Earth, and the question ``does Earth still have
rocketry?'' does not become false; it becomes unasked. Yet the argument
needs it asked, because the Imperfection Magnitude Drop it is trying to
exhibit is precisely a drop \emph{in the substituted planet's}
contribution to the Golden Point.

The difficulty is sharpened by the framework's own commitments. Under an
Aethic ontology there are no completed determinate histories across
which an object could retain a categorical substance, so nothing plays
the role that a rigid designator's referent plays in the possible-worlds
picture (\S\ref{to:sub:kripke}). The individuation must therefore come
from somewhere else, and the proposal is that it comes from the root the
analysis is already conditioned on.

% ---------------------------------------------------------------------
\subsection{The definition}\label{to:sub:definition}

\begin{definition}[Task-object; \textbf{[$\Phi$]}]\label{to:def:taskobject}
Fix a root Aethus $\mathcal R$. A \emph{task-object} relative to
$\mathcal R$ is an object individuated not by its perceived independence
from other objects but by the contribution it makes to
$\mathcal R$. Where the root is the Golden Point, a \emph{Task-Earth} is
--- for instance --- the planet on which the makers of the Golden Point
speciated, and the corresponding \emph{task-humans} and
\emph{task-events} are individuated likewise.\footnote{%
\emph{Provenance.} The concept and the Task-Earth example are the
author's, recorded 1 July 2022 and reproduced in the natural-historical
development of August 2022 \cite{GoldenScience22}, which makes this the
earliest dated item in the present corpus --- predating the extrusion
principle by some weeks and the Golden Point material it was written to
serve by a month. The formalisation below, the placement against Kripke,
and the pricing of the identity objection are later.}
\end{definition}

\begin{remark}[Root-relativity is by design, not by accident]
\label{to:rem:rootrelative}
The root is a parameter. Task-object individuation is root-relative by
construction, and the Golden Point is the root this programme selects
rather than one the framework requires; a different inquiry with a
different root would individuate differently and legitimately. This
matters for how the concept should be read outside this paper: nothing
in the machinery is committed to prosperity, survival, or any other
particular optimand. What is committed is that \emph{some} root must be
designated before individuation is well-defined, which is the
substantive claim.
\end{remark}

\begin{remark}[Why this is not a redescription of ordinary objects]
\label{to:rem:notredescription}
Ordinary objects and task-objects come apart in both directions, and the
examples are the author's. Earth $3.9$ billion years ago is called
Earth, despite differing from the present planet in continent
configuration, atmospheric composition, and much else. A perfect
duplicate in another star system would not be called Earth, despite
differing in none of them. Neither verdict tracks contribution to a
root, and both are stable features of ordinary usage --- which is
precisely why ordinary usage cannot serve where the analysis needs
comparability across perturbation.
\end{remark}

% ---------------------------------------------------------------------
\subsection{Relation to Kripke, and the departure}\label{to:sub:kripke}

The distinction between a designator that picks out its object robustly
under counterfactual variation and one that picks it out by a
description which the object might not have satisfied is Kripke's
\cite{kripke1980naming}. The present concept was arrived at independently, by a
different route --- asking what continuity there is between Earth's
persistence in the past and the prospects of a biosphere's near future
--- and the coincidence is worth recording as a convergence rather than
claimed as a discovery.

The departure is twofold.

\emph{First, in what is done with the distinction.} Kripke's conclusion
is prohibitive: do not use contingent designators where rigid ones are
required. The present move is instead to adopt the awkward category as
both ontologically apt and operationally useful --- the pattern by which
physics adopted Newton's first law over the intuition of motion
requiring a mover. That pattern is not general licence, and should not
be read as one: counterintuitiveness is usually a warning, conspicuously
so in moral matters, and the claim is only that this is one of the
uncommon cases in which the received intuition is out of step with the
ontology, plausibly because what served human societies through
evolution was never selected for tracking it.

\emph{Second, and more radically, in a denial of the metaphysics.} The
Aethic framework does not accept the objects to which rigid designators
are supposed to point. With Aethae rather than possible worlds as the
base unit, there is no categorical substance persisting across worlds
for such a designator to track: by the Aethic reversal principle, the
underlying determinate reality required to constitute such an object
cannot be well-defined in the first place \cite{Aethus}. The scope of
this denial should be stated carefully --- \emph{what is denied is
Kripke's metaphysical presupposition, not his linguistics.} The modal
argument survives untouched as a claim about how names function in
language; the present claim lives a level below, about what that
functioning tracks.

\begin{remark}[Robustness recovered without substance]
\label{to:rem:robustness}
What made rigid designation attractive was robustness under
perturbation, and that is recoverable without the metaphysics. Placing
several description classes of an object in Aethic superposition yields
a designation that survives perturbation by landing in another class
when one fails --- robustness as a property of the superposition rather
than of a substance beneath it. The strong form of the resulting
position is that an Aethic ontology must make task-objects fundamental
and admit rigid-designator talk only as non-arbitrary shorthand for
them.
\end{remark}

% ---------------------------------------------------------------------
\subsection{The two-way fan, and the selection problem}
\label{to:sub:selection}

\begin{proposition}[Many-to-many in both directions; \textbf{[S]}]
\label{to:prop:fan}
The correspondence between objects and task-objects is many-to-many.
One object admits indefinitely many task-objects: Earth may be
individuated as a spherical body composed largely of iron, or as a
region within a star's habitable zone, and so on without limit. One
task-object admits indefinitely many objects: fix Task-Earth as the
speciation site of the Golden Point's makers, and every planet
satisfying it qualifies --- including one on which Hawaii never formed,
which is not our Earth and is a Task-Earth.
\end{proposition}

The fan is the concept's principal technical liability, since it means a
task-object is not determined by the object it individuates, and some
selection is required. The author's own assessment is that the
selection carries irreducible speculation, but speculation that might be
linguistically systematised; the criterion offered is to prefer the
task-object yielding the simplest or most informative interplay with the
wider question the root encodes.

\begin{remark}[The criterion is a gauge choice, and the paper already
has one]\label{to:rem:gauge}
That criterion is recognisably a gauge fixing: one probes the freedom
until the description aligns most economically, exactly as one selects
coordinates in which a physical problem is cleanest. The comparison is
the author's and is more than decorative here, because
\S\ref{gr1:sub:coords} performs the same operation on the survival
model --- initial data there form a gauge orbit invisible to every
asymptotic exponent, and the analysis proceeds by identifying what the
freedom cannot touch. The two are the same methodological move at
different levels: find the quantities the arbitrary choice leaves
invariant, and state the theory in those.

This suggests where a systematisation would come from, and the framework
supplies a candidate it does not yet apply here. The construction for
degree-of-existence over description classes \cite{Aethus} selects
canonical classes as \emph{fixed points} --- those whose further
refinements no longer improve description quality --- which is exactly
the shape of criterion the task-object selection problem needs. Whether
that construction transfers is open \textbf{[O]}, and it is the first
thing we would attempt.
\end{remark}

% ---------------------------------------------------------------------
\subsection{The identity objection, and what answering it costs}
\label{to:sub:identity}

The standard objection is Kripke's own. Aristotle might not have taught
Alexander; he would have been Aristotle regardless; so the person is
what the name tracks, and no task-human individuation is required.
Under the present proposal the corresponding task-human --- the man who
taught Alexander in the sense of equipping him for some specified
subsequent achievement --- would simply not have been instantiated, and
the objection says this gets the identity facts wrong.

The reply is not to deny the objection's premise but to run it out.

\begin{remark}[The reply, as a reductio of the objection's own criterion]
\label{to:rem:reductio}
Grant that identity is settled by continuity of consciousness through a
common point of divergence, which is what the objection's force rests
on. Under centric unfolding, any other person's self-describing Aethus
is reachable by ascending to a sufficiently early ancestor Aethus ---
plausibly in infancy or before, wherever a substantial portion of the
universe enters disagreeing superposition --- and perturbing there. The
objection's own criterion therefore identifies not merely the untutored
Aristotle with Aristotle, but \emph{every} person with every other: the
set of human-pointing rigid designators collapses to a single
equivalence class.

The structure of the reply should be read carefully, since its force
depends on it. It is not a positive thesis that all persons are
identical, offered for acceptance. It is a demonstration that the
criterion the objection relies on, applied consistently under centric
unfolding, cannot deliver the fine-grained identity facts the objection
needs --- so the objection cannot be used against task-object
individuation without first supplying a criterion that survives its own
generalisation.

Two costs are recorded. The reply commits to reading continuity of
consciousness through the framework's own treatment of centric unfolding
\cite{Aethus}, which is a substantive commitment and not a neutral one.
And a reader may take the collapse as a reductio of the framework rather
than of the objection; the honest reply to that is that the collapse
follows from the objection's criterion together with centric unfolding
as stated, so contesting it while retaining Aethic reasoning would
require an extra-centric unfolding proposal, which nothing here
supplies \textbf{[O]}.
\end{remark}

% ---------------------------------------------------------------------
\subsection{What this licenses}\label{to:sub:licenses}

With individuation in hand, three things earlier in the paper become
well-formed rather than merely persuasive.

\emph{The perturbation arguments.} \S\ref{gc:subsub:pastalign-argument}
may compare our record with histories in which the dinosaurs persist,
because both contain a Task-Earth and a task-lineage against which a
missing capacity registers as missing. Without task-objects the
comparison has no common term and the argument does not run.

\emph{The substitution property.} The generalisation across rungs
\cite{Nexus} places a trait of ours in the conditioning seat and
deduces that alternative histories fail its occurrence. That manoeuvre
requires the trait to be identifiable in histories where it does not
occur, which is a task-object identification and not an object one.

\emph{Imperfection Magnitude Drops.} A drop is a fall in contribution to
the root under substitution \cite{Ben26Optimal}. Since a task-object
just is an object individuated by that contribution, the drop is a
change in a task-object's value rather than a comparison between
incommensurable objects --- which is what makes it a quantity at all.

\begin{remark}[Order of development]\label{to:rem:order}
The dating is worth one line, because it inverts the order of
presentation. Task-objects are dated 1 July 2022; the Golden Point
material follows in August \cite{GoldenScience22}; the extrusion
principle and with it tense-free blankness follow weeks after that. The
individuation scheme therefore came first and the ontology that
vindicates it came second --- the concept was constructed to serve a
question about a biosphere's prospects, and only later found itself
entailed by an ontology built for unrelated reasons. We record this as
provenance rather than as evidence, but it does bear on the charge that
task-objects are a device fitted to the argument they support.
\end{remark}

% ---------------------------------------------------------------------
\subsection{Ledger}\label{to:sub:ledger}

\paragraph{Costs.} \emph{Task-object individuation}
(\ref{to:def:taskobject}) --- \textbf{[$\Phi$]};
\emph{falsifier:} a perturbation argument of the kind
\S\ref{gc:sec} requires, shown to be well-formed under ordinary object
individuation, which would make the apparatus idle. \emph{The denial of
rigid designators' referents} --- \textbf{[$\Phi$]}, inherited from the
Aethic ontology \cite{Aethus} and not argued here;
\emph{falsifier:} the ontology's. \emph{The selection principle} ---
\textbf{[O]}, owed and not supplied; the fan of
Proposition~\ref{to:prop:fan} is real and a criterion beyond
``simplest or most informative'' is not given. \emph{The identity reply}
--- \textbf{[$\Phi$]}, resting on continuity of consciousness read
through centric unfolding, with the singleton collapse acknowledged as
a cost a reader may decline to pay.

\paragraph{Buys.} Well-formedness for the perturbation arguments, the
substitution property, and Imperfection Magnitude Drops, none of which
survives ordinary object individuation intact. A principled account of
why ordinary usage fails here rather than an appeal to its being
inconvenient. And a concept whose earliest dated statement precedes both
the material it serves and the ontology that vindicates it.

\paragraph{Not claimed.} That Kripke's modal argument is mistaken; that
all persons are identical; that the Golden Point is a privileged root;
or that the selection problem is solved.


\section{The model and its survivorship structure}\label{gr1:sec}

\begin{quote}\small
This section is mathematics. It establishes the object the philosophical
section reasons about, and it is organised by what that reasoning
consumes rather than by order of discovery: the object
(\S\ref{gr1:sub:object}), its coordinates (\S\ref{gr1:sub:coords}), the
verdict on survival (\S\ref{gr1:sub:verdict}), and the structure of
conditioning on survival (\S\ref{gr1:sub:cond}). Material that maps the
terrain around the model rather than establishing it --- horizon theory,
the free-energy background, the waiting room and gain profile, the
transfer computations, and the numerics --- is placed in appendices,
where it is available without standing in the way.

Claims are tagged \textbf{[P]} proved here, \textbf{[C]} classical or
cited, \textbf{[A]} conditional on the soft-wall Ansatz
(Ansatz~\ref{sh:ansatz}), \textbf{[S]} scaling-level, \textbf{[O]}
open. The tag \textbf{[C]} is load-bearing and frequent: the persistence
exponent, the harmonic function of the conditioned process, the
free-energy asymptotics of the background and the boundary
classification of one-dimensional diffusions are all imported, and this
section claims none of them. What it claims is the model, the
coordinates, and the assembly.
\end{quote}

\subsection*{Relation to the literature}

Four bodies of work border this section, and the reader deserves a
precise map of what is classical, what is adjacent, and what is claimed
as new.

\emph{The classical layer.} The state process is Kolmogorov's 1934
diffusion \cite{Kol34}, the primordial hypoelliptic example
\cite{Hor67}. Its persistence theory is a mature subject with a clean
lineage --- McKean \cite{McK63}, Goldman \cite{Gol71}, Lachal
\cite{Lac91}, Sinai \cite{Sin92}, Isozaki--Watanabe \cite{IW94},
Denisov--Wachtel \cite{DW15} --- surveyed in \cite{AS15} and, as the
\emph{random acceleration process}, possessing an extensive physics
literature reviewed in \cite{BMS13}. The conditioned (never-hitting)
process, our limit object in \S\ref{sh:sub:interp}, is constructed by
Groeneboom, Jongbloed and Wellner \cite{GJW99}. None of this is ours;
we import it.

\emph{The adjacent frameworks.} Our survivorship interpolation belongs
to two established genres. First, Theorem~\ref{sh:thm:Q} is a
\emph{penalisation} statement in the sense of Roynette--Vallois--Yor
\cite{RVY06}: a limit of Wiener-type measures reweighted by a
time-dependent functional, identified as a martingale change of measure
--- here the weight is survival itself, and the polynomial tail places
it in the clockless class. Second, the mode classification of
Theorem~\ref{sh:thm:modes} sits against the theory of quasi-stationary
distributions and $Q$-processes surveyed in \cite{MV12}: classical QSD
theory is at home in the exponential-tail class (our modes (I) and
(III)), while the critical polynomial mode admits no proper QSD in fixed
coordinates, its Yaglom object being self-similar. We use both
frameworks as positioning, not as inputs; the proofs here are
self-contained given the cited persistence theory.

\emph{The applied layer: what this is a model of.} Stochastic hazard
rates are not new, and the wrapper $S_t=\exp(-\int_0^th)$ with $h$ a
diffusion is entirely standard: it is the reduced-form, or
intensity-based, framework of credit risk \cite{Lan04} and of stochastic
mortality, where the force of mortality is routinely modelled by a
diffusion, the mean-reverting Brownian Gompertz specification of Milevsky
and Promislow \cite{MP01} being a canonical instance. What distinguishes
the present model is not the wrapper but \emph{where the noise is
placed}: standard practice diffuses $h$ or $\log h$, whereas we diffuse
$-\tfrac{d}{dt}\log h$, one integration further out. The consequences are
structural rather than cosmetic --- $\log h$ becomes $C^1$, the state
becomes degenerate and $3/2$-self-similar, and the governing persistence
exponent moves from $\tfrac12$ to $\tfrac14$ (Remark~\ref{sh:rem:iwp},
\S\ref{sh:subsub:tail}). We stress the point because it locates the
novelty precisely, and disclaims the rest.

\emph{The conceptual ancestor.} The thesis of \S\ref{sh:sub:interp} ---
that selection on survival makes the observed ensemble diverge
systematically from the unbiased one --- is not new in substance. It is
the frailty phenomenon of Vaupel, Manton and Stallard \cite{VMS79} and
Vaupel and Yashin \cite{VY85}: when the frail are removed first the
population hazard falls below the individual hazard, the gap widening
with age, producing mortality deceleration at the population level with
no corresponding deceleration for any individual. We claim no priority
over that insight, and regard this section as exhibiting its
\emph{critical dynamical} analogue. The differences are structural and
are the point. Frailty is a static random variable fixed at birth and the
analysis is a mixture computation over a population; here the randomness
is a driver evolving in time, and the object is a limit of path measures
singular with respect to the unbiased law
(Theorem~\ref{sh:thm:singular}). Where the frailty literature computes
how an observed hazard curve bends, the interpolation of
\S\ref{sh:sub:interp} identifies the limiting law of the survivor's
entire trajectory, together with the drift \eqref{sh:eq:G} that
conditioning manufactures.

\emph{What is claimed as new.} To our knowledge, the following are new
at least in assembly, and we flag them as this section's contribution:
the log-derivative placement of the noise, with its gauge analysis
(\S\ref{sh:subsub:dimensional}) and the two-sided soft-wall analysis
whose remaining gap is located exactly (Ansatz~\ref{sh:ansatz}); the
scale-function horizon trichotomy with mortal critical case
(\S\ref{sh:subsub:scale}); the adverse-drift survival rate
$3\nu^2/8\sigma^2$ with its front-load-and-coast optimal profile
(Proposition~\ref{sh:prop:advdrift}); and the survivorship
interpolation read as a monotone flow with the tail index as its
linear-response coefficient, classified across drivers into saturation,
transformation and concentration (\S\ref{sh:sub:interp}). Where a
statement merely re-expresses known results in this frame, the tags and
citations say so.

\begin{convention}\label{sh:conv}
For positive functions, $f\sim g$ means $f/g\to1$; $f\asymp g$ means
$0<\liminf f/g\le\limsup f/g<\infty$; $f\lesssim g$ means $f\le Cg$.
Hatted quantities are nondimensionalized by the intrinsic timescale of
Definition~\ref{sh:def:units}. $\Phi$ is the standard normal distribution
function; $\one_A$ is an indicator; $\Leb$ is Lebesgue measure. All processes
are defined on a filtered probability space satisfying the usual conditions,
and $W$ denotes a standard one-dimensional Brownian motion.
\end{convention}

% ---------------------------------------------------------------------
\subsection{The object}\label{gr1:sub:object}\label{tc:sub:model}\label{sh:sub:model}

% ---------------------------------------------------------------------

\subsubsection{Setup: hazard, credit, and the Kolmogorov diffusion}
\label{sh:subsub:setup}

\begin{definition}[Model]\label{sh:def:model}
Fix $\sigma>0$ and initial data $(h_0,\mu_0,D_0)\in(0,\infty)\times\R\times\R$.
Define the \emph{improvement rate}, \emph{credit}, \emph{hazard}, and
\emph{cumulative hazard} by
\begin{equation}\label{sh:eq:model}
\mu_t=\mu_0+\sigma W_t,\qquad
D_t=D_0+\int_0^t\mu_s\,ds,\qquad
h_t=h_0e^{-(D_t-D_0)},\qquad
\Lambda_t=\int_0^t h_u\,du,
\end{equation}
and write $x_t=\log h_t$, so that $\mu_t=-\dot x_t$. We also set
$I_t=\int_0^t W_s\,ds$, so that $D_t=D_0+\mu_0t+\sigma I_t$.
\end{definition}

\begin{definition}[Departure time; Cox construction]\label{sh:def:cox}
Let $E\sim\mathrm{Exp}(1)$ be independent of $W$ and define
$T:=\inf\{t\ge0:\Lambda_t\ge E\}\in(0,\infty]$. Then
$\PP(T>t\mid\Fc^W_\infty)=e^{-\Lambda_t}$ and
$\PP(T=\infty\mid\Fc^W_\infty)=e^{-\Lambda_\infty}$, with the convention
$e^{-\infty}=0$. The \emph{survival function} is
$S_t=\PP(T>t\mid\Fc^W_\infty)=e^{-\Lambda_t}$.
\end{definition}

\begin{definition}[Viability]\label{sh:def:viab}
The \emph{viability event} is $\Imm:=\{\Lambda_\infty<\infty\}$. On $\Imm$
the conditional survival probability $S_\infty=e^{-\Lambda_\infty}$ is
strictly positive; off $\Imm$ it vanishes. Thus
\begin{equation}\label{sh:eq:twotier}
\PP(T=\infty)\;=\;\E\!\left[e^{-\Lambda_\infty}\one_\Imm\right]
\;\le\;\PP(\Imm),
\end{equation}
and we distinguish throughout between \emph{viability} ($\PP(\Imm)$, the
probability that immortality is achievable) and \emph{immortality proper}
($\PP(T=\infty)$).
\end{definition}

The pair $(\mu,D)$ is the \emph{Kolmogorov diffusion} \cite{Kol34}:
\begin{equation}\label{sh:eq:kolm}
d\mu_t=\sigma\,dW_t,\qquad dD_t=\mu_t\,dt,\qquad
\Lc=\tfrac{\sigma^2}{2}\partial_\mu^2+\mu\,\partial_D .
\end{equation}

\begin{lemma}[Regularity]\label{sh:lem:hypo}
The free process $(\mu_t,D_t)$ is Gaussian with mean
$(\mu_0,\,D_0+\mu_0t)$ and covariance
$\bigl(\begin{smallmatrix}\sigma^2t&\sigma^2t^2/2\\
\sigma^2t^2/2&\sigma^2t^3/3\end{smallmatrix}\bigr)$;
in particular it admits the explicit smooth transition density written down
by Kolmogorov \cite{Kol34}. More generally, $\Lc$ is hypoelliptic: with
$X_1=\sigma\partial_\mu$ and $X_0=\mu\partial_D$ one has
$[X_1,X_0]=\sigma\partial_D$, so H\"ormander's bracket condition holds at
first order \cite{Hor67}. Consequently the process killed at rate $e^{x}$
retains smooth densities, which is the regularity relevant to the killed
problem of \S\ref{sh:subsub:reduction}.
\end{lemma}

\begin{proof}
Gaussianity and the moments follow from
Lemma~\ref{sh:lem:secondorder} below and linearity of
\eqref{sh:eq:model} in $W$; the bracket computation is immediate.
\end{proof}

Three structural remarks frame everything that follows.

\begin{remark}[No boundary in $h$]\label{sh:rem:noboundary}
The logarithmic parametrization makes $h>0$ automatic: there is no absorbing
or reflecting boundary in the hazard itself, and the entire problem concerns
the growth of the integrated driver.
\end{remark}

\begin{remark}[Credit is spent only through negative rate]\label{sh:rem:monot}
$D$ is nondecreasing exactly while $\mu\ge0$. Accumulated credit cannot be
destroyed except by first driving the improvement rate negative; this
monotone coupling is the source of all horizon intuition, and its
quantitative teeth are extracted in \S\ref{sh:subsub:softhorizon}.
\end{remark}

\begin{remark}[Nonparametric-Bayes reading of the model]\label{sh:rem:iwp}
Definition~\ref{sh:def:model} places an integrated Wiener process prior on
$-\log h$ --- equivalently, the classical cubic smoothing-spline prior of
Wahba \cite{Wah78}, since a prior whose second derivative is white noise
is precisely one whose first derivative is Brownian. The model is
therefore the canonical smoothness prior on the log-hazard rather than an
exotic construction, and the results below may be read as the survival
asymptotics that prior implies.
\end{remark}

\begin{lemma}[Exact gauge role of the initial hazard]\label{sh:lem:hgauge}
Write $\Lambda^{(h_0)}$ for the cumulative hazard with initial value $h_0$.
Then $\Lambda^{(h_0)}_t=h_0\Lambda^{(1)}_t$ pathwise; hence $\Imm$ does not
depend on $h_0$, while
$\PP(T=\infty)=\E[e^{-h_0\Lambda^{(1)}_\infty}\one_\Imm]$ is continuous and
nonincreasing in $h_0$, strictly decreasing whenever $\PP(\Imm)>0$, with
$\lim_{h_0\downarrow0}\PP(T=\infty)=\PP(\Imm)$ and
$\lim_{h_0\uparrow\infty}\PP(T=\infty)=0$.
\end{lemma}

\begin{proof}
The identity is \eqref{sh:eq:model}; the limits follow by monotone and
dominated convergence.
\end{proof}

Thus $h_0$ can never affect a dichotomy, only a value: it is the first and
most complete gauge degree of freedom. The same will be shown for
$(\mu_0,D_0)$ at the level of asymptotic exponents.

\subsubsection{Primitives and the hazard}\label{tc:subsub:primitives}

\begin{definition}[Two-channel model]\label{tc:def:model}
Fix $\sigma>0$, $\epsilon>0$, $n\in\mathbb N$, $\Gamma>0$ and
$\rho>1$. Define
\begin{align}
\text{driver:}&\quad \mu_t=\mu_0+\sigma W_t,
&&[\mu]=\mathsf T^{-1},\ [\sigma]=\mathsf T^{-3/2},
\label{tc:eq:driver}\\
\text{muffling exponent:}&\quad x_t=x_0+\int_0^t\mu_s\,ds,
&&[x]=1,
\label{tc:eq:x}\\
\text{background:}&\quad B_t=\sum_{i=1}^{n}e^{\epsilon W^{(i)}_t},
&&[B]=\mathsf T^{-1},\ [\epsilon]=\mathsf T^{-1/2},
\label{tc:eq:B}
\end{align}
and let the total hazard be the sum of two channels,
\begin{equation}\label{tc:eq:hazard}
h^{\mathrm{tot}}_t
=\underbrace{B_t\,\varsigma_b(x_t)}_{\text{crest}}
+\underbrace{\Gamma\,e^{-\rho x_t}}_{\text{trough}},
\qquad
\varsigma_b(x):=\frac{1+e^{-b}}{1+e^{x-b}}
=\frac{1+e^{-b}}{2}\Bigl(1-\tanh\tfrac{x-b}{2}\Bigr),
\end{equation}
with \emph{offset} $b>0$. The normalising factor $1+e^{-b}$ is chosen so
that $\varsigma_b(0)=1$ exactly, whence $h^{\mathrm{crest}}|_{x=0}=B_t$:
the void limit is the background itself, for every $b$. Setting $b=0$
recovers the unshifted form $B_t(1-\tanh\frac{x}{2})$
Survival is $\bar F_t=e^{-\Lambda_t}$ with
$\Lambda_t=\int_0^t h^{\mathrm{tot}}_u\,du$, and the departure time $T$
is given by the Cox construction of the driftless case. There is \emph{no}
capacity parameter and \emph{no} leak: $x$ is the undamped integral of
$\mu$, exactly as in the driftless case.
\end{definition}

\begin{remark}[Regimes]\label{tc:rem:regimes}
The three regimes of \eqref{tc:eq:hazard}:
\begin{center}\small
\begin{tabular}{@{}lll@{}}
\toprule
regime & total hazard & reading\\
\midrule
$x\gg b$ & $\approx B_te^{\,b-x}$ & shielding; background passes through
multiplicatively\\
$-k/\rho\ll x\ll b$ & $\approx B_t$ & \emph{void}: the system
contributes nothing\\
$x\ll0$ & $\approx\Gamma e^{\rho|x|}$ & internal damage; $B$ has dropped
out entirely\\
\bottomrule
\end{tabular}
\end{center}
The middle line is exact at $x=0$: the normalisation gives
$h^{\mathrm{crest}}|_{x=0}=B_t$, so the void limit is the background
itself. The offset $b$ places $x=0$ \emph{inside} the void rather than
at its upper edge --- see Remark~\ref{tc:rem:whyoffset}.
\end{remark}

\subsubsection{The four constraints}\label{tc:subsub:constraints}

The model is answerable to four structural requirements. We state them
as the author posed them and record where each is met.

\begin{constraint}[Background reaches the total at positive $x$]
\label{tc:C1}
Met for $\rho>1$, and \emph{only} then: for $x\gg b$ the crest channel
is $B_te^{\,b-x}$, multiplicative in $B$, and by
Proposition~\ref{tc:prop:dropout}(ii) it dominates the trough at all
large positive $x$ precisely when $\rho>1$. This --- not
Constraint~\ref{tc:C2} --- is what the brittleness condition buys.
\end{constraint}

\begin{constraint}[Background drops out at negative $x$]\label{tc:C2}
Met by Proposition~\ref{tc:prop:dropout}(i), for \emph{every} $\rho>0$.
Since $\Gamma$ is a primitive constant the trough contains no $B$
whatsoever, so at $x\ll0$ the background disappears from the total
entirely, not merely in its fluctuations.
\end{constraint}

\begin{constraint}[the driftless case undisturbed]\label{tc:C3}
Met as a theorem, \S\ref{tc:sub:asymptotics}: both walls are
$o(t^{3/2})$ and hence gauge.
\end{constraint}

\begin{constraint}[No smuggled primitives]\label{tc:C4}
Met by Theorem~\ref{tc:thm:gain}: with $B$ constant the total hazard is
$h^{\mathrm{tot}}(0)e^{-\Psi(x)}$ for a single soft-kink gain $\Psi$, so
the model is the driftless case's with the credit passed through a kink.
\end{constraint}

\begin{proposition}[Which channel dominates, and where; \textbf{P}]
\label{tc:prop:dropout}
\emph{(i) Negative $x$.} For every $\rho>0$,
$\Gamma e^{-\rho x}/\bigl(B\varsigma_b(x)\bigr)\to\infty$ as
$x\to-\infty$, since $\varsigma_b$ is bounded while $e^{\rho|x|}$ is
not. Hence $h^{\mathrm{tot}}\sim\Gamma e^{\rho|x|}$, free of $B$:
Constraint~\ref{tc:C2} needs no condition on $\rho$ beyond positivity.

\emph{(ii) Positive $x$.} As $x\to+\infty$ the crest is
$\sim Be^{\,b-x}$ and the trough is $\Gamma e^{-\rho x}$, so the crest
dominates for all large $x$ if and only if $\rho>1$. For $\rho<1$ the
two cross at
\begin{equation}\label{tc:eq:crossover}
x_{\times}=\frac{k+b}{1-\rho},\qquad k:=\log\frac{B}{\Gamma},
\end{equation}
beyond which the trough dominates and the background drops out
\emph{even for survivors}.
\end{proposition}

\begin{proof}
(i) $\varsigma_b\le1+e^{-b}$ by Proposition~\ref{tc:prop:bounded}, while
$\Gamma e^{\rho|x|}\to\infty$ for any $\rho>0$. (ii) Compare exponents
$b-x$ and $-k-\rho x$; they are equal at \eqref{tc:eq:crossover}, and
for $\rho>1$ the crest exponent is the larger for all $x$ exceeding a
finite threshold.
\end{proof}

\begin{remark}[What $\rho=1$ is, and is not]\label{tc:rem:rhorole}
Three separate claims must be kept apart, and the first two are easily
run together.
\emph{(a)} $\rho=1$ is \emph{not} the threshold for the background to
drop out at negative $x$; that happens for every $\rho>0$, by
(i).
\emph{(b)} $\rho=1$ \emph{is} the threshold at which the background
ceases to reach the survivor's hazard: for $\rho<1$ internal damage
overtakes the external channel beyond $x_{\times}$ and the system is
effectively closed to the external world, violating
Constraint~\ref{tc:C1}. This is a sharp threshold in a structural
parameter, and it concerns \emph{openness}, not survival.
\emph{(c)} $\rho$ does \emph{not} affect the survival exponent. For any
$\rho>0$ the dominant channel is an exponential in $x$, so the killing
wall is a level of $x$ and the persistence problem is the driftless case's up
to a rescaling of the credit by a constant. The exponent is $\tfrac14$
throughout, which strengthens Constraint~\ref{tc:C3} rather than
qualifying it.
\end{remark}

\subsubsection{Boundedness and the causal asymmetry}
\label{tc:subsub:asymmetry}

\begin{proposition}[The crest channel is bounded; \textbf{P}]
\label{tc:prop:bounded}
$0<\varsigma_b(x)<1+e^{-b}$ for all $x\in\R$, the supremum being
approached only as $x\to-\infty$. Hence the external channel can
\emph{reduce} the background hazard without bound, but can
\emph{increase} it by at most the factor $1+e^{-b}$, which tends to $1$
as $b$ grows.
\end{proposition}

\begin{proof}
$(1+e^{-b})/(1+e^{x-b})$ is decreasing in $x$ with limits $1+e^{-b}$ and
$0$.
\end{proof}

\begin{remark}[Why this matters]\label{tc:rem:asymmetry}
The slogan is the author's: \emph{strength can shield from the
background but never amplify it}, and it is now enforced structurally
rather than assumed. The earlier multiplicative form $B_te^{-x}$
violated it --- at $x<0$ it multiplied the background without bound,
which asserts that a failure of the system's own muffling magnifies the
external world, a causal claim the ontology does not license.
Amplification now arises only through the trough channel, which is
internal damage: a distinct mechanism, correctly separated from the
external one. Proposition~\ref{tc:prop:dropout} is the formal expression
of that separation. Note how much the offset strengthens the slogan: at
$b=0$ the background could still be doubled, whereas for $b\gtrsim6$ the
permitted amplification is under a quarter of one percent. Shielding is
unbounded; amplification is, for practical purposes, forbidden.
\end{remark}

% ---------------------------------------------------------------------
\subsubsection{The construction as an iterated operator}
\label{gr1:subsub:operator}

The placement of the noise --- on $\mu$, the negative logarithmic
derivative of the hazard, rather than on the hazard or its logarithm ---
is the model's single most consequential choice, and this subsection
states the sense in which it is not a choice at all.

\begin{definition}[The hazard operator]\label{gr1:def:operator}
For a positive function $f$ of time write
\begin{equation}\label{gr1:eq:operator}
\mathcal D[f]\;:=\;-\frac{d}{dt}\log f .
\end{equation}
\end{definition}

\noindent Applied to the survival function it returns the hazard, and
applied to the hazard it returns this model's driver:
\begin{equation}\label{gr1:eq:ladder}
\mathcal D[\bar F]=h,\qquad \mathcal D[h]=\mu .
\end{equation}
The driver is therefore the \emph{second iterate of the operation that
defines the hazard in the first place}, and the model's noise sits at
level two of a recursion whose level one is the hazard itself. This is a
different statement from ``we adopt a smoothness prior on the log-hazard
and choose the integrated-Wiener member of the family,'' which is how
the same object appears from the nonparametric-Bayes side, and it is a
stronger one: no member of the family is being selected, because the
ladder \eqref{gr1:eq:ladder} has rungs and the model occupies one of
them.\footnote{%
\emph{Provenance, with its stages kept apart.} The operator
\eqref{gr1:eq:operator} is the author's, derived in May 2021 by a route
worth recording because it is the physical one rather than the textbook
one: examine $\frac{d}{dt}(1-\bar F)$ at $t=0$, read the resulting
quantity as an instantaneous \emph{felt danger} propagating through the
system, rewrite the survival function with that quantity as a parameter,
replace the product $ht$ in the exponent by the time integral
$\int_0^t h$, and solve for $h$ --- which returns
\eqref{gr1:eq:operator}. The hazard thus arrives as a rate one can point
at, not as a definition one is handed, which is the same difference that
\S\ref{gr1:subsub:operator} turns on. The \emph{second} application is
later and its stages should not be run together: the author's record
places the present stochastic use of $\mu$ on the hazard in autumn 2025,
with an earlier and as yet unlocated use of a doubled operator on the
survival function, in connection with the Golden Point rather than with
hazard rates, somewhere in 2021--2023 and under a name of its own. That
earlier use is author-attested and not verified here; if the journal
entries are recovered, this footnote should be revised to cite them.
The logistic deformation of \S\ref{tc:subsub:asymmetry} is later still,
arising during the preparation of the present work from the causal
constraint recorded there.}

\begin{remark}[Where the ladder stops]\label{gr1:rem:ladderstops}
The rungs are not equally hospitable, and the model sits on the last one
that is analytically open. Noise at level one --- $\log h$ Brownian ---
makes the credit Brownian and the persistence exponent $\tfrac12$. Noise
at level two --- $\mu$ Brownian, the present model --- makes the credit
integrated Brownian and the exponent $\tfrac14$
(Theorem~\ref{sh:thm:persist}). At level three the credit is a twice
integrated Brownian motion, and no closed form for the exponent is
known: existence and monotonicity are available for fractionally
integrated processes \cite{AD13,AS15}, and the persistence of iterated
partial sums has been studied, but the value is open. Level one is close
to a definition and level three is close to an open problem; level two
is the rung on which the object is both nontrivial and solvable, and it
is the rung the hazard construction reaches by being applied twice.
\end{remark}

\begin{remark}[The two channels, in operator terms, and why they differ]
\label{gr1:rem:channelsoperator}
Applying $\mathcal D$ to the two channels of \eqref{tc:eq:hazard}
separately, at fixed background, gives
\begin{equation}\label{gr1:eq:channelops}
\mathcal D\bigl[h^{\mathrm{trough}}\bigr]=\rho\,\mu,
\qquad
\mathcal D\bigl[h^{\mathrm{crest}}\bigr]=\varsigma(x-b)\,\mu,
\qquad
\varsigma(u)=\frac{e^{u}}{1+e^{u}} .
\end{equation}
The internal channel therefore preserves the operator relation
\emph{exactly}, with constant gain $\rho$; the external channel
preserves it only up to a logistic weight running from $0$ to $1$.

The asymmetry is the causal constraint of
\S\ref{tc:subsub:asymmetry} in its sharpest form, and it is worth
stating as such. An exact relation $h=h_0e^{-x}$ is unbounded above as
$x\to-\infty$: it says that a failure of the system's own muffling
multiplies the background hazard without limit, which is a causal claim
the ontology declines --- a biosphere may fail to shield itself from the
external world, but it cannot thereby make that world more dangerous
than it is. No such constraint applies inward. \emph{One may harm
oneself without bound; one may not amplify the world.} Hence the trade
recorded in \eqref{gr1:eq:channelops}: the internal channel keeps the
exact relation because nothing forbids it, and the external channel must
be deformed because something does. The logistic is the minimal
deformation purchasing boundedness, and the offset $b$ is how much of it
is bought.

Three consequences are already in the text and are named here rather
than derived again. The technoanatomy $\rho$ of
Definition~\ref{gr1:def:anatomy} \emph{is} the internal channel's
operator gain, $\mathcal D[h^{\mathrm{int}}]/\mu$, which is a third and
perhaps the cleanest answer to why that exponent deserves the name: it
is the rate at which effort converts into change in self-inflicted
hazard, protective when positive and wounding when negative. The void
of \S\ref{gr1:sub:anatomy} is the region where $\varsigma(x-b)\approx0$,
so that $\mathcal D[h^{\mathrm{crest}}]\approx0$ even for $\mu\neq0$ ---
the operator is applied and nothing comes out, which is what
Proposition~\ref{tc:prop:mingain} measures from the other side. And the
concealed gain $\Psi'$ of Corollary~\ref{tc:cor:concealed} is exactly
the hazard-weighted blend of the two relations in
\eqref{gr1:eq:channelops}, which is why its range is $(0,\rho)$ and why
it tends to $1$ under shielding and to $\rho$ under damage.
\end{remark}


\subsubsection{Reduction to a persistence problem}\label{sh:subsub:reduction}

\begin{proposition}[Feynman--Kac]\label{sh:prop:FK}
Let $u(t,x,\mu):=\PP_{x,\mu}(T>t)=\E_{x,\mu}\bigl[e^{-\Lambda_t}\bigr]$.
Then $u$ is the (classical, by Lemma~\ref{sh:lem:hypo}) solution of
\begin{equation}\label{sh:eq:FK}
\partial_tu=\frac{\sigma^2}{2}\partial_\mu^2u-\mu\,\partial_xu-e^{x}u,
\qquad u(0,\cdot,\cdot)\equiv1 .
\end{equation}
\end{proposition}

\begin{proof}
It\^o's formula applied to $s\mapsto
u(t-s,x_s,\mu_s)\exp(-\int_0^se^{x_r}dr)$ exhibits the right side as a
martingale if and only if \eqref{sh:eq:FK} holds; regularity is supplied by
hypoellipticity of the killed generator.
\end{proof}

The killing rate $e^x$ vanishes rapidly for $x\ll0$ and diverges for
$x\gg0$: at exponent level it is an absorbing wall at $x=0$, of thickness
$O(1)$ in $x$ --- a factor $e$ in the hazard --- which is negligible against
excursions of size $t^{3/2}$. We elevate this to a labeled hypothesis, since
it is the single non-rigorous reduction on which the tail asymptotics of
\S\ref{sh:subsub:tail} rest.

\begin{ansatz}[Soft-wall reduction]\label{sh:ansatz}
Let $\tau:=\inf\{t\ge0:D_t\le0\}$ denote the hard-wall passage time of the
Kolmogorov diffusion. Then for every admissible initial datum there exist
constants $0<c\le C<\infty$ and $t_0<\infty$ with
\[
c\,\PP(\tau>t)\ \le\ \PP(T>t)\ \le\ C\,\PP(\tau>t),\qquad t\ge t_0,
\]
after an $O(1)$ adjustment of the initial credit. Equivalently,
$\log\PP(T>t)=\log\PP(\tau>t)+O(1)$.
\end{ansatz}

\begin{remark}[Status of the Ansatz]\label{sh:rem:ansatzstatus}
One half is rigorous modulo a cited input. \emph{Lower bound:} on the event
$\{D_u\ge f(u)\ \forall u\le t\}$ with $f(u)=2\log\frac{2+u}{2}$ one has
$\Lambda_t\le h_0e^{-D_0}\int_0^\infty\bigl(\tfrac{2+u}{2}\bigr)^{-2}du<\infty$,
whence $\PP(T>t)\ge e^{-C}\,\PP(D\ge f\text{ on }[0,t])$; that a logarithmic
moving boundary does not change the persistence exponent of an integrated
process is the standard moving-boundary phenomenon (see the survey
\cite{AS15} and references therein; cf.\ also the robustness expressed
by the universality results \cite{AD13,DW15}). \emph{Upper bound:} from
$\PP(T>t)\le e^{-K}+\PP(\Lambda_t\le K)$ and
$\Lambda_t\ge h_0e^{-D_0}\Leb\{u\le t:D_u\le0\}$, the event
$\{\Lambda_t\le K\}$ forces $D$ to remain positive off a set of bounded
Lebesgue measure. What is missing is the (plausible, unproved here)
statement that \emph{persistence with a bounded occupation tolerance has the
same exponent as strict persistence}. Everything tagged \textbf{[A]} below
depends on exactly this gap and nothing else; \S\ref{sh:sub:numerics}
describes its direct numerical test.
\end{remark}

\subsubsection{The transfer, stated}\label{tc:subsub:transferstated}

\begin{theorem}[Transfer of the driftless case; \textbf{A}]
\label{tc:thm:transfer}
In the two-channel model, for every admissible parameter set,
\begin{equation}\label{tc:eq:transfer}
\PP(T>t)\;\asymp\;\underbrace{\hm\bigl(\mu_{0}-\tfrac{\epsilon^{2}}{2},
\;D_{0}-\log n-b\bigr)}_{\text{shifted initial data}}\;\cdot\;
\underbrace{e^{-\Theta\left(n(b/\sigma)^{2/3}\right)}}_{\text{room
passage}}\;\cdot\;t^{-1/4},
\end{equation}
with $\hm$ the persistence-harmonic function of the driftless case. In
consequence: $\E[T]=\infty$ while $T<\infty$ a.s.; residual life is
asymptotically Pareto$(\tfrac14)$; quantiles grow as the fourth power of
rarity; the effective population hazard is $1/4t$; and the survivorship
interpolation, its pure Doob limit, its $s/4t$ linear response and its
critical mode all hold verbatim, the limit being the same $Q$-process
computed at the shifted data.
\end{theorem}

\begin{proof}
Above the room the hazard is the driftless case's up to bounded constants
(Lemma~\ref{tc:lem:aboveroom}); the wall is linear and absorbed exactly
into the shifted data (Theorem~\ref{tc:thm:absorb}, with the offset $b$
contributing a further constant); the band collapses
(Theorem~\ref{tc:thm:collapse}); the room contributes a bounded additive
term to $\Lambda$, hence a multiplicative prefactor
(Proposition~\ref{tc:prop:roomcost}). The listed consequences are those
of the driftless case, transported through
\eqref{tc:eq:transfer}.
\end{proof}

\begin{remark}[How close, in one sentence]\label{tc:rem:howclose}
Exactly as close as it is possible to be: every difference between the
two models is a shift of initial data or a bounded multiplicative
constant, and we proved both to be gauge. The two-channel model
buys its ontology --- a positive fluctuating background, a bounded
external channel obeying the causal asymmetry, an unbounded internal
channel, an ordered pair of onsets and a waiting room --- at a cost of
precisely zero asymptotic exponents, and it repays part of the price by
shrinking the one gap the earlier section had to assume.
\end{remark}

% ---------------------------------------------------------------------
\subsection{Coordinates: gauge, strength, anatomy}\label{gr1:sub:coords}\label{tc:sub:polar}

\subsubsection{Dimensional structure, self-similarity, and the gauge orbit}
\label{sh:subsub:dimensional}

\begin{definition}[Units and intrinsic scales]\label{sh:def:units}
Assigning $[t]=\mathsf T$ gives $[\mu]=\mathsf T^{-1}$,
$[\sigma]=\mathsf T^{-3/2}$, $[D]=1$, $[h]=\mathsf T^{-1}$. The unique
timescale constructible from $\sigma$ alone is $\sigma^{-2/3}$; the
time-equivalents of the initial data are
\begin{equation}\label{sh:eq:timescales}
t_\mu:=\frac{\mu_0^2}{\sigma^2},\qquad
t_D:=\Bigl(\frac{D_0}{\sigma}\Bigr)^{2/3}\quad(D_0>0);
\end{equation}
and the unique dimensionless state coordinate on the quadrant
$\{\mu>0,D>0\}$ is the \emph{horizon coordinate}
\begin{equation}\label{sh:eq:kappa}
\kappa:=\frac{\sigma^2D}{\mu^3},\qquad\text{satisfying}\qquad
\frac{t_D}{t_\mu}=\kappa_0^{2/3}.
\end{equation}
Hatted variables are $\hat\mu=\mu\sigma^{-2/3}$, $\hat D=D$,
$\hat h=h\sigma^{-2/3}$, $\hat t=t\sigma^{2/3}$.
\end{definition}

\begin{lemma}[Second-order structure]\label{sh:lem:secondorder}
$\Var(I_t)=t^3/3$, $\Cov(W_t,I_t)=t^2/2$, and for $u\ge1$,
$\E[I_1I_u]=\tfrac u2-\tfrac16$.
\end{lemma}

\begin{proof}
By Fubini, $\E[I_aI_b]=\int_0^a\!\int_0^b(\alpha\wedge\beta)\,d\beta\,d\alpha$.
For $a=b=t$ this is $2\int_0^t\!\int_0^\beta\alpha\,d\alpha\,d\beta=t^3/3$.
Next, $\E[W_tI_t]=\int_0^t(\alpha\wedge t)\,d\alpha=t^2/2$. Finally, for
$u\ge1$, $\int_0^u(\alpha\wedge\beta)\,d\beta
=\alpha u-\alpha^2/2$ for $\alpha\le1\le u$, and integrating over
$\alpha\in[0,1]$ gives $u/2-1/6$.
\end{proof}

\begin{proposition}[Exact self-similarity]\label{sh:prop:selfsim}
Fix $\sigma$ and $c>0$. If $(\mu,D)$ starts from $(\mu_0,D_0)$, then
\[
\bigl(\mu_{c^2t},\,D_{c^2t}\bigr)_{t\ge0}\ \eqd\
\bigl(c\,\mu'_t,\,c^3D'_t\bigr)_{t\ge0},
\]
where $(\mu',D')$ starts from $(\mu_0/c,\,D_0/c^3)$. Equivalently,
$(I_{ct})_{t\ge0}\eqd(c^{3/2}I_t)_{t\ge0}$.
\end{proposition}

\begin{proof}
Brownian scaling $(W_{c^2t})_t\eqd(cW_t)_t$, integrated once in time.
\end{proof}

\begin{proposition}[The gauge orbit]\label{sh:prop:gauge}
The initial data $(h_0,\mu_0,D_0)$ affect asymptotics only through
prefactors and timescales, never through exponents. Precisely:
\emph{(i)} $h_0$ is exactly gauge in the sense of
Lemma~\ref{sh:lem:hgauge};
\emph{(ii)} the deterministic contribution
$x_0-\mu_0t$ to $x_t$ is affine in $t$, hence
$(x_0-\mu_0t)/t^{3/2}\to0$, while $\sigma I_t/t^{3/2}\eqd\sigma I_1$ is
nondegenerate for every $t$: the initial data live in the subspace that the
$t^{3/2}$ fluctuations cannot resolve;
\emph{(iii)} under the scaling of Proposition~\ref{sh:prop:selfsim} the
initial data transform as $(\mu_0,D_0)\mapsto(c\mu_0,c^3D_0)$, so any
scale-invariant asymptotic quantity is constant on the two-parameter orbit
$\{(c\mu_0,c^3D_0):c>0\}$ and, when finite and positive, is a function of
$\kappa_0$ alone.
\end{proposition}

\begin{proof}
(i) is Lemma~\ref{sh:lem:hgauge}; (ii) is immediate from
Proposition~\ref{sh:prop:selfsim} and Lemma~\ref{sh:lem:secondorder};
(iii) holds because the orbits of the scaling action on the open quadrant
are exactly the level curves of $\kappa$.
\end{proof}

The sharp form of the slogan --- that even the \emph{prefactors} of survival
asymptotics depend on the data only through a single equivalent timescale ---
is Proposition~\ref{sh:prop:degrees} below.

% ---------------------------------------------------------------------
\subsubsection{The sign of the credit, and an exact half}
\label{gr1:subsub:sign}

One fact about the driftless model is used repeatedly below and is worth
isolating, because it is the engine of a proof and not merely a
description.

\begin{theorem}[Adverse configuration occupies exactly half of
logarithmic time; \textbf{P}]\label{tc:thm:half}
For the driftless driver, almost surely
\begin{equation}\label{tc:eq:half}
\lim_{S\to\infty}\frac1S\,\Leb\bigl\{s\in[0,S]:x_{e^{s}}<0\bigr\}
=\frac12 ,
\end{equation}
so the set $\{t:x_t<0\}$ --- the times at which the system's own
muffling sits \emph{above} rather than below its background hazard ---
has logarithmic density exactly $\tfrac12$, for every parameter value
and every initial condition.
\end{theorem}

\begin{proof}
$Z_s:=x_{e^{s}}/(\sigma e^{3s/2})$ is stationary, centred Gaussian and
ergodic (the Lamperti argument of Theorem~\ref{sh:thm:noimm}), so
Birkhoff gives density $\PP(Z_0<0)=\tfrac12$ by symmetry.
\end{proof}

\begin{remark}[What the exact half is for, and what it is not]
\label{gr1:rem:halfuse}
Two uses and one disclaimer.

\emph{Use.} It is the engine of Theorem~\ref{tc:thm:noviab2}: where the
driftless argument needed an unbounded integrand, the two-channel model
needs only that the system spends a positive fraction of logarithmic
time at hazard above background --- and the fraction is exactly one
half, for every parameter value. Boundedness of the hazard does not buy
immortality; it changes which argument rules it out.

\emph{Sharpening.} The no-viability proof of
\S\ref{gr1:sub:verdict} used only the qualitative input
$\PP(Z_0<-\varepsilon)>0$; here the constant is exact.

\emph{Disclaimer, and it is the important one.} This is a statement
about a \emph{state coordinate}, not about technoanatomy. Every
biosphere, however well constituted, spends half of logarithmic time
with $x<0$; the sign of the credit is a condition the system passes
through and re-enters, not a property it has. Technoanatomy in the sense
of Definition~\ref{gr1:def:anatomy} is the exponent $\rho$, which governs
what capacity \emph{does} when acquired and does not fluctuate at all.
Nothing about the frequency with which a biosphere is adversely
positioned bears on whether its capacity is adversely directed, and the
theorem should not be read as saying otherwise.
\end{remark}


% ---------------------------------------------------------------------
\subsection{The verdict: certain departure, infinite expectation}
\label{gr1:sub:verdict}\label{sh:sub:verdict}

\begin{remark}[What $T$ is the time of]\label{gr1:rem:whatTis}
Throughout this section $T$ is the time of the \emph{Departure Point}
(Definition~\ref{gc:def:departure}) --- the moment past which no
continuation advances --- not of biological death, and
$h^{\mathrm{tot}}$ is that event's hazard. Death is one mechanism of
foreclosure and irreversible loss of the capacity to advance is
another; the analysis does not distinguish them, because for
advancement they are the same event. It does, however, distinguish both
from a mere lull: a biosphere that has stopped progressing but could
still resume has not reached $T$. Where the classical results imported below speak of survival
and first passage, the vocabulary is theirs and the referent here is
departure. Read that way, $\PP(\Lambda_\infty<\infty)$ is the
probability of \emph{advancing}, not of not dying --- which is what the
Golden Point was always a claim about.
\end{remark}

\subsubsection{Impossibility of immortality}\label{sh:subsub:noimm}

\begin{theorem}[No viability; \textbf{P}]\label{sh:thm:noimm}
For every $(h_0,\mu_0,D_0,\sigma)$ in the driftless model
\eqref{sh:eq:model}, $\Lambda_\infty=\infty$ almost surely. Consequently
$S_\infty=0$ a.s., $\PP(\Imm)=0$, $\PP(T=\infty)=0$, and $T<\infty$ a.s.
\end{theorem}

\begin{proof}
Since $\Lambda_\infty=h_0e^{-D_0}\int_0^\infty e^{-\mu_0u-\sigma I_u}du$
and the prefactor is a fixed positive constant, it suffices to prove that
the integral diverges a.s.

\emph{Step 1 (Lamperti transform).} Define $Z_s:=e^{-3s/2}I_{e^s}$,
$s\ge0$. As a linear functional of $W$, $Z$ is a centered Gaussian process.
By Proposition~\ref{sh:prop:selfsim}, for each $r\ge0$,
$(Z_{s+r})_{s\ge0}=(e^{-3(s+r)/2}I_{e^re^s})_{s\ge0}
\eqd(e^{-3s/2}I_{e^s})_{s\ge0}=(Z_s)_{s\ge0}$: $Z$ is stationary. By
Lemma~\ref{sh:lem:secondorder}, for $s\ge0$,
\begin{equation}\label{sh:eq:lampcov}
\E[Z_0Z_s]=e^{-3s/2}\Bigl(\frac{e^{s}}{2}-\frac16\Bigr)
=\frac12e^{-s/2}-\frac16e^{-3s/2}\ \longrightarrow\ 0 .
\end{equation}

\emph{Step 2 (ergodicity).} A stationary Gaussian process whose covariance
vanishes at infinity is mixing \cite{Mar49}, hence ergodic.

\emph{Step 3 (Birkhoff).} Fix $\varepsilon>0$. Since
$Z_0\sim N(0,\tfrac13)$, $p_\varepsilon:=\PP(Z_0\le-\varepsilon)>0$. By the
continuous-time Birkhoff theorem,
$\tfrac1S\Leb\{s\in[0,S]:Z_s\le-\varepsilon\}\to p_\varepsilon$ a.s., so
$\Leb\{s\ge0:Z_s\le-\varepsilon\}=\infty$ a.s.

\emph{Step 4 (change of variables).} Let
$A:=\{u\ge1:I_u\le-\varepsilon u^{3/2}\}$. Under $u=e^s$ the set in Step~3
maps into $A$ with
$\Leb\{s:Z_s\le-\varepsilon\}=\int_A\frac{du}{u}\le\Leb(A)$,
the inequality because $u\ge1$. Hence $\Leb(A)=\infty$ a.s.

\emph{Step 5 (divergence).} Choose $u_1=u_1(\varepsilon,\mu_0,\sigma)$ so
that $\sigma\varepsilon u^{3/2}-\mu_0u\ge0$ for $u\ge u_1$. Then
\[
\int_0^\infty e^{-\mu_0u-\sigma I_u}\,du
\ \ge\ \int_{A\cap[u_1,\infty)}e^{\sigma\varepsilon u^{3/2}-\mu_0u}\,du
\ \ge\ \Leb\bigl(A\cap[u_1,\infty)\bigr)\ =\ \infty
\qquad\text{a.s.}\qedhere
\]
\end{proof}

\begin{remark}[The asymmetry of exponential response]\label{sh:rem:asym}
The driver is symmetric; the response is not. A favorable excursion
suppresses the integrand $e^{-D}$ multiplicatively toward zero ---
\emph{bounded} benefit for $\Lambda$ --- while an unfavorable excursion of
relative size $\varepsilon$ inflates it without bound. Exponential response
to a symmetric driver is structurally hostile to survival; the mechanism is
the sign-oscillation that makes low-dimensional random walks recurrent,
aggravated by this asymmetry. Identifying exactly what a horizon must break
--- the symmetry of the driver, not the response --- is the burden of
\S\ref{sh:subsub:scale}.
\end{remark}

\subsubsection{No viability, by a different route}\label{tc:subsub:noviab}

the driftless case's no-viability theorem was proved by exhibiting an integrand
that \emph{blows up}: on a set of infinite measure the credit satisfies
$D_u\le-\varepsilon u^{3/2}$, whence $e^{-D_u}\ge
e^{\sigma\varepsilon u^{3/2}}\to\infty$. That argument is unavailable
here whenever the crest channel acts alone, since $\varsigma_b$ is
\emph{bounded}. The conclusion nevertheless survives, by a weaker and
more robust mechanism: a positive floor on an infinite set.

\begin{theorem}[No viability in the two-channel model; \textbf{P}]
\label{tc:thm:noviab2}
For every $\sigma,\epsilon>0$, $n\ge1$, $b\ge0$, $\rho>0$ and every
$\Gamma\ge0$ --- including $\Gamma=0$, where the total hazard is
bounded --- one has $\Lambda_{\infty}=\infty$ almost surely. Hence
$\bar F_{\infty}=0$ a.s.\ and $T<\infty$ a.s.
\end{theorem}

\begin{proof}
Discard the trough, which is nonnegative. On $\{x_u\le0\}$ we have
$\varsigma_b(x_u)\ge\varsigma_b(0)=1$, so
\[
\Lambda_{t}\;\ge\;\int_{0}^{t}B_u\varsigma_b(x_u)\,du
\;\ge\;\Bigl(\inf_{u\le t}B_u\Bigr)\,
\mathrm{Leb}\{u\le t:x_u\le0\}.
\]
By Theorem~\ref{tc:thm:half} the set $\{u:x_u\le0\}$ has logarithmic
density $\tfrac12$; writing $u=e^{s}$,
$\mathrm{Leb}\{u\le t:x_u\le0\}=\int_{0}^{\log t}\one\{Z_s\le0\}e^{s}ds
\to\infty$ a.s., since the integrand is positive on a set of density
$\tfrac12$ and $e^{s}\to\infty$. As $\inf_{u\le t}B_u>0$ a.s.\ for each
$t$ and $B$ does not tend to zero, $\Lambda_{\infty}=\infty$.
\end{proof}


\subsubsection{The survival tail and its consequences}\label{sh:subsub:tail}

\begin{theorem}[Persistence of integrated Brownian motion; \textbf{C}]
\label{sh:thm:persist}
Let $(\mu,D)$ start from $(m,0)$ with $m>0$ and let
$\tau=\inf\{t:D_t\le0\}$. There is a constant $c_1>0$ such that
\begin{equation}\label{sh:eq:persist}
\PP_{(m,0)}(\tau>t)\ \sim\ c_1\Bigl(\frac{m^2}{\sigma^2t}\Bigr)^{1/4},
\qquad t\to\infty .
\end{equation}
The exponent $\tfrac14$ --- half the Brownian value, halved by one
integration --- is the persistence exponent of integrated Brownian
motion, and its provenance is a chain worth recording precisely. McKean
\cite{McK63} computed the first-passage (``winding'') law of the
Kolmogorov diffusion; Goldman \cite{Gol71} derived the passage-rate
asymptotic from McKean's law; Lachal \cite{Lac91} obtained the exact
first-passage distributions; Sinai \cite{Sin92} initiated and settled,
up to constants, the discrete counterpart $n^{-1/4}$ for the integrated
random walk; Isozaki--Watanabe \cite{IW94} supplied the sharp constants;
and Denisov--Wachtel \cite{DW15} proved the exact $n^{-1/4}$ asymptotics
universally for centered integrated random walks with $2+\delta$
moments, via a discrete potential theory (see also the universality
result of Aurzada--Dereich \cite{AD13} for integrated processes).
Starting states in the open quadrant obey the same power with a
prefactor proportional to the harmonic function of
Proposition~\ref{sh:prop:harmonic}; this general-start form is
established in \cite{GJW99}, which recovers \cite{Gol71,IW94} as
special cases. For the surrounding theory see the survey \cite{AS15}
and, on the physics side --- where the model is the \emph{random
acceleration process} and has an extensive literature of its own ---
the review \cite{BMS13}.
\end{theorem}

\begin{remark}[The $m^{1/2}$ law is exact given the $m=1$ asymptotic]
By Proposition~\ref{sh:prop:selfsim} with $c=1/m$ (and $\sigma=1$),
$\PP_{(m,0)}(\tau>t)=\PP_{(1,0)}(\tau>t/m^2)$, so once
\eqref{sh:eq:persist} is known for one $m$ it holds for all, prefactor
scaling as $m^{1/2}$ exactly. The initial rate enters through its
time-equivalent $t_\mu$ and nothing else.
\end{remark}

\begin{claim}[Survival tail; \textbf{A}]\label{sh:claim:tail}
Under Ansatz~\ref{sh:ansatz}, for every initial datum there is
$C=C(\hat h_0,\hat\mu_0,\hat D_0)\in(0,\infty)$ with
\begin{equation}\label{sh:eq:tail}
\PP(T>t)\ \asymp\ C\,t^{-1/4},\qquad t\to\infty .
\end{equation}
\end{claim}

All of the following consequences are elementary integrations of
\eqref{sh:eq:tail}; each inherits the tag \textbf{[A]}. Write
$\alpha=\tfrac14$.

\begin{proposition}[Consequences of the tail; \textbf{A}]
\label{sh:prop:tailconseq}\label{sh:cor:moments}\label{sh:cor:mean}%
\label{sh:cor:trunc}\label{sh:cor:quant}\label{sh:cor:residual}%
\label{sh:cor:lindy}
All of the following are elementary consequences of \eqref{sh:eq:tail}
by regular variation, with $\alpha=\tfrac14$; we state them together
rather than prove them severally, since each is a one-line integration
and none is original.

\begin{center}\small
\begin{tabular}{@{}lll@{}}
\toprule
quantity & value & reading\\
\midrule
$\E[T^{p}]$ & finite iff $p<\tfrac14$ &
even $\E[\sqrt T]=\infty$; heavier-tailed than Cauchy\\
$\E[T]$ & $\infty$, every initial condition &
null recurrence: certain departure, infinite mean\\
$\E[T\wedge\mathcal T]$ & $\asymp\tfrac{4C}{3}\mathcal T^{3/4}$ &
an observed mean is an artefact of the window\\
$t_q$ with $\PP(T>t_q)=q$ & $\asymp(C/q)^{4}$ &
quantiles grow as the fourth power of rarity\\
$\PP(T>\lambda t\mid T>t)$ & $\to\lambda^{-1/4}$ &
residual life Pareto$(\tfrac14)$, mean infinite at every age\\
$r(t)=-\tfrac{d}{dt}\log\PP(T>t)$ & $\approx 1/4t$ &
a Lindy law; $\int^\infty r=\infty$, so departure is certain\\
\bottomrule
\end{tabular}
\end{center}

\noindent Two remarks are worth more than the derivations. The truncated
mean shows that any finite dataset reports a finite mean lifetime
growing as the $3/4$ power of its own duration, so the quantity is not
estimable. And the effective hazard $1/4t$ sits far from the mortality
boundary --- any positive exponent gives certain departure --- but deep
inside the infinite-mean regime, whose boundary is exponent $1$.
\end{proposition}

% ---------------------------------------------------------------------
\subsection{Technoanatomy: the sign of the trough exponent}
\label{gr1:sub:anatomy}

Section~\ref{gr1:sub:object} introduced the trough channel
$\Gamma e^{-\rho x}$ under the standing restriction $\rho>1$, imposed so
that the background would reach the survivor's hazard
(Constraint~\ref{tc:C1}). That restriction has never been examined, and
lifting it turns out to supply the model's missing structural
parameter. We lift it now: $\rho$ ranges over $\R$.

\begin{definition}[Technoanatomy]\label{gr1:def:anatomy}
The exponent $\rho$ is the \emph{technoanatomy}\footnote{%
\emph{Provenance.} The notion, its square--cube framing, and the
exponential-against-linear scaling that
\S\ref{gr:subsub:boundary} later derives as $\rho=-\alpha$ all appear in
the author's journal three years before the present model, and the entry
is reproduced here in full because it states the mechanism more
compactly than the formal development does. Dated 23 July 2023,
1:42~pm, under the heading \emph{Golden Biosphere Analogue to the Square
Cube Law}:
\begin{quote}\small
``Just thought of the craziest idea! You know the quote about how
advanced aliens might attack us like how boys throw stones at frogs? The
implication I get from this is that such an alien biosphere is portrayed
as just a scaled up version of our own so to speak in every aspect of
society, such that we just get a biosphere which is far more advanced
technologically, but interacts about itself in similar ways, such as
they are precisely are to humans as humans are to ants.

So, here's why I think this won't work:

The idea is basically an analogue of the square cube law to
Astrobiology, working in a similar way to how you can't just scale an
ant up to the size of an elephant. There has to be a transformation in
the process instead.

The idea, essentially, is that the capacity for us to destroy ourselves
grows exponentially with technology, whereas the defense mechanisms put
in place with our current society, given the tragedy of the commons and
all, grows only linearly with technology. If we assigned each human
comparatively limitless power, such that the advanced analogue of
`boys' might attack humans with the same mindset as human boys attack
frogs, then such an advanced civilization will implode as soon as it
spawns. Why? Because the linear defense mechanisms of a human society
would be immediately overtaken by the vast exponential issues propagated
by such a technology.

Anyways, so this demonstrates essentially two things: A) we will not
remain constant in societal mindset with time, and B) we will be
essentially forced through the keyhole of the Golden Point mindset long
before we reach the capacity of supposedly treating humans like frogs to
boys.''
\end{quote}
The entry must be read as its author wrote it, which is
\emph{conditionally on survivorship} --- the conditioning having been the
paradigmatic heart of the programme since at least 2022, on the author's
record, and therefore predating this entry as well. Read that way, every
one of its claims is recovered below; read unconditionally, two of them
would appear to fail. The distinction bears directly on how much the
entry anticipates, and read correctly the answer is all of it.

Six things are worth marking, since each became a separate part of the
apparatus. The exponential-against-linear contrast is
$\rho<0$: destruction capacity growing as $e^{\alpha x}$ against
enforcement growing more slowly is exactly
\eqref{gr:eq:rhoextract}, and the entry states the mechanism without the
surface-versus-volume derivation that \S\ref{gr:subsub:boundary}
supplies. The square--cube analogy is the source of the names
\emph{technostrength} and \emph{technoanatomy}, so
Remark~\ref{gr:rem:namesclosure}'s observation that the joke turned out
to be the mechanism is a retrospective reading of this entry rather than
a later discovery. ``Implode as soon as it spawns'' is
Theorem~\ref{gr1:thm:trapped}. And ``Golden Point'' is the earlier name
for what the present work calls the Golden Biosphere.

The entry's two numbered conclusions are the remaining pair, and both
are theorems here once the conditioning is respected. Item~(B) --- that
a biosphere is forced through the keyhole \emph{before} attaining the
capacity in question --- is
Proposition~\ref{gr1:prop:tolerance}: the affordable duration of adverse
anatomy is $\delta\lesssim\Gamma^{-1}e^{-|\rho|x}$, so conditional on
being alive at large technostrength, adverse anatomy is essentially
absent, and a biosphere cannot have arrived there without having already
passed through. The keyhole is not a force acting on biospheres --- 
unconditionally, Corollary~\ref{gr1:cor:recurrentanatomy} allows a
biosphere simply to fail --- but a property of the survivors, which is
what the entry claims. Item~(A) --- that societal arrangement will not
remain constant --- is the anatomy selection measured in
\S\ref{gr:subsub:anatsel}, where the conditioned distribution of $\rho$
shifts away from its prior. Neither is a prediction the model
contradicts; both are among the results it was later built to state.}
of the system: the
degree to which increasing muffling capacity fails to be turned back
against the system that acquired it. Reading the trough channel
$\Gamma e^{-\rho x}$ as internal hazard:
\begin{center}\small
\begin{tabular}{@{}lll@{}}
\toprule
& internal hazard as $x$ grows & reading\\
\midrule
$\rho>0$ & falls, $\propto e^{-\rho x}$ & capacity is patched as it is
built\\
$\rho=0$ & constant at $\Gamma$ & capacity neither patched nor
weaponised\\
$\rho<0$ & rises, $\propto e^{|\rho|x}$ & \emph{the system eats itself}\\
\bottomrule
\end{tabular}
\end{center}
\end{definition}

\begin{remark}[Why this is the right home for the notion]
\label{gr1:rem:anatomyhome}
Two things recommend $\rho$ over the phase coordinate $\Theta$ of
\S\ref{gr1:sub:coords}, with which it is easily conflated.
First, $\Theta$ is a \emph{state coordinate}: it fluctuates, is uniform
under $\PP$, and every system occupies every value of it half of
logarithmic time (Theorem~\ref{tc:thm:half}). Anatomy in the intended
sense is not something a system passes through but something it
\emph{is}. Second, and decisively, the intended notion is a claim about
what \emph{strength does} --- whether acquiring capacity raises or
lowers one's own hazard --- and that is a statement about a derivative,
$\partial_x\log h^{\mathrm{int}} = -\rho$, not about a position. The
phase coordinate is retained under its own name; technoanatomy is
$\rho$.
\end{remark}

\begin{example}[The asteroid, and why the channels must be two]
\label{gr1:ex:asteroid}
The reason internal and external hazard cannot be folded into one term
is best seen in a case where the same capability serves both. Consider a
biosphere that becomes progressively better at deflecting asteroids.
Every increment of that capability lowers the external hazard: the
crest channel falls, and by \eqref{tc:eq:hazard} it falls
multiplicatively in the background, exactly as it should.

But an apparatus that can redirect a body large enough to end a
biosphere is, by construction, an apparatus that can \emph{aim} such a
body. The same increment therefore also creates a capability that did
not previously exist and that is directed at the biosphere itself. What
happens to the total hazard depends on nothing about asteroids and
everything about whether the biosphere is so arranged that the new
capability can be turned inward --- which is precisely $\rho$.

With $\rho>0$ the deflection apparatus is built into an arrangement that
patches its own misuse faster than it creates the possibility of it, and
the net effect of the capability is protective. With $\rho<0$ it is not,
and the biosphere has purchased a reduction in external hazard at the
price of a larger increase in internal hazard --- the trade being worse
the more capable the apparatus. The example is not incidental to the
model: it is why \eqref{tc:eq:hazard} carries two channels rather than
one, and why the second is unbounded above while the first is not
(Proposition~\ref{tc:prop:bounded}). A single-channel model cannot
represent a capability that cuts both ways, and every capability worth
having does.
\end{example}

\begin{proposition}[The anatomy trichotomy; \textbf{P}]
\label{gr1:prop:trichotomy}
Write $h(x)=B\varsigma_b(x)+\Gamma e^{-\rho x}$ with $B,\Gamma>0$.
\begin{enumerate}[label=\emph{(\roman*)},itemsep=1pt]
\item If $\rho>0$, $h$ is strictly decreasing with $h(x)\to0$ as
$x\to+\infty$: capacity is unboundedly protective.
\item If $\rho=0$, $h$ is strictly decreasing with
$h(x)\to\Gamma>0$: capacity is protective but against a floor.
\item If $\rho<0$, $h\to B(1+e^{-b})$ as $x\to-\infty$ and
$h\to\infty$ as $x\to+\infty$, with a unique interior minimum. Writing
$A:=Be^{b}$ and $|\rho|=-\rho$, the minimum sits at
\begin{equation}\label{gr1:eq:xstar}
x_\ast=\frac{1}{1+|\rho|}\log\frac{A}{|\rho|\Gamma},
\qquad
h_\ast=A^{\frac{|\rho|}{1+|\rho|}}\,\Gamma^{\frac{1}{1+|\rho|}}
\Bigl(|\rho|^{\frac{1}{1+|\rho|}}+|\rho|^{-\frac{|\rho|}{1+|\rho|}}\Bigr),
\end{equation}
in the regime $x_\ast\gtrsim b$ where the crest is exponential.
\end{enumerate}
\end{proposition}

\begin{proof}
$\varsigma_b$ is strictly decreasing with limits $1+e^{-b}$ and $0$;
$e^{-\rho x}$ is decreasing for $\rho>0$, constant for $\rho=0$, and
increasing for $\rho<0$. (i) and (ii) are immediate. For (iii),
approximating the crest by $Ae^{-x}$ above the offset, $h'=0$ gives
$Ae^{-x}=|\rho|\Gamma e^{|\rho|x}$, whence $x_\ast$; substituting back
and collecting the common factor
$A^{|\rho|/(1+|\rho|)}\Gamma^{1/(1+|\rho|)}$ gives $h_\ast$. Uniqueness
follows since $h''>0$ wherever $h'=0$. (Verified numerically to three
significant figures at $b=3$, $k=6$, $\rho=-0.2$: predicted
$x_\ast=8.84$, $h_\ast=1.744\times10^{-2}$; measured $8.88$ and
$1.757\times10^{-2}$.)
\end{proof}

\begin{theorem}[Non-positive anatomy is fatal, and fast; \textbf{P}]
\label{gr1:thm:trapped}
If $\rho\le0$ then $h_\ast:=\inf_x h(x)>0$, and therefore
\begin{equation}\label{gr1:eq:trappedtail}
\PP(T>t)\;\le\;e^{-h_\ast t}\qquad\text{for all }t,
\end{equation}
so that $\E[T]<\infty$ and \emph{no power-law structure survives}:
neither the $t^{-1/4}$ tail, nor infinite life expectancy, nor Pareto
residual life, nor the Lindy reading of
\S\ref{gr1:sub:verdict}. The bound holds pathwise and uses no property
of the driver whatever.
\end{theorem}

\begin{proof}
For $\rho=0$, $h\ge\Gamma>0$. For $\rho<0$, $h_\ast$ is the interior
minimum of Proposition~\ref{gr1:prop:trichotomy}(iii), strictly positive
since both channels are. Then $\Lambda_t\ge h_\ast t$ pathwise and
$\PP(T>t)=\E[e^{-\Lambda_t}]\le e^{-h_\ast t}$.
\end{proof}

\begin{remark}[What the trapped regime looks like]
\label{gr1:rem:trapped}
The reading of $\rho<0$ is worth stating plainly, because it inverts the
intuition the rest of the model builds. With negative technoanatomy the
\emph{safest available configuration is not the strongest but the
weakest}: $h\to B(1+e^{-b})$ as $x\to-\infty$, so a system with bad
anatomy is safest at background hazard with no capacity at all, and
every increment of capacity beyond $x_\ast$ is net harmful. The system
is not merely at risk; it is trapped between an external hazard it
cannot shed without capacity and an internal hazard that capacity
creates. Equation~\eqref{gr1:eq:xstar} prices the trap: the best
achievable hazard $h_\ast$ is a geometric interpolation between the
external scale $A$ and the internal scale $\Gamma$, weighted by anatomy.
\end{remark}


% ---------------------------------------------------------------------
\subsection{Transience, and the two-condition theorem}
\label{gr1:sub:twocond}

Anatomy governs whether capacity protects. It does not by itself deliver
survival, because protection must also be \emph{acquired}, and whether
it is depends on the driver. This subsection assembles the two into the
model's central structural result.

\subsubsection{The feedback condition}\label{gr1:subsub:feedback}

By Proposition~\ref{gr1:prop:trichotomy}, $\rho>0$ makes $h\to0$ as
$x\to+\infty$; but $\Lambda_\infty<\infty$ additionally requires
$x_t\to+\infty$ fast enough for $\int^\infty e^{-\min(1,\rho)x_t}dt$ to
converge, and $x_t=\int_0^t\mu_s\,ds$ grows without bound only if the
driver does not return. That is a transience condition on $\mu$, and
the boundary classification of one-dimensional diffusions settles it
exactly.

\begin{remark}[What the theorem computes]\label{gr1:rem:twocondreading}
By Remark~\ref{gr1:rem:whatTis}, $\PP(\Lambda_\infty<\infty)$ is the
probability that the biosphere \emph{advances} --- that some
continuation available to it reaches the Golden Point. The
two-condition theorem below therefore does not say what it takes to
avoid dying; it says what it takes for a route to exist at all, and the
two factors are the two things a route requires. That reading is what
makes the trapped regime intelligible: with $\rho<0$ the probability is
identically zero, so a biosphere there has departed in the strict
sense --- not because it has stopped, but because nothing it can do
advances.
\end{remark}

\begin{remark}[Transience is a feedback loop]\label{gr1:rem:feedback}
The canonical transient driver is
\begin{equation}\label{gr1:eq:uou}
d\mu_t=\nu\,\mu_t\,dt+\sigma\,dW_t,\qquad \nu>0,
\end{equation}
whose drift is proportional to the state: \emph{the rate of improvement
grows with the improvement already achieved}. This is the
self-reinforcing dynamic that the technological-singularity literature
posits, written as a diffusion, and it is not an addition to the model
but the member of the driver family whose scale function is bounded.
The connection made in \S\ref{gr1:subsub:twocondthm} is therefore not an
analogy: the condition under which the model admits immortality
\emph{is} the condition under which the driver exhibits feedback.
\end{remark}

\subsubsection{The theorem}\label{gr1:subsub:twocondthm}

\begin{theorem}[Two conditions for viability; \textbf{P}]
\label{gr1:thm:twocond}
Let the driver be a regular diffusion on $\R$ with scale function $s$,
and let $\rho\in\R$ be the technoanatomy. Then
\begin{equation}\label{gr1:eq:twocond}
\PP\bigl(\Lambda_\infty<\infty\bigr)
\;=\;
\underbrace{\one\{\rho>0\}}_{\text{anatomy}}
\;\cdot\;
\underbrace{\frac{s(\mu_0)-s(-\infty)}{s(+\infty)-s(-\infty)}}
_{\text{transience}},
\end{equation}
the second factor read as $0$ when $s(+\infty)=\infty$ and as $1$ when
$s$ is bounded above and $s(-\infty)=-\infty$. Both factors are
necessary: neither good anatomy without feedback, nor feedback without
good anatomy, admits viability.
\end{theorem}

\begin{proof}
On $\{\mu_t\to+\infty\}$: eventually $\mu_t\ge1$, so $x_t\ge x_{t_0}+
(t-t_0)$ and, for $\rho>0$, $h(x_t)\lesssim e^{-\min(1,\rho)(t-t_0)}$,
whence $\Lambda_\infty<\infty$. For $\rho\le0$,
Theorem~\ref{gr1:thm:trapped} gives $\Lambda_t\ge h_\ast t\to\infty$
regardless of the driver, so the first factor is necessary. On
$\{\mu_t\to-\infty\}$, $x_t\to-\infty$ and $h\to B(1+e^{-b})>0$ for
every $\rho$, so $\Lambda_\infty=\infty$; on the recurrent event the
argument of Theorem~\ref{sh:thm:noimm} applies. The hitting
probabilities of $\pm\infty$ are the classical scale-function formula
\cite{KT81,RY99}.
\end{proof}

\begin{corollary}[The solvable case; \textbf{P}]\label{gr1:cor:solvable}
For the feedback driver \eqref{gr1:eq:uou},
$s'(v)=\exp(-\nu v^{2}/\sigma^{2})$ is integrable at both ends, and
\begin{equation}\label{gr1:eq:sigmoid}
\PP\bigl(\Lambda_\infty<\infty\bigr)
=\one\{\rho>0\}\cdot\Phi\!\Bigl(\frac{\mu_0\sqrt{2\nu}}{\sigma}\Bigr).
\end{equation}
The viability probability is thus a \emph{product of two thresholds of
different character}: a sharp one in the technoanatomy, at $\rho=0$,
which is a discontinuity; and a smooth Gaussian sigmoid in the initial
improvement rate, of width $\sigma/\sqrt{2\nu}$, which sharpens to a
step only in the small-noise limit.
\end{corollary}

\begin{remark}[Against the singularity literature]
\label{gr1:rem:againstsing}
Equation \eqref{gr1:eq:twocond} states, in the model's own terms, an
objection to the singularity and techno-optimist literatures which is
sharper than the usual one. Those literatures supply
the second factor and are silent on the first: feedback is argued for at
length, and whether capacity is turned back on its holder is not a
parameter they carry. The model's verdict on that omission is not that
it leaves the conclusion unsupported but that it inverts it. With
$\rho<0$, feedback does not merely fail to secure survival; it
\emph{accelerates the failure}, since $\mu_t\to+\infty$ drives
$x_t\to+\infty$ and hence, by
Proposition~\ref{gr1:prop:trichotomy}(iii), $h\to\infty$
superexponentially. Supplying feedback alone is strictly worse than
supplying neither: a system with bad anatomy and no feedback sits near
$h_\ast$, while one with bad anatomy and feedback runs to unbounded
hazard as fast as its improvement compounds.
\end{remark}

\subsubsection{If anatomy itself fluctuates}\label{gr1:subsub:fluct}

Technoanatomy has been treated as a constant, which is an idealisation
and the model's most obvious place to give. Suppose instead
$\rho=\rho_t$ is a process. Two consequences follow at scaling grade,
and together they sharpen rather than blunt the event-horizon reading.

\begin{proposition}[The tolerance for adverse anatomy collapses with
strength; \textbf{S}]\label{gr1:prop:tolerance}
Let anatomy dip to $\rho_t=-|\rho|<0$ for a duration $\delta$ at a time
when technostrength is $x$. The contribution to the cumulative hazard
is $\asymp\delta\,\Gamma e^{|\rho|x}$, so that such an excursion is
survivable only if
\begin{equation}\label{gr1:eq:tolerance}
\delta\;\lesssim\;\Gamma^{-1}e^{-|\rho|x}.
\end{equation}
The affordable duration of adverse anatomy therefore decays
\emph{exponentially in technostrength}.
\end{proposition}

\begin{corollary}[Recurrent anatomy is fatal; \textbf{S}]
\label{gr1:cor:recurrentanatomy}
If $\rho_t$ is recurrent on $\R$ --- if anatomy returns to negative
values infinitely often --- then $\Lambda_\infty=\infty$ almost surely,
\emph{whatever the driver}. Each negative excursion occurring at
technostrength $x_t$ contributes $e^{|\rho|x_t}$, and along any path on
which strength accumulates this diverges superexponentially. Viability
therefore requires $\rho_t$ to be eventually positive almost surely ---
transient upward, or absorbed above some $\rho_c>0$ --- and the
two-condition theorem becomes a \emph{two-transience} condition, with
$\PP(\text{viable})$ a product of two scale functions rather than an
indicator times one.
\end{corollary}

\begin{remark}[This is the event horizon the fixed model could not
state]\label{gr1:rem:horizonproper}
Three readings, and the third is the one to keep.

\emph{The threshold becomes an escape rather than a level.} With $\rho$
fixed, the anatomy condition is a sign test one either passes or fails
from the outset. With $\rho$ fluctuating it is an event --- anatomy
escaping upward and not returning --- which is what an event horizon
is: not a place in the state space but a passage after which the past
is unreachable.

\emph{``Hard to regress past'' becomes quantitative.} Regression is
never impossible; what \eqref{gr1:eq:tolerance} says is that the
affordable \emph{duration} of it collapses exponentially in strength. A
biosphere may be badly configured for a long while when weak and for
almost no time at all when strong.

\emph{An ordering constraint falls out.} Since the cost of anatomical
failure grows with strength while the benefit of capacity does not
grow with anatomy, anatomy must be secured \emph{before} strength
accumulates. A biosphere that acquires capacity first and attends to
its own arrangement afterwards is running the two in the fatal order,
and \eqref{gr1:eq:tolerance} prices exactly how fatal: by the time it
turns to the problem, its margin for error has shrunk by
$e^{|\rho|x}$.
\end{remark}

\noindent The dynamics of $\rho_t$ are not supplied here, and supplying
them is the model's principal open direction: an anatomy that degrades
when capacity outruns the patching of it, $d\rho=(\text{patching})-
\lambda\,\dot x\,dt$ or similar, would make technoanatomy emergent
rather than parametric. Whether the asymptotics of the resulting
two-process system remain tractable is \textbf{[O]}.

\begin{remark}[Where the earlier results now sit]\label{gr1:rem:wherenow}
The trichotomy locates the rest of this section. The driftless model of
\S\S\ref{gr1:sub:verdict}--\ref{gr1:sub:cond} --- certain departure,
infinite life expectancy, the $t^{-1/4}$ tail, the soft horizon and the
whole survivorship interpolation --- is the case $\rho>0$ with a
\emph{recurrent} driver: good anatomy, no feedback. It is the boundary
between the trapped regime, where life expectancy is finite
(Theorem~\ref{gr1:thm:trapped}), and the viable regime, where
immortality has positive probability
(Corollary~\ref{gr1:cor:solvable}). That is the same structural position
the driftless case occupies within the scale-function classification of
Appendix~\ref{app:horizon}, and it is not a coincidence: both
trichotomies are the same statement resolved along different axes ---
one along anatomy, one along the driver.
\end{remark}


% ---------------------------------------------------------------------
\subsubsection{The internal channel as a perturbation}
\label{gr1:subsub:perturb}

The two channels admit a perturbative treatment in the internal
coupling $\Gamma$, and it is worth recording, both because it supplies
the model's one series expansion and because the series behaves in an
instructive way at the anatomy threshold.

\begin{proposition}[First-order internal correction; \textbf{P} to
first order, \textbf{S} for the convergence claim]
\label{gr1:prop:perturb}
Write $\Lambda_t=\Lambda^{\mathrm{crest}}_t+\Gamma\int_0^te^{-\rho
x_s}ds$. Expanding the survival functional in $\Gamma$,
\begin{equation}\label{gr1:eq:perturb}
\PP(T>t)\;=\;\PP_{\mathrm{crest}}(T>t)\left[\,1-\Gamma\,
\E_{Q_{\mathrm{crest}}}\!\Bigl[\int_0^t e^{-\rho x_s}\,ds\Bigr]
+O(\Gamma^{2})\,\right],
\end{equation}
where $Q_{\mathrm{crest}}$ is the crest-only survivorship measure of
\S\ref{gr1:sub:cond}. The leading correction is therefore $\Gamma$ times
the \emph{survivor ensemble's} expected exposure to the internal
channel, not the unconditioned ensemble's --- the conditioning enters
the correction rather than sitting outside it.
\end{proposition}

\begin{proof}[Proof of the first-order term]
$e^{-\Lambda}=e^{-\Lambda^{\mathrm{crest}}}(1-\Gamma\int
e^{-\rho x}+O(\Gamma^{2}))$ pointwise; take expectations and divide by
$\E[e^{-\Lambda^{\mathrm{crest}}}]$, which converts the expectation of
the bracket into a $Q_{\mathrm{crest}}$-expectation by the definition of
that measure.
\end{proof}

\begin{remark}[The expansion fails exactly at the anatomy threshold]
\label{gr1:rem:perturbfail}
Under $Q_{\mathrm{crest}}$ the credit grows as $x_s\asymp\sigma
s^{3/2}$, so the correction term behaves as
$\int^\infty e^{-\rho\sigma s^{3/2}}ds$, which converges if and only if
$\rho>0$. Numerically at $\sigma=1$: the integral over $[0,60]$ is
$0.90$ at $\rho=1$ and $6.7$ at $\rho=0.05$, but $2.2\times10^{10}$ at
$\rho=-0.05$ and $1.0\times10^{60}$ at $\rho=-0.3$, with essentially all
the mass in the far tail once $\rho<0$.

The reading is worth stating, because it is a check on the model rather
than a decoration. A perturbation series whose radius of convergence
collapses at a parameter value is the standard signature of a genuine
phase transition there, as opposed to a modelling artefact or a
definitional boundary: the expansion fails because the object it is
expanding around ceases to exist, not because the algebra becomes
awkward. That $\Gamma$-perturbation theory breaks down precisely at
$\rho=0$ is therefore independent evidence that the trichotomy of
\S\ref{gr1:sub:anatomy} is a real division of the model and not a
convenient way of describing it. It also explains why no uniform
treatment of the two channels is available: for $\rho\le0$ the internal
channel is not a perturbation of anything, being the dominant term at
every large $x$.
\end{remark}

\begin{remark}[What is available in closed form, and what is not]
\label{gr1:rem:closedform}
A brief inventory, since the model is now large enough that the question
arises. \emph{Exact:} the viability product \eqref{gr1:eq:twocond}; the
Gaussian sigmoid \eqref{gr1:eq:sigmoid}; the trapped-regime optimum
$x_\ast$ and floor $h_\ast$ of \eqref{gr1:eq:xstar}; the gain minimum
$\kappa(\rho)e^{-(\rho b+k)/(1+\rho)}$; the first-order term of
\eqref{gr1:eq:perturb}. \emph{Exact as an exponent, unknown as a
constant:} the survival tail $Ct^{-1/4}$, whose prefactor involves the
Groeneboom--Jongbloed--Wellner harmonic function and is not evaluated
here. \emph{Series:} the bias field expansion
$(1-s/t)^{-1/4}=1+s/4t+O(s^2/t^2)$ of \eqref{gr:eq:rate}, in the
interpolation parameter $s/t$; and \eqref{gr1:eq:perturb}, in $\Gamma$.
\emph{Variational rather than closed:} the adverse-drift rate
$3\nu^{2}/8\sigma^{2}$ and its optimal profile, obtained from the
Cameron--Martin action by Cauchy--Schwarz --- the model's analogue of an
Euler--Lagrange computation, and the natural processing step whenever a
conditioned path law is wanted at large deviation scale.
\emph{Not available:} the trapped regime's survivor tail beyond the
exponential upper bound, and the asymptotics of any model in which
$\rho$ carries its own dynamics.
\end{remark}


% ---------------------------------------------------------------------
\subsection{Conditioning on survival}\label{gr1:sub:cond}\label{sh:sub:interp}

\subsubsection{The interpolation is a monotone flow of upward tilts}
\label{sh:subsub:flow}

\begin{definition}[Survivorship interpolation; bias field]
\label{sh:def:interp}
For $t>0$ set $\PP^{(t)}:=\PP(\,\cdot\mid T>t)$, and analogously
$\PP^{(t)}_\tau:=\PP(\,\cdot\mid\tau>t)$ for the hard-wall functional;
$\PP^{(0)}=\PP$ is the unbiased law. For $0\le s<t$ the \emph{bias field}
on the window $[0,s]$ is the Radon--Nikodym derivative
\begin{equation}\label{sh:eq:bias}
\beta_t(s):=\frac{d\PP^{(t)}}{d\PP}\Big|_{\Fc_s}
=\one_{\{T>s\}}\,
\frac{\PP_{X_s}(T>t-s)}{\PP_{X_0}(T>t)},
\qquad X_s=(\mu_s,D_s),
\end{equation}
the second equality by the Markov property; likewise for
$\PP^{(t)}_\tau$ with $\tau$ in place of $T$. The interpolation at finite
$t$ is thus globally absolutely continuous with respect to $\PP$, with
density $\one_{\{T>t\}}/\PP(T>t)$ on $\Fc_\infty$.
\end{definition}

\begin{lemma}[Survivorship tilts upward; \textbf{P}]\label{sh:lem:monotone}
In the driftless model, couple two initial data
$(\mu_0,D_0,h_0)\le(\mu_0',D_0',h_0)$ (componentwise in the first two
arguments) synchronously, with the same $W$. Then pathwise
$\mu'_t-\mu_t=\mu_0'-\mu_0\ge0$,
$D'_t-D_t=(D_0'-D_0)+(\mu_0'-\mu_0)t\ge0$, hence $h'\le h$,
$\Lambda'\le\Lambda$, and both $\PP_x(T>t)$ and $\PP_x(\tau>t)$ are
nondecreasing in $(\mu_0,D_0)$ (and nonincreasing in $h_0$). Consequently
every bias field \eqref{sh:eq:bias} is, on $\{T>s\}$, a nondecreasing
function of the state $X_s$: at every conditioning horizon, survivorship
bias is an upward tilt in rate and credit, and the interpolation is a
monotone flow of such tilts.
\end{lemma}

\begin{proof}
Immediate from the displayed pathwise inequalities and
Definition~\ref{sh:def:cox}.
\end{proof}

\subsubsection{The critical limit is the pure Doob transform}
\label{sh:subsub:Q}

\begin{theorem}[Existence and identification of the limit;
\textbf{C}, soft version \textbf{A}]\label{sh:thm:Q}
Assume the pointwise refined tail asymptotics
$\PP_{(\mu,d)}(\tau>t)\sim c\,\hm(\mu,d)\,t^{-1/4}$ of
Theorem~\ref{sh:thm:persist} and
Proposition~\ref{sh:prop:harmonic}, together with the martingale property
$\E_{x}[\hm(X_s)\one_{\{\tau>s\}}]=\hm(x)$ furnished by \cite{GJW99}.
Then for every $s\ge0$ and every bounded $\Fc_s$-measurable $\Psi$,
\[
\E\bigl[\Psi\mid\tau>t\bigr]\ \longrightarrow\ \E_Q[\Psi],
\qquad t\to\infty,
\]
where $Q$ is the time-homogeneous Markov law with
\begin{equation}\label{sh:eq:Qdensity}
\frac{dQ}{d\PP}\Big|_{\Fc_s}
=\one_{\{\tau>s\}}\,\frac{\hm(X_s)}{\hm(X_0)} :
\end{equation}
the \emph{pure} Doob $\hm$-transform, at eigenvalue one. Under
Ansatz~\ref{sh:ansatz} the same limit holds for
$\PP^{(t)}=\PP(\,\cdot\mid T>t)$.
\end{theorem}

\begin{proof}
Write $f_t:=\one_{\{\tau>s\}}\PP_{X_s}(\tau>t-s)/\PP_{X_0}(\tau>t)$, so
that $\E[\Psi\mid\tau>t]=\E[\Psi f_t]$. By the Markov property
$\E[f_t]=1$ for every $t$. Factorizing,
\[
f_t=\one_{\{\tau>s\}}\,
\frac{(t-s)^{1/4}\,\PP_{X_s}(\tau>t-s)}{t^{1/4}\,\PP_{X_0}(\tau>t)}
\,\Bigl(1-\frac st\Bigr)^{-1/4}
\ \longrightarrow\
f:=\one_{\{\tau>s\}}\frac{\hm(X_s)}{\hm(X_0)}
\quad\text{a.s.},
\]
using the assumed asymptotics at the (fixed, random) starting point
$X_s$ and at $X_0$, and $(1-s/t)^{-1/4}\to1$. The martingale property
gives $\E[f]=1=\lim\E[f_t]$, so by Scheff\'e's lemma $f_t\to f$ in
$L^1(\PP)$, and $\E[\Psi f_t]\to\E[\Psi f]=\E_Q[\Psi]$ for bounded
$\Psi$. That \eqref{sh:eq:Qdensity} defines a consistent, time-homogeneous
Markov family is the standard theory of Doob transforms by invariant
harmonic functions; time-homogeneity with no eigenvalue factor is
possible precisely because $\hm$ is invariant (eigenvalue one) for the
killed semigroup.
\end{proof}

\begin{remark}[Why survivorship has no clock here; \textbf{P}]
\label{sh:rem:noclock}
The mechanism deserves isolation. For \emph{regularly varying} tails the
kinematic factor $\PP_x(\tau>t-s)/\PP_x(\tau>t)\to1$ for each fixed $s$:
the age of the conditioning drops out at leading order, and the limit is
an un-tilted $h$-transform. For \emph{exponential} tails
$\PP_x(T>t)\approx C(x)e^{-\theta t}$ the same factor converges to
$e^{\theta s}$, and survivorship acquires a clock --- the eigen-tilting
of mode (III) in Theorem~\ref{sh:thm:modes}. The polynomial class is
exactly the class in which survivorship bias is, asymptotically, a
change of measure with no time-dependence. In the taxonomy of
penalisations of Wiener measure \cite{RVY06}, the survival weight
$e^{-\Lambda_t}$ thus lands in the clockless (martingale-normalisable)
class precisely because its expectation is regularly varying.
\end{remark}

\begin{proposition}[Interpolation rate; the tail index as linear-response
coefficient; \textbf{P/C}]\label{sh:prop:rate}
Under the hypotheses of Theorem~\ref{sh:thm:Q}, for fixed $s$ as
$t\to\infty$,
\begin{equation}\label{sh:eq:rate-interp}
\beta_t(s)\;=\;\one_{\{\tau>s\}}\,\frac{\hm(X_s)}{\hm(X_0)}
\Bigl(1-\frac st\Bigr)^{-1/4}\bigl(1+o(1)\bigr),
\qquad
\Bigl(1-\frac st\Bigr)^{-1/4}=1+\frac{s}{4t}+O\!\Bigl(\frac{s^2}{t^2}\Bigr).
\end{equation}
The natural interpolation parameter is therefore $s/t$, and the
persistence exponent $\tfrac14$ is the linear-response coefficient of
survivorship bias: the leading finite-$t$ correction to the survivor's
universe is $s/4t$ per window of length $s$. Approach to the limit is
polynomial --- the critical signature, contrasting with the exponential
approach of mode (I) and the exponentially growing tilt of mode (III)
below.
\end{proposition}

\begin{proof}
Insert the assumed asymptotics into \eqref{sh:eq:bias} and expand.
\end{proof}

\subsubsection{The limit object: survivorship drift and promotion to
transience}\label{sh:subsub:promotion}

\begin{theorem}[Survivorship drift; \textbf{P} given \textbf{C} inputs]
\label{sh:thm:drift}
Under $Q$ the state solves
\begin{equation}\label{sh:eq:Qsde}
d\mu_t=\sigma^2\,\partial_\mu\log\hm(\mu_t,D_t)\,dt+\sigma\,dW^Q_t,
\qquad dD_t=\mu_t\,dt,
\end{equation}
for a $Q$-Brownian motion $W^Q$: conditioning adds drift only in the
noisy coordinate. With the representation \eqref{sh:eq:harm},
\begin{equation}\label{sh:eq:G}
b_Q(\mu,D)=\frac{\sigma^2}{\mu}\,G(\kappa),
\qquad
G(\kappa):=\frac12-3\kappa\,\frac{F'(\kappa)}{F(\kappa)},
\end{equation}
and: \emph{(i)} $b_Q\ge0$ everywhere (survivorship never pushes the rate
down), by Lemma~\ref{sh:lem:monotone} applied to $\hm$;
\emph{(ii)} $G(0^+)=\tfrac12$, since
$\hm(\mu,0^+)=c_1'\mu^{1/2}$ by Theorem~\ref{sh:thm:persist};
\emph{(iii)} $G(\kappa)\to0$ as $\kappa\to\infty$, since
$\hm(0^+,d)=c_2 d^{1/6}$ is independent of $\mu$ at leading order.
Statements (ii)--(iii) use the regularity of the profile $F$ furnished by
the explicit formulas of \cite{GJW99}.
\end{theorem}

\begin{proof}
For invariant harmonic $\hm$, the transformed generator is
\[
\Lc^{\hm}f=\hm^{-1}\Lc(\hm f)
=\Lc f+\sigma^2\bigl(\partial_\mu\log\hm\bigr)\,\partial_\mu f,
\]
since only $\partial_\mu^2$ is second order in
\eqref{sh:eq:kolm}; this is \eqref{sh:eq:Qsde}. Formula \eqref{sh:eq:G}
is the logarithmic derivative of \eqref{sh:eq:harm} using
$\partial_\mu\kappa=-3\kappa/\mu$. Positivity: synchronous coupling gives
$\PP_{(\mu,d)}(\tau>t)$ nondecreasing in $\mu$, hence
$\partial_\mu\hm\ge0$ wherever the derivative exists. The boundary values
follow from the quoted prefactor laws: $F(0)=c_1'\in(0,\infty)$ and
$\kappa F'(\kappa)/F(\kappa)\to0$ as $\kappa\to0$ by smoothness of $F$ at
the origin; as $\kappa\to\infty$, $\hm\sim c_2d^{1/6}$ forces
$\mu\,\partial_\mu\log\hm\to0$.
\end{proof}

\begin{corollary}[Wall drift; promotion to transience; \textbf{P/C}]
\label{sh:cor:bessel}
Near zero credit ($\kappa\to0^+$) the survivorship drift reduces to
$\sigma^2/2\mu$, so that in this regime the conditioned rate solves
\[
d\mu_t=\frac{\sigma^2}{2\mu_t}\,dt+\sigma\,dW^Q_t .
\]
Two cautions attach to this formula.
\emph{(i)} This is the Bessel drift of index $\delta=2$, not $\delta=3$:
the exponent $\tfrac12$ in $\hm(\mu,0^+)\propto\mu^{1/2}$ is half the
exponent $1$ of the classical harmonic function $x$ for Brownian motion
killed at the origin, the halving being exactly that of the persistence
exponent from $\tfrac12$ to $\tfrac14$ under one integration.
\emph{(ii)} Under $Q$ the coordinate $\mu$ is not autonomous --- by
\eqref{sh:eq:G} its drift depends on $\kappa$, hence on both coordinates
--- so the comparison is a statement about the \emph{form} of the drift in
a limiting regime, not an identification of processes.

The correct P\'olya statement is therefore not a Bessel index but a
recurrence class, and at that level it is unconditional: under $\PP$ the
credit is null recurrent, returning to its initial level a.s.\
(Theorem~\ref{sh:thm:noimm}) with infinite expected return time
(Corollary~\ref{sh:cor:mean}), whereas under $Q$ it is transient,
$D_s\to\infty$ a.s.\ (Theorem~\ref{sh:thm:singular}). Conditioning on
survival promotes the credit from the recurrent to the transient class ---
the P\'olya $2\mathrm{D}\to3\mathrm{D}$ transition --- with no literal
change of dimension anywhere. Dimension remains, as in
\S\ref{sh:subsub:scale}, the wrong order parameter for the state dynamics;
what survivorship alters is the recurrence class of the \emph{ensemble}.
\end{corollary}

\begin{theorem}[Locally ours, globally elsewhere; \textbf{P/C}]
\label{sh:thm:singular}
\emph{(i)} $Q(\tau=\infty)=1$.
\emph{(ii)} $\hm(X_s)\to\infty$ $Q$-a.s.\ as $s\to\infty$; moreover
$D_s\to\infty$ $Q$-a.s.\ \cite{GJW99}.
\emph{(iii)} $\liminf_sD_s=-\infty$ $\PP$-a.s.
\emph{(iv)} Hence $Q|_{\Fc_s}\ll\PP|_{\Fc_s}$ for every finite $s$, with
the explicit density \eqref{sh:eq:Qdensity}, while $Q\perp\PP$ on
$\Fc_\infty$: the survivor's universe is locally absolutely continuous
with respect to ours on survivors, and globally disjoint from it. The
exit from the measure class is completed only at $t=\infty$; at every
finite conditioning the bias is the bounded tilt of
Definition~\ref{sh:def:interp}.
\end{theorem}

\begin{proof}
(i) $Q(\tau>s)=\E_\PP[\one_{\{\tau>s\}}\hm(X_s)]/\hm(X_0)=1$ for every
$s$ by the martingale property. (ii) Under $Q$, $1/\hm(X_s)$ is a
nonnegative supermartingale:
$\E_Q[1/\hm(X_{s+r})\mid\Fc_s]
=\PP_{X_s}(\tau>r)/\hm(X_s)\le1/\hm(X_s)$; it therefore converges
$Q$-a.s., and since
$\E_Q[1/\hm(X_s)]=\PP_{X_0}(\tau>s)/\hm(X_0)\to0$, Fatou forces the a.s.\
limit to vanish, i.e.\ $\hm(X_s)\to\infty$; the sharper statement
$D_s\to\infty$ is part of the description of the conditioned process in
\cite{GJW99}. (iii) is contained in the proof of
Theorem~\ref{sh:thm:noimm} (the set $A$ there is unbounded, and on it
$D_u\le D_0+\mu_0u-\sigma\varepsilon u^{3/2}\to-\infty$). (iv) Local
absolute continuity is \eqref{sh:eq:Qdensity}; mutual singularity on
$\Fc_\infty$ follows since $\{\lim_sD_s=+\infty\}$ has $Q$-probability
one and $\PP$-probability zero.
\end{proof}

\subsubsection{The three modes of survivorship}\label{sh:subsub:modes}

\begin{theorem}[Modes of the interpolation; \textbf{P/C/S}]
\label{sh:thm:modes}
As $t\to\infty$ the survivorship interpolation exhibits exactly three
asymptotic modes across the driver classes of \S\ref{sh:sub:horizon}:
\begin{enumerate}[label=\emph{(\Roman*)},itemsep=2pt]
\item \emph{Saturation} \textbf{[P]} (favorable drivers,
$\PP(T=\infty)>0$): since $\{T>t\}\downarrow\{T=\infty\}$,
\[
\bigl\|\PP^{(t)}-\PP(\,\cdot\mid T=\infty)\bigr\|_{\mathrm{TV}}
\ \le\ \frac{2\,\PP(t<T<\infty)}{\PP(T>t)}\ \longrightarrow\ 0
\quad\text{on }\Fc_\infty,
\]
and every bias is bounded by $1/\PP(T=\infty)$: survivorship saturates at
a bounded, globally absolutely continuous reweighting. No new universe.
\item \emph{Transformation} \textbf{[C/A]} (critical drivers; the
driftless model): the limit is the pure Doob transform of
Theorem~\ref{sh:thm:Q}, locally absolutely continuous and globally
singular (Theorem~\ref{sh:thm:singular}), approached at the polynomial
rate of Proposition~\ref{sh:prop:rate}. The new universe is entered, but
only at infinity.
\item \emph{Concentration} \textbf{[S]} (adverse drivers, exponential
tails; e.g.\ Proposition~\ref{sh:prop:advdrift}): on fixed windows the
limit is the eigen-tilted transform of
Remark~\ref{sh:rem:noclock} with clock
$\theta=3\nu^2/8\sigma^2$, while at macroscopic scale the conditioned
path law collapses onto the deterministic optimal profile of
Proposition~\ref{sh:prop:advdrift}(iii):
$\sup_{v\in[0,1]}\bigl|\mu_{tv}/t-\tfrac{|\nu|}4v(2-3v)\bigr|\to0$ in
$\PP^{(t)}$-probability. Survivorship becomes a law of large numbers for
the rare: the survivor ensemble is asymptotically deterministic at
leading order.
\end{enumerate}
The three modes are in bijection with the horizon trichotomy
(Theorem~\ref{sh:thm:trichotomy}, Proposition~\ref{sh:prop:recurrent}),
and they constitute a strictly finer invariant: modes (II) and (III)
share $\PP(\Imm)=0$ but are separated by the interpolation --- by whether
the survivor's universe is approached polynomially, without a clock, or
exponentially, with one.
\end{theorem}

\begin{proof}
(I) is continuity of measure from above plus the displayed
total-variation estimate, whose numerator is
$2\PP(\{t<T<\infty\})$ and whose denominator is bounded below by
$\PP(T=\infty)>0$. (II) collects Theorems~\ref{sh:thm:Q}
and~\ref{sh:thm:singular} and Proposition~\ref{sh:prop:rate}. (III) is
the large-deviation sketch of Proposition~\ref{sh:prop:advdrift}(iii)
read under the conditional measure, together with
Remark~\ref{sh:rem:noclock}; elevating it to a full proof requires the
sharp prefactor asymptotics $\PP_x(T>t)\sim C(x)e^{-\theta t}$, which we
do not establish here.
\end{proof}

\begin{remark}[The verdict]\label{sh:rem:verdict}
The section's question can now be answered in one paragraph. The
survivorship interpolation $t\mapsto\PP^{(t)}$ is a monotone flow of
upward tilts (Lemma~\ref{sh:lem:monotone}), globally absolutely
continuous at every finite $t$, whose asymptotic character is the
sharpest invariant of the system and takes exactly three forms
(Theorem~\ref{sh:thm:modes}). The driftless model realizes the critical
form: the limit exists, is the pure Doob transform by the
persistence-harmonic function, carries no clock
(Remark~\ref{sh:rem:noclock}), is approached at rate $s/4t$ with the tail
exponent as linear-response coefficient
(Proposition~\ref{sh:prop:rate}), is locally equivalent on survivors to
the unbiased law yet globally singular
(Theorem~\ref{sh:thm:singular}), and differs in generator by the
universal survivorship drift \eqref{sh:eq:G}, equal to $\sigma^2/2\mu$
at the wall and carrying the credit from the recurrent to the transient
class (Corollary~\ref{sh:cor:bessel}). The gauge structure of
\S\ref{sh:sub:model} persists to the end: the initial data cancel from
the limit density \eqref{sh:eq:Qdensity} up to the normalization
$\hm(X_0)$, and the survivorship drift is universal. The event horizon
absent from state space is therefore real, but it is a feature of the
ensemble, not of the state: it is crossed by the conditioning parameter
$t$ --- which is to say, by survival itself --- and never by the system.
\end{remark}

% ---------------------------------------------------------------------
\subsection*{Summary of the model section}

% ---------------------------------------------------------------------


The system is a Kolmogorov diffusion in $(\mu,D)$ with exponential
killing; at exponent level it is the persistence problem for integrated
Brownian motion. Initial data are gauge: they act only through an
equivalent timescale $t_*$, and the one intrinsic coordinate is
$\kappa=\sigma^2D/\mu^3$ (\S\ref{sh:sub:model}). For the driftless model,
viability is impossible for every initial condition
(Theorem~\ref{sh:thm:noimm}), yet life expectancy is infinite:
$\PP(T>t)\asymp(t_*/t)^{1/4}$, with Pareto($\tfrac14$) residual life,
quantiles growing as the fourth power of rarity, and an effective Lindy
hazard $1/4t$ --- null recurrence, the two-dimensional case of the P\'olya
analogy (\S\ref{sh:sub:verdict}). No horizon exists there: the crossover
is a maximally soft quarter-power law smeared over orders of magnitude,
and recovery cost is sublinear in credit. Where genuine horizons exist,
the order parameter is the boundedness of the driver's scale function,
the horizon profile \emph{is} the normalized scale function
(Theorem~\ref{sh:thm:trichotomy}), the recurrent class is decided by the
stationary mean with mortal critical case
(Proposition~\ref{sh:prop:recurrent}), and the solvable unstable-OU
driver realizes the profile exactly as a Gaussian sigmoid
$\Phi(\mu_0\sqrt{2\nu}/\sigma)$ of width $\sigma/\sqrt{2\nu}$, a step
function only in the small-noise limit (\S\ref{sh:sub:horizon}). Under
adverse drift the survival tail is exponential with rate
$3\nu^2/8\sigma^2$, carried by a front-load-then-coast optimal
fluctuation. The destination is the survivorship interpolation
(\S\ref{sh:sub:interp}): a monotone flow of upward tilts $\PP^{(t)}$
whose critical limit is the pure Doob transform by the
persistence-harmonic function, approached without a clock at rate
$s/4t$ --- the tail exponent as linear-response coefficient --- locally
absolutely continuous on survivors yet globally singular, and generated
by the universal survivorship drift, equal to $\sigma^2/2\mu$ at the
wall: conditioning carries the credit from the recurrent to the transient
class, the P\'olya promotion effected by survival rather than by
dimension. Across
drivers the interpolation has exactly three modes --- saturation,
transformation, concentration --- in bijection with the horizon
trichotomy and strictly finer than the survival dichotomy. A horizon is
not a place. Where it exists in state space it is a ratio of scales; and
where it does not, it is real nonetheless --- located in the ensemble,
crossed not by the state but by the conditioning parameter, which is to
say by survival itself.



% ---------------------------------------------------------------------
\subsection*{Where these results are used}\label{gr1:sub:consumers}

This section has been developed as a self-contained account of a hazard
process, and is readable as one. It is not, however, freestanding within
this work: each of its principal results has a consumer elsewhere, and
naming them is the quickest way to see why the section is here rather
than in a paper of its own.

\emph{The two-condition theorem} (\ref{gr1:thm:twocond}) is the paper's
most-used result. It supplies the regime trichotomy on which
\S\ref{gr:subsub:regimes} scopes the General/Golden distinction --- and
which, as recorded there, restricts the measure-theoretic claim to the
critical case. It supplies the dominance argument of
\S\ref{hr:rem:concession}, where the extinctionist option and the
settled-at-$x_\ast$ option are shown to carry viability probability zero
alike. And its anatomy factor is what \S\ref{ad:sec} exists to compute.

\emph{Technoanatomy} (Definition~\ref{gr1:def:anatomy}) is consumed
twice over. \S\ref{gr:sub:extraction} identifies extractive organisation
with $\rho<0$, which is what converts the Extraction Hypothesis from a
vocabulary match into a claim about the sign of an exponent; and
\S\ref{ad:sub:toward} takes the derivation of $\rho$ from institutional
dynamics as its object, together with an account of why the machinery
there does not yet reach it.

\emph{The trapped-regime theorem} (\ref{gr1:thm:trapped}) does the work
in \S\ref{hr:sub:extinction}, where adverse anatomy is shown fatal
independently of the driver --- the step that makes the argument against
extinctionism turn on forfeiture rather than on comparative hazard.

\emph{The tolerance collapse} (\ref{gr1:prop:tolerance}) yields the
ordering constraint invoked in \S\ref{hr:subsub:stakes}: since the
affordable duration of adverse anatomy decays as $e^{-|\rho|x}$, anatomy
must be secured before capability accumulates, and a system that
attends to its arrangement late does so with the least margin it will
ever have.

\emph{The gauge results} of \S\ref{gr1:sub:coords} recur in a different
register in \S\ref{to:sub:selection}, where the selection of a
task-object among the many that individuate a given object is a gauge
fixing of the same shape: find the quantities the arbitrary choice
leaves invariant, and state the theory in those.

\emph{And the survivorship interpolation} of \S\ref{gr1:sub:cond} is
what \S\ref{gr:sec} is about in its entirety --- the Doob limit there
being the Golden Biosphere, and the local absolute continuity of
Theorem~\ref{sh:thm:singular} being the epistemic ceiling that
\S\ref{gr:sub:ceiling} turns against this work's own conclusions as
firmly as against its targets.

% ---------------------------------------------------------------------
\section{Golden reasoning}\label{gr:sec}

\begin{quote}\small
The preceding sections were mathematics. This one is philosophy, and
says so. Its method is to let the mathematics \emph{limit the
possibility space} and then to reason philosophically inside what
remains. Two streams of load-bearing claims therefore run through it,
and are tagged accordingly: mathematical results carry the tags of
\S\ref{gr1:sec}, while philosophical premises carry \textbf{[$\Phi$]} and
are numbered $\Phi1,\dots,\Phi5$, each graded in the corpus's own
vocabulary --- stipulation, declaration, conjecture-grade, premise,
imported result --- and each priced with its falsifier in the ledger of
\S\ref{gr:sub:ledger}. Nothing is smuggled: where the argument turns on
a premise rather than a theorem, the premise is named, argued for, and
the thing whose exhibition would break it stated.

The thesis is that the fundamental categorical division among
world-scale living systems is not Pre-Civilization versus Civilization
but \emph{General} versus \emph{Golden} --- $\PP$ versus $Q$ --- and
that this division, unlike the other, is not a threshold anyone chose.
The ideal category is empty, but its practical form is not: a biosphere
of age $t$ is $\varepsilon$-Golden over windows of length $4\varepsilon
t$, so that Goldenness is bought by having lasted and by nothing else.
From this follow: the rejection of the Kardashev scale, not because its
thresholds are arbitrary (though they are) but because the model admits
a terminal stage gated by a variable the ladder cannot index; the rejection of the
violent-aliens thesis, conditional on one named premise; and a strict
epistemic ceiling which limits this section's own conclusions as much as
its targets'.
\end{quote}

\subsection*{Two grading vocabularies}
\label{gr:sub:grades}

Two tag systems meet in this section and should not be confused.
\S\ref{gr1:sec} grade mathematical claims by how well established they
are: \textbf{[P]} proved, \textbf{[C]} classical or cited, \textbf{[A]}
conditional on the soft-wall Ansatz, \textbf{[S]} scaling-level,
\textbf{[O]} open. The companion program grades commitments instead by
\emph{what sort of thing they are} --- exposed postulate,
conjecture-grade, assumption, standing license, schema premise --- and
prices each with the thing whose exhibition would break it
\cite{Nexus, Ben26Found, Ben26Optimal}. The two are orthogonal
rather than rival, and this section uses both: the mathematical tags for
what it inherits from \S\ref{gr1:sec}, the companion's vocabulary for what
it adds. Each premise below accordingly carries its kind in its
statement, and \S\ref{gr:sub:ledger} assembles the whole into a ledger
in the companion's two-column format, on the companion's own principle
that a priced cost which could never be wrong is not a cost but a
decoration.

% ---------------------------------------------------------------------
\subsection{The referent problem}\label{gr:sub:referent}

\begin{premise}[The biospheric referent; \emph{methodological
stipulation}]\label{gr:phi1}
The appropriate unit of analysis for questions about the long-run
trajectory of world-scale living systems is the \emph{biosphere}: the
totality of the living and technical activity associated with a world,
considered as one system.
\end{premise}

The argument is one of referential stability across the relevant
timescale. ``Species'' fails backwards: the system that became us was
not a species for most of its history, and speciation is not the joint
along which the relevant quantity --- total hazard-bearing organisation
--- varies. ``Civilization'' fails backwards for the same reason and
worse, since it names a phase rather than a system, so that using it as
the unit builds a phase transition into the vocabulary before any
analysis has occurred. ``Planet'' fails forwards: it is a substrate, not
a system, and there is no reason the system should remain confined to
one. ``Biosphere'' is the only available term that refers to the same
thing at every point of the trajectory, which is exactly what an
asymptotic analysis requires.

This is not a terminological preference. A unit that changes referent
partway through cannot appear in a limit theorem. \S\ref{gr1:sec} concern
the $t\to\infty$ behaviour of a single system; the vocabulary must
survive the limit or the theorems cannot be applied at all.

\begin{remark}[What this already costs Kardashev]\label{gr:rem:kardvocab}
Kardashev's scale \cite{Kar64} classifies \emph{civilizations} by energy
capture. Whatever its other merits, the choice of unit forecloses the
move this section makes: one cannot ask what the universe looks like
from the vantage of a long-survived biosphere if the vocabulary presumes
a civilization, since ``civilization'' carries semantic content ---
agency, society, purpose --- that the analysis is supposed to derive
rather than assume. The unit is doing work that ought to be done by
argument.
\end{remark}

% ---------------------------------------------------------------------
\subsection{General and Golden}\label{gr:sub:generalgolden}

\begin{definition}[General and Golden biospheres]\label{gr:def:gg}
A \emph{General Biosphere} is a biosphere considered under $\PP$, the
unconditioned law of \S\ref{gr1:sec}. A \emph{Golden Biosphere} is a biosphere
considered under $Q$, the survivorship limit --- the Doob $\hm$-transform
of \S\ref{gr1:sec}, the $t\to\infty$ limit of conditioning on $\{T>t\}$. The
family $\PP^{(t)}=\PP(\,\cdot\mid T>t)$ interpolating them is the
\emph{goldening} of the biosphere.
\end{definition}

Three features of this definition do work that no threshold-based
division can do.

\emph{First, the split is categorical and unchosen --- in the critical
regime.} When $\PP(T=\infty)=0$, as it is under a recurrent driver with
$\rho>0$, the laws $Q$ and $\PP$ are \emph{mutually singular} on the
infinite-horizon $\sigma$-algebra: there is an event --- the credit
tending to infinity --- of $Q$-probability one and $\PP$-probability
zero. General and Golden are then not two regions of one population
separated by a line someone drew but different measures, the difference
established by a theorem rather than a stipulation. Pre-Civilization
versus Civilization has no such status, being a threshold in a
continuous variable with both the variable and the value contested.

The scoping is not a technicality and \S\ref{gr:subsub:regimes} states
it in full: outside the critical regime the singularity fails, and with
it the unchosenness. This is the section's strongest claim and it is
regime-specific.

\emph{Second, in the critical regime nothing occupies the Golden
category --- as an ideal.} By the mode classification of
\S\ref{sh:subsub:modes}, that regime is the critical one: the survivor's
law is approached but entered only at infinity. Every actual biosphere is somewhere on the interpolation
$\PP^{(t)}$, at finite $t$, and none is ideally Golden. This would make
the category inapplicable were it not for the practical form of
\S\ref{gr:subsub:practical}, which is occupiable, and which is what
Golden reasoning actually assumes of anything.

\emph{Third, the approach is quantified.} The bias field of \S\ref{gr1:sec}
gives, on a window of length $s$ at conditioning horizon $t$,
\begin{equation}\label{gr:eq:rate}
\beta_t(s)=\one_{\{T>s\}}\frac{\hm(X_s)}{\hm(X_0)}
\Bigl(1+\frac{s}{4t}+O(s^{2}/t^{2})\Bigr),
\end{equation}
so that goldening proceeds at rate $s/4t$ with the persistence exponent
$\tfrac14$ as its linear-response coefficient. ``The further out you
look, the more Golden it gets'' is not a slogan here; it is
\eqref{gr:eq:rate}.

\subsubsection{Golden reasoning in the three regimes}
\label{gr:subsub:regimes}

The two-condition theorem partitions the model, and Golden reasoning has
a different character in each cell. Since the apparatus of this section
was built inside one of them, the partition is stated before anything is
built on it.

\begin{center}\small
\begin{tabular}{@{}lll@{}}
\toprule
regime & $\PP(T=\infty)$ & what General/Golden becomes\\
\midrule
\emph{trapped}: $\rho\le0$ & $0$, $\E[T]<\infty$ &
concentration; $Q$ collapses onto one path\\
\emph{critical}: $\rho>0$, recurrent & $0$, $\E[T]=\infty$ &
\textbf{mutually singular; the split is categorical}\\
\emph{viable}: $\rho>0$, transient & $>0$ &
$Q\ll\PP$; an ordinary sub-population\\
\bottomrule
\end{tabular}
\end{center}

Three consequences, and the middle one costs this section something.

\emph{Trapped.} With adverse anatomy, life expectancy is finite
(Theorem~\ref{gr1:thm:trapped}) and survivors concentrate on the
optimal-fluctuation path: mode (III) of \S\ref{sh:subsub:modes}, where
the conditioned ensemble is asymptotically deterministic. There is a
General/Golden distinction, but the Golden ensemble has no randomness
left to study and no biosphere reaches it.

\emph{Critical.} This is where the rest of this section lives, and it is
the only cell in which the categorical claim holds. Death is certain but
life expectancy infinite; $Q$ is entered only at infinity and is
mutually singular with $\PP$; the interpolation $\PP^{(t)}$ is
non-trivial, and everything from \eqref{gr:eq:rate} through the
epistemic ceiling is a statement about it.

\emph{Viable.} With good anatomy and feedback, $\PP(T=\infty)>0$ and
conditioning on immortality is an \emph{ordinary conditional
probability}: $Q\ll\PP$ globally, total-variation convergence holds, and
the Golden Biosphere is a sub-population of the General one --- a set of
paths one could in principle be a member of, with no measure-theoretic
novelty and no unchosen categorical split. Mode (I), saturation.

The consequence for this section's central claim should be stated
without softening, because it is a real limitation. \emph{The
measure-theoretic recasting of the survivor category is available in
exactly one of the three regimes.} Where survival is genuinely
achievable, General and Golden differ by a reweighting rather than in
kind, and the categorical division reduces to the ordinary distinction
between a population and a favoured part of it. The framework's most
distinctive structural claim thus holds precisely where immortality does
\emph{not}, and a reader who believes we inhabit the viable regime
should regard \S\S\ref{gr:sub:generalgolden}--\ref{gr:sub:ceiling} as
describing a case they think we are not in.

Two things blunt the loss without removing it. Which regime obtains is
decided by $\rho$ and by the driver's scale function, neither of which
is observable (\S\ref{gr:sub:ceiling}), so no one is presently entitled
to the belief that we inhabit the viable regime. And the critical regime
is not a knife-edge in the pejorative sense but the boundary the other
two are defined against: it is where the apparatus is sharpest precisely
because it is where the two failure modes meet.

\subsubsection{Practical and Ideal Golden Biospheres}
\label{gr:subsub:practical}

The ideal category is empty, but the questions one asks of a biosphere
are not questions about its infinite horizon. They concern windows ---
the next millennium, the next galactic orbit. On a window the
distinction collapses to a tolerance.

\begin{definition}[Practical and Ideal]\label{gr:def:practical}
The \emph{Ideal Golden Biosphere} is $Q$. A biosphere of age $t$ is a
\emph{Practical Golden Biosphere to tolerance $\varepsilon$ on windows
of length $s$} if $\PP^{(t)}$ and $Q$ agree to within $\varepsilon$ on
the $\sigma$-algebra $\mathcal F_s$.
\end{definition}

\begin{proposition}[Practical Goldenness is a ratio, not a level;
\textbf{S}]\label{gr:prop:ratio}
By the bias field \eqref{gr:eq:rate}, the density of $\PP^{(t)}$ with
respect to $Q$ on $\mathcal F_s$ is $(1-s/t)^{-1/4}(1+o(1))$, so the two
laws deviate at order $s/4t$. Hence
\begin{equation}\label{gr:eq:practical}
\text{$\varepsilon$-Golden on windows of length } s
\qquad\Longleftrightarrow\qquad
t\;\gtrsim\;\frac{s}{4\varepsilon},
\end{equation}
and a biosphere of age $t$ is $\varepsilon$-Golden for every question
concerning a window shorter than the \emph{Golden horizon}
$s_\varepsilon(t)=4\varepsilon t$. The tag is \textbf{[S]}: the $s/4t$
scale follows from the leading term of \eqref{gr:eq:rate}, but the exact
total-variation constant depends on the $o(1)$ structure and is not
established here.
\end{proposition}

\begin{remark}[The invariant]\label{gr:rem:invariant}
The Golden horizon is a fixed fraction of age:
$s_\varepsilon(t)/t=4\varepsilon$, independent of $t$. \emph{Every
surviving biosphere is Golden over a constant fraction of its own
history, and the fraction is four times one's tolerance.} There is no
threshold in strength, in energy, in technology, or in any state
variable whatever; the only thing that buys Goldenness is having lasted,
and what it buys is proportional. To be $1\%$-Golden over the coming
millennium requires twenty-five thousand years behind one.

This sharpens \S\ref{gr:subsub:kardashev}: the objection to Kardashev is
not merely that its thresholds are placed arbitrarily, but that the
operative threshold does not live in the space it measures.
\end{remark}

\begin{remark}[The premise of Golden reasoning is weaker than it looks]
\label{gr:rem:weakpremise}
This answers the natural suspicion that Golden reasoning assumes
something optimistic. It does not. By \eqref{gr:eq:practical}, practical
Goldenness to any fixed tolerance follows from age alone; and age is
exactly what the question already supposes. To ask what our descendants
are like across any substantial future is already to condition on there
\emph{being} descendants across that future, and that conditioning by
itself confers practical Goldenness over windows proportional to the
horizon supposed. The assumption is not that the future goes well. It is
that there is one --- and the model converts that supposition, and
nothing further, into the whole apparatus.

The corollary for the present is deflationary and should be stated
plainly: we are not practically Golden to any interesting tolerance, and
Golden reasoning does not claim we are. It describes the ensemble our
descendants would be drawn from in those futures where they exist.
\end{remark}

\begin{premise}[Golden reasoning is a vantage, not an obligation;
\emph{scope declaration}]\label{gr:phi2}
Golden reasoning is the practice of re-analysing questions about
biospheres by taking $Q$ as the reference measure rather than $\PP$. It
is \emph{inferential}: it licenses conditional expectations, not duties.
\end{premise}

\begin{remark}[Why this is not longtermism]\label{gr:rem:notlongtermism}
The objection that any future-weighted reasoning collapses into ``act in
the interest of the future, as you happen to conceive it'' does not
apply, and the reason is structural rather than rhetorical. A normative
longtermism requires future beings whose interests ground obligations.
Golden reasoning has no such beings available: $Q$ is a measure that
nothing occupies in the critical regime
(Definition~\ref{gr:def:gg}), where moreover $\PP(T=\infty)=0$. In the
viable regime the objection would have purchase and is met differently,
by $\Phi2$ alone rather than by the vacuity of the reference class. What $Q$ supplies is a
\emph{conditional distribution}, from which one may infer what
long-survived systems are like without thereby acquiring any duty toward
them. The distinction is the ordinary one between a likelihood and a
value, and it is preserved here by the mathematics rather than by
stipulation. Whether one ought to care about the far future is a
question this section does not address and does not need.
\end{remark}

% ---------------------------------------------------------------------
\subsection{Technostrength and technoanatomy}\label{gr:sub:coords}

\begin{definition}[The two coordinates]\label{gr:def:coords}
\emph{Technostrength} is the muffling exponent $x_t$: the capacity the
biosphere has accumulated, positive when its own organisation holds its
hazard below background and negative when it holds it above.
\emph{Technoanatomy} is the exponent $\rho$ of
Definition~\ref{gr1:def:anatomy}, the sign of
$-\partial_x\log h^{\mathrm{int}}$: whether acquiring capacity lowers or
raises the biosphere's hazard \emph{to itself}.

The two are of different kinds, and keeping them apart is the whole of
what follows. Technostrength is a \emph{state coordinate}: it
fluctuates, and by Theorem~\ref{tc:thm:half} every biosphere ---
however well constituted --- spends exactly half of logarithmic time
with $x<0$. Technoanatomy is a \emph{structural property}: it does not
fluctuate at all, it says what capacity does when acquired, and
$\rho<0$ names the biosphere whose own strength is turned against it,
which is what Definition~\ref{gr1:def:anatomy} means by eating oneself.
\emph{Adverse strength} is $x<0$; \emph{adverse anatomy} is $\rho<0$,
and only the second is a defect.
\end{definition}

Two results of \S\ref{gr1:sec} are worth restating in this vocabulary because
they are counterintuitive.

\emph{Technostrength is not cumulative --- in the critical regime.}
Under a recurrent driver $x$ oscillates on scale $t^{3/2}$ and returns
to every level infinitely often: it is not a stock that accumulates but
a magnitude that fluctuates, so a framework treating it as monotonically
increasing has the wrong kind of object. Under a \emph{transient} driver
it does accumulate, and the ladder picture is locally right --- which is
the concession \S\ref{gr:subsub:regimes} makes precise and prices.

\emph{Every biosphere is adversely configured half the time.} By
Theorem~\ref{tc:thm:half} the set $\{t:x_t<0\}$ has logarithmic density
exactly $\tfrac12$, for every parameter value and every biosphere.
Adverse \emph{strength} is therefore not a defect distinguishing some
biospheres from others but a condition all of them occupy half of
logarithmic time --- which is precisely why it cannot be what
``bad anatomy'' means, and why the two must be kept apart.

\subsubsection{What conditioning does to each coordinate}
\label{gr:subsub:measured}

The following are simulation results for the two-channel model of
\S\ref{gr1:sec} at $\sigma=1$, $\epsilon=0.05$, $n=3$, $b=3$, $\rho=2.5$,
$k=6$, conditioning horizon $T=4$, $N=1.2\times10^{5}$ particles.

\begin{center}\small
\begin{tabular}{@{}lrr@{}}
\toprule
& General ($\PP$) & Golden ($Q$)\\
\midrule
$\PP(x<0)$ (adverse strength) & $0.501$ & $0.0037$\\
median $x/\sigma T^{3/2}$ & $-0.001$ & $+1.085$\\
$\PP(\rho>0)$, $\rho\sim\mathrm{Unif}(-1,3)$ & $0.750$ & $0.839$\\
$5\%$ quantile of $\rho$ & $-0.80$ & $-0.44$\\
\bottomrule
\end{tabular}
\end{center}

\begin{remark}[Truncation versus rescaling]\label{gr:rem:truncresc}
Two numbers, and they differ in kind rather than in force.
Adverse strength is suppressed by a factor of $137$ --- abruptly, and
within a few multiples of the intrinsic timescale. Technoanatomy is
selected too, but \emph{gradually}: $\PP(\rho>0)$ moves from $0.75$ to
$0.839$ at the same horizon, reaching $0.909$ only by $T=6$, with the
negative tail truncating steadily ($5\%$ quantile $-0.80\to-0.44\to
-0.14$).

The asymmetry is structural rather than incidental, and it follows from
what the two objects are. Strength is a state coordinate, so
conditioning pushes it out of a half-space it can and does re-enter ---
fast, and reversibly. Anatomy is a fixed structural parameter, so it can
be selected only through the accumulated hazard cost of never changing
--- slow, and irreversibly. Anatomy is the slower and the more decisive:
by Theorem~\ref{gr1:thm:trapped}, adverse anatomy is eventually fatal
with certainty, while adverse strength is a condition every biosphere
occupies half of logarithmic time and survives.


\end{remark}

\subsubsection{Does conditioning select technoanatomy?}
\label{gr:subsub:anatsel}

The measurement above concerns technostrength, a state coordinate; the
claim the argument needs is about technoanatomy, a structural parameter,
which therefore cannot be selected \emph{within} a trajectory. It can be
selected across an ensemble, and that is directly computable: place a
prior on $\rho$, evolve, and reweight by survival. Taking
$\rho\sim\mathrm{Uniform}(-1,3)$, so that $\PP(\rho>0)=0.75$ a priori:

\begin{center}\small
\begin{tabular}{@{}rrrr@{}}
\toprule
$T$ & $\PP(\text{survive})$ & $\PP(\rho>0\mid\text{survived})$ &
$5\%$ quantile of $\rho$\\
\midrule
$2$ & $2.8\times10^{-3}$ & $0.745$ & $-0.80$\\
$4$ & $2.2\times10^{-4}$ & $0.839$ & $-0.44$\\
$6$ & $1.9\times10^{-4}$ & $0.909$ & $-0.14$\\
\bottomrule
\end{tabular}
\end{center}

\noindent Survivorship does select for positive technoanatomy, and the
adverse tail truncates steadily. But the selection is \emph{gradual}
where the strength truncation was abrupt: adverse strength falls by a
factor of $137$ at $T=4$, while $\PP(\rho>0)$ has moved only from
$0.75$ to $0.84$ at the same horizon. Anatomy is the slower and the more
decisive: by Theorem~\ref{gr1:thm:trapped} an adverse $\rho$ is
eventually fatal with certainty, while adverse strength is a condition
every biosphere occupies half of logarithmic time and survives.

\begin{remark}[The coordinates stay separable]\label{gr:rem:separable}
Strength and anatomy cannot fuse under conditioning, and the reason is
structural rather than measured: $\rho$ is a parameter of the model and
$x$ a functional of the path, so no amount of conditioning makes one a
function of the other. They therefore remain distinct diagnostic
questions rather than two aspects of a single ``advancement'' --- which
is exactly what a one-dimensional scale cannot represent.
\end{remark}

% ---------------------------------------------------------------------
\subsection{What Golden reasoning refutes}\label{gr:sub:refutes}

\subsubsection{The Kardashev scale}\label{gr:subsub:kardashev}

The usual complaint against Kardashev's scale \cite{Kar64} is that its
thresholds --- planetary, stellar, galactic energy capture --- are
arbitrary. That complaint is correct but weak, since a scale may be
useful with arbitrary units. Golden reasoning supports three stronger
objections, of which the first is a theorem.

\emph{(i) The terminal stage exists, and is gated by a variable the
scale cannot see.} It is tempting to argue instead that no terminal
stage exists at all, since $\Lambda_\infty=\infty$ almost surely --- but
that holds of the driftless model and not in general, and the stronger
claim is available. By Theorem~\ref{gr1:thm:twocond},
\[
\PP(\text{viable})=\one\{\rho>0\}\cdot
\frac{s(\mu_0)-s(-\infty)}{s(+\infty)-s(-\infty)},
\]
so there \emph{is} an absorbing state of safety --- and reaching it
requires two conditions, of which Kardashev's scale tracks at most one.
The transience factor is the accumulation of capacity, which energy
capture does index, however crudely. The anatomy factor is not on the
scale at all: $\rho$ does not appear in energy capture, is not
monotone in it, and cannot be inferred from it.

The objection is therefore not that the ladder has no top but that
\emph{climbing it is neither necessary nor sufficient for reaching the
top}. Worse, by Remark~\ref{gr1:rem:againstsing}, a biosphere with
$\rho<0$ that climbs fast does not merely fail to arrive: its hazard
diverges superexponentially in the very capacity the scale is
measuring. A scale indexed on capture alone cannot distinguish the
ascent that arrives from the ascent that is a fall. That last point is
the vulnerable-world hypothesis \cite{Bostrom19} restated with a
continuous parameter in place of a discrete draw, and
\S\ref{gr:subsub:xrisk} sets out what is and is not shared.

\emph{(ii) The indexed quantity is not monotone.} Technostrength is
stationary and recurrent in logarithmic time
(Definition~\ref{gr:def:coords}). A ladder presumes a quantity that
ratchets. Energy capture may of course ratchet while technostrength does
not --- but then energy capture is not tracking the quantity that
governs survival, which returns us to (iii).

\emph{(iii) A magnitude index cannot represent categorical selection.}
By Remark~\ref{gr:rem:truncresc}, survivorship truncates technoanatomy
while merely rescaling technostrength. A one-dimensional index of
magnitude has no coordinate in which to express a deleted half-space. It
is not that Kardashev measures the less important variable --- that
claim is false, and was corrected above --- but that it measures the one
whose selection is quantitative, and is structurally unable to express
the one whose selection is categorical.

\begin{remark}[On Dyson spheres]\label{gr:rem:dyson}
Dyson's proposal \cite{Dys60} is often read as the natural
Type~II signature. Under Golden reasoning it is better read as a
\emph{hypothesis about technoanatomy} disguised as one about
technostrength: it asserts not merely that a biosphere could capture
stellar output but that a long-lived one would be organised so as to
want to. That assertion is exactly what \S\ref{gr:sub:extraction}
disputes, and it should be recognised as an assertion rather than an
extrapolation.
\end{remark}

\subsubsection{The violent-aliens thesis}\label{gr:subsub:violence}

The thesis in view holds that intelligence evolves preferentially in
predators, hence that advanced biospheres are likely to be predatory or
violent. We engage the position rather than any proponent.

\begin{quote}\small
``It was the saying of Bion, that though the boys throw stones at frogs
in sport, yet the frogs do not die in sport but in earnest.''
\hfill --- Plutarch, \emph{Moralia} \cite{Plutarch}
\end{quote}

\noindent The line is Bion of Borysthenes', preserved by Plutarch, who
deploys it against hunting for sport; it has since been borrowed by
astrobiological writing as the image of what a technologically superior
biosphere might do to an inferior one. It is the right epigraph for this
subsection because it names the assumption at issue more exactly than
the conquest analogies usually offered. The serious form of the
violent-aliens thesis is not that an advanced biosphere would hate us.
It is that we would be \emph{beneath its consideration} --- that its
power could be spent across the boundary between itself and us
casually, without the expenditure ever rising to the level of a
decision. The stones are thrown in sport. That is the claim, and any
reply pitched at malice misses it.

The image also fixes what the reply must not do. It must not argue that
such a biosphere would be kind, since kindness is not what the thesis
denies; nor that it would have outgrown violence, since sport is not
violence in the relevant sense. What it must do is address the
\emph{scaling}: whether a biosphere can be arranged so that power is
spendable in that manner and still arrive at the scale where the
spending would matter. The author's own route to the present framework
began from exactly this observation, three years before the model
existed, and \S\ref{gr1:sub:anatomy} records the entry.

Before the model is brought to bear, a further placement is owed,
because the \emph{conclusion} here is the one most likely to be mistaken
for a new one and it is not. The conclusion --- that anything old enough to encounter must be
peaceable --- is a standing argument in SETI, sometimes called the
\emph{technological adolescence} argument, and it long predates this
paper: a civilisation reaching the capacity to destroy itself must
either become socially responsible or fail to persist, so whatever
survives has already made that transition. Cooper documents the
argument and its currency, and also supplies the standing rebuttal: that
it conflates altruism with amiability, that it is typically asserted
rather than derived, and that it presumes a trajectory of moral
improvement for which there is no evidence \cite{Cooper19}. Nothing
below claims the conclusion as new. What is at issue is whether a
mechanism can be supplied that survives Cooper's objection.

The inference under dispute has the form: predatory origin
$\Rightarrow$ predatory maturity. Golden reasoning denies the
implication, on the grounds that it neglects the conditioning that must
be applied before any encounter can be contemplated. Any biosphere one could meet is old; any biosphere
that is old is drawn from $\PP^{(t)}$ for large $t$, not from $\PP$; and
$\PP^{(t)}\to Q$, under which technoanatomy is truncated. Whatever
selection operated at origin, a second and stronger selection operates
over the trajectory, and it is the second that determines what survives
to be encountered.

\begin{premise}[Origin does not transfer; \emph{premise}]
\label{gr:phi4}
Facts about the selective regime that produced intelligence in a
biosphere do not, without further argument, transfer to the selective
regime that governs that biosphere's persistence at world scale. The two
regimes act on different units, at different scales, over different
timescales.
\end{premise}

$\Phi4$ is not exotic; it is the ordinary observation that a trait
adaptive for an organism in a population need not be adaptive for a
biosphere in a universe. What the model adds is that the second regime
is not merely different but \emph{stronger} in the relevant sense: it
deletes a half-space of configurations, and it does so with a force
measured in \S\ref{gr:subsub:measured}.

What the model contributes to the frogs case in particular is a reply
that never mentions disposition. A biosphere that can spend power
casually across a self/other boundary is one that maintains such a
boundary and enforces it --- and by \S\ref{gr:subsub:boundary} the
enforcement of a boundary is interface-like while the domain it encloses
is volume-like, so the arrangement carries $\rho<0$ by
\eqref{gr:eq:rhoextract}. By Corollary~\ref{gr:cor:extractcost} its
viability probability is then exactly zero. The stone-throwers are not
ruled out because they would have grown kind; they are ruled out
because the organisation that permits stone-throwing at that scale is
the organisation the scaling argument says cannot reach that scale. The
frogs survive not by the boys' mercy but by the boys' having imploded
before attaining the requisite size --- which is, in a form Bion would
have recognised, the square--cube law again.

Three further features distinguish this from the folk argument, and they
are exactly the three points Cooper's rebuttal presses on. \emph{It is not
a moral claim.} What is selected is $\rho$, the sign of
$-\partial_x\log h^{\mathrm{int}}$ --- a structural property of how a
biosphere is arranged, carrying no implication that its members are
kind, and applying identically to a biosphere with no members in the
relevant sense at all. The conflation of altruism with amiability has
nothing to attach to. \emph{It is not asserted but measured.}
\S\ref{gr:subsub:anatsel} exhibits the selection on $\rho$ and finds it
\emph{gradual} --- $0.75$ to $0.91$ over the horizons computed ---
where the folk argument implicitly treats the filtering as complete.
\emph{And it is not a trajectory.} Nothing here says biospheres improve;
$\rho$ does not change at all in the fixed model, and what conditioning
does is reweight an ensemble, not move a system along a path. The
selection is Darwinian in form rather than developmental, which is
what the folk version gets wrong and Cooper is right to object to.

This much, however, only establishes that predatory maturity is not
\emph{inherited}. To conclude that it is positively \emph{selected
against}, one needs the identification of predation with negative
technoanatomy. That identification is a philosophical premise, and it is
the subject of the next subsection. Everything in
\S\ref{gr:sub:consequences} is downstream of it.

\subsubsection{Trajectory arguments for internal neglect}
\label{gr:subsub:neglect}

A third target belongs here, and it lies inside the biosphere rather
than out among the aliens. Longtermism in its strong form holds that
what most matters about our actions is their effect on the very long
run \cite{GM21}; and a recurring corollary within that literature is
that resources directed at the welfare of some populations do less for
the long-run future than the same resources directed elsewhere, so that
broad investment in the worse-off is not what a future-oriented agent
should fund. We engage the corollary as a position and name no one, as
with the violent-aliens thesis.

The position is already under sustained criticism, and it is worth being
exact about how the present objection differs, because the difference is
what makes it worth adding. The standing critiques are \emph{external}:
they fault longtermism on ethical and political grounds --- as an
ideology with damaging real-world effects and a flawed underlying moral
theory \cite{Crary23}, and as a doctrine whose cosmic accounting
licenses the discounting of present suffering \cite{Torres21}. Those
objections deny the framework's premises. The objection developed here
does the opposite: it \emph{grants} the long-run criterion entirely, and
shows the position failing by that criterion, on the framework's own
terms and in its own currency. If it succeeds it is therefore not an
alternative to the ethical critiques but a complement of a different
kind --- an internal result rather than an external one.

Stated at its strongest it is not a callous argument but an
optimisation: interventions differ in long-run leverage, a
future-oriented agent should allocate to high-leverage interventions,
and if the leverage of some interventions is low then funding them is a
loss measured against the future. Nothing in that reasoning is invalid.
The question is whether the leverage calculation is complete.

Golden reasoning says it is not, and the missing term is technoanatomy.
No new premise is required, and this is the point worth making: a
position of this kind is $\Phi3$ with the boundary drawn \emph{inside}.
Writing off part of a system is the maintenance of a self/resource
distinction across an internal interface --- a region of the biosphere
to which protection is not extended --- which is structurally the same
object as the extractive boundary of \S\ref{gr:subsub:boundary}, and it
inherits the same scaling. Enforcement of an internal boundary is
interface-like; the interior it excludes is volume-like; so by
\eqref{gr:eq:rhoextract} the arrangement carries $\rho<0$, and
Corollary~\ref{gr:cor:extractcost} applies verbatim.

The consequence is that the position is self-undermining in its own
currency, and sharply rather than rhetorically. By
Corollary~\ref{gr:cor:extractcost}(ii) a biosphere with $\rho<0$ has
viability probability \emph{exactly zero} however much capacity it
acquires; by (i) its life expectancy is finite where a
$\rho>0$ biosphere's is infinite; and by
Remark~\ref{gr1:rem:againstsing} acquiring capacity faster makes its
position worse, since $x\to\infty$ with $\rho<0$ drives hazard to
infinity superexponentially. An argument that trades internal cohesion
for external capacity, offered in the name of the long-run future, is
therefore an argument for a shortened one --- not because neglect is
distasteful, which the model cannot see, but because the boundary it
requires is the structure the model identifies as fatal.

The thought experiment that makes this concrete is worth stating,
because it isolates the variable. Hold a society's technoanatomy fixed
at whatever value its present arrangements realise, and scale its
technostrength to galactic magnitude. If $\rho\ge0$, nothing goes
wrong. If $\rho<0$, the hazard does not rise gently but spikes as
$e^{|\rho|x}$: at every scale the same arrangement supplies more
micro-opportunities for the system to act against itself, and each
opens further ones. The elephant-sized ant of
Remark~\ref{gr:rem:namesclosure} is exactly this, and the square--cube
law is exactly why.

Four fences, because this is the section's most contestable
application. It refutes a \emph{class} of positions --- any that trades
internal cohesion for external capacity --- and not any author's wider
view, nor longtermism as such: a longtermism that valued internal
cohesion would be untouched, and the argument is in that sense a
constraint on longtermist reasoning rather than a refutation of it. It is downstream of $\Phi3$ and inherits its conjecture grade;
reject $\Phi3$ and this goes with it. And it establishes nothing about
what anyone ought to do, by $\Phi2$: what it establishes is that a
particular argument for neglect, evaluated by its own long-run
criterion, fails on that criterion.

% ---------------------------------------------------------------------
\subsection{A prediction outside cosmology}\label{gr:sub:crossdomain}

Most of what this section has established is a re-derivation. The
conclusions of \S\ref{gr:sub:refutes} were available before it --- the
Kardashev scale has long been thought crude, the peaceability of
long-lived civilisations is a standing SETI argument
\cite{Cooper19}, and the mechanism behind $\Phi3$ has precedents in
the collapse and existential-risk literatures
\cite{Tainter88,Bostrom19}. What the section offers is that they follow
from one premise with a named exponent, which is a unification rather
than a discovery.

Unification is a respectable contribution, but it carries an obligation
that a discovery does not. A framework which only reproduces what it was
built to reproduce is a reformulation, and the way to tell the two apart
is whether anything falls out that was not among the inputs. This
subsection states the clearest such consequence, in a domain the
framework was not designed for and where it can be checked.

\subsubsection{The claim}\label{gr:subsub:crossclaim}

The asymmetry measured in \S\ref{gr:subsub:measured} --- adverse
strength truncated abruptly, technoanatomy selected gradually --- is not
a fact about biospheres. Nothing in its derivation mentions one. It is a
consequence of conditioning a two-coordinate system on survival, where
one coordinate is a functional of the path and the other a parameter of
the model, and it should therefore appear wherever that structure does.

\begin{quote}
\emph{Cross-domain prediction.} In any population subject to
survivorship selection, carrying a fluctuating state coordinate
alongside a fixed structural parameter that governs whether the state's
growth raises or lowers hazard, conditioning on survival should
\textbf{truncate} the state coordinate --- deleting a region it can
otherwise re-enter --- and \textbf{shift} the structural parameter,
progressively and without deleting anything, the two effects being
comparable in magnitude and opposite in kind.
\end{quote}

Candidate domains are not exotic: firms with a fluctuating balance-sheet
position and a fixed governance structure; institutions with a
fluctuating resource base and a fixed constitutional arrangement;
lineages with a fluctuating population size and a fixed life-history
strategy. In each the two coordinates are separately measurable, old
members are selected from many, and the comparison the prediction asks
for is between surviving and full cohorts rather than between a model
and the sky.

\subsubsection{What would count against it}
\label{gr:subsub:crossfalsify}

Three outcomes would tell against the framework, and they are worth
stating plainly, since a prediction whose failure conditions are vague
is not doing the work this subsection is meant to do.

If the selection acted \emph{categorically on the structural parameter
and quantitatively on the state coordinate} --- the asymmetry inverted
--- the account of why the two behave differently would be wrong, since
that account turns on which of them can be re-entered.

If \emph{neither} coordinate showed selection in a population where
survivorship is known to operate, the framework would be describing a
mechanism that does not act at accessible scales, whatever its formal
standing.

And if the structural parameter showed selection \emph{as abrupt} as
the state coordinate, the distinction drawn in
\S\ref{gr:sub:coords} between a state and a structure would not be doing
the work claimed for it, and the argument of
\S\ref{gr:subsub:violence} --- which rests on that distinction to answer
the charge of positing a moral trajectory --- would lose its reply.

\subsubsection{Two honest qualifications}
\label{gr:subsub:crossqual}

The prediction is weaker than it may appear, in two respects that should
be recorded rather than left for a reader to find.

It is a claim about \emph{shape}, not magnitude. The framework does not
predict how large either effect should be in any particular domain,
since that depends on the hazard scale, the horizon, and the prior
spread of the structural parameter --- none of which the framework
supplies. What it predicts is that one effect is a truncation and the
other a shift, which is a qualitative signature and correspondingly
easier to satisfy by accident.

And it is not \emph{unique} to this framework. Any account on which
survival depends on a structural property will predict some selection on
that property; what is distinctive here is the predicted \emph{contrast}
in kind between the two coordinates, and only the contrast is
diagnostic. A domain exhibiting selection on both, with no difference in
character, would be consistent with much and would confirm little.

With those qualifications the prediction stands as the section's answer
to the question of whether its unification is more than a restatement.
We do not know the answer; the point of stating the prediction in this
form is that someone could find out.

% ---------------------------------------------------------------------
\subsection{The Extraction Hypothesis}\label{gr:sub:extraction}

\begin{premise}[Extraction Hypothesis; \emph{conjecture-grade}]
\label{gr:phi3}
Organisation built on unbounded extraction --- consuming the resources
of a locale and moving on --- has \emph{negative technoanatomy}:
$\rho<0$ in the sense of Definition~\ref{gr1:def:anatomy}, so that
internal hazard rises with acquired capacity rather than falling with
it.
\end{premise}

The premise is now a claim about the sign of an exponent rather than an
identification between a social form and a vocabulary, and this is worth
pausing on, because it is what the anatomy parameter buys. Before, one
had to be told what ``bad technoanatomy'' meant and then persuaded that
extraction was an instance. Now $\Phi3$ says something with a truth
value fixed by the model: that for extractive organisation,
$-\partial_x\log h^{\mathrm{int}}<0$. The consequences follow
mechanically rather than interpretively, and they are severe.

\begin{corollary}[What $\Phi3$ costs an extractive biosphere; \textbf{P}
given $\Phi3$]\label{gr:cor:extractcost}
If extractive organisation has $\rho<0$ then, by
Theorem~\ref{gr1:thm:trapped} and Theorem~\ref{gr1:thm:twocond}:
\emph{(i)} its life expectancy is finite, with
$\PP(T>t)\le e^{-h_\ast t}$ for the floor $h_\ast$ of
\eqref{gr1:eq:xstar};
\emph{(ii)} its viability probability is \emph{exactly zero}, whatever
feedback it achieves --- the anatomy factor in \eqref{gr1:eq:twocond}
annihilates the transience factor;
\emph{(iii)} it inhabits the trapped regime of
\S\ref{gr:subsub:regimes}, where none of the Lindy structure holds and
the survivor ensemble concentrates rather than diversifying;
and \emph{(iv)} its safest available configuration is not its strongest
but its \emph{weakest}, since $h\to B(1+e^{-b})$ as $x\to-\infty$ while
$h\to\infty$ as $x\to+\infty$.
\end{corollary}

\noindent Item (ii) is the one to hold onto. An extractive biosphere
does not merely have poor prospects; it has \emph{no} prospect of
viability, and acquiring capacity faster makes its position worse rather
than better. Item (iv) states the same thing from the other end and is
the sharpest form of the thesis: for a system that turns its strength on
itself, every increment of strength beyond $x_\ast$ is a net loss, and
the counsel of the model is to stop.

This is the section's single load-bearing philosophical premise. It is
not a theorem and is not presented as one. We argue for it.

\subsubsection{The argument from boundary cost}\label{gr:subsub:boundary}

Extraction requires a maintained distinction between the extracting
system and the extracted resource: a boundary across which value flows
one way. The boundary is not free. It must be enforced --- against
diversion, against defection, against the tendency of a system's own
subunits to treat the system as their resource, which is precisely what
a tumour does.

Now consider how the two sides of this ledger scale with technostrength,
and note that the anatomy parameter lets the scaling be written down
rather than gestured at. Let capacity be indexed by $x$. The domain over
which extraction occurs grows with reach and is volume-like,
$V\sim e^{\alpha x}$; the boundary across which enforcement is exercised
is interface-like, $S\sim e^{\beta x}$. The claim is
$\beta<\alpha$: volume outgrows surface. The unenforced fraction of the
domain is then $1-e^{-(\alpha-\beta)x}\to1$, so the unprotected
extractive interior grows as $e^{\alpha x}$, and internal hazard with
it:
\begin{equation}\label{gr:eq:rhoextract}
h^{\mathrm{int}}\;\sim\;e^{\alpha x}
\qquad\Longrightarrow\qquad
\rho\;=\;-\alpha\;<\;0 .
\end{equation}
So the argument does not merely conclude that extraction is bad
anatomy; it \emph{names the exponent}. Technoanatomy is minus the
volumetric growth rate of the extractive domain, and the more reach an
extractive organisation acquires per unit of capacity, the worse its
anatomy is. Hence:

\begin{quote}
\emph{As a biosphere's technostrength grows, the cost of maintaining the
self/resource distinction grows faster than the capacity to maintain
it, and the shortfall is the exponent $\alpha-\beta>0$.}
\end{quote}

\noindent This mechanism is not original to the present work, and
\S\ref{gr:subsub:xrisk} records its ancestry: Tainter's account of
collapse from the declining marginal returns of complexity
\cite{Tainter88} reaches the same interior optimum from archaeological
evidence, and Bostrom's vulnerable-world hypothesis \cite{Bostrom19}
reaches a compatible conclusion from technological risk. What the
present treatment adds is the exponent \eqref{gr:eq:rhoextract} and the
hazard-theoretic consequences that follow from its sign.

Past some scale, the cheapest available target of an extractive
subsystem is the biosphere itself --- not through malice but through
gradient, exactly as a tumour exploits the body's own supply mechanisms
rather than building its own. Extraction at scale therefore turns
inward, and in the model's vocabulary this is precisely the
\emph{trough} channel of \S\ref{gr1:sec}: internal damage, unmediated by the
environment, unbounded above, with brittleness exponent $\rho>1$
ensuring that its costs outrun the crest channel's gains per unit.

\begin{remark}[The names were the mechanism]\label{gr:rem:namesclosure}
The argument just given is the square--cube law. That is the law from
which ``technostrength'' and ``technoanatomy'' were named --- initially
as a joke about giants. The joke turns out to be the mechanism: a
biosphere whose extractive reach has outgrown its capacity to remain one
system is an elephant-sized ant, and collapses for the same reason. The
model's brittleness exponent $\rho$ is the quantitative form of ``how
badly does the mismatch hurt''.
\end{remark}

\subsubsection{Objections}\label{gr:subsub:objections}

\begin{objection}[Bounded extraction]\label{gr:obj:bounded}
A biosphere might extract only from what is genuinely external and never
grow past the scale at which boundary cost bites.
\end{objection}

\noindent This is granted, and it is not a counterexample but a
concession of the point. A biosphere that halts its extractive growth at
a sustainable scale is not the unboundedly expanding entity whose
possibility is at issue; it is not a Kardashev trajectory. $\Phi3$
concerns \emph{unbounded} extraction, and the objection concedes that
unbounded extraction is unavailable.

\begin{objection}[Cheap boundaries]\label{gr:obj:cheap}
Sufficiently advanced technique might make boundary enforcement scale
with volume rather than surface.
\end{objection}

\noindent This is the serious objection and we do not refute it. In the
notation of \eqref{gr:eq:rhoextract} it is the claim $\beta\ge\alpha$:
enforcement scaling at least volumetrically, so that $\rho\ge0$ and the
whole corollary lapses. That is an empirical claim about the scaling of
coordination costs, it is where $\Phi3$ is falsifiable, and stating it
as an inequality between two exponents is the chief gain of the anatomy
parameter --- the objection and the premise now disagree about a
measurable number rather than about a description. We note only that
every known instance of large-scale organisation --- organisms,
ecosystems, institutions --- exhibits $\beta<\alpha$, and that the
burden of the exception lies with whoever claims it.

\begin{objection}[Assuming the conclusion]\label{gr:obj:circular}
Defining bad technoanatomy as ``that which raises hazard'' and then
asserting extraction raises hazard makes the argument circular.
\end{objection}

\noindent The charge is now easy to answer precisely, and the anatomy
parameter is what makes it so. The model defines $\rho$ as
$-\partial_x\log h^{\mathrm{int}}$ and proves what follows from its
sign; that is a theorem about an exponent and mentions extraction
nowhere. $\Phi3$ claims that a particular form of organisation realises
a negative value of that exponent, and argues for it from boundary
scaling. Two distinct objects, one proved and one argued, and neither
borrowed from the other. The identification is supplied from outside,
by \S\ref{gr:subsub:boundary}, and could be false --- it would be false
if $\beta\ge\alpha$. The division of labour is exactly the one this section's method
requires: mathematics limits the possibility space, philosophy says
which region of it we occupy.

\begin{remark}[What would falsify $\Phi3$, and what would not]
\label{gr:rem:falsify}
Three falsifiers, and they are \emph{operational} in the companion's
sense --- checkable in domains that exist, rather than requiring a
capability the argument itself says no one has \cite{Ben26Optimal}.
$\Phi3$ fails if $\beta\ge\alpha$ in \eqref{gr:eq:rhoextract} --- if
enforcement is shown to scale at least volumetrically in reach; if an
extractive organisation at scale is exhibited stable over timescales
long compared with its own doubling time; or if internal hazard is
observed to \emph{decrease} with technostrength in an extractive system,
which is $\rho>0$ measured directly and refutes the premise outright. The first
is the live question, and it is live in institutional, ecological and
organismal data rather than only in cosmology.

What would \emph{not} falsify it: the bare possibility that extraction
might turn out fine, or a hypothetical technology stipulated to make
boundaries cheap. The claim is about scaling, scaling is measurable
wherever large-scale organisation exists, and a falsifier must therefore
exhibit the scaling rather than assert its negation.
\end{remark}

% ---------------------------------------------------------------------
\subsection{Consequences}\label{gr:sub:consequences}

The following are downstream of $\Phi3$ and inherit its status. Each is
falsifiable in principle; none is currently testable.

\emph{(a) Encountered biospheres are not extractive.} Not because
predation is impossible but because extractive organisation is a
transient: it may be instantiated, but the conditioning under which any
encounter occurs selects against it. A predatory biosphere
``copy-pasted'' into the galaxy would be a configuration of the
$\PP$-ensemble, not the $Q$-ensemble, and would not persist to be met.

\emph{(b) Absence of megastructures is not evidence of absence.} If
$\Phi3$ holds, stellar-scale extraction is a signature of biospheres on
their way out, not of biospheres that lasted. Searches keyed to such
signatures \cite{Dys60} are therefore keyed to the wrong tail of the
distribution, and their null results carry less evidential weight
against the existence of old biospheres than is usually assumed.

\emph{(c) Old biospheres are Lindy --- and here the novelty is narrow,
so it is stated narrowly.} From the survival tail
$\PP(T>t)\asymp Ct^{-1/4}$ (\S\ref{gr1:sec}): residual life is
asymptotically Pareto$(\tfrac14)$ in units of attained age, and the
effective population hazard is $1/4t$. The longer a biosphere has
persisted, the longer it may be expected to persist, in exact
proportion --- while still departing with probability one.

That conclusion is \emph{not} new, and it would be a misrepresentation
to present it as such. Lindy longevity for technological civilizations
is an active subfield: Balbi and \'Cirkovi\'c model the termination rate
as a Weibull hazard and note that a decreasing rate yields a
Pareto-like fat tail of exactly Lindy character \cite{BC21}; Balbi and
Grimaldi investigate Lindy's law for technosignature longevity directly,
and show it can overturn the expectation that a first detection be
long-lived \cite{BG24}; Kipping and collaborators reach related
conclusions about the age of a first contact from lifetime distributions
\cite{Kip20}; and Ord surveys the mechanisms from which a Lindy effect
can arise at all, including one that conjures it from an uninformative
prior \cite{Ord23,Tal12}.

What this section adds to that literature is one thing only. In each of
those treatments the power law is \emph{parameterized}: the exponent is
a free shape parameter, chosen for tractability or fitted, and the
question studied is what follows from a given value of it. Here the
exponent is an \emph{output} --- $\tfrac14$, and it would have been
$\tfrac12$ had the noise entered the log-hazard rather than its
derivative (\S\ref{gr1:sec}). The claim is therefore not that biospheres are
Lindy, which is others', but that the Lindy exponent is fixed by the
smoothness of whatever drives the hazard, and is in that sense a
measurable structural fact rather than a modeling choice. Whether that
addition is worth much depends entirely on whether the driver's
smoothness is ever ascertainable, which this section does not claim.

\emph{(d) The age distribution of extant biospheres is heavy-tailed.}
Quantiles scale as the fourth power of rarity. Any sample of biospheres
found at a given epoch is dominated by the old, and moments of order
$\ge\tfrac14$ diverge: the sample mean age of such a population is an
artifact of the observation window, growing as $\mathcal T^{3/4}$.

\subsubsection{Relation to anthropic reasoning}\label{gr:subsub:anthropic}

Golden reasoning stands in an unusual relation to the observation
selection literature \cite{Car83,Got93,Bos02}, and the relation is worth
stating precisely because it is the section's clearest methodological
contribution.

Gott's Copernican argument \cite{Got93} derives a predictive
distribution for a system's remaining duration from the premise that the
observation time is uniformly distributed over the system's lifetime.
The Doomsday argument \cite{Car83,Bos02} proceeds similarly from a
reference class and an indifference principle. In both, the age
distribution is an \emph{input}, justified by a principle of
indifference.

In Golden reasoning the age distribution is an \emph{output}. The tail
$t^{-1/4}$ is not postulated; it is the persistence exponent of
integrated Brownian motion, which follows from where the noise is placed
in the hazard (\S\ref{gr1:sec}) and nothing else. The reference class is not
chosen; it is the conditioned ensemble $\PP^{(t)}$, which the model
constructs. The exponent could have been $\tfrac12$ --- it is, for a
mean-reverting driver --- and which it is depends on the smoothness of
the driver rather than on any epistemic principle.

The same contrast holds, and holds more sharply, against the
technosignature-longevity literature, which is the nearer neighbour and
is already model-based rather than indifference-based
\cite{BC21,BG24,Kip20}. Those treatments carry a hazard function with a
free shape parameter --- Weibull $k$, or a Pareto exponent $\alpha$ ---
and study the consequences of its value. The disagreement is therefore
not about method but about where a parameter comes from: there it is
chosen, here it is forced by the driver. Ord's survey is the useful
corrective in both directions \cite{Ord23}, since it shows how many
distinct mechanisms deliver a Lindy effect --- which is a warning that
observing Lindy behaviour would badly underdetermine the mechanism
producing it, and hence that the exponent, not the effect, is where any
inferential weight can sit.

This does not vindicate the Copernican arguments; it relocates the
disagreement. Where they require a defence of indifference, Golden
reasoning requires a defence of the model. That is a better place for
the burden to sit, because a model can be wrong in identifiable ways.

\begin{premise}[Rarity of origins; \emph{imported result}]
\label{gr:phi5}
Biospheres are rare: the difficulty lies in origination, not in
persistence past origination.
\end{premise}

$\Phi5$ is the rare-earth position \cite{WB00}, and it is not imported
from the general literature but from the author's companion program,
where it is derived rather than assumed: the Nexic framework
\cite{Nexus} argues that the astrobiological Copernican principle
is inconsistent with the mediocrity commitments that were supposed to
support it \cite{Ben26Found}, and re-reads the natural-historical record
accordingly \cite{Ben26Optimal}. Its interaction with the rest of this
section is worth making explicit, because the combination is more
coherent than either part alone. If origins are rare and survival is
Lindy, then the observed silence requires no filter ahead of us: it is
explained by the scarcity of starts, not by the destruction of the
started. Golden reasoning and rare-earth are complementary --- the first
says that what persists is not extractive and therefore not loud, the
second says there is little of it. Neither requires the Great Filter to
lie in our future \cite{Han98,SDO18}.

\subsubsection{Relation to the Nexic framework}
\label{gr:subsub:nexic}

The relation to that companion program is closer than a shared
conclusion, and stating it precisely is worth a subsection, because the
two halves turn out to partition the same problem.

\emph{The same move, in opposite temporal directions.} The Accordance
Principle constrains a realized natural-historical attribute $A$ against
its complement $A'$, with $H$ the production of observers, by
$P(H\mid A')/P(H\mid A)\lesssim P(A)/P(A')$: realized improbability must
be compensated by observer-relevance \cite{Ben26Found}. Golden reasoning
performs the same reweighting on a trajectory rather than an attribute,
the Doob density $\hm(X_s)/\hm(X_0)$ tilting the path measure by the
probability of the conditioning event downstream. Both are Bayes applied
to a selection event; they differ in what is conditioned on and in which
direction time is read. Accordance conditions on observers \emph{having
arisen} and reads backward over the record; Golden reasoning conditions
on survival \emph{to a future horizon} and reads forward over the
trajectory. The two are complementary rather than redundant:
retrodictive and predictive conditioning on the same kind of event.

\emph{The companion names the degenerate case.} This is not a
resemblance noticed from outside. The foundation paper itself identifies
survivorship reasoning as the coarsest limiting case of its own
constraint: restrict the state space to a single history's population
and degrade the graded observer-condition $H$ to a binary survival event
$S$, and the identity behind the principle reads
$P(A\mid S)/P(A'\mid S) = \bigl(P(S\mid A)/P(S\mid A')\bigr)\cdot
\bigl(P(A)/P(A')\bigr)$ --- the textbook survivorship correction
\cite{Ben26Found}. What the present section supplies is that limit
developed \emph{dynamically}: not a correction applied at one attribute
but a measure on paths, with a limit object $Q$, an exponent
$\tfrac14$, and an interpolation between the two.

\emph{Golden reasoning supplies what the companion declares itself
silent about.} The Optimal-Earth paper states as a principled limitation
that its framework is structurally silent on filters ahead: its
constraint conditions on observer-production, which our own case has
already secured, so no bound it issues bears on what becomes of
observers afterward \cite{Ben26Optimal}. That is precisely the gap this
section fills, because conditioning on $\{T>t\}$ is conditioning on a
\emph{future} event. The absence of a terminal stage, the Lindy residual
life, the $t^{-1/4}$ tail and the goldening interpolation are all
statements about what becomes of observers afterward. The two wings
therefore partition the timeline at the observer: Accordance behind,
Golden reasoning ahead.

\emph{One identification, flagged rather than assumed.} The connection
is a shared logical form plus one substantive identification, and the
identification is not free. The companion's cosmic weighting is a
primitive measure over partial information states; the $\PP$ of
\S\ref{gr1:sec} is a path measure induced by a
specific model. Reading the two as one object requires taking the
model's trajectory space as a family of partial states and its $\PP$ as
the cosmic weighting restricted to them --- a move neither paper
performs. Absent it, what is shared is the logical form, which is worth
stating and is not an entailment. Nothing in this section's results
depends on the identification being made.

\emph{An apparent clash, dissolved.} The Optimal-Earth Conjecture
\cite{Ben26Optimal} holds our record to be observer-optimal on
effectively any attribute one can interrogate, while \S\ref{gr:subsub:practical} holds
that we are not practically Golden to any interesting tolerance. These
do not conflict, and the reason is the temporal split just drawn: the
first conditions on observer-production and grades the realized past,
the second conditions on survival to a future horizon and measures
distance from $Q$. A record may be optimal in every respect Accordance
interrogates while its bearer sits at small $t$ on the goldening
interpolation. Optimality behind is not Goldenness ahead.

\subsubsection{Relation to the existential-risk and collapse
literatures}\label{gr:subsub:xrisk}

Two established positions occupy ground adjacent to $\Phi3$ and the
two-condition theorem, and both were found after those results were
written. They are recorded here because they narrow what is claimed as
new --- and because, on inspection, they support the premise rather than
threatening the paper.

\emph{Collapse from declining returns.} Tainter's account of societal
collapse \cite{Tainter88} holds that complexity is adopted as a
problem-solving strategy, that its marginal returns decline, and that
past a threshold further complexity is pure cost, at which point
collapse becomes likely. Structurally this is the U-shape of
Proposition~\ref{gr1:prop:trichotomy}(iii) arrived at from archaeology
thirty-six years earlier: a benefit-minus-cost curve in capacity with an
interior optimum, beyond which acquisition is net harmful. The
correspondence is close enough that the boundary-cost argument of
\S\ref{gr:subsub:boundary} should be read as a re-derivation of
Tainter's mechanism in hazard-theoretic terms rather than as an
independent discovery of it.

Three differences remain, and they are what this section adds. Tainter's
declining quantity is \emph{economic and energetic return}; ours is
\emph{hazard}, so the two curves are about different things even where
their shapes agree. Tainter's threshold is qualitative; ours is
$x_\ast$ of \eqref{gr1:eq:xstar}, in closed form and with a computed
floor $h_\ast$. And Tainter's account contains no analogue of the
transience factor, so it cannot state the two-sided result of
Theorem~\ref{gr1:thm:twocond} --- that capacity growth is simultaneously
\emph{necessary} for viability and, absent anatomy, fatal.

\emph{The vulnerable world.} Bostrom's hypothesis \cite{Bostrom19} holds
that continued technological development will at some point attain
capabilities making civilizational devastation extremely likely unless
civilization exits what he calls the semi-anarchic default condition,
with stabilization requiring greatly amplified capacities for preventive
policing and global governance. The mapping into the present vocabulary
is direct and should be stated without hedging: \emph{the semi-anarchic
default condition is low technoanatomy, and Bostrom's stabilization
conditions are the raising of $\rho$}. Quantitative development of the
hypothesis has begun, using survival-function methods on the urn model
\cite{EF24}, which is the same class of object as
\S\ref{gr1:sub:verdict}.

What survives the comparison is narrower than the anti-acceleration
result of Remark~\ref{gr1:rem:againstsing} might suggest on first
reading, and is stated narrowly. The vulnerable-world hypothesis is an
\emph{existence claim about a discrete draw} --- that the urn contains a
black ball --- whereas $\rho$ is a \emph{continuous scaling relation}
with an exponent derived from boundary geometry. And the vulnerable-world
hypothesis is one-sided: development is dangerous. The two-condition
theorem is two-sided, requiring the same quantity that it makes lethal,
and the product structure of \eqref{gr1:eq:twocond} --- in which one
factor annihilates the other --- has no counterpart there.

\emph{Why this strengthens rather than weakens the section.} $\Phi3$ is
conjecture-grade and argued from a single line of reasoning. It now has
two independent precedents reaching a compatible conclusion from
unrelated evidence: one archaeological and economic, one
technological and policy-facing. A premise supported by convergence
across literatures that share no premises is in better standing than one
argued from scratch, and the ledger's grading of $\Phi3$ should be read
in that light --- still conjecture-grade, since none of the three
derivations is a proof, but no longer solitary.

% ---------------------------------------------------------------------
\subsection{The epistemic ceiling}\label{gr:sub:ceiling}

We close with a limitation which constrains this section's own
conclusions at least as much as its targets'.

\begin{theorem}[No finite observation certifies Goldenness]
\label{gr:thm:ceiling}
$Q$ is locally absolutely continuous with respect to $\PP$: on every
finite window the two measures are mutually absolutely continuous on the
survival event, with the explicit bounded density of \S\ref{gr1:sec}. They
become singular only on the infinite-horizon $\sigma$-algebra.
Consequently no observation of bounded duration can establish that a
biosphere is Golden rather than General, and no bounded evidence can
place a system on the interpolation $\PP^{(t)}$ at one $t$ rather than
another beyond the resolution the bias field permits.
\end{theorem}

\begin{remark}[What this costs each party]\label{gr:rem:evenhanded}
Against the violent-aliens thesis: no finite observation of a biosphere,
however extensive, could establish that it is predatory in the
asymptotically relevant sense, so the thesis cannot be confirmed by
encounter.

Against \emph{this section}: no finite observation could establish that
an encountered biosphere is \emph{ideally} Golden either. The claims of
\S\ref{gr:sub:consequences} are claims about a conditional
distribution, not predictions about any individual system, and they must
not be read as the latter.

The ceiling is, however, specific to the ideal. Practical Goldenness
(Definition~\ref{gr:def:practical}) is a finite-window property, and
Theorem~\ref{gr:thm:ceiling} does not forbid establishing it: mutual
singularity is an infinite-horizon phenomenon, and on $\mathcal F_s$ the
two laws are equivalent with bounded density. What obstructs the
practical version is not structure but sample size --- one would need an
ensemble of biospheres, and we have one. That is a contingent barrier
rather than a theorem, and the two kinds of limitation are worth keeping
apart: the ideal category is unknowable in principle, the practical one
merely out of reach.

This symmetry is not a weakness of the position but a consequence of the
structure that the position is built on, and we would regard its absence
as a warning sign. An argument from survivorship that permitted
confident identification of survivors would have proved too much.
\end{remark}

% ---------------------------------------------------------------------
\subsection{What the section costs and what it buys}
\label{gr:sub:ledger}

The apparatus is priced here in the companion's format
\cite{Ben26Optimal}: costs enumerated with grade and falsifier, buys
enumerated after, and the whole kept short enough to audit.

\paragraph{Costs.} Six load-bearing commitments, no more, and their
grades are not uniform.

\begin{enumerate}[itemsep=2pt]
\item \emph{The biospheric referent} ($\Phi1$) ---
\textbf{methodological stipulation}. \emph{Falsifier:} exhibit a unit of
analysis referring to the same object at every point of the trajectory
while carrying less semantic content than ``biosphere''. The stipulation
is that none exists; one counterexample retires it, and the theorems
would then be applied in that unit instead.

\item \emph{Vantage, not obligation} ($\Phi2$) --- \textbf{scope
declaration}, and a self-imposed restriction rather than an assumption
doing work. \emph{Falsifier:} exhibit a normative conclusion derivable
from the apparatus without adding a value premise. Its failure would be
an expansion of the section's reach, not a defeat --- but it would
defeat the claim that Golden reasoning is not longtermism, which is why
it is priced here rather than assumed.

\item \emph{The Extraction Hypothesis} ($\Phi3$) ---
\textbf{conjecture-grade}, argued from boundary cost but not derived,
and the section's principal exposure --- though no longer a solitary
one: \S\ref{gr:subsub:xrisk} records two independent precedents
\cite{Tainter88,Bostrom19} reaching compatible conclusions from
unrelated evidence, which is convergence rather than proof but improves
the premise's standing. In its present form it is the
claim $\rho<0$ for extractive organisation
(Definition~\ref{gr1:def:anatomy}), which is why
Corollary~\ref{gr:cor:extractcost} follows from it mechanically and why
\S\ref{gr:subsub:neglect} needs no further premise. Its position in the argument is
exact and worth stating in the ledger: it is not a premise of any
theorem, but the \emph{identification} that connects the theorems to
their application, so its failure does not falsify anything in
\S\ref{gr1:sec} --- it \emph{empties}
\S\ref{gr:sub:consequences}, leaving the selection results intact and
silent about predation. \emph{Falsifiers:} the three of
Remark~\ref{gr:rem:falsify}, all operational.

\item \emph{Origin does not transfer} ($\Phi4$) --- \textbf{premise}.
\emph{Falsifier:} a mechanism by which the selective regime that
produced intelligence in a lineage transfers to the regime governing a
biosphere's persistence at world scale. Note this is the weaker of the
two anti-violence commitments and can stand alone: $\Phi4$ alone
establishes that predatory maturity is not \emph{inherited}, and only
$\Phi3$ adds that it is selected \emph{against}.

\item \emph{Rarity of origins} ($\Phi5$) --- \textbf{imported result},
carried at the companion's grade and not re-argued
\cite{Nexus, Ben26Found, Ben26Optimal}. \emph{Falsifier:} the
companion's, and nothing this section adds.

\item \emph{The inherited mathematical apparatus} --- \textbf{[A]} and
below. Every quantitative claim here --- the $t^{-1/4}$ tail, the Lindy
residual life, the $s/4\varepsilon$ criterion, the measured truncation
--- inherits the tags of \S\ref{gr1:sec}, in
particular the soft-wall Ansatz. \emph{Falsifier:} the numerical tests
specified there, chiefly a fitted exponent robustly away from
$-\tfrac14$.
\end{enumerate}

\noindent One further entry belongs in this column, though it is not a
premise: \emph{the section's novelty is narrower than its apparatus}.
The mathematics of \S\ref{gr1:sec} is classical or routine --- the
$\tfrac14$ exponent is McKean's, the background construction is an
application of known results on sums of random exponentials --- the Lindy
conclusion of \S\ref{gr:sub:consequences} is established prior work
\cite{BC21,BG24,Kip20,Ord23}; and the mechanism behind $\Phi3$ has
precedents in the collapse and existential-risk literatures
\cite{Tainter88,Bostrom19,EF24}; and the conclusion of
\S\ref{gr:subsub:violence} is the standing technological-adolescence
argument of SETI, already current and already rebutted
\cite{Cooper19}. What is claimed there is a mechanism that survives the
rebuttal, not the conclusion. What is claimed as new is the assembly
and three specific items: the measure-theoretic recasting of the
survivor category, the practical/ideal ratio criterion, and the
epistemic ceiling. A reader who finds any of those three anticipated
should discount the section accordingly, and the surrounding apparatus
provides no independent support for them.

One accounting note closes the column, and it governs what all
six costs are costs \emph{of}. The unconditional list is short, and the
ledger states it in full: that conditioning selects, that adverse
anatomy is fatal (Theorem~\ref{gr1:thm:trapped}), and that viability
requires two conditions rather than one
(Theorem~\ref{gr1:thm:twocond}). The mutual singularity of $Q$ and
$\PP$ is \emph{not} on that list; it is critical-regime only. What is priced is every sentence connecting that structure to
biospheres, aliens, or the future. No rhetorical weight anywhere in this
section is permitted to exceed the grade of $\Phi3$, and where a
consequence is stated in \S\ref{gr:sub:consequences} it is stated as
downstream of it.

\paragraph{Buys.} Against those six, the section returns the following.

The categorical division is \emph{unchosen}, in the critical regime:
General and Golden are mutually singular measures there, so the split is
established by a theorem rather than by a threshold someone placed,
which is what Pre-Civilization/Civilization and every Kardashev boundary
cannot say. \S\ref{gr:subsub:regimes} records the price --- the claim is
unavailable in the viable regime --- and the entry is worth exactly what
it is worth after that subtraction.
Practical Goldenness is \emph{computable}: $t \gtrsim s/4\varepsilon$,
with the invariant $s_\varepsilon(t)/t = 4\varepsilon$, so the question
``how Golden is this'' has an answer with quantities in it. The premise
is \emph{weaker than it looks}: age alone confers practical Goldenness,
and age is what the question already supposes, so the apparatus assumes
there is a future rather than that the future goes well. The age
Lindy \emph{exponent} is derived rather than parameterized --- a
narrower buy than it first appears, and stated as such in
\S\ref{gr:sub:consequences}, since Lindy longevity for civilizations is
established prior work \cite{BC21,BG24,Kip20,Ord23} and what is added is
that the exponent is fixed by the driver's smoothness rather than
chosen. Against the indifference-based anthropic tradition
\cite{Car83,Got93,Bos02} the contrast is larger, the age distribution
there being postulated outright; but the nearer neighbour is already
model-based, and the honest gain is one parameter rather than a
paradigm.
The epistemic ceiling is \emph{symmetric}, constraining the section
exactly as tightly as its targets. And the section \emph{fills a
declared gap} in the companion: conditioning on $\{T>t\}$ is
conditioning forward, which is what \cite{Ben26Optimal} states its own
constraint cannot do.

One entry is a positive expectation rather than a dissolution, and it is
entered deliberately, since an apparatus whose entire return is relief
from other people's difficulties is answerable to nothing. The
truncation/rescaling asymmetry of \S\ref{gr:subsub:measured} is not a
fact about biospheres; it is a fact about survivorship conditioning on a
two-coordinate state. It should therefore appear in \emph{any} domain
carrying survivorship selection and an orientation-like coordinate
alongside a magnitude-like one --- firms, institutions, lineages, any
population where the old are selected from the many. The prediction is
specific: the orientation coordinate should show a deleted region among
survivors, the magnitude coordinate a shifted distribution of roughly
preserved shape, and the two divergences should be comparable in
magnitude while opposite in kind. That is checkable outside cosmology,
on data that exist, and it would embarrass the apparatus if the
selection turned out to act on magnitude categorically and orientation
quantitatively instead. It is entered here as an expectation and not a
result.

% ---------------------------------------------------------------------
\subsection{What this section does not claim}\label{gr:sub:fence}

The fence is stated hard, on the companion's discipline that an
apparatus which prices what it claims must also refuse what it does not
\cite{Ben26Optimal}.

It is not claimed that any observed system is Golden, ideally or
practically; by Theorem~\ref{gr:thm:ceiling} the ideal version is
unknowable in principle and the practical version is out of reach for
want of an ensemble. It is not claimed that the companion's cosmic
weighting and the path measure $\PP$ of
\S\ref{gr1:sec} are the same object; that
identification is flagged as unperformed in
\S\ref{gr:subsub:nexic} and nothing here depends on it. It is not
claimed that any individual biosphere, ours included, will behave as
$\S\ref{gr:sub:consequences}$ describes; those are statements about a
conditional distribution. No normative conclusion is claimed
($\Phi2$). And the companion's natural-historical program --- its
conjectures, its bounds, its paleontological predictions --- is not
claimed, imported, or re-argued; $\Phi5$ enters as a stated result and
nowhere else.

One thing is owed on the far side of the fence, because a fence is
legible only if what it excludes can be seen. A program built on this
section would inherit three commitments, stated here as refusable rather
than as claims. It would be committed to the boundary-cost scaling of
$\Phi3$ showing up in institutional and ecological data, not merely in
argument. It would be committed to the truncation/rescaling asymmetry
appearing in survivorship-selected populations generally. And it would
be committed to the exponent's dependence on noise placement --- $\tfrac14$
for integrated drivers, $\tfrac12$ for Brownian ones --- being visible
in any domain where hazard dynamics can be measured directly. None of
the three is established here, all three could be found false without
touching a line of \S\ref{gr1:sec}, and that is the
point of naming them.

% ---------------------------------------------------------------------
\subsection*{Section summary}

The unit is the biosphere ($\Phi1$), because no other term survives the
limit the theorems are taken in. The fundamental division is General
versus Golden --- $\PP$ versus $Q$ --- and it is categorical without
being chosen, since in the critical regime the two are mutually
singular; but nothing occupies
the Golden category \emph{as an ideal}, which is a measure to reason
from rather than a state to be in, and this is what distinguishes Golden
reasoning from longtermism ($\Phi2$). The practical form is occupiable
and is what the section actually assumes: a biosphere of age $t$ is
$\varepsilon$-Golden over windows of length $4\varepsilon t$, a constant
fraction of its own history, so that the premise of Golden reasoning is
not that the future goes well but merely that there is one. The two coordinates are technostrength, a
state coordinate, and technoanatomy, a structural parameter, and
conditioning treats them differently in kind: adverse strength is
truncated abruptly, $137$-fold, while anatomy is selected gradually,
$\PP(\rho>0)$ moving only from $0.75$ to $0.84$ over the same horizon.
The second is the slower and the more decisive, since adverse anatomy is
eventually fatal with certainty while adverse strength is a condition
every biosphere occupies half of logarithmic time and survives. The Kardashev scale fails not because its
thresholds are arbitrary but because the terminal stage it points at is
gated by a technoanatomy the scale cannot index, so that climbing is
neither necessary nor sufficient for arriving --- and, with adverse
anatomy, climbing fast is itself the failure. The violent-aliens thesis fails at the
inference from origin to maturity ($\Phi4$), and fails further if
extraction is bad technoanatomy ($\Phi3$) --- which we argue from
boundary cost, recovering the square--cube law from which the names
came, and which we mark as the section's one load-bearing premise, with
its falsification conditions. Downstream: encountered biospheres are not
extractive, megastructure nulls are weak evidence, and old biospheres
are Lindy with a derived exponent rather than an assumed one --- which
is where Golden reasoning improves on Copernican anthropics, by making
the age distribution an output rather than an input. With rarity of
origins ($\Phi5$), the silence needs no filter ahead. And by the
epistemic ceiling, none of this can be confirmed by encounter --- a
limitation that binds this section exactly as tightly as it binds what
the section argues against.

% ---------------------------------------------------------------------

 ============================

\section{The human role: two errors about it}\label{hr:sec}

\begin{quote}\small
Two positions occupy the ends of a spectrum concerning what humans are
for. At one end, techno-optimism: capability is the thing, and more of
it is the answer. At the other, a strain of environmentalism holding
that humanity is a net harm to the biosphere and that its voluntary
disappearance would serve the living world. Golden reasoning finds both
mistaken, and mistaken in ways that are structurally opposite --- each
neglects exactly the factor the other over-weights.

The argument against techno-optimism is already in hand and is
consolidated here: it neglects technoanatomy, and by
Remark~\ref{gr1:rem:againstsing} acceleration without anatomy is not
merely insufficient but actively worse than stasis
(\S\ref{hr:sub:optimism}). The argument against extinctionism is new and
is the section's substance: on the extinctionist's \emph{own} valuation
--- the biosphere's persistence --- removing the biosphere's only
hazard-lowering agent lowers the thing they mean to protect
(\S\ref{hr:sub:extinction}). Both arguments are conditional, and the
conditions are stated rather than assumed.

What emerges is a position on environmental value that neither end
occupies: nature just \emph{is} the biosphere, the Golden Point is the
only route by which it acquires an unbounded expectancy, and human
existence is on that accounting not an accident at odds with the living
world but structurally implicated in its only route to persisting
(\S\ref{hr:sub:environment}). The threat is not thereby diminished ---
$\rho<0$ remains the live danger, and this section is as much an
argument for urgency about anatomy as against despair about presence.
\end{quote}

% ---------------------------------------------------------------------
\subsection{Against extinctionism}\label{hr:sub:extinction}

\subsubsection{The position, in its own terms}\label{hr:subsub:position}

The target is a specific and explicitly biosphere-valuing position, and
it must be stated in its own terms or the argument does not engage it.
The Voluntary Human Extinction Movement holds that humanity should
abstain from reproduction so as to bring about its own gradual
disappearance, on the ground that this would prevent environmental
degradation \cite{VHEMT}. The reasoning proceeds \emph{from} the value
of the living world: the movement's own framing weighs the greater value
of Earth's entire biosphere against the lesser value of one species, and
holds that once one envisages the biosphere entire and grants the
intrinsic value of every life form, voluntary extinction begins to make
sense.

That framing is what makes the position addressable here. Its premise
--- that the biosphere's persistence and flourishing is what matters ---
is one this framework can evaluate, because persistence of a biosphere
is the quantity \S\ref{gr1:sec} models.

\begin{remark}[What this argument does not touch]\label{hr:rem:scope}
It is essential to distinguish the target from a neighbouring position
it is often assimilated to. Benatar's antinatalism holds that coming
into existence is always a serious harm, and concludes that human
extinction would be preferable, on an asymmetry between the badness of
pain and the goodness of pleasure \cite{Benatar06}. That argument does
not proceed from biosphere value at all; its currency is suffering, and
extinction is a consequence rather than the objective.

\emph{Nothing below engages it.} An argument that removing humans lowers
the biosphere's expectancy is no reply to someone who never valued the
expectancy. The suffering-focused position must be met on its own ground
and this section does not attempt it. What follows applies to
extinctionism \emph{in its environmental form} --- and that restriction
is not a weakening, since the environmental form is the one that appeals
to a premise the framework shares.
\end{remark}

\subsubsection{The argument}\label{hr:subsub:extinctionargument}

\begin{proposition}[Removing the muffling agent forfeits viability;
\textbf{[P]} given the model]\label{hr:prop:extinction}
In the model of \S\ref{gr1:sec}, a biosphere with no hazard-lowering
agent is the case $x\equiv0$: no muffling is accumulated, and the total
hazard reduces to the background, $h_t=B_t$. Then
$\Lambda_t=\int_0^t B_u\,du\to\infty$, so
\[
\PP(\Lambda_\infty<\infty)=0
\]
identically --- not small, but zero, and for every parameter value.

Compare the alternative. By Theorem~\ref{gr1:thm:twocond}, a biosphere
that does accumulate muffling has
$\PP(\Lambda_\infty<\infty)=\one\{\rho>0\}\cdot
\bigl(s(\mu_0)-s(-\infty)\bigr)/\bigl(s(+\infty)-s(-\infty)\bigr)$,
which is strictly positive whenever anatomy is positive and the driver
transient. The extinctionist proposal is therefore a choice of an option
with viability probability exactly zero over one whose viability
probability may be positive.
\end{proposition}

The magnitude of what is forfeited is not left abstract. The companion
treats the external hazard as roughly constant over long timescales,
since the cosmic background does not change appreciably from era to era,
giving a survival probability $e^{-\lambda t}$ and an external life
expectancy of $1/\lambda$ \cite{Nexus}. The framework's own reading of
the Rare Earth and High-Contrast Earth hypotheses is that
$\lambda t_0>1$, where $t_0$ is the length natural history actually
took --- which is to say that \emph{Earth's remaining expectancy under
the external hazard alone is shorter than the span already elapsed.}

\begin{remark}[Undisturbed nature is not preserved nature]
\label{hr:rem:undisturbed}
This is the point on which the position turns, and it is worth stating
without ornament. The extinctionist picture is of a living world
continuing undisturbed once the disturbance is removed. But
``undisturbed'' is not a stable state; it is the case $x\equiv0$, in
which the biosphere has no means of lowering its own hazard and the
background runs unopposed.

The reframing of Definition~\ref{gc:def:departure} sharpens what is
being chosen, and sharpens it against the position. Under a reading on
which the relevant event is biological death, ``leave nature alone''
sounds like survival, and the argument must show that death follows
anyway. Under the reading on which the Departure Point is the closure of
possibility, no such demonstration is needed. At $x\equiv0$,
Proposition~\ref{hr:prop:extinction} gives
$\PP(\Lambda_\infty<\infty)=0$ identically: \emph{no} continuation
available to that biosphere advances. The modal quantifier of
Definition~\ref{gc:def:departure} is therefore satisfied at once, and
the Departure Point is not a fate awaiting the undisturbed world but
its present condition. In the terms of
\S\ref{gc:subsub:cascade}: $x\equiv0$ lies in
$\mu G\setminus\mathrm{Base}$ --- it is a departure at which no hazard
has yet fired, entered by configuration rather than by event. The
extinctionist proposal is a proposal to enter the cascade
deliberately. A biosphere left alone has not been preserved;
it has been placed in the one configuration from which nothing
leads anywhere. The extinctionist is not choosing a world that
survives without us. They are choosing a world that cannot go
anywhere. Destabilisation of the star, a sufficiently
close hypernova, orbital perturbation by a passing mass --- these do not
become less likely because humanity has withdrawn. They become the only
things that happen.

And the choice is not between a damaged biosphere and an undamaged one.
It is between a biosphere with a positive probability of persisting
indefinitely and one with none. On the extinctionist's own valuation,
the second is worse.
\end{remark}

\subsubsection{The concession, and why it does not rescue the position}
\label{hr:subsub:concession}

The argument has a real limit, and stating it is what keeps it honest.

\begin{remark}[Adverse anatomy does not vindicate the position]
\label{hr:rem:concession}
The natural concession at this point is that under $\rho<0$ the
extinctionist is closer to right, since $x\equiv0$ leaves hazard merely
at background while adverse anatomy drives it upward. The concession
should not be made, and seeing why changes which options the argument is
really between.

\emph{First, the local comparison goes the other way.} A biosphere with
adverse anatomy that \emph{halts} at the optimum $x_\ast$ of
Proposition~\ref{gr1:prop:trichotomy}(iii) sits at hazard $h_\ast$, and
$h_\ast<h(0)$ for every $\rho<0$. Numerically, at $B=1$, $b=3$, $k=6$:
\begin{center}\small
\begin{tabular}{@{}rrrrr@{}}
\toprule
$\rho$ & $h(0)$ & $x_\ast$ & $h(x_\ast)$ & expectancy gain\\
\midrule
$-0.05$ & $1.00$ & $11.47$ & $0.0046$ & $217\times$\\
$-0.20$ & $1.00$ & $8.88$ & $0.0176$ & $57\times$\\
$-1.00$ & $1.00$ & $4.28$ & $0.407$ & $2.5\times$\\
$-2.00$ & $1.00$ & $1.82$ & $0.898$ & $1.1\times$\\
\bottomrule
\end{tabular}
\end{center}
Even at $\rho=-2$ --- anatomy adverse enough that capability doubles
internal hazard twice over --- halting still beats having no muffling at
all. So $x\equiv0$ is strictly dominated for every $\rho<0$. Adverse
anatomy is worse than absence only if the biosphere \emph{fails to
halt}, which is a failure to stop rather than a property of the anatomy.

\emph{Second, and this is the point the local comparison obscures, both
options forfeit.} By Theorem~\ref{gr1:thm:twocond},
$\PP(\text{viable})=\one\{\rho>0\}\cdot(\text{transience factor})$. At
$x\equiv0$ there is no muffling agent, hazard sits at background, and
$\Lambda_\infty=\infty$. At $\rho<0$ halted at $x_\ast$, hazard sits at
$h_\ast>0$, and $\Lambda_\infty=\infty$. \emph{Both are zero.} The
extinctionist and the biosphere that settles at $x_\ast$ are not
choosing between a worse and a better outcome; they are choosing between
two ways of dying, and the model does not distinguish them in the only
respect that admits an answer other than death.

The response to adverse anatomy is therefore not to select among the
zeros. It is to move $\rho$, which is the one intervention that changes
the indicator rather than the multiplicand --- and which
\S\ref{ad:sec} exists to model. A reader persuaded of the danger and
unpersuaded that $\rho$ can be moved still has a coherent position, and
it remains the strongest form of the view opposed here; but it is a
claim about the tractability of anatomy, not a defence of extinction.
\end{remark}

\subsubsection{Why forfeiting is not the same as never having had}
\label{hr:subsub:stakes}

The preceding leaves a gap that arithmetic alone will not close. If both
options carry $\PP(\text{viable})=0$, why prefer either to the other,
and why regard the choice as weighty rather than indifferent?

The answer is that the two are not symmetric with respect to what has
already been spent --- and note that under
Definition~\ref{gc:def:departure} what is at stake is the trajectory
rather than the lives on it, which is why the accumulation is the right
object to weigh. A biosphere at $x\equiv0$ has not merely a low
viability probability; it has \emph{surrendered} an accumulation, and
the accumulation is the subject of the companion's Miracle sequence
\cite{Nexus}. Whatever its correct magnitude, the record is a chain of
events each of which had to succeed for the next to be attempted: an
origination, a transfer if panspermia obtained, the eukaryotic
transition, the Dark Era, the end-Cretaceous impact, the
$\alpha$-gal epidemic. The companion's estimate for a single such
step --- the interstellar transfer --- is of order $10^{-21}$
\cite{Nexus}.

\begin{remark}[The stakes rise monotonically]\label{hr:rem:monotone}
One structural feature of that chain is worth extracting, because it is
what makes the argument an argument rather than an appeal to sunk cost.
Each Miracle in the sequence is \emph{more} consequential than the last
--- in the companion's formulation, because the gap to the Golden Point
was by then that much closer \cite{Nexus}. The reason is not sentimental
but combinatorial: a step late in the chain carries all the improbability
of every prior step behind it, so the expected cost of failing at step
$n$ is the whole product $\prod_{i<n}p_i$ rather than $p_n$ alone.

Two consequences follow. The stakes attaching to a decision rise
monotonically along the trajectory, so the present moment carries the
highest stakes the lineage has yet faced --- not by self-importance but
by construction. And sunk cost is not the structure here: the classical
sunk-cost fallacy consists in letting past expenditure govern a choice
whose future payoffs are unchanged by it, whereas here the past
expenditure has \emph{purchased} the option that remains, and forfeiting
is forfeiting that option rather than honouring a debt. The distinction
is that between a bad reason to persist and a live asset one may
discard.
\end{remark}

\begin{remark}[Forfeiture is not observable]\label{hr:rem:silentforfeit}
The cascade of \S\ref{gc:subsub:cascade} adds a feature to the stakes
argument that is worth stating separately, because it is the feature
that makes the argument urgent rather than merely large.

A biosphere enters $\mu G$ at the moment the last advancing route
closes, and by \eqref{gc:eq:departurefix} that moment precedes the
firing of any hazard --- possibly by a great deal. Nothing marks it.
There is no event to observe, no loss to register, and by
\S\ref{gr:sub:ceiling} no finite observation that would distinguish a
biosphere which has just crossed from one which has not. The
accumulation is not forfeited in some future catastrophe one might see
coming and reconsider; it is forfeited silently, at whatever point the
routes ran out, and the catastrophe is only the execution.

This does not make the stakes argument more permissive --- see the
limits below --- but it does relocate what the argument is about. It is
not a claim that a visible disaster would be very bad. It is a claim
that the relevant transition is invisible, which is why deferring
attention to anatomy until the danger is legible is deferring it past
the point where deferral is possible.
\end{remark}

\begin{remark}[What this does and does not license]\label{hr:rem:stakeslimit}
The stakes argument establishes an asymmetry between forfeiting and
never having had, and nothing beyond it. In particular it does not
license: that any present cost is acceptable, since the argument is
silent on rates of exchange; that the Golden Point will be reached,
which \S\ref{gr:sub:ceiling} forbids anyone from knowing; or that
persons alive now bear a duty derived from the accumulation, which
$\Phi2$ excludes and which the framework has no machinery to generate.

What it does license is narrow and, we think, sufficient: that the two
options which forfeit viability are not the neutral defaults they are
often taken for, and that between an option which changes the indicator
and options which do not, the model has a determinate preference.
\end{remark}

% ---------------------------------------------------------------------
\subsection{Against techno-optimism}\label{hr:sub:optimism}

The opposite error requires less new argument, since the result it
founders on is already proved.

Techno-optimism holds that continued growth in capability is the answer
to the risks capability creates. In the model's terms it supplies the
transience factor of Theorem~\ref{gr1:thm:twocond} and is silent on the
anatomy factor --- $\rho$ does not appear in any accounting of energy
capture, capability, or rate of progress, and cannot be inferred from
them.

The verdict is not that this leaves the conclusion unsupported but that
it inverts it. With $\rho<0$, feedback does not fail to secure survival;
it \emph{accelerates the failure}, since $\mu_t\to+\infty$ drives
$x_t\to+\infty$ and hence, by
Proposition~\ref{gr1:prop:trichotomy}(iii), $h\to\infty$
superexponentially. Supplying capability alone is strictly worse than
supplying neither: a poorly-constituted biosphere that stagnates sits
near its hazard floor $h_\ast$, while one that accelerates runs at its
own cliff.

\begin{remark}[The cascade forces anatomy to be a process]
\label{hr:rem:forcesprocess}
Both arguments of this section, taken with
\S\ref{gc:subsub:cascade}, deliver a verdict on the fixed-$\rho$ model
that is worth confronting rather than absorbing.

If $\rho<0$ and $\rho$ is a \emph{constant}, then
Theorem~\ref{gr1:thm:twocond} gives
$\PP(\Lambda_\infty<\infty)=0$ identically --- at every time, including
$t=0$. By the modal definition the biosphere is then in $\mu G$ from the
outset: no continuation available to it ever advances, and it departed
before anything happened. That is a degenerate verdict. It says the
question of what such a biosphere should do was settled before it was
asked, which is not a claim about biospheres but a symptom of the
model.

The degeneracy has exactly one exit, and it is not a matter of taste.
It is avoided if and only if $\rho$ can change --- if adverse anatomy
now leaves open a continuation in which anatomy is not adverse later.
That is the fluctuating-anatomy setting of
\S\ref{gr1:subsub:fluct}, where
Corollary~\ref{gr1:cor:recurrentanatomy} and
Proposition~\ref{gr1:prop:tolerance} take over, and where the question
becomes whether the escape is completed in time rather than whether it
was ever available.

So the cascade converts emergent anatomy from a desirable extension
into a \emph{requirement of coherence}. A framework that treats
technoanatomy as a parameter and departure as modal closure will
pronounce every adverse-anatomy biosphere already lost; only a
framework in which anatomy has its own dynamics can say the thing this
section is trying to say, which is that the matter is open and the
window is closing. The arguments above should be read as holding under
the fluctuating reading, and the fixed-$\rho$ statements as the
limiting case in which the window has already shut.
\end{remark}

\begin{remark}[The two errors are the same error]\label{hr:rem:sameerror}
Placed side by side, the positions are not opposites but a single
mistake taken in two directions. Both treat capability as the only
variable. The techno-optimist maximises it; the extinctionist minimises
it; neither carries a parameter for whether capability, once acquired,
is turned inward. Once $\rho$ is in the model the spectrum they occupy
is revealed as one axis of a two-dimensional space, and the interesting
region --- positive anatomy with accumulating capability --- lies on
neither end of it. That is why the disagreement between them has proved
unresolvable by argument along their shared axis: the quantity that
settles it is not on that axis.
\end{remark}

% ---------------------------------------------------------------------
\subsection{What this implies for environmental value}
\label{hr:sub:environment}

Three consequences follow, of decreasing formality and increasing
interest. They are stated as consequences of the foregoing and inherit
its conditions.

\emph{Nature is the biosphere.} The unit \S\ref{gc:sub:vocab} adopts for
methodological reasons --- the totality of living and technical activity
associated with a world, taken as one system --- coincides with what
environmental thought means by nature when it speaks of intrinsic value
in the whole rather than in particular organisms. The framework does not
need to argue for that identification; it arrives at the same unit from
the requirement that a conditioning target apply across the whole record
(\S\ref{gc:sub:continuity}).

\emph{The Golden Point is the only route to an unbounded expectancy.}
Every alternative has a finite one. A biosphere without a
hazard-lowering agent has expectancy set by the external background; a
biosphere with adverse anatomy has a shorter one still. Only the
conjunction of positive anatomy and accumulating capability yields
$\PP(\Lambda_\infty<\infty)>0$. Whatever else environmental value
recommends, if it recommends the living world's persistence then it
recommends passage through that conjunction, since no other route is
available.

\emph{Human existence is not ontologically at odds with the living
world.} This is the claim most likely to be over-read, so its content
and its limits are given precisely.

\begin{remark}[What is and is not established]\label{hr:rem:notmistake}
\emph{Established, conditional on the foregoing:} the outcome in which
the biosphere persists indefinitely runs through the existence of an
agent capable of lowering its hazard. Humanity is, so far as the record
shows, that agent's occurrence in this biosphere. So on this accounting
human existence is not contingent noise afflicting an otherwise
self-sufficient nature; it is structurally implicated in the only
trajectory on which nature persists. Under Golden-conditioning
(\S\ref{gc:sec}) the implication is stronger still, since the record is
then read as produced under a conditioning that the passage occur.

\emph{Not established:} that the harm is small, that any particular
human activity is vindicated, or that the present trajectory is good.
The threat is real, $\rho<0$ is the live danger, and
Remark~\ref{hr:rem:concession} concedes that a biosphere with adverse
anatomy is worse than none. The claim concerns humanity's
\emph{structural position}, not its conduct, and the two are
independent: an agent can be necessary to an outcome and be failing at
it.

\emph{Also not established, and worth disclaiming explicitly:} any
conclusion about the standing of human meaning-making at large. It is
tempting to read the above as supplying, from a source outside human
self-interpretation, a verdict that humanity is not a mistake --- and as
answering a broader cultural pessimism thereby. The temptation should be
resisted at this grade. What the framework delivers is a conditional
structural claim about one biosphere's route to persistence, resting on
$\Phi3$ and on the Golden-conditioning argument, both of which are
argued rather than proved. That is a substantial thing to have. It is
not a philosophy of history, and presenting it as one would trade a
defensible result for an indefensible one.
\end{remark}

% ---------------------------------------------------------------------
\subsection{Ledger}\label{hr:sub:ledger}

\paragraph{Costs.} \emph{The valuation premise} --- the argument against
extinctionism is conditional on valuing the biosphere's persistence, and
is no argument at all against suffering-focused antinatalism
\cite{Benatar06}; \emph{falsifier:} none required, the restriction is
stated. \emph{The identification of humanity as the hazard-lowering
agent} --- \textbf{[$\Phi$]}; \emph{falsifier:} a mechanism by which a
biosphere lowers its own hazard without anything playing that role.
\emph{$\Phi3$ and Golden-conditioning}, inherited by
\S\ref{hr:sub:environment} and carrying their grades unchanged.

\paragraph{Buys.} An argument against environmental extinctionism that
proceeds \emph{on its own valuation} rather than by denying it. A
diagnosis on which the two ends of the spectrum share one error rather
than opposing each other. And a position on environmental value that is
neither anthropocentric nor misanthropic, obtained without adding a
premise for the purpose.

\paragraph{Not claimed.} That human conduct is vindicated; that the
danger is overstated; that suffering-focused antinatalism is answered;
or that anything here bears on the standing of human meaning-making
beyond the structural claim stated.


\section{Adversivity: institutional balance and the microfoundation of
technoanatomy}\label{ad:sec}

\begin{quote}\small
Technoanatomy entered \S\ref{gr1:sec} as a parameter: the exponent
$\rho$ governing whether capacity acquired is capacity turned inward.
Nothing there says where it comes from. This section builds the
apparatus intended to supply it --- a model of institutions that
inflict hazard on one another and reorganise under a single principle
--- and reports how far that apparatus gets on its own. It does not yet
reach $\rho$; \S\ref{ad:sub:toward} states precisely what remains.

The principle is that every institution seeks \emph{net adversivity
balance}: its weighted alliances should cancel its weighted enmities.
Its two motivating cases are historical, and the section's numerical
work reproduces both. Two literatures border it closely and are placed
in \S\ref{ad:sub:placement}; the principle is not a refinement of
either, and against structural balance in particular it has
\emph{disjoint} ground states.
\end{quote}

% ---------------------------------------------------------------------
\subsection{The principle}\label{ad:sub:principle}

\begin{premise}[Net adversivity balance; \emph{governing principle}]
\label{ad:phi:balance}
Each institution assigns every other a signed weight, positive for
alliance and negative for enmity, and seeks to make the sum of those
weights --- taken against the others' sizes --- vanish. It is pressed
away from having only allies exactly as it is pressed away from having
only enemies.\footnote{%
\emph{Provenance.} The principle is recorded in the author's journal of
August 2025, with both cases below and with the accompanying diagnosis
that political theory has attended to its various doctrinal categories
where the underlying process is one of stable and unstable equilibria:
``if there are all alliances, the weights shift and institutions start
targeting each other, and if there are more adversaries than not,
alliances form simply to counter that weight and for no particular other
reason.'' The formalisation, the placement against the two literatures
of \S\ref{ad:sub:placement}, and everything numerical below are later
and are the assistant's.}
\end{premise}

Two cases motivate it, and both are recovered numerically in
\S\ref{ad:sub:experiments}.

\emph{Removal of a common enemy.} When the Welsh drove back the
Anglo-Saxons, the alliance among Welsh tribes did not survive the
threat that had produced it. On the principle, the alliances were never
about one another: they were the positive weight required to offset a
large negative one, and their justification vanished with it.

\emph{Arrival of a common enemy.} Conversely, the only circumstance
under which Earth's states would plausibly unify is the arrival of
something opposed to all of them --- whereupon, on the principle, they
would unify not by persuasion but by the same arithmetic running the
other way.

\begin{remark}[The load-bearing departure]\label{ad:rem:departure}
The principle's distinctive content is that \emph{all-allied is
unstable}. That is what the removal case requires, and it is precisely
where the principle parts from structural balance, in which all-allied
is a ground state. Everything in \S\ref{ad:sub:placement} turns on this
one disagreement.
\end{remark}

% ---------------------------------------------------------------------
\subsection{Placement: two literatures, and a disjointness}
\label{ad:sub:placement}

\subsubsection{Structural balance}\label{ad:subsub:heider}

Signed-graph balance theory begins with Heider \cite{Hei46}, and was
formalised for graphs by Cartwright and Harary \cite{CH56}. A triad is
\emph{balanced} when the product of its three signs is positive, so that
$(+,+,+)$ and $(+,-,-)$ are stable while $(+,+,-)$ and $(-,-,-)$ are
not. The structure theorem states that a complete signed graph is
balanced if and only if it splits into at most two camps, positive
within and negative between. Continuous-time dynamics were given by
Marvel, Kleinberg, Kleinberg and Strogatz \cite{MKKS11} and a
statistical-physics treatment by Antal, Krapivsky and Redner
\cite{AKR05}.

The Welsh case is instructive in that vocabulary. While the
Anglo-Saxons are present, every pair of tribes sits in a triad with two
negative edges, which forces the third positive: the tribes are allied
by triad closure, and nothing in their own relations holds them
together. Remove the shared enemy and those props vanish, leaving
unbalanced triads that must repartition.

\begin{remark}[Disjoint ground states; \textbf{P}]
\label{ad:rem:disjoint}
The two principles are not variants. Writing $\hat E_{\mathrm H}$ for
the normalised Heider energy and $\hat E_{\mathrm A}$ for mean-square
net adversivity, at $n=5$ with unit weights:
\begin{center}\small
\begin{tabular}{@{}lrr@{}}
\toprule
configuration & net adversivity & Heider energy\\
\midrule
all-allied (Heider ground state) & $(4,4,4,4,4)$ & $-10$\\
two blocs $3/2$ (Heider ground state) & $(0,0,0,-2,-2)$ & $-10$\\
pentagon$^{+}$/pentagram$^{-}$ & $(0,0,0,0,0)$ & $0$\\
\bottomrule
\end{tabular}
\end{center}
The adversivity optimum is Heider's \emph{worst} case --- maximally
frustrated --- and both Heider optima fail adversivity balance. Nor is
two-bloc structure available even approximately: with unit weights a
majority bloc of size $p$ and a minority of size $q$ would need
$p-1=q$ and $q-1=p$ simultaneously. The principle's stable states are
frustrated ones, which fits the observation that real international
systems do not settle into two clean camps.
\end{remark}

\subsubsection{Balance of power}\label{ad:subsub:realism}

The second neighbour is a different tradition, with a name confusingly
close to the first. Realist theory in international relations holds
that states balance against concentrations of power \cite{Wal79},
refined by Walt into balancing against perceived \emph{threat} rather
than raw capability \cite{Wal87}. Both motivating cases above are textbook for
that tradition, and the author's stipulation that perceived threat
tracks an adversary's present unity rather than its capacity to
recombine is close to Walt's central move.

The present model sits between the two: the signed-network formalism of
the first, the node-local balancing condition of the second. Whether
that combination is already available in the literature we have not
established, and it is the first thing a reader should check.

% ---------------------------------------------------------------------
\subsection{The energy}\label{ad:sub:energy}

Institutions are nodes; $w_{ij}\in[-1,1]$ is $i$'s weight on $j$;
$s_j>0$ is $j$'s size.

\subsubsection{The model}\label{ad:subsub:constrained}

The model is a two-term energy on a constrained set:
\begin{equation}\label{ad:eq:model}
\min_{W}\ \ a\,E_1(W)+b\,E_2(W)
\qquad\text{subject to}\qquad
W=W^{\!\top},\quad \sum_{j\ne i}|w_{ij}|=\beta ,
\end{equation}
with
\begin{align}
E_1&=\frac1n\sum_i\Bigl(\sum_{j\ne i}w_{ij}s_j\Bigr)^{2}
&&\text{adversivity balance; node-local, quadratic}\label{ad:eq:E1}\\
E_2&=\frac{-1}{2\binom n3}\!\!\sum_{i<j<k}\!\!
\bigl(w_{ij}w_{jk}w_{ki}+w_{ji}w_{kj}w_{ik}\bigr)
&&\text{Heider frustration; triad-local, cubic}\label{ad:eq:E2}
\end{align}
and $\beta=n-1$. Equation \eqref{ad:eq:E2} is computed from
$\operatorname{tr}(W^{3})$, which equals three times the triad sum when
the diagonal vanishes.

There is \emph{one} free parameter, the ratio $b/a$, and the two
constraints are substantive rather than technical: relations are
\emph{mutual}, and attention is \emph{finite}. Both are claims about
what an institution is, and \S\ref{ad:subsub:relaxation} records what
goes wrong if either is dropped.

\subsubsection{The relaxation used for computation}
\label{ad:subsub:relaxation}

Hard constraints are awkward to anneal, so the numerical work below
minimises a four-channel relaxation in which each constraint is
replaced by a penalty:
\begin{align}
\hat E_3&=\frac{1}{8n(n-1)}\sum_{i\ne j}(w_{ij}-w_{ji})^{2}
&&\text{reciprocity; relaxes } W=W^{\!\top}\label{ad:eq:E3}\\
\hat E_{\mathrm b}&=\frac1n\sum_i\Bigl(\sum_{j\ne i}|w_{ij}|-\beta\Bigr)^{2}
&&\text{engagement budget; relaxes the norm}\label{ad:eq:Eb}
\end{align}
with $E=a\hat E_1+b\hat E_2+c\hat E_3+d\hat E_{\mathrm b}$. The two forms
agree on everything reported here. At $n=9$, $b=4$, over independent
seeds:
\begin{center}\small
\begin{tabular}{@{}lll@{}}
\toprule
& four-term relaxation & two-term constrained\\
\midrule
$a=0$ & lopsided $6/3$, $7/2$ & lopsided $6/3$, $6/3$\\
$a$ small & even $4/5$, $\hat E_1=1.78$ & even $4/5$, $E_1=1.78$\\
$a=1$ & frustrated web, $\hat E_1\approx0$ & frustrated web, $E_1=0.03$\\
\bottomrule
\end{tabular}
\end{center}
The relaxation is therefore a computational convenience and
\eqref{ad:eq:model} is the model. Three facts about the relaxation are
nonetheless worth recording, because each explains why the
corresponding constraint is in \eqref{ad:eq:model} at all.

\begin{remark}[The budget is not optional]\label{ad:rem:apathy}
Drop the norm constraint and the principle has a trivial optimum:
\emph{be indifferent to everyone}. Numerically, pure $\hat E_1$ reaches
exact balance at $n=6$ by leaving eight of fifteen ties weak. Balance by
apathy is not what the principle intends, and finite attention --- a
fixed total spent across one's ties --- is the substantive claim that
rules it out. This is why the norm appears in \eqref{ad:eq:model} as a
constraint rather than as a preference to be traded against others.
\end{remark}

\begin{remark}[Reciprocity is not emergent]\label{ad:rem:recip}
$E_1$ depends only on row sums and is therefore blind to the
antisymmetric part of $W$; $E_2$ constrains cyclic products and does not
force symmetry either. Removing $\hat E_3$ from the relaxation and
re-optimising produces asymmetries at the maximum,
$|w_{ij}-w_{ji}|=2$: total one-sided devotion. Symmetry must therefore
be \emph{imposed}, which is why it is a constraint in
\eqref{ad:eq:model}. Its near-zero value at relaxed optima indicates
that it is cheap to \emph{satisfy}, not that it is redundant --- an
inference the low value invites and which the test above refuses.
\end{remark}

\begin{remark}[Normalise by the mean, not the maximum]
\label{ad:rem:normalisation}
This concerns the relaxation, and is the one point at which a natural
choice is badly wrong. Dividing $\hat E_1$ by its maximum --- attained
at the all-allied state ---
crushes it, since near any configuration of interest it then varies over
only $O(n^{-2})$, and the Heider term dominates at any comparable
weight. The mean-square form \eqref{ad:eq:E1} instead assigns the
balanced two-bloc state a cost of exactly $1$ \emph{for every $n$},
since there $\sum_j w_{ij}=(p-1)-q=-1$ regardless of size, matching the
Heider floor of $-1$. Only under this normalisation is $b/a$ an
interpretable weight, and only then is the model's behaviour independent
of $n$. No re-powering of $\hat E_2$ achieves the same effect, and
taking roots of it would destroy the sign structure that makes it
frustration in the first place.
\end{remark}

% ---------------------------------------------------------------------
\subsection{Ground states}\label{ad:sub:ground}

\begin{proposition}[Perfect balance requires $n\equiv1\pmod4$;
\textbf{P}]\label{ad:prop:congruence}
Among symmetric configurations with all ties saturated at $\pm1$ and
equal sizes, exact adversivity balance is attainable if and only if
$n\equiv1\pmod4$. Each node requires $(n-1)/2$ positive ties, forcing
$n$ odd; and the total number of positive edges, $n(n-1)/4$, must be an
integer, forcing $n\equiv1\pmod4$. The admissible sizes are
$n=5,9,13,17,\dots$
\end{proposition}

\begin{proof}
Both conditions are immediate; their conjunction holds exactly on
$n\equiv1\pmod4$. Exhaustive search confirms the minimum of
$\sum_i(\text{net}_i)^2$ is zero at $n=5$ and positive at
$n=3,4,6$, the first case being realised by the pentagon of positive
ties with the complementary pentagram of negative ones.
\end{proof}

\begin{proposition}[Two failure modes, by parity; \textbf{P}]
\label{ad:prop:parity}
Write $P$ and $M$ for the numbers of alliances and enmities, so that
$P+M=\binom n2$. Summing every node's ledger counts each tie twice:
\begin{equation}\label{ad:eq:parity}
\sum_i \mathrm{net}_i \;=\; 2\,(P-M).
\end{equation}
Since $P-M$ and $P+M$ share a parity, the total imbalance is
\emph{forced} away from zero whenever $\binom n2$ is odd:
$\sum_i\mathrm{net}_i\equiv2\pmod 4$. Two regimes follow.

\emph{Even $n$: the imbalance spreads.} Here $n-1$ is odd, so every
$\mathrm{net}_i$ is a sum of an odd number of $\pm1$ terms and is
therefore odd; \emph{no} node can balance, and the cheapest option
distributes the residual --- at $n=4$ the optimum is $(-1,-1,-1,-1)$.

\emph{$n\equiv3\pmod4$: the imbalance concentrates.} Here individual
balance is possible but the total is not, and because the penalty is
\emph{quadratic} it is cheaper to load a mandatory total of $2$ onto one
node ($\sum\mathrm{net}_i^2=4$) than to spread it over three
($-2,-2,+2$ sums correctly but costs $12$). Squares reward
concentration.
\end{proposition}

\begin{remark}[The scapegoat, and what it does not license]
\label{ad:rem:scapegoat}
Proposition~\ref{ad:prop:parity} accounts for every computed residual:
\begin{center}\small
\begin{tabular}{@{}rrrrr@{}}
\toprule
$n$ & $\binom n2$ & forced $\sum\mathrm{net}_i$ & predicted min
$\sum\mathrm{net}_i^2$ & computed\\
\midrule
3 & 3 (odd) & $\equiv2$ & 4 & 4\\
4 & 6 (even) & $\equiv0$ & 4 & 4\\
5 & 10 (even) & $\equiv0$ & \textbf{0} & 0\\
6 & 15 (odd) & $\equiv2$ & 6 & 6\\
7 & 21 (odd) & $\equiv2$ & 4 & 4\\
\bottomrule
\end{tabular}
\end{center}
At $n=7$ the optimum leaves six nodes exactly balanced and one carrying
the whole residual; annealing under continuous weights reproduces it,
six nodes within $0.17$ of zero and one at $-1.64$. So where the
arithmetic forbids universal balance, the cheapest arrangement is not
shared discomfort but a single universally-disliked member --- and this
is forced by counting rather than emergent from dynamics.

Two limits on the reading. The result assumes unit weights and equal
sizes; with continuous weights every $n$ reaches balance by weakening
ties, so the congruence marks where balance is \emph{free} rather than
where it is possible. And the scapegoat mechanism, unlike the
congruence, survives that relaxation: the quadratic penalty rewards
concentrating an unavoidable residual on one member whatever the
weights, which is the part worth carrying forward.
\end{remark}

\begin{remark}[Where factions appear]\label{ad:rem:factions}
Under the normalisation of Remark~\ref{ad:rem:normalisation}, the two
channels are commensurate at $b\approx a$, and at $n=9$ the positive-tie
components remain a single connected web up to $b\approx2$, splitting
into two factions between $b=2$ and $b=4$. The interesting regime is
$b\approx4$, where the system exhibits \emph{both} a $5/4$ split
\emph{and} near-perfect balance, $\hat E_1\approx0.005$; only above
$b\approx8$ is balance abandoned, $\hat E_1\to1.6$, in favour of
saturated bipolarity.

At large $b$ the adversivity term takes on a further role, and it is
more than tie-breaking. Heider admits two ground states, one bloc and
two blocs, both at $\hat E_2=-1$; since all-allied costs
$\hat E_1=(n-1)^{2}$ while two-bloc costs $1$, any positive $a$ breaks
that degeneracy in favour of two blocs. But Heider is also indifferent
among two-bloc \emph{partitions}: every bipartition sits at $\hat E_2=
-1$ whatever the bloc sizes. Numerically at $n=9$, $b=4$, the effect of
$a$ across three seeds:
\begin{center}\small
\begin{tabular}{@{}rlrr@{}}
\toprule
$a$ & faction sizes found & $\hat E_1$ & $\max_i|\mathrm{net}_i|$\\
\midrule
$0$ & $6/3,\;3/6,\;7/2$ --- lopsided & $12.1$ & $4.67$\\
$0.001$--$0.1$ & $4/5,\;5/4$ --- even, every seed & $1.78$ & $2.00$\\
$1$ & no clean bipartition & $0.003$ & $0.10$\\
\bottomrule
\end{tabular}
\end{center}
Heider alone produces \emph{lopsided} blocs, since it cannot see their
sizes. A thousandth of $\hat E_1$ makes them \emph{even}, because
near-equal blocs balance far better --- which is a structural prediction
(bipolar systems should have near-equal camps) rather than a
degeneracy-breaking convenience. At $a\sim1$ the bipartition dissolves
into a frustrated balanced web. The adversivity term therefore selects
among Heider's ground states on a criterion Heider does not possess, at
every weight at which it is present.
\end{remark}

% ---------------------------------------------------------------------
\subsection{Two experiments}\label{ad:sub:experiments}

Both motivating cases are simulated by representing the common enemy as
an exogenous field: each institution carries a fixed tie $-h$ to it, so
its ledger reads $\sum_j w_{ij}s_j-h$. A first attempt that instead
\emph{initialised} an enemy node and let it optimise failed, the node
simply joining a faction; a threat must be exogenous rather than a
participant balancing its own books.

\subsubsection{Removal: the Welsh}\label{ad:subsub:welsh}

Seven institutions equilibrated at $h=1$, then relaxed at $h=0$ from
that state.

\begin{center}\small
\begin{tabular}{@{}lrrrr@{}}
\toprule
& mean tie & allied & opposed & $\hat E_1$\\
\midrule
threat present & $+0.146$ & 6 & 3 & $0.728$\\
threat removed, before relaxation & $+0.146$ & 6 & 3 & $1.476$\\
after relaxation & $+0.004$ & 4 & 4 & $0.004$\\
\bottomrule
\end{tabular}
\end{center}

The middle row is the mechanism displayed: nothing has moved, and the
same configuration has doubled in cost. The alliances were balancing the
field; without it they are surplus. Relaxation then sheds them, and the
two factions present under threat become \emph{three plus an isolate}.

\begin{remark}[Fragmentation, not repolarisation --- and no flips]
\label{ad:rem:welshreading}
Two features distinguish this from what either neighbouring theory would
predict, and the second is a limitation.

Balance theory predicts repolarisation into two fresh camps, since two
camps is what a balanced complete graph must be. The present model
predicts \emph{fragmentation}: with no external field, the cheapest
balance is many small mutual enmities rather than two large ones. The
historical record on the Welsh after the Anglo-Saxon reverse ---
increased factional subdivision without descent into general war --- is
consistent with the second and not the first. We record this as a
retrodiction and not a test: the simulation was run before the record
was consulted, but the parameters were not fixed in advance and a single
case cannot discriminate.

A third feature was first recorded here as a limitation and is better
read as a match. \emph{No alliance inverted}: no pair went from ally to
enemy. What happened instead is that two weak alliances lapsed to
neutrality while one previously neutral pair turned hostile, so
alliances fell from six to four and enmities rose from three to four.
That is coalition dissolution accompanied by fresh hostility among
parties who were already distinct --- not the betrayal of close allies.
The historical pattern is the former: the Welsh tribes were separate
before the reverse and became mutually hostile after it, rather than
turning on partners with whom they had been closely bound. A model
producing ally-to-enemy inversions would in fact overshoot the record.

The distinction is worth keeping because the two are easy to conflate
and make different predictions. Betrayal requires a commitment term
penalising intermediate $|w|$, so that ties cannot fade and must flip;
dissolution-plus-new-hostility requires no such term and is what the
present energy produces. Which of the two a given historical episode
exhibits is checkable, and the model is committed to the second.
\end{remark}

\subsubsection{Arrival: the aliens}\label{ad:subsub:aliens}

Nine institutions equilibrated alone, then a tenth appended with size
exceeding all nine combined and with its own weights \emph{frozen} at
$-1$ toward each --- it takes no view among them.

\begin{center}\small
\begin{tabular}{@{}lrrr@{}}
\toprule
& mean tie among the nine & allied pairs & $\hat E_1$\\
\midrule
alone & $+0.000$ & 10 of 36 & $0.001$\\
threat present & $+0.583$ & \textbf{28 of 36} & $8.1$\\
\bottomrule
\end{tabular}
\end{center}

Eight previously opposed pairs became allied --- inverted, not merely
softened. The reinforcement sustaining those enmities loses its force
once every ledger is dominated by a single enormous term.

\begin{remark}[Incomplete unification, and an asymmetry]
\label{ad:rem:aliensreading}
Two results here were not designed for.

\emph{Unification saturates short of completion.} Only $28$ of $36$
pairs allied, with $\hat E_1$ blown up to $8.1$: even universal alliance
supplies at most $+8$ against a threat term of $-18$, so the system
cannot balance and carries permanent strain. Overwhelming threat
produces near-total but incomplete unification, which is a more
qualified prediction than the verbal argument yields.

\emph{Arrival and removal are not symmetric.} Adding a threat produced
eight sign inversions; removing one produced none. The asymmetry follows
from size-weighting --- an adversary larger than everyone combined
dominates every ledger at once, where a unit-sized field leaves only a
small surplus to shed --- and yields the prediction that \emph{threat
arrival is a sharp transition while threat removal is a slow
relaxation}. Nothing was inserted to produce this.
\end{remark}

% ---------------------------------------------------------------------
\subsection{Toward technoanatomy}\label{ad:sub:toward}

The purpose of this apparatus is to supply what \S\ref{gr1:sec} assumes.
Internal hazard there is $\Gamma e^{-\rho x}$ with both constants
primitive; here internal hazard would be what institutions inflict on
one another, and $\rho$ would be how that scales with capability. In the
notation of \S\ref{gr:subsub:boundary}, where $\rho=-(\alpha-\beta)$
with $\alpha$ the growth of reach and $\beta$ of enforcement, the
adversivity model supplies referents: $\alpha$ is how fast the damage
one institution can do another grows with technology, and $\beta$ how
fast joint-agreement structure grows.

\begin{remark}[What is missing, stated exactly]\label{ad:rem:missing}
The section does not deliver $\rho$, and the obstruction is specific
rather than a matter of unfinished labour. This model contains conflict
dynamics and no coordination scaling: $\hat E_1$ through
$\hat E_{\mathrm b}$ describe how institutions arrange themselves given
their capacities, and nothing in them says how the \emph{cost of
maintaining joint agreement} grows as those capacities grow. That
quantity is $\beta$, and technoanatomy is entirely a claim about it.
Balance dynamics of any kind will supply $\alpha$ and be silent on
$\beta$.

The route from here is therefore twofold, and both parts are open
\textbf{[O]}. A coarse-graining must be specified --- merging
positively-tied clusters into single nodes at the next scope, with sizes
adding and weights averaging --- and the energy must be shown to keep
its form under it, so that every scope obeys the same dynamics and the
biosphere's $\rho$ is the top of a flow rather than a separate
posit. And a coordination cost must be added, whose scaling with
capability is the anatomy the whole construction is for.
\end{remark}

% ---------------------------------------------------------------------
\subsection{Ledger}\label{ad:sub:ledger}

\paragraph{Costs.} \emph{The principle} (\ref{ad:phi:balance}) ---
\textbf{governing premise}; \emph{falsifier:} an institutional system
observed to rest stably in universal alliance without an external
threat, or to be indifferent to the balance of its commitments.
\emph{The engagement budget} --- \textbf{structural premise}, without
which the principle is trivially satisfied by indifference;
\emph{falsifier:} institutions observed to hold arbitrarily many ties
at full strength. \emph{The weighting $b/a$} --- the model's \textbf{sole free
parameter}, controlling the multipolar-to-bipolar transition;
\emph{falsifier:} it is fitted rather than derived, and a system whose
faction structure is inconsistent with any single value across eras
would embarrass the model. \emph{Mutuality and finite attention} ---
\textbf{constraints}, not preferences; \emph{falsifiers:} stable
one-sided regard between institutions, and institutions holding
arbitrarily many ties at full strength, respectively. \emph{Completeness} --- the two-faction
structure theorem assumes a complete graph, and real institutional
networks are sparse; the incomplete case admits more than two factions
and the clean result is unavailable.

\paragraph{Buys.} A single principle reproducing both motivating cases,
with the removal case matching the historical pattern more closely than
either neighbouring theory. A counting constraint that produces
scapegoat structure without it being posited. An arrival/removal
asymmetry that was not designed for. A construction with one free
parameter and two interpretable constraints, the four-channel form used
for computation being a relaxation that agrees with it. And a stated,
checkable route to $\rho$, together with an exact account of why this
apparatus cannot reach it alone.

% ============================ 

% ---------------------------------------------------------------------
\appendix

\section{Horizon theory}\label{app:horizon}

% ---------------------------------------------------------------------
\label{sh:sub:horizon}

\subsection{The driftless horizon is a timescale, and it is maximally
soft}\label{sh:subsub:softhorizon}

In the driftless model the naive horizon quantity
$\Pi(\mu,D):=\PP_{(\mu,D)}(D\text{ never returns to }0)$ vanishes
identically (Theorem~\ref{sh:thm:noimm}): null recurrence leaves no escape
probability to be gradual about, and the horizon must be recast in the only
currency the model possesses --- time.

\begin{proposition}[Recovery cost is sublinear in credit; \textbf{P}]
\label{sh:prop:recovery}
Let $\rho(d)$ denote the median of the return time
$\inf\{t:D_t\le0\}$ from $(\mu,D)=(0,d)$, $d>0$. Then exactly
\begin{equation}\label{sh:eq:recovery}
\rho(d)=\rho(1)\,\Bigl(\frac d\sigma\Bigr)^{2/3}\sigma^{2/3}\cdot
\sigma^{-2/3}=\rho_1\Bigl(\frac d\sigma\Bigr)^{2/3},
\qquad \rho_1:=\rho(1)\big|_{\sigma=1}\in(0,\infty),
\end{equation}
and the return time is finite a.s.\ with tail index $\tfrac14$. Doubling
accumulated credit multiplies the recovery timescale by only
$2^{2/3}\approx1.59$: credit is far easier to undo than a linear intuition
suggests.
\end{proposition}

\begin{proof}
Finiteness a.s.\ is contained in the proof of
Theorem~\ref{sh:thm:noimm} (the set $A$ there is unbounded). The scaling
identity is Proposition~\ref{sh:prop:selfsim} with $c=d^{-1/3}$ applied to
the median, a scale-equivariant functional; positivity and finiteness of
$\rho_1$ follow since the return time is a nondegenerate positive random
variable. The tail index is Theorem~\ref{sh:thm:persist} read for the
returning coordinate.
\end{proof}

\begin{proposition}[Horizon level sets are $\kappa$-rays; \textbf{P}]
\label{sh:prop:rays}
Every scale-invariant function on the open quadrant is constant exactly on
the curves $D\propto\mu^3/\sigma^2$; if a horizon exists in any
scale-invariant sense, it is a ray of constant $\kappa$, and $\kappa$ is
its natural transverse coordinate.
\end{proposition}

\begin{proof}
Proposition~\ref{sh:prop:gauge}(iii).
\end{proof}

\begin{definition}[Deadline-relative horizon and its width]
\label{sh:def:deadline}
Fix a time budget $\mathcal T$. The \emph{deadline horizon} is the level
set $\{t_*(\mu,D)\asymp\mathcal T\}$, located at
$\mu^*\asymp\sigma\mathcal T^{1/2}$, $D^*\asymp\sigma\mathcal T^{3/2}$.
Its \emph{softness} is the logarithmic sensitivity of survival to state.
\end{definition}

\begin{corollary}[Maximal softness; \textbf{A}]\label{sh:cor:soft}
Under Claim~\ref{sh:claim:tail},
$\dfrac{\partial\log\PP(T>\mathcal T)}{\partial\log t_*}=\dfrac14+o(1)$.
Changing survival probability by a factor $e$ therefore requires changing
$t_*$ by $e^4\approx55$ and, since $t_*\propto D^{2/3}$ in the
credit-dominated regime, changing $D$ by $e^6\approx403$. The driftless
horizon is smeared over roughly two orders of magnitude in timescale and
nearly three in credit: there is no regime in which resistance has
arbitrarily little effect, only a $\tfrac14$-power degradation --- about as
gentle a degradation as a power law permits.
\end{corollary}

\subsection{The scale-function criterion: P\'olya's theorem, correctly
translated}\label{sh:subsub:scale}

Let the driver be a regular diffusion on $\R$,
\begin{equation}\label{sh:eq:generaldriver}
d\mu_t=b(\mu_t)\,dt+\varsigma(\mu_t)\,dW_t,\qquad
s(\mu):=\int_{0}^{\mu}\exp\Bigl(-\int_{0}^{v}\frac{2b(w)}{\varsigma^2(w)}
\,dw\Bigr)dv,
\end{equation}
with $\varsigma>0$ locally bounded away from $0$ and $\infty$; $s$ is the
scale function and Feller's test (see \cite[Ch.~15]{KT81} or
\cite[Ch.~VII]{RY99}) classifies the boundary behavior.

\begin{theorem}[Trichotomy; \textbf{P}]\label{sh:thm:trichotomy}
For the model \eqref{sh:eq:model} with driver
\eqref{sh:eq:generaldriver}:
\begin{enumerate}[label=\emph{(\roman*)},itemsep=1pt]
\item If $s(+\infty)<\infty$ and $s(-\infty)=-\infty$, then
$\mu_t\to+\infty$ a.s., $\Lambda_\infty<\infty$ a.s., $\PP(\Imm)=1$, and
$\PP(T=\infty)=\E[e^{-\Lambda_\infty}]\in(0,1)$.
\item If $s(-\infty)>-\infty$ and $s(+\infty)=\infty$, then
$\mu_t\to-\infty$ a.s., $\Lambda_\infty=\infty$ a.s., and $T<\infty$ a.s.
\item If both limits are finite, then
\begin{equation}\label{sh:eq:scaleformula}
\PP(\Imm)\;=\;\frac{s(\mu_0)-s(-\infty)}{s(+\infty)-s(-\infty)},
\end{equation}
with $\Lambda_\infty<\infty$ on $\{\mu_t\to+\infty\}$ and
$\Lambda_\infty=\infty$ on the complementary event
$\{\mu_t\to-\infty\}$.
\item If both limits are infinite, the driver is recurrent and the outcome
is \emph{not} determined by $s$; see
Proposition~\ref{sh:prop:recurrent}.
\end{enumerate}
\end{theorem}

\begin{proof}
(i) $\PP(\lim_t\mu_t=+\infty)=1$ is Feller's test \cite{KT81}. On this
event choose $t_0(\omega)$ with $\mu_t\ge1$ for $t\ge t_0$; then
$D_t\ge D_{t_0}+(t-t_0)$, so
$\Lambda_\infty\le\Lambda_{t_0}+h_0e^{-(D_{t_0}-D_0)}\int_0^\infty
e^{-v}dv<\infty$ a.s. Positivity of $\PP(T=\infty)$ follows since
$e^{-\Lambda_\infty}>0$ a.s.; it is $<1$ since $\Lambda_\infty\ge
\Lambda_1>0$ a.s. (ii) is the mirror image: eventually $\mu_t\le-1$, so
$e^{-(D_t-D_0)}\to\infty$ and $\Lambda_\infty=\infty$; death then occurs in
finite time a.s.\ by Definition~\ref{sh:def:cox}. (iii) The hitting
probabilities of $\pm\infty$ are the classical scale-function formula, and
the two events partition the space up to null sets; apply the pathwise
arguments of (i) and (ii) on each. (iv) Recurrence under unbounded scale is
Feller's test again; Example~\ref{sh:ex:recurrentcounter} shows $s$ alone
cannot decide.
\end{proof}

\begin{proposition}[Recurrent drivers; \textbf{P}]\label{sh:prop:recurrent}
Suppose the driver \eqref{sh:eq:generaldriver} is positive recurrent with
stationary law $\pi$ and $\int|\mu|\,\pi(d\mu)<\infty$; write
$\bar m:=\int\mu\,\pi(d\mu)$.
\begin{enumerate}[label=\emph{(\alph*)},itemsep=1pt]
\item If $\bar m>0$ then $D_t/t\to\bar m$ a.s., $\Lambda_\infty<\infty$
a.s., and $\PP(\Imm)=1$.
\item If $\bar m<0$ then $\Lambda_\infty=\infty$ a.s.\ and $T<\infty$ a.s.
\item If $\bar m=0$ and the cycle integral defined in the proof is
nondegenerate (automatic for nondegenerate diffusions), then
$\Lambda_\infty=\infty$ a.s.: \emph{the critical case falls on the mortal
side.}
\end{enumerate}
For null-recurrent drivers no general statement is offered
\textbf{[O]}; the Brownian driver, the fundamental null-recurrent case, is
settled unconditionally by Theorem~\ref{sh:thm:noimm}.
\end{proposition}

\begin{proof}
(a), (b): the ergodic theorem for positive recurrent one-dimensional
diffusions gives $D_t/t\to\bar m$ a.s.\ (see \cite[Ch.~X]{RY99}); if
$\bar m>0$ then eventually $D_t\ge D_0+\tfrac{\bar m}2t$ and
$\Lambda_\infty<\infty$; if $\bar m<0$ the integrand
$e^{-(D_t-D_0)}\to\infty$.

(c) Fix interior levels $a<b$ and regenerate at up-crossings: let
$R_0:=\inf\{t:\mu_t=b\}$ and inductively
$S_k:=\inf\{t>R_{k-1}:\mu_t=a\}$, $R_k:=\inf\{t>S_k:\mu_t=b\}$. By the
strong Markov property the cycles over $[R_{k-1},R_k)$ are i.i.d.; positive
recurrence gives $\E[L]<\infty$ for the cycle length
$L:=R_1-R_0$, and the cycle occupation identity yields, for the cycle
integral $\xi_k:=\int_{R_{k-1}}^{R_k}\mu_s\,ds$,
\[
\E[\xi_1]=\E[L]\int\mu\,d\pi=0,\qquad
\E|\xi_1|\le\E\Bigl[\int_{\text{cycle}}|\mu_s|ds\Bigr]
=\E[L]\int|\mu|\,d\pi<\infty .
\]
Thus $S_n:=D_{R_n}-D_{R_0}=\sum_{k\le n}\xi_k$ is a nondegenerate mean-zero
i.i.d.\ walk, so $\liminf_nS_n=-\infty$ a.s.\ by the classical
oscillation dichotomy for random walks (see, e.g., \cite{Dur19});
in particular the events $B_n:=\{D_{R_n}\le-1\}$ occur infinitely often
a.s. Choose $\delta>0$ and $M<\infty$ with
\[
q:=\PP_b\Bigl(L\ge\delta,\ \sup_{0\le u\le\delta}
\Bigl|\int_0^u\mu_s\,ds\Bigr|\le M\Bigr)>0,
\]
possible since $L>0$ a.s.\ and the supremum is finite a.s. Let
$G_n$ be the corresponding event for cycle $n{+}1$; by the strong Markov
property $\PP(G_n\mid\Fc_{R_n})=q$, so
$\sum_n\PP(B_n\cap G_n\mid\Fc_{R_n})=q\sum_n\one_{B_n}=\infty$ a.s., and
L\'evy's extension of the Borel--Cantelli lemma (see, e.g.,
\cite{Dur19}) gives
$B_n\cap G_n$ infinitely often a.s. On $B_n\cap G_n$, for
$u\in[0,\delta]$ one has $D_{R_n+u}\le-1+M$, hence
\[
\int_{R_n}^{R_n+\delta}e^{-(D_s-D_0)}\,ds\ \ge\ \delta\,e^{\,1-M+D_0}\,>\,0 ,
\]
and these windows are disjoint. Summing over infinitely many of them,
$\Lambda_\infty=\infty$ a.s.
\end{proof}

\begin{example}[Recurrence alone does not decide]\label{sh:ex:recurrentcounter}
The mean-reverting Ornstein--Uhlenbeck driver
$d\mu=\theta(\bar\mu-\mu)dt+\sigma dW$ with $\theta>0$ has Gaussian-growing
$s'$, hence $s(\pm\infty)=\pm\infty$: it is (positive) recurrent for every
$\bar\mu$. Yet Proposition~\ref{sh:prop:recurrent} gives
$\PP(\Imm)=1$ for $\bar\mu>0$ and certain death for $\bar\mu\le0$. A
horizon exists here --- sharp, at $\bar\mu=0$ --- but it lives in
\emph{parameter} space, not in the initial condition, which is again pure
gauge.
\end{example}

\begin{remark}[Two warnings]\label{sh:rem:warnings}
Two traps in this circle of ideas deserve explicit flags.
\emph{First}, it is tempting to summarize \eqref{sh:eq:scaleformula} as
``recurrent driver $\Rightarrow$ certain death''; Example
\ref{sh:ex:recurrentcounter} shows this is false --- in the recurrent class
the order parameter is the stationary mean $\bar m$, with the critical case
$\bar m=0$ mortal, and only in the null-recurrent Brownian case is
symmetry itself the whole story.
\emph{Second}, under adverse drift it is tempting to read the survival tail
off the typical path, which yields doubly-exponential decay; in fact the
tail is only \emph{exponential}, dominated by an optimal ``fighting-back''
fluctuation --- Proposition~\ref{sh:prop:advdrift} below --- a
Freidlin--Wentzell phenomenon in miniature.
\end{remark}

\begin{center}
\begin{tabular}{@{}ll@{}}
\toprule
\textbf{Random walk (P\'olya)} & \textbf{Hazard dynamics}\\
\midrule
dimension $d$ & scale function $s$ of the driver\\
$d\le2$: recurrent, no escape & $s(\pm\infty)=\pm\infty$: no viability from
data\\
$d\ge3$: escape probability $>0$ & $s(\pm\infty)$ finite: viability
probability in $(0,1)$\\
$d=2$: null recurrent & Brownian driver: certain death, infinite mean
life\\
escape probability from $x$ & normalized scale function at $\mu_0$\\
\bottomrule
\end{tabular}
\end{center}

\begin{remark}[Slogan]\label{sh:rem:slogan}
\emph{Existence} of a horizon in the initial condition
$=$ boundedness of the scale function;
\emph{sharpness} of the horizon $=$ steepness of the scale function.
The Brownian driver has $s(\mu)=\mu$, unbounded and of constant slope: its
horizon degenerates not by becoming infinitely sharp but by becoming
infinitely smeared --- a linear scale function distinguishes no location.
Finally, recall from \eqref{sh:eq:twotier} that
\eqref{sh:eq:scaleformula} computes \emph{viability}; immortality proper is
$\E[e^{-\Lambda_\infty}\one_\Imm]$, strictly smaller, and this is the
unique juncture at which the gauge parameter $h_0$ re-enters
(Lemma~\ref{sh:lem:hgauge}) --- affecting values, never dichotomies.
\end{remark}

\subsection{An exactly solvable gradual horizon}\label{sh:subsub:solvable}

\begin{theorem}[Unstable Ornstein--Uhlenbeck driver; \textbf{P}]
\label{sh:thm:uou}
Let the driver be
\begin{equation}\label{sh:eq:uou}
d\mu_t=\nu\,\mu_t\,dt+\sigma\,dW_t,\qquad \nu>0 .
\end{equation}
Then $e^{-\nu t}\mu_t\to Z$ a.s.\ and in $L^2$, where
$Z\sim N\!\bigl(\mu_0,\ \sigma^2/2\nu\bigr)$, and
\begin{equation}\label{sh:eq:sigmoid}
\PP(\Imm)\;=\;\PP(Z>0)\;=\;\Phi\!\Bigl(\frac{\mu_0\sqrt{2\nu}}{\sigma}\Bigr).
\end{equation}
On $\{Z>0\}$ the hazard decays doubly exponentially,
$\log h_t\sim-(Z/\nu)e^{\nu t}$, and $\Lambda_\infty<\infty$; on
$\{Z<0\}$ the hazard explodes doubly exponentially and $T<\infty$ a.s.
\end{theorem}

\begin{proof}
Variation of constants gives
$\mu_t=e^{\nu t}\bigl(\mu_0+\sigma\int_0^te^{-\nu s}dW_s\bigr)
=:e^{\nu t}M_t$. The martingale $M$ has
$\E\langle M\rangle_\infty=\sigma^2\!\int_0^\infty e^{-2\nu s}ds
=\sigma^2/2\nu<\infty$, hence converges a.s.\ and in $L^2$ to a Gaussian
limit $Z$ with the stated mean and variance. On $\{Z>0\}$ pick
$t_0(\omega)$ with $\mu_t\ge\tfrac Z2e^{\nu t}$ for $t\ge t_0$; then
$D_t\ge D_{t_0}+\tfrac{Z}{2\nu}(e^{\nu t}-e^{\nu t_0})$ and the integrand
$e^{-(D_t-D_0)}$ is dominated by a doubly-exponentially decaying function,
so $\Lambda_\infty<\infty$. On $\{Z<0\}$ symmetrically
$\mu_t\le\tfrac Z2e^{\nu t}$ eventually, $D_t\to-\infty$ doubly
exponentially, the integrand diverges, $\Lambda_\infty=\infty$, and
$T<\infty$ a.s. Finally $\PP(Z=0)=0$, and
$\PP(Z>0)=\Phi(\mu_0/\sqrt{\sigma^2/2\nu})$.
\end{proof}

\begin{remark}[Consistency with the scale-function formula]
\label{sh:rem:crosscheck}
For \eqref{sh:eq:uou}, $s'(v)=\exp(-\nu v^2/\sigma^2)$, so
$s(\pm\infty)$ are finite and the normalized scale function at $\mu_0$ is
the mass below $\mu_0$ of the centered Gaussian with variance
$\sigma^2/2\nu$ --- precisely \eqref{sh:eq:sigmoid}. The exactly solvable
case is thus the generic case of
Theorem~\ref{sh:thm:trichotomy}(iii) in its purest form.
\end{remark}

\begin{definition}[Horizon profile, location, width]\label{sh:def:width}
For a driver with bounded scale, the \emph{horizon profile} is
$\Pi(\mu_0):=\PP(\Imm\mid\mu_0)$, the \emph{location} its median point
$\Pi=\tfrac12$, and the \emph{width} the natural spread of the transition,
e.g.\ the interquartile range of the measure $d\Pi$. For
\eqref{sh:eq:uou}: location $\mu_0=0$; width
\begin{equation}\label{sh:eq:width}
w=\frac{\sigma}{\sqrt{2\nu}},\qquad\text{sharpness}\quad
\Sigma:=w^{-1}=\frac{\sqrt{2\nu}}{\sigma}.
\end{equation}
As $\Sigma\to\infty$ (small noise or strong instability),
$\Pi\to\one\{\mu_0>0\}$: a hard event horizon. As $\Sigma\to0$,
$\Pi\to\tfrac12$ everywhere: no horizon, a fair coin.
\end{definition}

\begin{remark}[Sharpness is a small-noise limit; general principle]
\label{sh:rem:FW}
The width \eqref{sh:eq:width} is exactly the scale of fluctuation of the
terminal charge $Z$ capable of overturning the deterministic verdict
$\operatorname{sign}(\mu_0)$: the sharp horizon is the Freidlin--Wentzell
limit \cite{FW12} of the soft family. We record the principle in the form
answering the question that motivates this part: \emph{a horizon is
never intrinsically sharp; its width is the fluctuation scale of the
driver, and it appears sharp exactly when that scale is small against the
spread of initial conditions.} Operationally, comparing $w$ to the
population spread of $\mu_0$ yields three regimes: $w\ll$ spread --- sharp
horizon, the population splits into the fated; $w\sim$ spread --- genuinely
gradual horizon, initial conditions matter but do not determine;
$w\gg$ spread --- no horizon, the outcome is noise.
\end{remark}

\subsection{Model atlas and the role of driver smoothness}
\label{sh:subsub:atlas}

\begin{proposition}[Adverse linear drift: the price of survival;
\textbf{P/S}]\label{sh:prop:advdrift}
Let $\mu_t=\mu_0+\nu t+\sigma W_t$ with $\nu<0$. Then:
\begin{enumerate}[label=\emph{(\roman*)},itemsep=1pt]
\item \textbf{[P]}
$\displaystyle\liminf_{t\to\infty}\frac1t\log\PP(T>t)\ \ge\
-\frac{\nu^2}{2\sigma^2}$;
\item \textbf{[P]} there exists $c>0$ with $\PP(T>t)\le e^{-ct}$ for all
large $t$;
\item \textbf{[S]} the exact rate is
\begin{equation}\label{sh:eq:rate}
\lim_{t\to\infty}\frac1t\log\PP(T>t)\ =\ -\frac{3\nu^2}{8\sigma^2},
\end{equation}
attained by the optimal survival profile
$\mu_{tv}\approx\tfrac{|\nu|t}{4}\,v\,(2-3v)$,
$D_{tv}\approx\tfrac{|\nu|t^2}{4}\,v^2(1-v)$, $v\in[0,1]$: the conditioned
survivor front-loads credit, lets the rate go negative at $v=\tfrac23$, and
coasts so as to exhaust its credit exactly at the deadline.
\end{enumerate}
\end{proposition}

\begin{proof}
(i) By the Markov property and a finite-time event of positive probability
we may assume $\mu_0\ge1$; constants below depend on the data. Let
$a:=|\nu|/\sigma$, $b_0:=(1-\mu_0)/\sigma\le0$, and
\[
A_t:=\Bigl\{W_u\ge b_0+au\ \ \forall u\in[0,t],\ \ W_t\le b_0+at+1\Bigr\}.
\]
On $A_t$, $\mu_u\ge1$ for all $u\le t$, hence $D_u\ge D_0+u$ and
$\Lambda_t\le h_0$, so $\PP(T>t)\ge e^{-h_0}\PP(A_t)$. Girsanov with drift
$a$ (density $\exp(-a\widetilde W_t-a^2t/2)$, $\widetilde W:=W-a\,\cdot$
a Brownian motion under the tilted measure $Q$) gives
\[
\PP(A_t)\ \ge\ e^{-a^2t/2}\,e^{-a(b_0+1)}\,
Q\bigl(\widetilde W_u\ge b_0\ \forall u\le t,\ \widetilde W_t\in[b_0,b_0+1]
\bigr),
\]
and the $Q$-probability is of order $t^{-3/2}$ by the reflection
principle: polynomial corrections do not affect the rate.
(ii) With $K:=\tfrac{h_0}4t$,
$\PP(T>t)\le e^{-K}+\PP(\Lambda_t\le K)$. On $\{\Lambda_t\le K\}$,
$\Leb\{u\le t:D_u\le D_0\}\le K/(h_0e^{-D_0}\wedge h_0)\lesssim t/4$ (adjust
constants by $D_0$), so there exists
$u_*\in[t/2,\,3t/4+1]$ with $D_{u_*}>D_0$, forcing
$\sigma I_{u_*}\ge|\nu|u_*^2/2-\mu_0u_*\ge|\nu|t^2/16$ for large $t$. By
the Borell--TIS inequality for the Gaussian process $I$ on $[0,t]$ (whose
maximal variance is $t^3/3$ and whose expected supremum is
$O(t^{3/2})=o(t^2)$),
$\PP(\sup_{u\le t}I_u\ge|\nu|t^2/16\sigma)\le
\exp(-c't^4/t^3)=e^{-c't}$; combine.
(iii) Rescale $W_{tv}=\sqrt t\,B_v$, so
$I_{tv}=t^{3/2}\int_0^vB$ and
$D_{tv}=\mu_0tv+\tfrac{\nu t^2v^2}2+\sigma t^{3/2}\!\int_0^vB$. Keeping
$\Lambda_t=O(1)$ requires, at leading order, $D_{tv}\ge0$ for all
$v\in[0,1]$; with $B=\sqrt t\,\gamma$ this reads
$\int_0^v\gamma\ge\beta v^2$, $\beta:=|\nu|/2\sigma$, at
Cameron--Martin cost $\exp(-t\cdot\tfrac12\int_0^1\dot\gamma^2)$. The
variational problem
$\min\tfrac12\int_0^1\dot\gamma^2$ subject to
$\int_0^v\gamma\ge\beta v^2\ \forall v$, $\gamma(0)=0$, is solved by
relaxing to the endpoint constraint
$\int_0^1\gamma=\int_0^1(1-s)\dot\gamma(s)\,ds\ge\beta$: Cauchy--Schwarz
gives $\int\dot\gamma^2\ge3\beta^2$ with equality for
$\dot\gamma(s)=3\beta(1-s)$, i.e.\ $\gamma(v)=3\beta(v-\tfrac{v^2}2)$,
which satisfies the \emph{full} constraint,
$\int_0^v\gamma=\tfrac{\beta v^2}2(3-v)\ge\beta v^2$ on $[0,1]$ with
equality only at $v=1$. Hence the minimum is $\tfrac32\beta^2
=3\nu^2/8\sigma^2$, and unwinding the substitutions yields the profiles
stated: $\mu_{tv}=\mu_0+t(\nu v+\sigma\gamma(v))
\approx\tfrac{|\nu|t}4v(2-3v)$ and
$D_{tv}\approx t^2\bigl(\sigma\int_0^v\gamma-\tfrac{|\nu|v^2}2\bigr)
=\tfrac{|\nu|t^2}4v^2(1-v)$. Elevating this computation to a full
large-deviation proof (upper bound with matching constant, treatment of the
soft constraint) is routine but lengthy and is not carried out here.
\end{proof}

\begin{remark}\label{sh:rem:advdrift}
Three comments. \emph{(1)} The naive lower-bound strategy ``hold
$\mu\ge1$ throughout,'' used in (i), costs $\nu^2/2\sigma^2$ per unit time
and is strictly suboptimal: the true optimizer spends the final third of
the budget with $\mu<0$, spending credit rather than defending the rate ---
survival to a deadline is cheaper than survival in perpetuity, and the
discount is exactly the factor $\tfrac34$. \emph{(2)} Conditioned on the
rare event $\{T>t\}$, the path concentrates on the stated profile: once
again the appearance of purposeful resistance is manufactured by
conditioning. \emph{(3)} The analogue for a mean-reverting driver with
$\bar\mu<0$ is an exponential rate given by the corresponding
Freidlin--Wentzell action; we do not compute it here.
\end{remark}

\begin{center}
\small
\begin{tabular}{@{}lllll@{}}
\toprule
\textbf{Driver} $\mu_t$ & $\PP(\Imm)$ & \textbf{tail of }$\PP(T>t)$ &
$\E[T]$ & \textbf{horizon}\\
\midrule
constant $\mu_0>0$ & $1$ & $\to e^{-h_0/\mu_0}>0$ & $\infty$ (atom) &
trivial step\\
constant $\mu_0<0$ & $0$ & $\exp\bigl(-\tfrac{h_0}{|\mu_0|}
e^{|\mu_0|t}\bigr)$ & finite & ---\\
$\mu_0+\sigma W_t$ & $0$ & $t^{-1/4}$ \textbf{[A]} & $\infty$ & none
(null rec.)\\
$\mu_0+\nu t+\sigma W_t$, $\nu>0$ & $1$ & $\to$ const $>0$ & $\infty$ &
dichotomy in $\nu$\\
$\mu_0+\nu t+\sigma W_t$, $\nu<0$ & $0$ &
$e^{-\frac{3\nu^2}{8\sigma^2}t(1+o(1))}$ \textbf{[S]} & finite & ---\\
OU, mean $\bar\mu>0$ & $1$ & $\to$ const $>0$ & $\infty$ & sharp in
$\bar\mu$\\
OU, mean $\bar\mu=0$ & $0$ & $t^{-1/2}$ \textbf{[S]} & $\infty$ & null
rec., new exponent\\
OU, mean $\bar\mu<0$ & $0$ & $e^{-ct}$ & finite & ---\\
unstable OU \eqref{sh:eq:uou} & $\Phi\bigl(\mu_0\sqrt{2\nu}/\sigma\bigr)$
& mixture & $\infty$ & \textbf{gradual, width }$\sigma/\sqrt{2\nu}$\\
general bounded scale & \eqref{sh:eq:scaleformula} & mixture & $\infty$ &
shape $=$ scale fn.\\
\bottomrule
\end{tabular}
\end{center}

\begin{remark}[Exponents measure smoothness, not recurrence class]
\label{sh:rem:smoothness}
The mean-zero mean-reverting driver is mortal
(Proposition~\ref{sh:prop:recurrent}(c)) with expected tail $t^{-1/2}$, not
$t^{-1/4}$: at large scales its credit is Brownian-like (an additive
functional obeying an invariance principle \cite{Bha82}), so the relevant
persistence exponent is that of Brownian motion, $\tfrac12$, not of its
integral, $\tfrac14$. The survival exponent is a \emph{persistence exponent
of the credit process} and is set by the smoothness of the driver rather
than by its recurrence class: smoother driver, smaller exponent, heavier
survival tail. For a fractional Brownian driver of Hurst index $H$ the
persistence exponent of the driver itself is $1-H$ \cite{Mol99}; the
exponent for its integral --- the relevant one here --- is not known
rigorously \textbf{[O]}: the numerical evidence of Molchan--Khokhlov
\cite{MK04} supports the conjecture $\theta(H)=H(1-H)$ for persistence
from the origin, which correctly returns $\tfrac14$ at $H=\tfrac12$;
see also \cite{AS15}.
\end{remark}

% ---------------------------------------------------------------------
\section{The background as a free energy}\label{app:background}

% ---------------------------------------------------------------------
\label{tc:sub:background}

The background must be positive, variable, and never near zero.
Construction \eqref{tc:eq:B} achieves all three: a sum of $n$ positive
terms is small only if \emph{all} are small, an event of probability
$\Phi(-a/\epsilon\sqrt t)^{n}$, exponentially small in $n$. This is the
mechanism the author identified. What was not anticipated is that the
construction is a familiar object.

\subsection{Freezing and Gompertz}\label{tc:subsub:freezing}

\begin{remark}[The background is a partition function]
\label{tc:rem:partition}
Writing $\epsilon W^{(i)}_t=\epsilon\sqrt t\,g_i$ with $g_i$ i.i.d.\
standard normal,
\[
B_t=\sum_{i=1}^{n}e^{\beta_t g_i},\qquad \beta_t:=\epsilon\sqrt t,
\]
the partition function of $n$ i.i.d.\ Gaussian energy levels at inverse
temperature $\beta_t$. Since $\beta_t$ \emph{increases} with $t$, the
background is a disordered system that \emph{cools as it ages}: it is
Derrida's random energy model \cite{Der81,Der80} with time as inverse
temperature. The corresponding limit theory for sums of random
exponentials is \cite{BBM05}; the mathematical treatment of the REM is
\cite{Bov06}.
\end{remark}

\begin{theorem}[Freezing; \textbf{C}]\label{tc:thm:freezing}
With $\theta=\epsilon\sqrt{t/\log n}$, as $n\to\infty$ at fixed
$\theta$,
\begin{equation}\label{tc:eq:freezing}
\frac{\log B_t}{\log n}\;\longrightarrow\;
\begin{cases}
1+\theta^{2}/2, & \theta<\sqrt2\quad(\text{delocalised}),\\
\theta\sqrt2, & \theta\ge\sqrt2\quad(\text{frozen}),
\end{cases}
\end{equation}
in probability, with freezing time
$t_{\mathrm{fr}}=2\log n/\epsilon^{2}$.
\end{theorem}

\begin{proof}[Proof and status]
The two candidates are the annealed free energy
$\log\E B_t=\log n+\epsilon^{2}t/2$ and the extremal contribution
$\max_i\epsilon W^{(i)}_t\approx\epsilon\sqrt{2t\log n}$; normalised
these are $1+\theta^{2}/2$ and $\theta\sqrt2$, and
$(1+\theta^{2}/2)-\theta\sqrt2=\tfrac12(\theta-\sqrt2)^{2}\ge0$ with
equality only at $\theta=\sqrt2$ --- the annealed value dominates
everywhere and the curves meet tangentially, the REM's characteristic
contact. The second moment gives
$\E B_t^{2}/(\E B_t)^{2}=e^{\epsilon^{2}t}/n+(n-1)/n\to1$ iff
$\epsilon^{2}t<\log n$, so concentration is immediate for
$t<t_{\mathrm{fr}}/2$; the full delocalised range needs the truncated
second moment, and the frozen phase is the standard extreme-value
computation. Both are classical \cite{Der81,Bov06,BBM05}.
\end{proof}

\begin{remark}[The author's two mechanisms are the two phases]
\label{tc:rem:twophases}
The brief describes the background as held up both by ``the high ones
dominating'' and by ``the fuzz of all the intermediate-sized ones adding
together.'' These are not two descriptions of one regime but exact
descriptions of the two: the frozen phase is carried by the single
largest term, the delocalised phase by the bulk.
\end{remark}

\begin{corollary}[Gompertz; \textbf{P}]\label{tc:cor:gompertz}
For $t<t_{\mathrm{fr}}$, $B_t=n\,e^{\epsilon^{2}t/2}(1+o(1))$ in
probability: the background grows exponentially in age at rate
$\epsilon^{2}/2$. This is the Gompertz law \cite{Gom25}, arriving
unbidden --- the construction was chosen for positivity, not to
reproduce it.
\end{corollary}

\subsection{The background is gauge}\label{tc:subsub:gauge}

\begin{lemma}[Channels define moving walls; \textbf{P}]
\label{tc:lem:walls}
Set $D_t:=x_t-x_0$. Each channel becomes $O(1)$ --- i.e.\ contributes
appreciably to $\Lambda$ --- at a \emph{killing wall}:
\begin{equation}\label{tc:eq:killwalls}
x^{\mathrm{crest}}_t=\log B_t+b,
\qquad
x^{\mathrm{trough}}=\frac{\log\Gamma}{\rho},
\end{equation}
the first random and slowly moving, the second a fixed constant. the driftless case's soft-wall reduction applies with its fixed wall replaced by
$\max$ of these.
\end{lemma}

\begin{proof}
Solve $B_te^{\,b-x}=1$ and $\Gamma e^{-\rho x}=1$.
\end{proof}

\begin{theorem}[Exact absorption of a linear background; \textbf{P}]
\label{tc:thm:absorb}
Suppose $\log B_u=a+cu$ exactly --- the Gompertz phase, $a=\log n$,
$c=\epsilon^{2}/2$. Then persistence above the crest wall
\eqref{tc:eq:killwalls} is \emph{identical} to the driftless case's fixed-wall
problem with initial data
\begin{equation}\label{tc:eq:shift}
\mu_0\longmapsto\mu_0-c,\qquad D_0\longmapsto D_0-a-b .
\end{equation}
By that section's gauge proposition the background therefore moves
prefactors and timescales alone, and no asymptotic exponent.
\end{theorem}

\begin{proof}
With $D_u=D_0+\mu_0u+\sigma I_u$, $I_u=\int_0^u W_s\,ds$, the constraint
$D_u\ge a+b+cu-x_0$ reads
$\sigma\bigl(I_u-\tfrac{c-\mu_0}{\sigma}u\bigr)\ge a+b-x_0-D_0$, and
$I_u-\kappa u=\int_0^u(W_s-\kappa)ds$ is the integral of a Brownian
motion started at $-\kappa$: the same Kolmogorov diffusion from shifted
data.
\end{proof}

\begin{remark}[A headwind]\label{tc:rem:headwind}
An exponentially growing background is a \emph{constant headwind} of
size $\epsilon^{2}/2$ on the improvement rate: the system must improve
at that rate merely to hold its hazard fixed, and its initial credit is
debited by $\log n+b$. Ageing, here, is not a new mechanism but an
offset.
\end{remark}

\begin{theorem}[Sublinear walls are invisible; \textbf{A}]
\label{tc:thm:sublinear}
If both walls are $o(u^{3/2})$ almost surely --- which holds here, the
crest wall being $O(u)$ in the delocalised phase and $O(\sqrt u)$ in the
frozen phase, and the trough wall constant --- then, granting the
moving-boundary robustness of the driftless case's Ansatz class, the
persistence exponent is $\tfrac14$, unchanged.
\end{theorem}

\begin{proof}[Sketch and status]
After rescaling by $t^{3/2}$ the boundary converges to zero uniformly on
compacts, so the constraint converges to the fixed-wall constraint.
Converting this to an exponent identity is the moving-boundary
robustness discussed in the driftless case and surveyed in \cite{AS15}.
\end{proof}

% ---------------------------------------------------------------------
\section{The waiting room and the gain profile}\label{app:room}

% ---------------------------------------------------------------------
\label{tc:sub:room}

\subsection{Two kinds of wall}\label{tc:subsub:twowalls}

A distinction must be drawn that the informal account elides. There are
two different senses in which a channel ``has a wall'', and both are
needed.

\begin{definition}[Shape and killing walls]\label{tc:def:walls}
The \emph{shape walls} are where a channel rises above the background
floor, i.e.\ where it becomes comparable to $B_t$:
\begin{equation}\label{tc:eq:shapewalls}
x^{\mathrm{crest}}_{\mathrm{shape}}\approx b,
\qquad
x^{\mathrm{trough}}_{\mathrm{shape}}\approx-\frac{k_t}{\rho},
\qquad
k_t:=\log\frac{\E[B_t]}{\Gamma},
\end{equation}
so that the room is the \emph{asymmetric} interval $[-k_t/\rho,\,b]$,
containing the origin in its interior.
The \emph{killing walls} are where a channel makes $\Lambda$ accumulate,
i.e.\ where it becomes $O(1)$ in absolute terms; these are
\eqref{tc:eq:killwalls}. The shape walls govern what the hazard
\emph{looks like} as a function of $x$; the killing walls govern
the driftless case's persistence geometry. They are different objects and are
not interchangeable.
\end{definition}

\begin{remark}[$\Gamma$ primitive, $k$ derived]\label{tc:rem:kderived}
$\Gamma$ is a primitive of Definition~\ref{tc:def:model}; the width
$k_t$ is \emph{not} a parameter but a bookkeeping quantity recording
where $\Gamma$ sits relative to the background. Two consequences follow
immediately and are worth separating, since they run in opposite
directions:
\begin{itemize}[itemsep=1pt]
\item \emph{The room widens.} In the delocalised phase
$\E[B_t]=ne^{\epsilon^{2}t/2}$, so $k_t=\log(n/\Gamma)+\epsilon^{2}t/2$
grows linearly and the shape-wall separation $k_t/\rho$ widens at rate
$\epsilon^{2}/2\rho$: \emph{grace grows as the background rises}.
\item \emph{The killing walls separate faster.} The crest killing wall
$\log(2B_t)$ rises at $\epsilon^{2}/2$ while the trough killing wall
$\rho^{-1}\log\Gamma$ is constant, so their separation grows at
$\epsilon^{2}/2$.
\end{itemize}
Both are harmless asymptotically, being $O(t)=o(t^{3/2})$.
\end{remark}

\subsection{The void, and when it is flat}\label{tc:subsub:void}

\begin{proposition}[The room, and when it is flat; \textbf{P}]
\label{tc:prop:flat}
The crest is a sigmoid of transition width $O(1)$ centred at $x=b$:
$\varsigma_b\approx1$ for $x\lesssim b-2$ and $\varsigma_b\approx
e^{\,b-x}$ for $x\gtrsim b+2$. Hence the total hazard is flat, at
$\approx B_t$, on
\begin{equation}\label{tc:eq:flat}
-\frac{k_t}{\rho}\;\ll\;x\;\lesssim\;b-2 ,
\end{equation}
a genuine plateau when $k_t/\rho+b\gg1$. It contains $x=0$ in its
interior precisely when $b\gtrsim2$, which is the structural role of the
offset.
\end{proposition}

\begin{proof}
$\varsigma_b(x)\in(0.88,1.05)\cdot$ its plateau value for
$x\le b-2$, and the trough exceeds the crest below $-k_t/\rho$ by
Proposition~\ref{tc:prop:dropout}(i).
\end{proof}

\begin{constraint}[Internal onsets before external]\label{tc:C5}
As the amplitude grows, $x$ reaches the trough wall before the crest
wall if and only if
\begin{equation}\label{tc:eq:C5}
\boxed{\;b\;>\;\frac{k_t}{\rho}\;}
\end{equation}
--- the offset must exceed the trough's half-width. Under
\eqref{tc:eq:C5} a system of unfavourable anatomy suffers internal
damage first, and only later does its external hazard begin to fall.
\end{constraint}

\begin{remark}[Why the offset is needed]\label{tc:rem:whyoffset}
At $b=0$ the origin sits at the crest's inflection, where the local gain
is $\tfrac12$: a system with no muffling at all would already be half
effective, contradicting the reading of the void as the regime in which
the system contributes nothing. The offset moves the origin into the
plateau, where by Theorem~\ref{tc:thm:gain}(iii) the gain is $O(e^{-b})$.
The void then contains its own centre, and \eqref{tc:eq:C5} orders the
two onsets.
\end{remark}

\begin{remark}[Void as mutedness]\label{tc:rem:mutedness}
The void is where the system contributes nothing. That is stated
qualitatively by Remark~\ref{tc:rem:regimes} and quantitatively by
Proposition~\ref{tc:prop:mingain}: the gain $\Psi'$, which measures how
much the driver moves the hazard, falls to
$\kappa(\rho)e^{-(\rho b+k)/(1+\rho)}$ on the void, exponentially small
in the room's width. A system in the waiting room is not merely at low hazard;
its efforts have almost no purchase, and the wider the room the less.
\end{remark}

\subsection{Activation: time supplies the interpolation}
\label{tc:subsub:activation}

\begin{proposition}[Activation time; \textbf{S}]\label{tc:prop:tact}
In log-time $s=\log t$,
\begin{equation}\label{tc:eq:logamp}
\log\Ac_{e^{s}}=\log\sigma+\tfrac32s+\log R_{e^{s}},
\qquad R\ \text{stationary},
\end{equation}
so the amplitude is exponentially small at early log-times and grows
exponentially, with a stationary fluctuation riding on a deterministic
ramp of slope $\tfrac32$. The void therefore ends, and the channels
engage, at
\begin{equation}\label{tc:eq:tact}
t^{\mathrm{int}}_{\mathrm{act}}\asymp\Bigl(\frac{k}{\rho\,\sigma}
\Bigr)^{2/3},
\qquad
t^{\mathrm{ext}}_{\mathrm{act}}\asymp\Bigl(\frac{b}{\sigma}\Bigr)^{2/3},
\end{equation}
the times at which the typical amplitude reaches the trough and crest
walls respectively. Under Constraint~\ref{tc:C5} these are ordered,
$t^{\mathrm{int}}_{\mathrm{act}}<t^{\mathrm{ext}}_{\mathrm{act}}$: the
internal channel engages first and the external one only later.
\end{proposition}

\begin{proof}[Scaling argument]
$\Ac_t=\sigma t^{3/2}R_t$ with $R_t=O_P(1)$ stationary; setting
$\sigma t^{3/2}\asymp k/\rho$ and $\sigma t^{3/2}\asymp b$ give
\eqref{tc:eq:tact}. Dimensions:
$[(k/\rho\sigma)^{2/3}]=(\mathsf T^{3/2})^{2/3}=\mathsf T$. A sharp
constant would require the first-passage law of the amplitude, which we
do not compute.
\end{proof}

\begin{remark}[Why no capacity parameter is needed]
\label{tc:rem:nocapacity}
Revision~1 installed a capacity (strength) parameter with a leak, whose
purpose was to produce a default regime in which the system contributes
nothing and $\mu$ matters only rarely. Equation \eqref{tc:eq:logamp}
supplies that regime for free: the amplitude of $x$ is exponentially
small in log-time at early times, so the void is not an assumption to be
installed but \emph{what the model already does}, and it is left at
\eqref{tc:eq:tact} --- in two stages, internal then external --- without
any additional primitive. The interpolation
parameter is $k$ --- large $k$, long void, background-dominated; small
$k$, immediate engagement --- and $k$ is itself derived
(Remark~\ref{tc:rem:kderived}).
\end{remark}

\begin{remark}[The cost of dropping capacity, recorded]
\label{tc:rem:cost}
This is a real loss and should not be glossed.
\begin{center}\small
\begin{tabular}{@{}lll@{}}
\toprule
& capacity (revision~1, dropped) & $k$ (retained)\\
\midrule
what it moves & an \emph{exponent}, $\tfrac14\to\tfrac12$ &
a \emph{timescale}\\
sharpness & genuine threshold at $S_c$ & crossover at
$t_{\mathrm{act}}$\\
\bottomrule
\end{tabular}
\end{center}
Revision~1's sharp-threshold theorem is therefore gone, and its ``sharp
in parameter space'' register weakens to a crossover. What replaces it
is the threshold at $\rho=1$
(Proposition~\ref{tc:prop:dropout}(ii)), beyond which the background
still reaches the survivor's hazard and below which the system closes
off from the external world. That is a threshold in a structural
parameter with a mechanical meaning rather than one installed to produce
a threshold --- but, as Remark~\ref{tc:rem:rhorole}(c) records, it moves
no exponent. The exchange is favourable and honest, not free.
\end{remark}

% ---------------------------------------------------------------------
\subsection{The kinked gain, and what the driver was concealing}
\label{tc:sub:gain}

\subsection{The exact gain}\label{tc:subsub:gain}

\begin{theorem}[Soft-kink gain; \textbf{P}]\label{tc:thm:gain}
Fix $B_t\equiv B$ constant and define the \emph{gain}
\begin{equation}\label{tc:eq:Psi}
\Psi(x):=-\log\frac{h^{\mathrm{tot}}(x)}{h^{\mathrm{tot}}(0)},
\qquad\text{so that}\qquad
h^{\mathrm{tot}}(x)=h^{\mathrm{tot}}(0)\,e^{-\Psi(x)}
\end{equation}
identically, with $h^{\mathrm{tot}}(0)=B+\Gamma$. Then:
\begin{enumerate}[label=\emph{(\roman*)},itemsep=1pt]
\item $\Psi(0)=0$ and $\Psi$ is smooth and strictly increasing;
\item $\Psi(x)=x+O(e^{-x})$ as $x\to+\infty$ and
$\Psi(x)=\rho x+O(1)$ as $x\to-\infty$: a \emph{soft kink} of slopes
$1$ and $\rho$;
\item the local gain is a hazard-weighted average of the two channels'
slopes,
\begin{equation}\label{tc:eq:Psiprime}
\Psi'(x)=\frac{h^{\mathrm{crest}}(x)\,\frac{e^{x}}{1+e^{x}}
+h^{\mathrm{trough}}(x)\,\rho}
{h^{\mathrm{crest}}(x)+h^{\mathrm{trough}}(x)}\;\in\;(0,\rho),
\end{equation}
with $\Psi'(0)\approx e^{-b}+\rho e^{-k}$, $\Psi'\to1$ as
$x\to+\infty$, $\Psi'\to\rho$ as $x\to-\infty$, and an interior minimum
on the void quantified in Proposition~\ref{tc:prop:mingain}.
\end{enumerate}
Consequently, with constant background, the model is exactly the driftless case's with the credit passed through $\Psi$: Constraint~\ref{tc:C4} is
met.
\end{theorem}

\begin{proof}
(i) is immediate from \eqref{tc:eq:Psi} and positivity of
$h^{\mathrm{tot}}$; monotonicity from (iii). (ii): as $x\to+\infty$,
$h^{\mathrm{tot}}=2Be^{-x}(1+O(e^{-(\rho-1)x}))$, so $\Psi=x+O(\cdot)$;
as $x\to-\infty$, $h^{\mathrm{tot}}=\Gamma e^{-\rho x}(1+O(e^{(\rho-1)
x}))$, so $\Psi=\rho x+O(1)$. (iii): for a sum $h=h_1+h_2$ one has
$-h'/h=(h_1s_1+h_2s_2)/(h_1+h_2)$ with $s_i=-h_i'/h_i$; here
$s_2=\rho$ and, differentiating $2B/(1+e^{x})$,
$s_1=e^{x-b}/(1+e^{x-b})\in(0,1)$. The stated limits follow by
evaluating the weights, using $\Gamma/B=e^{-k}$ at $x=0$.
\end{proof}

\begin{proposition}[The void is muted, exponentially in $k$;
\textbf{P}]\label{tc:prop:mingain}
On the void the gain \eqref{tc:eq:Psiprime} attains an interior minimum
at
\begin{equation}\label{tc:eq:xstar}
x^{*}=\frac{b-k+2\log\rho}{1+\rho},
\qquad
\min\Psi'=\kappa(\rho)\,e^{-\frac{\rho b+k}{1+\rho}},
\qquad
\kappa(\rho):=\rho^{\frac{2}{1+\rho}}+\rho^{\frac{1-\rho}{1+\rho}} .
\end{equation}
The gain therefore does not merely dip but is \emph{exponentially small
in $(\rho b+k)/(1+\rho)$}: the wider the room --- on either side --- the
more completely the driver is muted within it. For $\rho=\tfrac32$,
$\kappa=2.045$; for $\rho=\tfrac52$, $\kappa=2.443$.
\end{proposition}

\begin{proof}
On the void the crest is saturated, so its weight is $\approx1$ and its
slope is $s_{1}\approx e^{x-b}$, while the trough weight is
$\approx(\Gamma/B)e^{-\rho x}=e^{-k}e^{-\rho x}$. Hence
$\Psi'\approx f(x):=e^{x-b}+\rho e^{-k}e^{-\rho x}$, and $f'(x)=0$ gives
$e^{(1+\rho)x}=\rho^{2}e^{\,b-k}$, which is \eqref{tc:eq:xstar}.
Substituting back, both terms carry the common factor
$e^{-(\rho b+k)/(1+\rho)}$ with coefficients $\rho^{2/(1+\rho)}$ and
$\rho^{(1-\rho)/(1+\rho)}$, summing to the stated value. (The
approximation is verified numerically to four significant figures over
$\rho\in\{1.5,2.5\}$, $b\in[6,14]$, $k\in[8,20]$.)
\end{proof}

\begin{remark}[Correction to the reference specification]
\label{tc:rem:speccorr}
Two points where the informal specification requires amendment, both
recorded in the ledger.

\emph{(a)} The closed form $\psi_\rho(x)=-\log(e^{-x}+e^{-\rho x})$ is
\emph{not} an identity for \eqref{tc:eq:hazard}: at $x=0$ the model
gives $B+\Gamma$ while $\psi_\rho$ gives $2h_0$. It is the clean
representative of the soft-kink class, not the model's own gain;
\eqref{tc:eq:Psi} is exact by construction and has the same asymptotic
slopes, so nothing structural is lost.

\emph{(b)} The gain range is $(0,\rho)$ and not $[1,\rho]$. The crest
channel's local slope is $e^{x}/(1+e^{x})$, which tends to $0$, not $1$,
as $x\to-\infty$. This is not a blemish but the sharpest available
statement of what the void \emph{is}: on the plateau the gain falls to
$\kappa(\rho)e^{-(\rho b+k)/(1+\rho)}$
(Proposition~\ref{tc:prop:mingain}), so
the driver is muted and its efforts scarcely move the hazard. The gain
profile thus encodes all three regimes --- exponentially small in the
void, $\to1$ under shielding, $\to\rho$ under damage. The mutedness is a
property of a \emph{wide} room: if $\rho b+k$ is $O(1)$ the minimum is
$O(1)$ and there is no muted regime, consistently with
Proposition~\ref{tc:prop:flat}.
\end{remark}

\subsection{What the original driver was concealing}
\label{tc:subsub:concealed}

\begin{corollary}[State-dependent gain; \textbf{P}]
\label{tc:cor:concealed}
Along a path, the effective negative logarithmic derivative of the total
hazard is
\begin{equation}\label{tc:eq:effmu}
\tilde\mu_t=-\frac{d}{dt}\log h^{\mathrm{tot}}_t
=\Psi'(x_t)\,\mu_t,
\qquad \Psi'\in(0,\rho).
\end{equation}
That is: \emph{the driftless case's driver was the present driver times a
state-dependent gain}, and the old Brownian motion was absorbing that
gain as extra noise. The revision therefore extracts structure from the
original noise term rather than adding a primitive to it.
\end{corollary}

\begin{proof}
Chain rule on \eqref{tc:eq:Psi} with $\dot x=\mu$.
\end{proof}

\begin{remark}[Reading of the gain]\label{tc:rem:gainreading}
The new $\mu$ is correspondingly narrower than the old, and the
difference is now explicit rather than hidden. Where the old model
attributed all variation in the log-hazard's slope to the driver, the
present one attributes part of it to \emph{where the system is}: the
same effort buys a full unit of log-hazard reduction under shielding,
$\rho$ units of damage in a trough, and nothing at all in the void.
\end{remark}

% ---------------------------------------------------------------------
\section{Transfer, in detail}\label{app:transfer}

\subsection{Reduction to the driftless case}\label{tc:subsub:reduction}

\begin{proposition}[Reduction; \textbf{P}]\label{tc:prop:reduction}
On the slice $\epsilon=0$ (so $B\equiv n$), $\Gamma\to0$, and $x\gg0$,
the model is the driftless case's with $h_0=2n$ and credit $x-x_0$. More
usefully: for $\epsilon=0$ and any $\Gamma$, the total hazard is a
deterministic function of $x$ alone, and Theorem~\ref{tc:thm:gain}
identifies it as the driftless case's hazard composed with a fixed gain.
\end{proposition}

\begin{proof}
At $\epsilon=0$, $B\equiv n$; for $x\gg0$, $h^{\mathrm{tot}}\approx
2ne^{-x}$ since the trough is $O(e^{-\rho x})=o(e^{-x})$.
\end{proof}

% ---------------------------------------------------------------------
\subsection{Asymptotics: the driftless case survives intact}
\label{tc:sub:asymptotics}

\begin{theorem}[the driftless case is undisturbed; \textbf{A}]
\label{tc:thm:survives}
In the two-channel model, for every $\rho>0$, $\Gamma>0$ and $b\ge0$:
\begin{enumerate}[label=\emph{(\roman*)},itemsep=1pt]
\item the persistence exponent is $\tfrac14$, exactly;
\item $\PP(T>t)\asymp Ct^{-1/4}$;
\item $\E[T]=\infty$ while $T<\infty$ a.s.\ (null recurrence);
\item the survivorship interpolation of the driftless case, its pure Doob
limit, its $s/4t$ linear response and its critical mode, are unchanged.
\end{enumerate}
\emph{No truncation window and no $\gamma=0$ caveat}:
Constraint~\ref{tc:C3} holds as a theorem.
\end{theorem}

\begin{proof}
By Lemma~\ref{tc:lem:walls} the killing walls are $\log(2B_t)$ and
$\rho^{-1}\log\Gamma$. The latter is constant, hence an $O(1)$ shift,
which we proved to be gauge. The former is $O(t)$ in the
delocalised phase and $O(\sqrt t)$ in the frozen phase, hence
$o(t^{3/2})$, and is absorbed exactly in the linear case by
Theorem~\ref{tc:thm:absorb} and invisibly in general by
Theorem~\ref{tc:thm:sublinear}; the offset $b$ enters
\eqref{tc:eq:shift} as a further $O(1)$ shift and is likewise gauge.
Survivors live at large positive $x$, where by
Proposition~\ref{tc:prop:dropout}(ii) the dominant channel is the crest
if $\rho>1$ and the trough if $\rho<1$; in either case the dominant
channel is an exponential in $x$, so the killing wall is a level of the
credit and the persistence problem is the driftless case's up to a constant
rescaling. The conclusions are then the driftless case's, transported --- and,
notably, for every $\rho>0$ rather than only $\rho>1$.
\end{proof}

\begin{remark}[This repairs revision~1's defect]\label{tc:rem:repair}
Revision~1's additive internal hazard contributed $\gamma t$ to
$\Lambda$, destroying the $t^{-1/4}$ tail and forcing a truncation
window with finite life expectancy $\asymp g^{-3/4}$. The trough channel
does not, and the reason is structural rather than fortunate:
$\Gamma e^{-\rho x}$ is small exactly where survivors live, whereas a
constant is not small anywhere. Both walls are $O(1)$-or-gauge shifts,
and we proved such shifts to be gauge.
\end{remark}

\begin{remark}[What was given up]\label{tc:rem:givenup}
Both channels are decreasing in $x$, so the total hazard is monotone and
\emph{more strength is always better}. Revision~1's U-shaped hazard, and
with it the optimal finite strength and the $(S,A)$ optimisation, are
gone. This is the price of Constraint~\ref{tc:C3}: an interior optimum
requires a term that grows with $x$, and any such term destroys the
tail. An optimum and an undisturbed the driftless case were never going to
coexist. What survives of the earlier optimisation story is the
threshold at $\rho=1$ and the trade recorded in
Remark~\ref{tc:rem:cost}.
\end{remark}

% ---------------------------------------------------------------------
\subsection{Transfer: how much of the driftless case survives, and where the
differences go}\label{tc:sub:transfer}

Theorem~\ref{tc:thm:survives} asserted that the driftless case's conclusions
carry over. This subsection makes the transfer quantitative: it
enumerates the four ways the two-channel hazard differs from
$h_{0}e^{-x}$, shows that each lands in gauge or in the prefactor, and
isolates the one place where the two-channel model is genuinely
\emph{better} conditioned than the driftless case.

\subsection{The four differences}\label{tc:subsub:four}

\begin{center}\small
\begin{tabular}{@{}llll@{}}
\toprule
difference & size & where it goes & reference\\
\midrule
wall location $\log B_u+b$ & $\Theta(u)$ & gauge (exact absorption) &
Thm.~\ref{tc:thm:absorb}\\
graded boundary of width $W_u$ & $\Theta(u)$ & gauge (collapses under
rescaling) & Thm.~\ref{tc:thm:collapse}\\
waiting-room passage & $O(1)$ & prefactor & Prop.~\ref{tc:prop:roomcost}\\
background fluctuation & $O(\sqrt u)$ & gauge &
Thm.~\ref{tc:thm:sublinear}\\
\bottomrule
\end{tabular}
\end{center}

\noindent Nothing reaches an exponent. The transfer is therefore exact
at the level at which the driftless case states its results.

\begin{lemma}[Above the room the model is the driftless case's; \textbf{P}]
\label{tc:lem:aboveroom}
For $x\ge b$,
\begin{equation}\label{tc:eq:sandwich}
\tfrac12\bigl(1+e^{-b}\bigr)e^{\,b-x}\;\le\;\varsigma_b(x)\;\le\;
\bigl(1+e^{-b}\bigr)e^{\,b-x},
\end{equation}
so $h^{\mathrm{crest}}\asymp B_ue^{\,b-x}$ with constants in
$[\tfrac12,1]\cdot(1+e^{-b})$. If moreover $\rho>1$ then
$h^{\mathrm{trough}}/h^{\mathrm{crest}}\to0$ as $x\to\infty$, so
$h^{\mathrm{tot}}\asymp B_ue^{\,b-x}$: above the room the hazard is
the driftless case's $h_{0}e^{-x}$ with $h_{0}=B_ue^{b}$, up to bounded
multiplicative constants.
\end{lemma}

\begin{proof}
For $x\ge b$ one has $1\le e^{x-b}$, hence
$e^{x-b}\le1+e^{x-b}\le2e^{x-b}$; invert. The trough comparison is
Proposition~\ref{tc:prop:dropout}(ii).
\end{proof}

\subsection{The graded boundary collapses}\label{tc:subsub:collapse}

The substantive structural difference is that the killing boundary is
neither hard nor soft but \emph{graded}. Above $\log B_u+b$ the hazard
is negligible; across a band below it the hazard is $\asymp B_u$,
bounded; below $(\log\Gamma)/\rho$ it diverges.

\begin{theorem}[The band is asymptotically invisible; \textbf{A}]
\label{tc:thm:collapse}
Let
\begin{equation}\label{tc:eq:band}
W_u:=\underbrace{\log B_u+b}_{\text{crest wall}}
-\underbrace{\frac{\log\Gamma}{\rho}}_{\text{trough wall}}
=\Bigl(1-\frac1\rho\Bigr)\log B_u+b+\frac{k}{\rho}
\end{equation}
be the width of the graded band. In the delocalised phase
$W_u=\Theta(u)$ and in the frozen phase $W_u=\Theta(\sqrt u)$; in both,
\[
\frac{W_u}{u^{3/2}}\longrightarrow0 .
\]
Hence under the driftless case's rescaling the whole band contracts to a single
point, and the graded boundary is asymptotically the hard wall of that
section. Granting the moving-boundary robustness of its Ansatz class,
the persistence exponent is $\tfrac14$.
\end{theorem}

\begin{proof}
Insert $\log B_u=\log n+\epsilon^{2}u/2$ (delocalised,
Corollary~\ref{tc:cor:gompertz}) or
$\log B_u=\epsilon\sqrt{2u\log n}$ (frozen,
Theorem~\ref{tc:thm:freezing}) into \eqref{tc:eq:band}; both are
$o(u^{3/2})$. The rescaling and the robustness step are as in
Theorem~\ref{tc:thm:sublinear}.
\end{proof}

\begin{remark}[Softening and re-hardening do not cancel each other ---
they both cancel against self-similarity]\label{tc:rem:netcancel}
It is tempting to read the situation as a competition: the logistic
bounds the hazard from above and would soften the wall, while the trough
is unbounded and re-hardens it, the two effects offsetting. That is not
what happens. The two effects act at \emph{different places} --- the
logistic at the top of the band, the trough at its bottom --- and
neither undoes the other. What renders both irrelevant is the third
thing: the band they enclose has width $o(u^{3/2})$, while the credit
they constrain fluctuates on scale $u^{3/2}$. A boundary structure of
any shape whatsoever, confined to a band of width $o(u^{3/2})$, is
invisible to the exponent. The net effect of the logistic and the trough
is therefore exactly nil asymptotically --- not by mutual cancellation
but because self-similarity outruns them both.
\end{remark}

\begin{proposition}[The tolerance shrinks: this model is better
conditioned than the driftless case's; \textbf{P}]\label{tc:prop:tolerance}
Survival to time $t$ requires
\begin{equation}\label{tc:eq:tolerance}
\mathrm{Leb}\{u\le t:x_u\lesssim b\}\;=\;O\!\bigl(B_t^{-1}\bigr),
\end{equation}
since the hazard on the band is $\asymp B_u$ and $\Lambda_t=O(1)$ is
required. In the delocalised phase $B_t^{-1}=n^{-1}e^{-\epsilon^{2}t/2}
\to0$, so the admissible occupation below the wall tends to zero: the
problem converges to \emph{strict} persistence.
\end{proposition}

\begin{proof}
$\Lambda_t\ge\bigl(\inf_{u\le t}B_u\bigr)\cdot c\cdot
\mathrm{Leb}\{u\le t:x_u\lesssim b\}$ with $c=\varsigma_b(b)>0$;
survival to $t$ with probability bounded below forces $\Lambda_t=O(1)$.
\end{proof}

\begin{remark}[The Ansatz gap narrows]\label{tc:rem:gapnarrows}
the driftless case's sole load-bearing heuristic was that persistence with
\emph{bounded occupation tolerance} shares the exponent of strict
persistence. The two-channel model does not evade that assumption, but
it weakens what is being assumed: by
Proposition~\ref{tc:prop:tolerance} the tolerance here is not $O(1)$ but
$O(B_t^{-1})\to0$. The richer, more ontologically committed model is
thus \emph{closer} to the rigorous case than the simpler one --- an
unusual direction for added structure to move a proof, and worth
recording as such.
\end{remark}

\subsection{The waiting room is a prefactor, and it is not a
shelter}\label{tc:subsub:roomcost}

\begin{proposition}[Cost of the room passage; \textbf{S}]
\label{tc:prop:roomcost}
The room $[-k/\rho,b]$ lies entirely below the crest killing wall
$\log B_u+b$, since $\log B_u>0$. A trajectory therefore traverses it
only before escaping upward, at time
$t^{\mathrm{ext}}_{\mathrm{act}}\asymp(b/\sigma)^{2/3}$
(Proposition~\ref{tc:prop:tact}), accumulating
\begin{equation}\label{tc:eq:roomcost}
\Lambda_{\mathrm{room}}\;\asymp\;
B_0\,t^{\mathrm{ext}}_{\mathrm{act}}
\;\asymp\;n\Bigl(\frac{b}{\sigma}\Bigr)^{2/3},
\end{equation}
an $O(1)$ quantity in $t$. Its entire effect is to multiply the
prefactor in $\PP(T>t)\asymp Ct^{-1/4}$ by
$e^{-\Theta\left(n(b/\sigma)^{2/3}\right)}$.
\end{proposition}

\begin{remark}[Grace, but not shelter]\label{tc:rem:notshelter}
The room deserves a more careful reading than ``grace interval''
suggests. Inside it the hazard is \emph{exactly the background}: neither
raised by damage nor lowered by shielding. It is therefore grace
relative to the trough --- unfavourable anatomy costs nothing while the
amplitude is small enough to keep $x$ above $-k/\rho$ --- but it is a
\emph{cost} relative to the crest, because survivors must be above the
wall and time spent in the room is time not spent shielding. Formula
\eqref{tc:eq:roomcost} prices that cost. The wider the room and the
larger the background, the more expensive the wait: the offset $b$ buys
the ordering of the two onsets (Constraint~\ref{tc:C5}) and pays for it
in the prefactor, at rate $n(b/\sigma)^{2/3}$ in the exponent.
\end{remark}

\subsection{The conditioning ladder}\label{sh:subsub:ladder}

We quantify how each conditioning enhances survival.

\begin{proposition}[Degrees of the initial data; \textbf{C/A}]
\label{sh:prop:degrees}
\emph{(i) Initial rate.} From $(m,0)$ with $m>0$: the rate stays positive
until $W$ reaches $-m/\sigma$, a time of order $t_\mu=m^2/\sigma^2$ during
which $D$ increases; thereafter, exactly by scaling,
$\PP_{(m,0)}(\tau>t)\sim c_1(t_\mu/t)^{1/4}\propto m^{1/2}$
\emph{(degree $\tfrac12$: a tenfold gain in survival odds costs a
hundredfold increase in $m$)}.
\emph{(ii) Initial credit.} From $(0,d)$ with $d>0$: by scaling with
$c=d^{-1/3}$, $\PP_{(0,d)}(\tau>t)=\PP_{(0,1)}(\sigma^{2/3}\tau>
t\,d^{-2/3}\sigma^{2/3})\asymp(t_D/t)^{1/4}\propto d^{1/6}$
\emph{(degree $\tfrac16$)}.
\emph{(iii) Unification.} All initial data act only through the equivalent
timescale they generate:
\begin{equation}\label{sh:eq:tstar}
\PP(T>t)\ \asymp\ \Bigl(\frac{t_*}{t}\Bigr)^{1/4},\qquad
t_*\ \asymp\ \sigma^{-2/3}\max\bigl(1,\ \hat\mu_0^{\,2},\ \hat
D_0^{\,2/3}\bigr),
\end{equation}
and which term dominates is decided by the horizon coordinate: by
\eqref{sh:eq:kappa}, $t_D/t_\mu=\kappa_0^{2/3}$, so rate dominates for
$\kappa_0\ll1$ and credit for $\kappa_0\gg1$.
\end{proposition}

\begin{proof}
The scaling identities are exact (Proposition~\ref{sh:prop:selfsim});
transfer to $\PP(T>t)$ is Ansatz~\ref{sh:ansatz}. That the exponent from
interior starting states is unchanged is part of
Theorem~\ref{sh:thm:persist}.
\end{proof}

\begin{proposition}[Homogeneity of the harmonic function; \textbf{P/C}]
\label{sh:prop:harmonic}
On the open quadrant define
\[
\hm(\mu,d)\ :=\ \lim_{t\to\infty}t^{1/4}\,\PP_{(\mu,d)}(\tau>t);
\]
the limit exists, is finite and positive \cite{IW94,GJW99}. Then $\hm$ is
invariant-harmonic for the killed Kolmogorov diffusion and homogeneous of
degree $\tfrac12$ under the scaling of
Proposition~\ref{sh:prop:selfsim}, $\hm(c\mu,c^3d)=c^{1/2}\hm(\mu,d)$.
Consequently
\begin{equation}\label{sh:eq:harm}
\hm(\mu,d)=\mu^{1/2}F\bigl(\kappa\bigr)=d^{1/6}\tilde F(\kappa),
\qquad \kappa=\sigma^2d/\mu^3,
\end{equation}
for profiles $F,\tilde F$ determined by the explicit
confluent-hypergeometric formulas constructed in \cite{GJW99}. (We flag
a provenance point: the special functions arising here are confluent
hypergeometric, not Airy; Airy functions belong to the neighbouring
parabolic-drift literature.)
\end{proposition}

\begin{proof}
Given existence of the limit, homogeneity is immediate from
Proposition~\ref{sh:prop:selfsim}:
\[
t^{1/4}\,\PP_{(c\mu,c^3d)}(\tau>t)
= c^{1/2}\,\bigl(t/c^2\bigr)^{1/4}\,\PP_{(\mu,d)}\bigl(\tau>t/c^{2}\bigr)
\ \longrightarrow\ c^{1/2}\,\hm(\mu,d).
\]
Harmonicity follows from the Markov property and the same limit taken one
step later. Representation \eqref{sh:eq:harm} is homogeneity restricted to
the orbit coordinates of Proposition~\ref{sh:prop:gauge}(iii).
\end{proof}

\begin{remark}[Conditioning on survival forever: the Doob transform;
\textbf{C}]\label{sh:rem:doob}
Conditioning on $\{T=\infty\}$ is vacuous (Theorem~\ref{sh:thm:noimm});
its canonical regularization is the $\hm$-transform, the
$t\to\infty$ limit of conditioning on $\{\tau>t\}$. Under the transformed
measure the process never returns to zero credit and, by self-similarity,
\[
\mu_s\ \asymp\ \sigma s^{1/2},\qquad D_s\ \asymp\ \sigma s^{3/2},
\qquad\text{hence}\qquad \log h_s\ \sim\ -\sigma s^{3/2}:
\]
conditioned survivors exhibit superexponentially decaying hazard and
appear, from the inside, unambiguously immortal --- although the
unconditioned process is certainly mortal and the model contains no drift.
The drift is manufactured by selection; this is the strongest true form of
the event-horizon intuition, and it is a statement about ensembles, not
about state space. The accompanying Yaglom-type statement --- convergence
in law of $(\mu_t/\sigma t^{1/2},\,D_t/\sigma t^{3/2})$ conditionally on
$\{\tau>t\}$ to a nondegenerate limit --- is the correct
``stationary profile of a long-lived system,'' stationary only in
self-similar coordinates; see \cite{GJW99,AS15}, and \cite{MV12} for
the general theory of quasi-stationary limits against which this
self-similar variant should be read.
Section~\ref{sh:sub:interp} promotes this remark to the section's central
result: the conditioning that manufactures the drift becomes itself the
object of study, and the transform described here is identified as the
exact limit of the survivorship interpolation.
\end{remark}

% ---------------------------------------------------------------------
\section{Reference sheet: the filter at a glance}\label{app:sheet}

% ---------------------------------------------------------------------
\label{sh:sub:sheet}

This subsection collects, without proof, the quantitative content of the
section in the form most useful for reference. Everything here is
established above; cross-references point to the source.

\subsection{The conditioned dynamics in closed form}
\label{sh:subsub:sheet:dyn}

Under the survivorship limit $Q$ of Theorem~\ref{sh:thm:Q} the system is
\begin{equation}\label{sh:eq:sheet}
d\mu_t=\frac{\sigma^2}{\mu_t}\,G(\kappa_t)\,dt+\sigma\,dW^Q_t,
\qquad
dD_t=\mu_t\,dt,
\qquad
\kappa=\frac{\sigma^2D}{\mu^3},
\end{equation}
with $G(\kappa)=\tfrac12-3\kappa F'(\kappa)/F(\kappa)$ read off
$\hm=\mu^{1/2}F(\kappa)$ as in \eqref{sh:eq:harm}--\eqref{sh:eq:G}.
Dimensionally $[\sigma^2/\mu]=\mathsf T^{-2}=[\mu]/[t]$, as required. The
shape function interpolates between two endpoints
(Theorem~\ref{sh:thm:drift}):
\[
G(0^+)=\tfrac12\ \ (\text{drift }\sigma^2/2\mu,\ \text{maximal}),
\qquad
G(\infty)=0\ \ (\text{drift }0,\ \text{none}),
\]
so that \emph{survivorship pushes hardest near the wall and not at all
once credit is abundant}, the strength along any ray of fixed $\kappa$
scaling as $1/\mu$.

\begin{remark}[The endpoints are ray limits, not coordinate limits]
\label{sh:rem:raylimit}
One may not simply let $\mu\to0$ in $\sigma^2/2\mu$: at fixed $D>0$ that
sends $\kappa\to\infty$ and hence $G\to0$, and the two factors compete.
Only the statement along rays of constant $\kappa$ is unambiguous, which
is precisely why $\kappa$ and not $(\mu,D)$ separately is the coordinate
in which the conditioned dynamics should be read
(Proposition~\ref{sh:prop:rays}).
\end{remark}

\subsection{The three levels, before and after}
\label{sh:subsub:sheet:levels}

\begin{center}\small
\begin{tabular}{@{}p{0.24\textwidth}p{0.34\textwidth}p{0.34\textwidth}@{}}
\toprule
& \textbf{Before} ($\PP$) & \textbf{After} ($Q$)\\
\midrule
\multicolumn{3}{@{}l}{\emph{Level 1: the driver $\mu$}}\\
law & Gaussian, driftless & non-Gaussian, drift
$(\sigma^2/\mu)G(\kappa)\ge0$\\
autonomy & $\mu$ Markov by itself & only the pair $(\mu,D)$ is Markov\\
quadratic variation & $\sigma^2\,dt$ & $\sigma^2\,dt$ \emph{(unchanged)}\\
increments & stationary, independent & neither\\
self-similarity & index $3/2$ & index $3/2$ \emph{(preserved)}\\
\addlinespace
\multicolumn{3}{@{}l}{\emph{Level 2: the hazard $h$}}\\
path regularity & $C^1$ & $C^1$ \emph{(unchanged)}\\
law at time $s$ & lognormal, $\log$-variance
$\sigma^2s^3/3$ & non-Gaussian, one-sided\\
range & exceeds $h_0$ infinitely often,
unbounded above & $h_s<h_0$ for \emph{all} $s$ (if $D_0=0$)\\
credit $D$ & null recurrent,
$\liminf_sD_s=-\infty$ & transient, $D_s\to\infty$,
$D_s\asymp\sigma s^{3/2}$\\
\addlinespace
\multicolumn{3}{@{}l}{\emph{Level 3: survival}}\\
$\Lambda_\infty$ & $=\infty$ a.s. & $<\infty$ a.s.,
$\asymp\hat h_0=h_0\sigma^{-2/3}$\\
quenched curve $e^{-\Lambda_t}$ & $\to0$; collapse
$\exp(-e^{\sigma t^{3/2}})$ & plateau at $e^{-\Lambda_\infty}>0$\\
population curve & $\asymp Ct^{-1/4}$ & $\equiv1$\\
lifetime & $T<\infty$ a.s., $\E[T]=\infty$ & $T=\infty$\\
\bottomrule
\end{tabular}
\end{center}

\noindent Sources: Level 1, Theorems~\ref{sh:thm:Q}
and~\ref{sh:thm:drift}; Level 2, Definition~\ref{sh:def:model} and
Theorem~\ref{sh:thm:singular}; Level 3, Theorem~\ref{sh:thm:noimm},
Claim~\ref{sh:claim:tail}, Corollary~\ref{sh:cor:mean}.

\begin{remark}[Notable structures at Level 2]\label{sh:rem:sheet:L2}
Three features are worth isolating. \emph{(i)} The hazard is one
derivative \emph{smoother} than the noise driving it, a direct
consequence of the log-derivative placement
(Remark~\ref{sh:rem:iwp}); conditioning does not change this, since
Girsanov alters drift and never path regularity. \emph{(ii)} Before
conditioning $\log h_s$ is exactly Gaussian with variance
$\sigma^2s^3/3$ --- \emph{cubic} in time, the integrated-Brownian
fingerprint --- so $h_s$ is lognormal. \emph{(iii)} After conditioning,
the survivor's hazard never once exceeds its initial value, whereas
before it does so infinitely often and by unbounded amounts: this is the
wall of \S\ref{sh:subsub:reduction} in its most concrete form. Finer
still, Groeneboom, Jongbloed and Wellner \cite{GJW99} identify
$t\mapsto t^{9/10}$ as a critical curve for the conditioned process, the
expected time spent below $t^\alpha$ being finite for $\alpha<9/10$ and
infinite for $\alpha\ge9/10$: survivors do lag their typical
$s^{3/2}$ growth, and $9/10$ is exactly how much.
\end{remark}

\begin{proposition}[The quenched--annealed gap; \textbf{P}]
\label{sh:prop:gap}
In the driftless model, almost surely
\begin{equation}\label{sh:eq:gap}
\log\Lambda_t=\sup_{u\le t}\bigl(D_0-D_u\bigr)+O(\log t),
\end{equation}
and $t^{-3/2}\sup_{u\le t}(D_0-D_u)\to\sigma\sup_{v\le1}(-I_v)$ in law, a
nondegenerate positive limit. Hence the quenched survival curve
collapses \emph{doubly exponentially}: writing
$c_t:=t^{-3/2}\log\Lambda_t$,
\[
e^{-\Lambda_t}=\exp\bigl(-e^{\,c_t\,t^{3/2}}\bigr),
\qquad
c_t\ \xrightarrow{\;d\;}\ \sigma\sup_{v\le1}(-I_v)\;>\;0,
\]
while by Claim~\ref{sh:claim:tail} the population curve decays only as
$Ct^{-1/4}$. The entire population tail is therefore carried by a set of
paths of vanishing probability.
\end{proposition}

\begin{proof}
Write $M_t=\sup_{u\le t}(D_0-D_u)$ and
$\Lambda_t=h_0\int_0^te^{D_0-D_u}du$. The upper bound
$\Lambda_t\le h_0te^{M_t}$ is immediate. For the lower bound, $D$ is
$C^1$ with $\dot D=\mu$ and $\sup_{u\le t}|\mu_u|=O(t^{1/2})$ a.s.\ up to
logarithmic factors, so on an interval of length $\asymp t^{-1/2}$ around
a near-maximiser the integrand exceeds $e^{M_t-1}$, giving
$\Lambda_t\gtrsim h_0t^{-1/2}e^{M_t}$. Taking logarithms yields
\eqref{sh:eq:gap}. The scaling limit is
Proposition~\ref{sh:prop:selfsim} applied to $I$, with
$\sup_{v\le1}(-I_v)>0$ a.s.
\end{proof}

\begin{remark}[The interpolation as a one-parameter family]
\label{sh:rem:sheet:family}
Between the two columns, the survival curve of a $t$-survivor is
\[
\PP(T>t+s\mid T>t)\;\approx\;\Bigl(1+\frac st\Bigr)^{-1/4},
\]
a scale-free family flattening from a power law at $t=0$ to the constant
$1$ as $t\to\infty$; equivalently the $t$-survivor carries effective
hazard $1/4t$ (Corollaries~\ref{sh:cor:residual}
and~\ref{sh:cor:lindy}). The filter does not deform the survival curve
arbitrarily: it slides it along a single one-parameter orbit.
\end{remark}

\subsection{What conditioning can and cannot do}
\label{sh:subsub:sheet:rigid}

\begin{proposition}[Rigidity of the survivorship transform; \textbf{P}]
\label{sh:prop:rigid}
Let $\widetilde\PP$ be any measure locally absolutely continuous with
respect to $\PP$ on each $\Fc_s$ --- in particular $Q$. Then:
\begin{enumerate}[label=\emph{(\roman*)},itemsep=1pt]
\item \emph{The kinematics survive exactly.} $dD_t=\mu_t\,dt$ holds
$\widetilde\PP$-a.s.\ unchanged: credit remains literally the integral of
rate. Conditioning cannot alter it.
\item \emph{Only the driver may be tilted.} Since the noise enters solely
the $\mu$ coordinate in \eqref{sh:eq:kolm}, Girsanov can add drift only
to $\mu$; the transform has exactly one functional degree of freedom.
\item \emph{Self-similarity fixes its form.} If in addition
$\widetilde\PP$ respects the scaling of
Proposition~\ref{sh:prop:selfsim}, that degree of freedom is forced into
the shape $(\sigma^2/\mu)G(\kappa)$ of \eqref{sh:eq:sheet}, a single
function of one dimensionless variable.
\item \emph{The cascade is destroyed.} Under $\PP$ the system is
triangular: $\mu$ is autonomous and $D$ is slaved to it. Under $Q$ the
credit re-enters the drift of $\mu$ through $\kappa$, so $\mu$ alone is
no longer Markov although the pair remains so.
\end{enumerate}
\end{proposition}

\begin{proof}
(i) and (ii) are Girsanov's theorem for \eqref{sh:eq:kolm}, whose
diffusion matrix is degenerate with range spanned by $\partial_\mu$; the
finite-variation coordinate admits no equivalent change. (iii) is
Proposition~\ref{sh:prop:harmonic} together with
Theorem~\ref{sh:thm:drift}. (iv) is the $\kappa$-dependence in
\eqref{sh:eq:G}, non-constant by Theorem~\ref{sh:thm:drift}(ii)--(iii).
\end{proof}

\begin{remark}[Survivorship coupling is collider bias in continuous time]
\label{sh:rem:collider}
The feedback of Proposition~\ref{sh:prop:rigid}(iv) is informational, not
mechanical: nothing pushes $\mu$. Under $\PP$ the credit accumulated by
time $s$ and the future increments of the driver are independent, but
both feed a common outcome, survival; conditioning on that outcome
renders them dependent. This is exactly collider bias --- Berkson's
paradox, or ``explaining away'' --- realised in continuous time, and
\eqref{sh:eq:sheet} is its quantitative form: little credit banked
implies the remaining path must work harder, and the drift says by how
much. The same object appears in large-deviation theory as the driven or
effective process associated with a rare event. The moral of
Proposition~\ref{sh:prop:rigid} is then sharp: survivorship cannot bend
the kinematics of a system, only the statistics of what drives it ---
and, when the system is self-similar, only along a one-parameter family
of shapes.
\end{remark}

% ---------------------------------------------------------------------
\section{Numerical verification}\label{app:numerics}

% ---------------------------------------------------------------------
\label{sh:sub:numerics}

The one load-bearing heuristic, Ansatz~\ref{sh:ansatz}, admits a direct
test, which we specify for completeness; this subsection may be read as an
appendix.

\begin{lemma}[Exact simulation of the pair; \textbf{P}]\label{sh:lem:exact}
For a step $\Delta>0$ let $\xi:=W_{t+\Delta}-W_t$ and
$\eta:=\int_t^{t+\Delta}(W_s-W_t)\,ds$. Then $(\xi,\eta)$ is independent of
$\Fc_t$, centered Gaussian with covariance
$\bigl(\begin{smallmatrix}\Delta&\Delta^2/2\\
\Delta^2/2&\Delta^3/3\end{smallmatrix}\bigr)$, and the update
\[
\mu_{t+\Delta}=\mu_t+\sigma\xi,\qquad
D_{t+\Delta}=D_t+\mu_t\Delta+\sigma\eta
\]
reproduces the law of \eqref{sh:eq:kolm} at the grid points with zero
discretization error.
\end{lemma}

\begin{proof}
$(\xi,\eta)$ is a measurable functional of the post-$t$ increments and is a
copy of $(W_\Delta,I_\Delta)$; apply Lemma~\ref{sh:lem:secondorder}. The
update identity is
$I_{t+\Delta}=I_t+W_t\Delta+\eta$ multiplied through by $\sigma$.
\end{proof}

The remaining error is the quadrature of
$\Lambda_t=h_0e^{-D_0}\int_0^te^{-(D_u-D_0)}du$, whose integrand varies on
the scale $|\Delta D|\lesssim1$, i.e.\ $\Delta\lesssim1/|\mu|\sim
t^{-1/2}$: the constraint tightens with time. Accumulate in logarithmic
space; the integrand ranges over hundreds of orders of magnitude. To reach
the tail --- probabilities of $10^{-3}$ to $10^{-4}$ require
$t/t_*$ up to $10^{12}$ and beyond --- use multilevel splitting along the
thresholds $D_k=\sigma(\theta^kt_*)^{3/2}$, or importance sampling tilted
by the harmonic function of Proposition~\ref{sh:prop:harmonic}, the
asymptotically optimal tilt. Validate the code first on the hard-wall
functional, where \eqref{sh:eq:persist} is a theorem; then run the
soft-killing functional and fit $\log\PP$ against $\log t$ over at least
three decades, discarding $t\lesssim10\,t_*$.
\emph{Falsifiers:} a fitted slope robustly different from $-\tfrac14$
(Ansatz fails: the killing term carries its own scale); any dependence of
the slope on $h_0$ (gauge violation, contradicting
Proposition~\ref{sh:prop:gauge}); non-convergence of the conditioned
rescaled credit $D_s/\sigma s^{3/2}$ (failure of the self-similar Yaglom
picture of Remark~\ref{sh:rem:doob}). The survivorship results of
\S\ref{sh:sub:interp} add three sharper tests: estimate the window bias
$\beta_t(s)$ under splitting and regress it against $\hm(X_s)$, whose
proportionality tests Theorem~\ref{sh:thm:Q}; regress conditioned
increments $\Delta\mu$ on $1/\mu$ at small $\kappa$, whose slope
$\sigma^2/2$ tests the wall asymptotics of
Corollary~\ref{sh:cor:bessel}; and verify the linear response
$\beta_t\approx1+s/4t$ at large $t/s$, whose coefficient $\tfrac14$ tests
Proposition~\ref{sh:prop:rate}.

% ---------------------------------------------------------------------
\label{tc:sub:numerics}

Four tests specific to this part, additional to the driftless case's.

\emph{(i) Freezing.} Simulate $B_t$ for $n\in\{10^{3},10^{4},10^{5}\}$
and plot $\log B_t/\log n$ against $\theta$; expect collapse onto
\eqref{tc:eq:freezing} with tangential contact at $\theta=\sqrt2$, and
sample-to-sample fluctuations $o(1)$ before the transition and $O(1)$
after, the frozen phase being non-self-averaging.

\emph{(ii) The two coordinates under conditioning.} Verify the sign
truncation and the anatomy selection separately, since they are the
section's two measured claims and they differ in kind. For the first,
estimate $\PP(x_T<0\mid T_{\text{death}}>T)$ by the particle filter and
compare with the unconditioned $\tfrac12$ of
Theorem~\ref{tc:thm:half}; the prediction is an abrupt collapse, to
order $10^{-3}$ within a few multiples of the intrinsic timescale. For
the second, place a prior on the phase --- $\rho\sim
\mathrm{Unif}(-1,3)$ serves --- and track the conditioned
$\PP(\rho>0)$ against horizon; the prediction is a \emph{gradual} rise
from the prior together with steady truncation of the negative tail. A
selection on $\rho$ as abrupt as the selection on $\operatorname{sign}
x$, or an absence of selection on $\rho$ altogether, would each
contradict \S\ref{gr1:sub:anatomy}.

\emph{(iii) Exponent robustness.} Estimate the survival tail for a grid
of $(\epsilon,n,\Gamma,\rho)$ and fit the local slope. Prediction:
$-\tfrac14$ throughout, with no dependence on any of the four ---
this is Theorem~\ref{tc:thm:survives} and the section's central
falsifiable claim.

\emph{(iv) The gain profile.} Compute $\Psi'$ from
\eqref{tc:eq:Psiprime} and check the limiting values of
Theorem~\ref{tc:thm:gain}(iii); then test
Proposition~\ref{tc:prop:mingain} by regressing $\log\min\Psi'$ on
$\rho b+k$ across a grid of $\rho$. The predicted slope is
$-1/(1+\rho)$, the
predicted intercept $\log\kappa(\rho)$, and the minimiser is at
\eqref{tc:eq:xstar}; all three are sharp, making this the most stringent
test in the section.

\emph{Falsifiers.} A fitted exponent differing from $-\tfrac14$, or
depending on $\Gamma$ or $\rho$ (contradicting
Theorem~\ref{tc:thm:survives}); a gain reaching a floor of $1$ rather
than $0$ on the plateau (contradicting
Remark~\ref{tc:rem:speccorr}(b)); an absence of anatomy selection under
conditioning (contradicting Theorem~\ref{gr1:thm:trapped}, which makes
adverse anatomy eventually fatal and hence selectable).


% =====================================================================
\PrintRefs

\end{document}
